% ============================================================
%  R. Nielsen, "Grundideernes Logik. II." [The Logic of the Fundamental Ideas, II]
%  Kjøbenhavn: Forlagt af den Gyldendalske Boghandel (F. Hegel);
%  Trykt hos J. H. Schultz, 1866.
%  TRANSLATION (English).
%  Source of truth: transcription.tex (Danish), transcribed in full and diffed
%  batch by batch against the BL scan's embedded text layer 2026-09-22 (see its
%  ACCURACY header). This file mirrors it 1:1: every printed page carries
%  \opage{N} at the corresponding point in the English; the front matter
%  carries roman \opage{[III]}..\opage{XVI}.
%  Conventions: ../../../TRANSLATION-PLAYBOOK.md, and the terminology of the
%  Volume I translation (texts/nielsen/grundideernes-logik/translation.tex),
%  which this volume continues -- Vol. II's page 1 opens under Vol. I's
%  lettering ("B. Logiske Functioner").
%  Latin, Greek, French and German quotations kept as printed; Danish work
%  titles in citations kept in Danish. Printer's errors in the Danish are not
%  imitated in the English, and the transcription's % PRINTER: and % UNSURE:
%  notes are not carried over; the book's own errata are applied and flagged
%  at the site by % Errata: (list below).
%  The Hebrew of the p. 55 footnote is set as Hebrew, with cjhebrew, as in
%  transcription.tex; an editorial English gloss follows it in brackets. The
%  pointing is the transcription's reading: the scan is 1-bit and the type has
%  broken up, so the consonants and word order are certain, the points are not.
%
%  ERRATA printed at p. VI ("Rettelser"), applied in this translation:
%   p. 79  l.15 f.o.  udvikes      -> udvikles
%   p. 125 l. 4 f.o.  en           -> er
%   p. 130 l.11 f.n.  des          -> das
%   p. 139 l. 7 f.n.  I den        -> 3) I den
%   p. 194 l.12 f.n.  høist        -> høiest
%   p. 195 l. 6 f.n.  belyste      -> belyse
%   p. 382 l. 7 f.n.  καλοκαγάθία  -> καλοκἀγαθία   [the 1866 list keys this
%                    correction to "282"; the word is on printed p. 382, and the
%                    entry sits out of order between the 315 and 398 entries]
%   p. 305 l. 2 f.n.  De           -> Det
%   p. 315 l.10 f.n.  Syllogistisk -> 1) Syllogistik
%   p. 398 l.10 f.n.  realiset     -> realiseret
%   p. 400 l.10 f.o.  blot Række   -> Række blot
% ============================================================
\documentclass[12pt]{book}
\usepackage[utf8]{inputenc}
\usepackage[T1]{fontenc}
\usepackage[english]{babel}
\usepackage[protrusion=true,expansion=true]{microtype}
\usepackage[margin=1.4in]{geometry}
\usepackage{setspace}
\usepackage{amsmath}
\usepackage{libertinus}
\usepackage{libertinust1math}
\usepackage{textalpha}
\usepackage{cjhebrew}
\usepackage{graphicx}
\DeclareUnicodeCharacter{0254}{\reflectbox{c}}
\usepackage{enumitem}
\usepackage{fancyhdr}
\usepackage[colorlinks=true,linkcolor=black,urlcolor=black,
  pdftitle={The Logic of the Fundamental Ideas II},
  pdfauthor={Rasmus Nielsen}]{hyperref}
\setlength{\headheight}{15pt}
\pagestyle{fancy}
\fancyhf{}
\fancyhead[LE,RO]{\thepage}
\fancyhead[RE]{\textit{The Logic of the Fundamental Ideas}}
\fancyhead[LO]{\textit{\leftmark}}
\setlength{\parindent}{1.5em}
\setlength{\parskip}{0pt}
\setstretch{1.3}
\title{\textbf{The Logic of the Fundamental Ideas}\\[4pt]\large II.\\[8pt]
  \normalsize (\textit{Grundideernes Logik}, translated from the Danish)}
\author{R.\ Nielsen}
\date{Copenhagen: Published by the Gyldendal Bookshop (F.\ Hegel).\\[2pt]
  \small Printed by J.\ H.\ Schultz. 1866.}
\usepackage{marginnote}
\newcommand{\opage}[1]{\marginnote{\footnotesize\textit{[#1]}}}
\newcommand{\apage}[1]{\marginnote{\footnotesize\textit{[#1]}}}
% \tocentry is used by the translated Indhold (the book's own analytic contents).
\newcommand{\tocentry}[3]{\par\noindent\hangindent=\dimexpr#1em+2.2em\relax\hangafter=1\hspace*{#1em}#2\dotfill\ #3\par}
\begin{document}
\maketitle
\tableofcontents

% FRONT MATTER (printed roman folios III--XVI), English.
% Three pieces: the book's own analytic table of contents (III--V), the
% errata (VI), and the unheaded, dated and signed preface (VII--XVI), whose
% \chapter* heading is editorial, as in the Danish.

% --- p. III ---
\chapter*{Contents.}
\addcontentsline{toc}{chapter}{Contents}
\opage{[III]}%
\begin{center}---\end{center}
{\setlength{\parskip}{2pt}
\hfill Page.\par
\tocentry{0}{B. \textbf{Logical Functions}}{1.}
\tocentry{1.5}{a. Functions of the Judgment: Formal Functions}{17.}
\tocentry{3}{α. The Form of the Judgment: Functions of the Word of Knowledge}{19.}
\tocentry{4.5}{αα) Language-forming Functions}{20.}
\tocentry{6.5}{Word-forming Functions}{30.}
\tocentry{6.5}{Form-forming Functions}{35.}
% Errata: Contents, p. III --- the print gives this page number as "42", without the closing full stop every other entry carries; set here with the stop.
\tocentry{4.5}{ββ) Thought-forming Functions}{42.}
\tocentry{4.5}{γγ) Idea-forming Functions}{70.}
\tocentry{6.5}{Mythical Functions}{73.}
\tocentry{6.5}{Gnostic Functions}{80.}
\tocentry{6.5}{Theosophical Functions}{86.}
\tocentry{6.5}{1) The Word and God in ambiguous Unity}{89.}
\tocentry{6.5}{2) The Word and God in ambiguous Opposition}{90.}
\tocentry{6.5}{3) The Word and God in developed Unity-in-Opposition}{92.}
\tocentry{3}{β. Judgment and Object: Functions of the Word of Power}{95.}
\tocentry{4.5}{αα) Knowledge of Judgment and Object: theoretical Functions.}{101.}
% Errata: Contents, p. III --- the print has "Porallelisme" for "Parallelisme"; translated as corrected.
\tocentry{6.5}{1) Parallelism of Judgment and Object: supposing Functions}{103.}
\tocentry{6.5}{2) The Object determines the Judgment: abstracting Functions}{108.}
\tocentry{6.5}{3) The Judgment determines the Object: determining Functions}{112.}
\tocentry{4.5}{ββ) Power in Judgment and Object: practical Functions.}{117.}
\tocentry{6.5}{1) In the logical Judgment the Object is preformed: anticipating Functions}{122.}
\tocentry{6.5}{2) With the antilogical Judgment the Object is realized: executive Functions}{128.}
\tocentry{6.5}{3) In the logical-antilogical Judgment the Object is explained: explicative Functions}{139.}
% --- p. IV ---
\opage{IV}\hfill Page.\par
\tocentry{4.5}{γγ) The Idea in Judgment and Object: ideal Power-Functions.}{147.}
\tocentry{6.5}{1) The Ideality of the Power-Word: imperative Functions}{148.}
\tocentry{6.5}{2) The Reality of the Power-Word: indicative Functions}{157.}
\tocentry{6.5}{3) The categorical Power-Language: interpreting Functions}{166.}
\tocentry{3}{γ. The Validity of the Judgment: the Modality of the Functions}{176.}
\tocentry{4.5}{αα) The assertoric Judgment: Functions of immediate Determinateness}{180.}
\tocentry{6.5}{1) The antilogical Element in the assertoric Judgment: the objective Side of Modality}{182.}
\tocentry{6.5}{2) The logical Moment in the assertoric Judgment: the subjective Side of Modality}{191.}
\tocentry{6.5}{3) The logical-antilogical Essentiality of the assertoric Judgment: two-sided Modality}{201.}
\tocentry{4.5}{ββ) The problematic Judgment: Functions of sublated Determinateness}{210.}
% lower-case "the" as printed (the other entries begin with a capital)
\tocentry{6.5}{1) the skeptical Possibility: the subjective Side of Modality}{212.}
\tocentry{6.5}{2) The dogmatic Possibility: the objective Side of Modality}{222.}
\tocentry{6.5}{3) The Ontology of the Possibilities: two-sided Modality.}{234.}
\tocentry{4.5}{γγ) The apodictic Judgment: Functions of reflected Determinateness}{247.}
\tocentry{6.5}{1) The Logical-Apodictic: essential Necessity.}{253.}
\tocentry{6.5}{2) The Antilogical-Apodictic: actual Necessity.}{264.}
\tocentry{6.5}{3) The Rational-Apodictic: absolute Necessity.}{275.}
\tocentry{1.5}{b. Functions of the Syllogism: Modal Functions}{282.}
\tocentry{3}{α. Subjective Syllogism: Modal Functions of Knowledge}{293.}
\tocentry{4.5}{αα) Immediate Syllogism: analytic Modal Functions.}{294.}
\tocentry{6.5}{1) Schematism: the Schema of the so-called immediate Syllogisms}{299.}
\tocentry{6.5}{2) The analytic Modal Functions}{302.}
\tocentry{6.5}{3) The Transition of the immediate Syllogism into the mediate}{309.}
\tocentry{4.5}{ββ) The mediate Syllogism: synthetic Modal Functions}{312.}
\tocentry{6.5}{1) Syllogistic}{315.}
\tocentry{6.5}{2) Deduction of syllogistic Forms}{319.}
% --- p. V ---
\opage{V}%
\tocentry{6.5}{3) The Dissolution of the Syllogistic}{324.}
\tocentry{4.5}{γγ) The immanent Syllogism: dialectical Modal Functions}{329.}
\tocentry{6.5}{1) The immanent Syllogism as Form of Subjectivity itself}{330.}
\tocentry{6.5}{2) The immanent Syllogism as Form of the Content of Subjectivity}{344.}
\tocentry{8.5}{Categorical Immanence-Syllogisms: the Problem of Subsumption}{347.}
\tocentry{8.5}{Disjunctive Immanence-Syllogisms: the Problem of Disjunction}{351.}
\tocentry{8.5}{The Totalization of the Syllogism-Systems: the Problem of Individuation}{356.}
\tocentry{6.5}{3) The immanent Syllogism as Form of free Necessity}{363.}
\tocentry{3}{β. Objective Syllogism: Modal Functions of Power}{384.}
\tocentry{4.5}{αα) The Formal Concept of the objective Syllogism}{390.}
\tocentry{4.5}{ββ) The Word of Power in the objective Syllogism}{408.}
\tocentry{6.5}{1) In the Word of Power Gravity is an objective Syllogism.}{409.}
\tocentry{6.5}{2) In the Word of Power Affinity is an objective Syllogism}{414.}
\tocentry{6.5}{3) In the Word of Power the organic Process is an objective Syllogism}{420.}
\tocentry{4.5}{γγ) The objective Syllogism and the categorical Power-Language.}{424.}
\tocentry{6.5}{1) The mechanical Totality in categorical Meaning}{426.}
\tocentry{6.5}{2) The second Fundamental Totality of Objectification determined by the categorical Power-Language}{429.}
\tocentry{6.5}{3) The third Fundamental Totality of Objectification determined by the categorical Power-Language}{432.}
\tocentry{3}{γ. Idea-Syllogism: Modal Functions of Truth}{434.}
\tocentry{4.5}{αα) Ideality-Syllogism: the Pre-existence of the Ideas}{435.}
\tocentry{4.5}{ββ) Reality-Syllogism: the Real Existence of the Ideas}{450.}
\tocentry{4.5}{γγ) Conclusion: the Completion of Idea-Existence}{456.}
}
% the printed double rule closes the Contents
\begin{center}---\end{center}

% --- p. VI ---
% The entries are kept in Danish because they correct Danish spellings; the corrections themselves are applied in the body of this translation.
\chapter*{Corrigenda.}
\addcontentsline{toc}{chapter}{Corrigenda}
\opage{[VI]}%
\begin{center}---\end{center}
{\setlength{\parindent}{0pt}
p.\ 79, l.\ 15 from top: udvikes, read: udvikles.\\
--- 125, --- 4 from top: en, read: er.\\
--- 130, --- 11 from bottom: des, read: das.\\
% the print has "læs;" for "læs:" in this entry; set here with the colon
--- 139, --- 7 from bottom: I den, read: 3) I den.\\
--- 194, --- 12 from bottom: høist, read: høiest.\\
--- 195, --- 6 from bottom: belyste, read: belyse.\\
--- 305, --- 2 from bottom: De, read: Det.\\
--- 315, --- 10 from bottom: Syllogistisk, read: 1) Syllogistik.\\
% the two Greek forms are given exactly as printed: the faulty reading carries two acutes, the correction a coronis on the second alpha
--- 282, --- 7 from bottom: καλοκαγάθία, read: καλοκἀγαθία.\\
--- 398, --- 10 from bottom: realiset, read: realiseret.\\
--- 400, --- 10 from top: blot Række, read: Række blot.\par}
\begin{center}---\end{center}
\begin{center}{\bfseries Corrigenda to The Logic of the Fundamental Ideas I.}\end{center}
{\setlength{\parindent}{0pt}
p.\ 327, l.\ 14 from bottom: frembringes, read: forbruges.\\
--- 328, --- 14 from top: (H O), read: (H\textsubscript{2} O).\\
--- 328, --- 8 from bottom: denne Kjendsgjerning, read: dette som Kjendsgjerning.\\
--- 239, --- 17 from top: dannede, read: dannende.\par}
\begin{center}---\end{center}

% --- p. VII ---
\chapter*{[Introduction]}
\addcontentsline{toc}{chapter}{[Introduction]}
\opage{[VII]}The Logic of the Fundamental Ideas is an attempt in pure philosophy, an attempt at a pronounced theism, a scientific protest against all theology.

\emph{Pure philosophy.} The task is to free thinking from extraneous representations, fancies and sense-images, so that it may be put in a position to develop concepts and, in the development of the concepts, to cognize and hold fast the necessary forms of the Idea. When, however, it is said, according to a usage once accepted, that in logic it is thinking alone as thinking that accomplishes everything, then here as in so many other cases it holds that: \textit{a potiori fit denominatio.} So-called pure thinking is, when taken by itself, an abstraction which, despite speculative sorceries with pure being, would not be able to stir from the spot without the help of intuition. It is for that reason a misunderstanding when ``pure philosophy'' is to be restricted to a simple grinding out of the old tune on being and thinking and thinking and being. Many and heterogeneous phenomena of spirit may be introduced into pure philosophy, and examples chosen from the most diverse spheres of life and of science alike; what everything essentially depends on in this respect is an observance of the following three fundamental demands: 1) The abstraction must be carried through, so that the concept together with its consequences can become sufficiently clear
% --- p. VIII ---
\opage{VIII}and transparent at every stage of development; 2) The critique of the phenomena of spirit in question and the working-out of the examples chosen must be so laid out that only the logical, and what serves in a peculiar degree to render ontological problems and their solution distinct, comes into consideration; 3) The development must advance in essential continuity, so that what follows can be seen, step by step, to proceed with logical necessity from what precedes it: only in a continuous stream can the Idea be present everywhere in the movement.

\emph{Pronounced theism.} Although a logic of the fundamental ideas can in no section coincide immediately with the philosophy of religion, it is nevertheless quite necessary that it --- unless it should be able to help itself with an atheistic principle --- must develop the logical form of the idea of God as well as of the idea of the world, must logically determine the relation between these ideas and thus furnish the foundation for a corresponding philosophy of religion. The impossibility of giving either a logic of the Idea or a logic of the ideas without at the same time indicating the fundamental features of a philosophical world-view must surely be evident to anyone who has any concept of the matter. But by what, then, is such a philosophical world-view actually determined? Psychologically seen, it might almost look as if the philosophizer had a personal view ready in advance, and then, instead of adjusting his philosophical world-view to the logic, on the contrary accommodated the logic to his personal view. In this way an atheist would lay out an atheistic, a pantheist a pantheistic, a deist a deistic, and a theist a theistic logic.

That the personal view of the philosophizer must in fact become a factor in the peculiar manner in which he apprehends and determines the logical essentiality of the Idea cannot
% --- p. IX ---
\opage{IX}be denied; but philosophy would not be a science if it were not able to procure for itself a quite exact and incontrovertible standard by which the greater or lesser scientific validity of the heterogeneous, mutually conflicting logical systems can be judged. In the testing of a logical system everything depends chiefly on investigating how its principle, beginning and method are originally determined. If the given system really deserves the name of a system, then a glance at the plan as a whole, together with a thorough acquaintance with the beginning and the first sections, will ordinarily be sufficient to decide to which world-view such a system must necessarily lead. Those who know the Hegelian philosophy, e.g., will be able to see already from the first part of the Logic, from the first sections of the first division on being, and especially from the beginning with pure being, that the world-view answering to such a logic must necessarily become a logical pantheism, a panlogism, which only by mystifications of ``the concept,'' the pure concept, the logical Proteus of the immanent method, will be able to give itself a theistic tinge.

That The Logic of the Fundamental Ideas, in the two attempts now before us, is and must be theism, unambiguous theism, in plan, principle and method, can likewise be seen from the beginning itself and the first sections. If it is true that philosophy cannot begin with a unity of thinking and being unless it expressly begins with knowledge; if it is true that the principle of knowledge must absolutely coincide with the principle of subjectivity; if it is true that subjectivity must be principle, not merely for itself, but for all objectivity, and can be the former only by virtue of absolute knowledge, the latter by virtue of absolute power, and both by virtue of that unity-in-opposition of knowledge and power which alone gives absolute truth: then it will already from this
% --- p. X ---
\opage{X}be clear that the logical determination of the fundamental relation between the idea of God and the idea of the world must become theistic.

\emph{Scientific protest against all theology.} That an atheistic philosophy protests against all theology is in order. ``But,'' one might say, ``how is it possible that a philosophy which announces itself as an attempt at a pronounced theism, which acknowledges the idea of God and consequently, from its standpoint, must set about developing the concept of God, can present itself as a scientific protest, not merely against theology in one determinate form, but unconditionally against all theology? A scientific development of the concept of God must, after all, always coincide with an objective doctrine of God. Since now any doctrine whatever of God --- one need only recall the old orthodox fourfold division: \textit{theologia accipitur sensu quadruplici: generalissime pro quavis de Deo doctrina, generaliter pro theologia vera, specialiter pro theologia revelata, specialissime pro doctrina de Deo uno et trino} --- must be called theology, the question is: how is it possible that a philosophy should be able to set up a doctrine of God which yet, in order not to be theology, dare not be any doctrine of God?''

In answering this question we prefer to begin with a reduction of the fourfold form of theology, comprehending the first two definitions under the designation: natural theology (\textit{theologia naturalis}), the last two under the category: revealed theology (\textit{theologia revelata}).

That natural theology, which is originally an abstraction from a positive church doctrine, must be very different from true philosophy is easy to see. Natural theology, namely, begins in its own theological fashion with an immediate representation of God. In such an immediate representation God himself is then apprehended as a supramundane objective being with certain determinate properties. The task is: to prove that such a
% --- p. XI ---
\opage{XI}supramundane objective being with such properties really exists, that this being really created the world, and so forth. With such an immediate representation of God, a representation that makes God into an object outward to contemplation, philosophy, if it is truly to be science, cannot possibly begin. Instead of getting involved with immediate representations of God and bungling with argumentative proofs of God's existence, philosophy begins by presupposing what is not something beyond or absent, but something essentially present for common knowledge, namely existence itself with its infinitely rich content. This content of knowledge is then determined, according to its content-rich unity with itself, as the Absolute, and the task becomes to develop and clarify the forms of knowledge adequate to the absolute content of being.

This task is scientific, for such forms can be developed, revised and corrected with an exactness which in dialectic is not inferior to what is exact in mathematics. For the development of the forms the method is so laid out that the content streams in of itself, as though under a pressure-head of the Idea: if only the form is true, then every doubt about the reality of the content is thereby at once removed. As the perfect form for the absolute content the concept of God is then developed; in the unity of the developed concept with its content the idea of God is cognized; in the idea of God, God is cognized; but this cognition is not theological.

As for revealed theology next, the ecclesiastically positive theology, in particular ``speculative dogmatics,'' the author of the present writing has already expressed his view so often and so unreservedly that he can here restrict himself to a stereotyped repetition of what he, in order to concentrate attention on the first thing needful, must regard as decisive points.

The first thing, then, is to indicate what is theological in
% --- p. XII ---
\opage{XII}theology by a conceptual determination whose validity every expert must concede. We hold fast to the following criterion: ``What is theological in theology is recognized by this, that the cause from which all effects in the realm of nature as in that of spirit are derived and explained, partly mediately, partly immediately, is sought not in the nature of things or the essence of reason, but in the absolute will, in the Almighty who creates out of nothing. But even the purest, most acute and most consistent application of such a principle (occasionalism, Malebranche) makes science unscientific.''

In consequence of the conceptual determination set up, the second thing is how the account between theology and philosophy is to be settled. For philosophy's part we choose the following settlement: In the Middle Ages it was theology that attacked philosophy, because philosophy was not believing enough; in our own day it is conversely philosophy that attacks theology, because theology is not scientific enough. It is then necessary in this respect that a sharp distinction be made between scientific in the wider and in the narrower sense. ``Scientific in the wider sense is everything that becomes an object of science's treatment: what is scientific lies in the treatment; scientific in the narrower sense, on the other hand, is only that whose validity science discerns, whose truth it can judge, test and acknowledge.''

By exploiting the ambiguity that lies in the category ``scientific,'' the doctrine of faith sometimes succeeds in giving itself the appearance of science; it is then the task of the theory of science to point out the most important sophisms belonging here.

\emph{First sophism.} ``If faith is not to rest on subjective notions and be abandoned to the whims of individuals, it must be possible to show expressly what one is to believe: faith must have
% --- p. XIII ---
\opage{XIII}its object. To investigate the objects of faith thoroughly is to investigate them scientifically. The doctrine of faith rests on such an investigation; thus the doctrine of faith is a science.''

This sophism gives the matter the appearance that the objects of faith are secured by science because they are tested in science. The matter stands the other way about. The objects of faith dissolve into vapor as soon as science is given leave to apply its critique. It is just as impossible for a scientific critique to acknowledge the gospel of the Risen One as it is impossible for science to acknowledge the myth of Hercules. To acknowledge a miracle by virtue of scientific critique is a contradiction, since the possibility of a miracle makes scientific critique impossible: where the wonder speaks, the understanding must be silent; where the miracle begins, science must end.

\emph{Second sophism.} ``Every distinctive science must necessarily have its presuppositions. Just as logic presupposes and justifies the validity of thinking, and aesthetics the validity of the beautiful, so must the doctrine of faith for its part presuppose and justify the validity of faith and of faith's content. Whoever proves that the doctrine of faith is not a science because it has its peculiar presuppositions proves at the same time that there is no science at all; but here it holds that: \textit{qui nimium probat, nihil probat.} Thus the doctrine of faith is a science.''

This sophism dazzles by the ambiguity that lies in bringing a presupposition along with one. What makes the doctrine of faith unscientific is by no means the circumstance that it brings presuppositions along, but precisely the circumstance that the presuppositions it brings along conflict with science. It presupposes the miracle, it presupposes the absolute validity of subjectivity and of subjective assurance; with these presuppositions
% --- p. XIV ---
\opage{XIV}it has, so to speak, thrown itself upon science and broken the scientific standard. If the consequences are to be scientific, the presuppositions must be so too; on unscientific ground no scientific edifice can be raised: a science of nature is possible, but a science of miracles is impossible.

\emph{Third sophism.} ``The supernatural presuppositions of the doctrine of faith have, each for itself, validity no doubt only for the believer; but among these propositions themselves there is a logical and to that extent objective connection. Since the connection is objective, and the propositions of faith find precisely their true interpretation and deeper grounding in this connection, the grounding is itself objective, and the doctrine of faith thus, by virtue of this grounding, a science.''

What is dazzling in this sophism rests on this, that a logical connection becomes without further ado an objective connection, and this in turn an objective grounding. ``All men are triangles; Cajus is a man, thus Cajus is a triangle. In every triangle the sum of the angles is equal to two right angles. Cajus is a triangle, thus the sum of the angles in Cajus is equal to two right angles.'' Among these propositions there is here undeniably a logical connection; but is there on that account an objective scientific grounding? From the standpoint of life there may be an eternity-difference between the holy and the profane, the evangelical and the mythological; but when once science cannot judge this holy, cannot master this evangelical, the difference is, in science's eyes, as though it were not at all. --- Since the propositions of faith are, each for itself, unconditionally subjective, they are in science just as arbitrary as the most fantastical fancies. If, then, there is a scientific grounding of doctrines of faith, there is also a scientific grounding of doctrines of myth, a scientific grounding of nonsense.

% --- p. XV ---
\opage{XV}\emph{Fourth sophism.} ``The spirit in the doctrine of faith, `the Holy Spirit,' is not merely in the subjective sense a spirit of faith; in the objective sense it is the absolute spirit of truth, the spirit that `searches the deep things of God' and `guides us into all truth.' If it now stands firm that every science presupposes a truth; if it stands firm that the lower truths must always seek their final explanation in a highest truth, then it also stands firm that the spirit which guides us into all truth must be the spirit which guides us into all science; the doctrine of faith is thus not merely a science, but the highest science, that in which all other science is explained.''

What is dazzling in this sophism rests on a double subreption. It is assumed that the spirit which runs through the theologians' objective doctrine of faith really is ``the Holy Spirit''; this assumption is false. Every church communion has, as is well known, a doctrine of faith of its own; the one doctrine of faith stands in conflict with the other; the adherents of the one contradict and condemn the adherents of the other. If the doctrine of faith were essentially a work of ``the Holy Spirit,'' then ``the Holy Spirit'' would have to be at variance with itself; but a spirit that is at variance with itself cannot be the spirit of truth. It is assumed that ``the Holy Spirit'' must in its principle be a scientific spirit. This assumption is false. If ``the Holy Spirit'' were a spirit of science, then Christ's and the apostles' cognition of the eternal truth would have had to be the highest scientific cognition; then Christ and his apostles, in other words, would have had to be the greatest men of science on earth. If the spirit of faith were in truth a spirit of science, then it would surely have to go forward with theology, inasmuch as on the whole it goes forward with scientific rigor. But how does it go? Scientific rigor advances, but theology goes backward.

% --- p. XVI ---
\opage{XVI}Now that the chief points to which a scientific protest against all theology can on the whole be referred have thus been set up, the third thing is to get these points \emph{either scientifically acknowledged or scientifically refuted}.

Those who now suppose that they have ``the testimony of the Spirit'' for theology's really being a science will find in the propositions set up a convenient and, as we hope, a welcome occasion to raise theology's repute and let its light shine, in that they illuminate point for point the errors and misunderstandings of which the author of the present writing, be it in the criterion set up or in the sophisms pointed out, may demonstrably have made himself guilty. What matters, be it noted, is to demonstrate \emph{determinate} errors in the \emph{determinate} propositions; for on vulgar evasive remarks, improvised views, alterations or rehashings of old and long since refuted ways of seeing, and the like, the undersigned will not reflect.

\bigskip
\noindent{\bfseries March 11, 1866.}

\begin{flushright}{\bfseries R.\ Nielsen.}\end{flushright}

% --- p. 1 ---
\begin{center}
{\large\bfseries B.}\\[3pt]
{\Large\bfseries Logical Functions.}
\end{center}
\addcontentsline{toc}{chapter}{B. Logical Functions}

\opage{1}It is an old objection that philosophy hangs upon words, that its explanations are explanations of words, its controversies disputes about words, its wisdom verbiage. This objection is not wholly without ground when one one-sidedly fixes one's attention upon the shadow-sides, upon the barren subtleties, the tasteless aberrations into empty word-cleverness, over which a Johan of Salisbury\footnote{\textit{Nihil est odibilius subtilitate, ubi nihil aliud est quam subtilitas.}}, an Erasmus of Rotterdam, a Baco of Verulam and among us a Holberg have, each in his own age, held proceedings against the clients of speculative knowledge. The objection becomes groundless and empty, on the other hand, when, in blind zeal for the real, one disparages the ideal and imagines that the understanding sees the thing best of all when it overlooks the word. That philosophy cannot possibly enter into such a way of regarding the matter goes without saying. Philosophy is the science of the word in a peculiar sense; it must necessarily hold to the word, so surely as it means to hold to thought and to hold fast the concept.

Or how should it indeed be possible to separate concepts from the corresponding words, developments of the concept from explanations of words? The concept being and the word being, the concept nothing and the word nothing, the concept something and the word something are surely so inseparable that one cannot alter the word without changing the concept, cannot give up the word without giving up the concept. Just as
% --- p. 2 ---
\opage{2}a soul cannot be without a body, so neither can the concept, as the concept determinate in itself, be without a determining word. It is to that extent quite significant that logic has its name from Logos, and that Logos is both word and thought; for if thought were not, then there would be no word, but if the word were not, then neither would there be any determinate thought.

What holds in this respect of philosophy holds in a corresponding way of the special science. It is unthinkable that any real science should really find itself able to cognize the thing otherwise than by means of the word, to set up concepts that did not rest upon words, to give us explanations of the matter that were not at the same time explanations of words. In every science a difference must necessarily be made between a doctrine of things and the things themselves; it is the doctrine of things, not the things themselves, that constitutes the content of the real science. If science could make the doctrine of things one with the things, then one would surely have to find again, e.g., the actual starry heaven in astronomy, the plant world in botany, the animal world in zoology, etc.; but since it is not the object, but the doctrine of the object, that science offers us, it follows of itself that the doctrine must hold to the word. Even in mathematics, where the exact designations do in so many ways take the place of the word, the word is necessary in order to raise the designations to determinations of concepts.

If it is explanations of words when the logician develops the difference between something and other, because the concept can contain nothing other than what lies in the word, then it is an explanation of words also when the mathematician makes a difference between angle and arc, the astronomer between parallax and aberration, the botanist between leaf and branch, the zoologist between paw and hoof; for arc, parallax, branch, hoof are words just as much as something and other. What is to be thought precisely by such words as paw and hoof must, even if we had paws and hoofs in a thousand, indeed ten thousand specimens
% --- p. 3 ---
\opage{3}before us, be explained by other words just as necessarily as it must be explained by other words what lies logically in the words something and other.

That disputes about words are no more to be avoided in the real science can then be illustrated by countless examples; we choose an example from botany. \textit{Taxus baccata} (the yew tree) plainly has its epithet from the fact that its fruit has been taken for a berry (\textit{bacca}). Now is it really a berry, or are we to use another designation? Such a question cannot be dismissed with a remark such as this, that it surely does not turn upon a name, a word, but upon the matter; for an explanation of the matter is here impossible without a corresponding explanation of words. --- That the alleged fruit has a likeness to a berry, is fleshy, rounded, colored and is the result of a flower, is not to be denied; but the alleged fruit is nonetheless no fruit. A fruit is the ovary developed after fertilization; but here there is no ovary. If there is no ovary here, then neither is there any fruit here; if there is no fruit here, then neither is there any berry here, for a berry is a fruit. Is the so-called berry then perhaps an open carpel, such a one as is generally found among the naked-seeded plants? Neither is this assumption tenable. A carpel must necessarily be a leaf of its own kind (\textit{folium sui generis}). But what we have to do with here is no leaf; it is a placenta, a funicle, hence an axis with a development quite analogous to that in the fig, the rose, etc. One sees from this how the dispute about the alleged berry becomes a dispute about words; for berry is a word, fruit is a word, ovary, carpel, placenta, etc., are nothing but words. The dispute about the thing thus becomes a dispute about words; without such a dispute about words a scientific development of the concept would be impossible.

``But,'' it is said, ``there is nonetheless, as regards dispute about words, an essential difference between the special sciences and philosophy. No doubt every special science must express its concepts in
% --- p. 4 ---
\opage{4}words, but the concepts expressed in words become, as the above example also sufficiently shows, tested and secured by indisputable facts to such a degree that in consequence it is not the concept that must conform to the word, but on the contrary the word that must conform to the concept. In philosophy, on the other hand, attention is so one-sidedly directed upon thought, and through thought upon the word, that the thing becomes dependent upon thought and thought upon the word.'' The opposition between philosophy and real science would then, from the standpoint hereby indicated, be apprehended thus: in philosophy it is the word that determines the thing, in the real science it is conversely the thing that determines the word. But how groundless this view is will be shown in what follows.

That things are determined by things through the mediation of words gives an explanation of the matter; that words are determined by words through the mediation of things, an explanation of words. It is now to be investigated how far one is able in a finite manner to hold the explanation of the matter outside the explanation of words and yet satisfy the demand of the concept. As regards the word, philology is that one of the special sciences which stands nearest to philosophy; we choose, then, an example from philology. ``For the consideration of that series of designations of \emph{the relations in the sentence}, those which concern the relations in intuition of the individual substantival representations,'' says Madvig\footnote{Om de grammatikalske Betegnelsers Tilblivelse og Væsen, 1ste Stykke. Indbydelsesskrift til Kjøbenhavns Universitets Fest i Anledning af Hans Majestæt Kongens Fødselsdag den 6 October 1856. p.~20--22.}, ``it is of the most decisive importance to make one's own and to hold fast the cognition that the grammatical designations have nothing whatever to do with the things' own real relations, but solely with those relations in intuition in which they are either referred to one another in contentless generality or show themselves as members in a determinate single image of action or state, from which the relation receives content. The master's relation
% --- p. 5 ---
\opage{5}to the slave is, according to its content, apprehended no less as the opposite of the slave's to the master, because I designate the connection in which the representations \emph{master} and \emph{slave} enter for intuition, when the one is determined by its relation to the other, by one and the same case from both sides: \emph{the master's slave} and \emph{the slave's master}. In grammar the slave is, if one will, just as much the master's owner as the master the slave's; that is to say, grammar has nothing whatever to do with the relation of ownership, which is a real relation, and \emph{possessive form} is quite as incorrect a name as the old Latin case-names, which are fetched from real relations; it brings, like those, one of the countless shapes before the mind which in intuition is subsumed under the relation of connection, and thereby it calls forth, indirectly and unconsciously, this true representation of the form's significance: belonging-together.'' --- With regard to certain German philologists' speculatively confused apprehensions of the significance of case, in particular of the genitive as the case which ``designates something (the representation whose name stands in the genitive) as ``Ursprung oder Bedingung des Werdens eines Andern''\,'',
% The Danish never closes the quotation opened before "designates"; closed here.
there is added in a note: ``Who does not now see, when he is reminded of it, that we are here in the midst of real relations and in the very crudest confusion? In ``des Vaters Sohn'' the father is no doubt Ursprung und Bedingung des Werdens des Sohnes, but in ``des Sohnes Vater'' is the son also Ursprung des Vaters? Perhaps the author would be shrewd enough to answer that the son is Bedingung des Werdens des Vaters, because the latter is father only through him. He would then merely confound the coming-to-be of an essence and that of a name. But how would he get free, e.g., of ``des Vaters Hand.'' In such obscurity one tumbles about, and by such stuff many let themselves be impressed, because there is once for all confusion in this domain and because one lacks the force of abstraction to grasp the matter as it is.''

It is a strong light that with these few but clarifying
% --- p. 6 ---
\opage{6}words falls upon the opposition between grammar and logic. That the grammatical designations have nothing whatever to do with the real relations of things means, then, in other words, that they have nothing to do with the real concepts of things. Explanation of words and explanation of the matter in this way step quite visibly apart from one another. Insofar as, in the talk of the master and the slave, the grammatical is here looked to, the slave is undeniably just as much the master's owner as the master the slave's; and why should he not be so as well? they both belong, after all, to the same part of speech, are both of equal rank, for the slave is surely a substantive just as much as the master. It is otherwise when in the talk the matter is expressly aimed at. In reality the master A is now an actual individual, the slave B likewise; but between the two individuals the real relation of master and slave is altogether contingent. For A as master could just as well be thought to have another slave, and B in the capacity of slave to have another master; indeed, so far as the immediate fact is concerned, it could just as well have been B who was master, and A slave, as the reverse. There is then, as is seen from this, neither in the pure explanation of words nor in the immediate explanation of the matter any question of comprehending proper. From the grammatical standpoint of the explanation of words the concept is lacking, because here every regard to the real relations of things is lacking; combinations such as ``the slave's master,'' ``the coat's servant,'' ``the moon's trousers'' give in that respect equally good sense. From the immediately real standpoint of the explanation of the matter the concept is lacking, because the real relation in its pure facticity lacks what is necessary to the concept, the universally valid: now ``A is master and B slave,'' no! ``now it is B who is master, and A who is slave.'' As the mere individual the slave is no concept, nor as the mere substantive either; if a trace of concept is to show itself here, the thing and the word, the individual and the substantive, must enter into process with one another.
% The Danish sentence lacks its verb ("Individet og Substantivet i Proces"); supplied in the English.

The development of the concept demands a reciprocal determination of the\opage{7}%
% --- p. 7 ---
explanation of words and the explanation of the matter, an objective reflection in which the thing determines the word, and the word in turn the thing: the thing is antilogical, only the word is logical. Insofar as the thing determines the word, the slave is dependent upon the master, not conversely; but insofar as the word determines the thing, the master as master is himself essentially dependent upon the slave. Insofar as the thing determines the word, the father is ``Ursprung und Bedingung des Werdens des Sohnes,'' not conversely; insofar, on the other hand, as the word determines the thing, the son is likewise essentially ``Ursprung des Vaters.'' What actuality parcels out, essence retains in reflected unity; in order to cognize the essential unity of the opposed determinations, the relation of actuality must continually be turned around. That the relation is turned around signifies, in the logical sense, by no means that the coming-to-be of an essence is merely exchanged for the coming-to-be of a name; it means on the contrary that the intuition of an actual is exchanged for the insight into an essential. It is by the effect that the cause not merely gets the name of, but essentially is, cause; it is by the slave that the master not merely gets the name of, but essentially is, master; it is by the son that the father not merely gets the name of, but essentially is, father.

In the reflection of the reversed relations dialectic is seen. Dialectic is essentially explanation of words, a mark of character that is strikingly designated by the Schleiermacherian expression that dialectic is an ``art of language.'' But the external, finite separation of explanation of words and explanation of the matter brings about confusion in science. In the feeling that dialectic is explanation of words, and, by the reversal of the relations, even a teasing explanation of words, one lumps it without more ado together with wordplay, word-splitting, etc. One speaks of the essential and means the actual; when dialectic then, in order to make the essential recognizable, turns the relation of actuality around, the verbally acute man can of course see in the word nothing other than the bare word, and thus protests,
% --- p. 8 ---
\opage{8}without knowing it, against the essential, and then the dialectician is at once to be a word-splitter.

A mitigated form of the same misunderstanding makes itself known in this, that one does indeed concede a dialectic of the formal and a priori concepts, but not of the real, the empirical ones. ``A priori concepts, such as: something and other, essence and seeming, ground and consequence and the like, can easily be brought forth from one another by the application of contradiction. He who posits the ground already tacitly understands the consequence; the ground is thus at bottom both itself and the consequence; but when the ground first becomes a ground proper on the ground of the consequence, then the ground is surely consequence, and the consequence ground. The same thing is not merely ground in one respect and consequence in another, but the two respects at the same time coincide in one respect: to dialecticize the concepts out of one another is, if I may use this expression, the same as to dialecticize \textit{diversi respectus} into one another. As regards the a posteriori concepts, the concepts of experience, on the other hand, it is impossible to dialecticize existence out of the concept; or who would indeed attempt to dialecticize the actual horse out of the horse's concept, or the actual snail out of the snail's concept? But if it is impossible to dialecticize the existent out of its concept, then it is presumably for the same reason also impossible to dialecticize the one concept of existence out of the other, the concept of horse, e.g., out of the concept of fly, and this in turn out of the concept of snail. In the world of experience \textit{diversi respectus} fall outside one another, and where \textit{diversi respectus} fall outside one another, dialectic plays a very subordinate role.''\footnote{Forelæsning over ``Phil.\ Propæd.'' Universitetsaar 1860--61 p.~25.}

The whole misunderstanding is removed when it is seen that a real concept can be formed and developed only on the condition that determinations of the antilogically actual and the logically essential articulate one another. The existing acid, e.g., is
% --- p. 9 ---
\opage{9}no doubt in actuality independent and to that extent indifferent toward the existing base. To want to maintain, from the standpoint of actuality, that it is the acid that makes the base a base, and conversely, just as in the relation of causality it is the cause that makes the effect an effect, and conversely; that the acid is thus itself base, and the base acid, just as the cause is itself effect, and the effect cause, is unreasonable and meaningless: in the antilogically actual \textit{diversi respectus} fall outside one another. But that in the essential sense, hence when the concept is to come along as concept, it must absolutely be the acid that makes the base a base, and conversely, is not merely incontrovertible in reflection, but is even indicated in the phenomenon of actuality by this, that the same substance can in one combination play the role of the acid, in another that of the base: in silicate of alumina the alumina is base, but in aluminate of magnesia it is acid; the acid as silicic acid thus makes the alumina a base, while the base magnesia makes the alumina an acid. Instead of reproaching philosophy with its dialectical explanations of words, one ought, for the benefit of science, at length to learn to see that the dialectical explanations of words are necessary for the all-sided development of the real concepts.

If in truth so finite a separation could be made between the word and the thing that the thing let itself be apprehended as the purely objective, the word as the purely subjective, dialectic would undeniably lose its higher validity; but what indeed is a word's essential significance, if it is not its objectivity, and what sort of thing is it whose objectivity does not let itself be apprehended and determined by a corresponding word? There may indeed be something subjective in usage, and to that extent it rests with ourselves what significance we will attach to a determinate word in a determinate connection. But in all usage there is also something objective; once a word has received a determinate significance, then the significance must be respected. New words can be formed, but not arbitrarily:
% --- p. 10 ---
\opage{10}the reciprocal action between the popular language and scientific language has therefore also long since been acknowledged.\footnote{\textit{Ce sont les opinions du peuple et le point de vûe sous lequel il envisage les objets, qui donnent la forme au Langage. A mesure que le savoir et la politesse se répandent dans une nation, et à proportion de la durée de leur regne ils portent leurs influences dans la Langue. C'est qu' alors le peuple s'approprie plusieurs expressions que les savans ont inventées; comme ceux-ci adoptent souvent à leur tour, des expressions populaires.}

\textit{De l'influence des opinions sur le langage. Par Mr.\ J.\ D. Michaëlis. Traduit de l'Allemand. à Breme. MDCCLXII. p.~9.}}
The thing is indeed independent of the word insofar as the word is an arbitrary designation; but a thing so objective that no word corresponded to its objectivity is unthinkable. The objective would then in this way transform itself into an unnameable, an in itself indeterminable, a something that nevertheless was not something, since something is after all a designation, and the unnameable falls outside every designation.

In order that the word shall be able to be determined by the thing, the thing must in turn be determined by the word; in order that the word's essence can be objective in the thing, the thing's essence must be subjective in the word. The thing and the word are thus in essence one and yet actually different: this is the meaning of dialectic.

Just as the word, dialectically regarded, is one with and yet different from the thing, so it is also one with and different from thought. If the word is posited against the thing, then the thing becomes the objective, the word the subjective; if the word is posited against thought, then the word becomes the objective, thought the subjective. Immediately thought and thing can never be united; only with the word does explanation come, and with explanation an understanding. As a general doctrine of the Logos, a doctrine of the pure Logos, logic is a doctrine of thought only insofar as it is at the same time a doctrine of the word in its essential
% --- p. 11 ---
\opage{11}unity with thought. Yet here, as everywhere, the essential unity is of course grounded upon difference and opposition; the doctrine of the word's unity with thought is precisely a doctrine of the opposition of word and thought. If it is misleading to want to hold thought fast in finite separation from the word, then it is no less misleading to want to identify the word outright with thought. The latter error characterizes speculation, which overlooks the word's immediacy and lays one-sided weight upon the fact that Logos is essentially thought.

That logic, the rational doctrine of the Logos, must from first to last use the word is then a matter of course; but it is one thing to use the word, another to reflect upon the word's use. This reflection presents itself precisely in the present connection, where the task is to develop the logical functions. If we abstract from what is peculiar in the relation between thought, word and thing, then it becomes impossible to subject the logical functions to any special treatment. One must then either presuppose the word's unity with thought and reflect upon subjectivity, or conversely the word's unity with the thing and reflect upon objectivity. In the first case the functions are so entirely absorbed into the universal that a special development of their own inner peculiarity and mutual relation becomes impossible. In the latter case the objective differences become so obtrusive that it cannot be the logical functions themselves, the subjective acts, but the logic of the functions, a system of objective determinations of relation, that comes to development. We have indeed in the doctrine of \emph{reflecting subjectivity} treated a plurality of peculiar forms of reflection: reflection of the understanding, essence-reflection, actuality-reflection with corresponding subdivisions; but these reflections, although from one side they must be apprehended as logical functions, are nevertheless in another respect so entirely fused with the objective as well as the subjective formations of the content that they cannot possibly come forward in their difference from
% --- p. 12 ---
\opage{12}the content itself. This difference, on the other hand, announces itself immediately when reflection is directed upon the word and the word's relation to thought and to the thing.

In the unity-in-difference of the word and thought there is seen a unity-in-difference of reflection and judgment: the reflecting one judges, the judging one reflects; but reflecting and judging are not on that account the same. That reflection, which always hovers between opposed determinations and has its rest in this hovering, must also rest in a hovering between the subjective and the objective, and to that extent be just as ready to give the determination of form an objective as a subjective stamp, is a matter of course. In the mirroring of the members, in the transparency of the moments, reflection is objective; but judgment cannot be objective in this way. In judgment as judgment the judging subjectivity posits an express unity-in-difference of objective and subjective, in that the object of the judgment is indeed objective, but the judgment-act subjective; the unity-in-difference is determined by the word: in the judgment's subject there is seen the unity of the word and the thing, in the predicate the unity of the word and thought, in the copula the unity of thing and thought by means of the word. Only by reflecting upon judgment as judgment can thought reflect upon what is peculiar in its own many-sided activity; but the development of what is peculiar in thought's formal activity coincides precisely with a development of the logical functions. The fundamental function is judgment as judgment: all determinate and determining thinking is essentially a judging; this judging expresses itself not merely in individual judgments, but in combinations of judgments, not merely in determinations of the judgment's formal, but also of its real validity. Judgment passes over with inner necessity into syllogism, syllogism into proof. What is general in the general fundamental function then articulates itself into three distinct fundamental functions: to \emph{judge}, to \emph{syllogize}, to \emph{prove}.

In each of these distinct fundamental functions the unity-in-\opage{13}%
% --- p. 13 ---
opposition of the subjective and the objective is easy to point out; for the unity-in-opposition is determined by the word.

``Formal logic'' takes the matter subjectively; it is right insofar as the functions judging, syllogizing and proving must necessarily be subjective, absolutely subjective activities. But formal logic is wrong insofar as it apprehends the subjective in so one-sided and undialectical a separation from the objective that the empty forms of thought are set up as though they were indifferent toward the content. That such an abstraction can be carried through only by the aid of the word is comprehensible; but in treating thinking as a psychologically conditioned activity of spirit, formal logic at the same time apprehends the word as a merely psychological product, a designation impotent in itself, which falls outside the thing's essence. By such a separation of the thing and the word and, what coincides with this, of the thing and thought, logic becomes idea-less. It is precisely the questionable thing about formal logic that, under the pretense of developing the ``nature and essence'' of thinking, it in reality carries it no further than to a development of the ``nature and essence'' of trivial abstractions, insipid remarks and endless ratiocinations.

``Speculative logic'' takes the matter objectively, indeed so objectively that the objective judgment must judge itself, the objective syllogism syllogize itself, and the objective proof of the method end at last by proving itself. Speculative logic has the immortal merit of having dissolved the finite thoughts in the idea's inner infinity and of having put into force the demand for a judging, syllogizing and proving that oversteps the narrow boundary of the psychological conditions. But when objectivity absorbs the subjective to such a degree that thought becomes word and the word thing, then subjectivity reacts and in turn dissolves the objective to such a degree that the thing becomes word and the word thought. Only a clarification of thought, word and thing in their original unity and opposition will
% --- p. 14 ---
\opage{14}be able to remove the mystification in which speculative logic has wanted to drown the dualism of the in itself subjective and the in itself objective; but such a clarification must then rest upon a special development of the logical functions.

The difference between thing and thought transforms itself in the logical judgment into a difference between subject and predicate. The subject is the singular, the singular in its immediacy, the object of the judgment, i.e. of the judgment-act, hence the thing; the predicate, the universal, namely the universal in reflection, is on the contrary thought, hence the subjective. In the immediate judgment thought now becomes so united with the thing that it looks as though the thing were thought, and thought thing; this comes of thought and the thing being united in the word. In judgments such as: this is a rose, this rose is red, --- the word rose as predicate is surely an expression for the concept, hence for thought, whereas as subject it is a designation of the object, hence of the thing. But according to the meaning the predicate rose must be just as objective as the subject this, and the predicate red just as objective as the subject rose; for the copula \emph{is} holds to objectivity.

It is thus not a part of the judgment, but the whole judgment, that states the objective; but precisely for the same reason it is also not a part of the judgment, but the whole judgment, that rests upon the judging subjectivity. The judging subjectivity determines the object as subject, in order that this subject may through its predicate become objectively determined, i.e. determined in the quality of object. The judgment-act is, seen from this side, an objectification-act. The whole judgment is objective, for the judging subjectivity is determined by the object, and expresses the determinateness categorically, in that it makes the object of the judgment the subject of the inhering predicate. The whole judgment is subjective, for the judging subjectivity must itself determine its object. The thing as thing is in its objective immediacy just as little a subject as the property existing in the thing
% --- p. 15 ---
\opage{15}is in its immediacy predicate: only in objectification does the existing thing become subject by means of the word; only through the word decisive for objectification can the existing property be transformed into predicate, into the universal under which the subject can be subsumed.

Insofar as the judgment, the whole judgment, is by its nature subjective, there is here, then, in the finite sense a possibility that the judging subjectivity, hovering above its object, may in an arbitrary manner connect subject with predicate: here, then, doubt about the judgment's validity can arise. If the doubt is to be removed, the validity must be grounded; this grounding has in turn a subjective and an objective side.

If the grounding is taken from the subjective side, the connection between subject and predicate must rest upon a development of thought, and the development of thought upon the formal connection of the factors in question. For full illumination of this connection it is required that subject and predicate step forth as opposing extremes; between the extremes a middle term is inserted which is so two-sided that it can show itself in one judgment united with the subject, in another with the predicate: as a judgment derived from the judgments thus formed, the validity of the disputed judgment is formally grounded; the grounding is a syllogism-act, and the judgment thereby grounded a syllogism. That the syllogism-act is a logical function, an act of thinking that rests just as originally upon the word as the immediate judgment does, is then a matter that goes without saying. If, as the Hegelians maintain, there could be objective syllogisms, syllogisms which, without any co-operation of a thinking, judging, syllogizing subjectivity, went on in the objective essence of objective reason: then the thoughts in objective reason would surely become judgments that could judge themselves; the self-thinking thoughts, the self-judging judgments, would then have to transform themselves into thinking and judging subjectivities:
% --- p. 16 ---
\opage{16}what a swarm of subjectivities in the objective essence of
objective reason!

If we take the matter from the objective side, the validity of the
judgment must be made good mediately by immediate demonstration.
Under the immediate demonstration the judging subjectivity, and
thereby the act of judgment, the judging function, becomes determined
by the object of the judgment: the judgment is objective. With this
objectivity the judgments in which each of the syllogism's extremes
is connected with the middle term now come forward separately, one
after the other, now coordinate, now in a relation of dependence; yet
always in such a way that the syllogism resulting therefrom is
determined as an objectively valid judgment. By this mode of
syllogizing the immediate is parceled out into separate members; the
demonstration is immediate: the immediate demonstration, mediate
nevertheless through its articulations, is the original mark of the
proof-function. Since the proof-function is, however, subjective as a
function of judgment and syllogism, the mediate cannot remain as a
parceling out and piecing together of the immediate, but must dissolve
itself in reflection. The mediate is here the objective by means of
the immediate; objectivity is recognizable precisely by that
alternation of mediate and immediate to which the demonstration
corresponds. Reflection is the subjective; reflection transforms the
immediate members of mediacy into reflexes in the movement of thought,
into a totality of transparent moments determinative for insight; by
this transformation the demonstration becomes grounding; in the
reflected unity of grounding and demonstration we have the proof.

If we bring the three logical functions together under one point of
view, it is clear that they designate a progressive objectification.
In the judgment the subjective alternates immediately with the
objective; in the syllogism the objective becomes a moment for a
formal development of the subjective; in the proof the subjective
becomes a moment for a development of what in knowledge is ob%
% --- p. 17 ---
\opage{17}jective; this doubling is for each function conditioned by
the word. That actual determinations of objectivity can be reflected
in subjective forms is due to the word as a word of \emph{knowledge};
that subjective determinations of form can pass over into content of
actual existence is due to the word as a word of \emph{power}; only
through the opposition between the word in knowledge and the word in
power, and through this opposition's content-rich unity in the word of
truth, do the logical functions allow themselves to be developed and
determined many-sidedly.

Since the judgment as judgment expressly shows how thought forms
itself by shaping its ideal material in the word, we shall in what
follows treat the functions of judgment under the designation:
\emph{formal functions}; since, further, the syllogism according to
its concept adds no new or peculiar determinations to the judgment in
its singularity, but precisely indicates the manner in which the one
judgment partly presupposes, partly entails and confirms the other, we
apprehend the functions of the syllogism as a whole as \emph{modal
functions}. The peculiar element that the proof as proof introduces
into the form of the syllogism, namely the real moment, the objective
reflex of the word, then leads finally to insight into the articulated
reciprocal determination of the matter and of its exposition; the
manner in which the one-sided, in the conduct of the proof, is
resolved into the many-sided, the incomplete supplemented, what is
defective in the grounding integrated, thus serves to show \emph{the
reality of the modal functions}.

\medskip
\noindent\hspace*{1em}a) \emph{Functions of the Judgment: Formal Functions.}
\addcontentsline{toc}{subsubsection}{a) Functions of the Judgment: Formal Functions}
\medskip

Just as thought is a function of the act of thinking, and conversely,
so the judgment is a function of the act of judgment: the act of
thinking from which the determinate thought results is originally
determined by the resulting thought, and likewise the judgment
determined by the act of judgment is determinative for the act of
judgment itself. For finite thinking and finite judging there must of
course be finite points of departure, finitely determined
presuppositions; but the finite acts of thinking and of judgment
% --- p. 18 ---
\opage{18}are only conditioned expressions of the essential judging
that is identical with self-objectification, of the universal thinking
that is identical with the being of subjectivity. For essential
judging, for infinite thinking, there is just as little a first fixed
point of departure as there is for absolute knowledge or for
ontological subjectivity.

If one separates the individual acts of thought from the principle of
thinking, the particular acts of judgment from the ``power of
judgment'' that expresses itself in all judgments, then one must
either live on in sheer immediacy, and to that extent in relative
thoughtlessness, or in the end run up against a great problem, the
problem that so characteristically torments the Licentiate in a
Student's Adventures: ``You see, my friend! A movement presupposes a
direction. The understanding cannot advance without moving in a
certain line, but before it follows this line it must surely have
thought that line beforehand. Thus one has always already thought
every thought before one thinks it. Every thought that seems to be the
work of a minute thus presupposes a whole eternity. It could almost
drive me mad. How, then, can any thought be born, since it must always
have existed before it comes into existence''.\footnote{Efterladte
Skrifter af Paul M. Møller. First volume. Third edition. p.~120.}

The infinite thought is made finite by the material of finitude over
which it thinks. It is not enough, however, that thought in the finite
has something to think about; it must at the same time have something
to think by, something to think in: that in which and by which
thinking takes place is language. Only when language transforms the
immediate images of sense-perception into representations does it
become possible for reflection to transform the representations into
thoughts; only by such a transformation can thinking become attentive
to itself and to its functions. The word, which by transforming image
into representation serves to de%
% --- p. 19 ---
\opage{19}termine thought, then serves, by converting representation
into image, to determine the object of thought as well. In the
proposition: this is a tree, the word does not primarily awaken a
representation of itself as word, but a representation of the existing
tree, and thus of the object.

Insofar as the judgment as judgment is now apprehended from the
subjective side, the word stands in the service of thought and is from
this standpoint to be apprehended as a word of knowledge; the
subjective essentiality of the judgment is form: all functions
belonging under the form of the judgment are \emph{functions of the
word of knowledge}. Insofar as the judgment is considered from the
objective side, that is, in relation to the object, the word stands in
the service of the object. But that the word stands in the service of
the object presupposes precisely that the object likewise stands in
the service of the word; the word is in this phase a word of power:
all functions objectively corresponding to the relation of judgment
and object are \emph{functions of the word of power}. A complete
development of the content of knowledge and of power in the word of
truth cannot, however, be given in the form of the judgment; as a
particular fundamental function the judgment still has its syllogism
outside itself. The development of the judgment must then end with a
determination of its subjective validity, a preliminary syllogism
which precisely introduces the syllogism itself, in that it leads us
to insight into \emph{the modality of the functions}.

\medskip
\noindent\hspace*{1em}\textit{α)} \emph{The Form of the Judgment: Functions of the Word of Knowledge.}
\addcontentsline{toc}{paragraph}{α) The Form of the Judgment: Functions of the Word of Knowledge}
\medskip

Both that about which judgment is passed and that in which and by
which judgment is passed are words; words are judged about by means of
words; here thought goes into one with language. The question how far
animals have a power of judgment and a ``natural logic'' is most
closely connected with the question how far they have a language
properly so called. What is at issue is a determination of the
relation between animal sense-consciousness and human
self-consciousness. If the difference between the animal's and
% --- p. 20 ---
\opage{20}the human being's consciousness is only a difference of
degree, then the difference between the animal's and the human being's
``natural logic'', between the animal's and the human being's
language, is also only a difference of degree. Those who, trusting in
the perfectibility of animals, take pains to develop the animal's
natural capacity for thought must then at the same time take pains to
develop and perfect the animal's natural language. If the latter
succeeds only in a very restricted sense, then the former endeavor too
will yield but little. The animal's language has sound enough, but no
words; from the most expressive sound in which the animal utters its
momentary condition to the logical word, human in the proper sense,
there is the same distance as from animal sense-consciousness to human
self-consciousness. For a closer understanding of the relation between
thought and language there serves a special development of the formal
functions as \emph{language-forming}, \emph{thought-forming} and
\emph{idea-forming functions}.

\textit{αα)} \emph{Language-forming Functions}. What is language? The
ideal corporeality of thought, the medium for the self-revelation of
subjectivity, the inner essential other in which the light of spirit
can refract its rays. As the first immediate revelation of thought,
language is indeed a system of signs, but not of arbitrary signs; here
it is impossible to begin with a specially conscious separation,
determined by agreement, of the sign itself and the thing signified.
If the formation of language were to begin with arbitrary
designations, then thought would have to exist, prior to every such
designation, as the determinate thought, for only the thought
determinate in itself can arbitrarily choose its sign; but how should
thought determine itself prior to language, when language is word, and
the word itself is thought's first immediate determinateness. In order
to form language, thought must act; in order to act, it must be; to
that extent the activity of thought certainly precedes language:
language is a function of thought. But thought can determine only
insofar as it is itself
% --- p. 21 ---
\opage{21}determined; in order to act determiningly upon language, the
active thought must be determined by language: thought is in this
sense a function of language. As a function of language, while
language is a function of thought, thought is a function of a function
of itself: by this circuit infinity is designated.

Yet we must not forget that the identity of thought and word thus
emphasized has, from the standpoint of philology, and not without a
hint of self-conscious philosophical assurance, been rather sharply
rejected. ``It is sometimes put forward,'' says Madvig,\footnote{Indbydelsesskrift
til Kjøbenhavns Universitets Fest i Anledning af
Hans Majestæt Kongens Fødselsdag, den 18de September 1842.
Om Sprogets Væsen, Udvikling og Liv. p.~5--6.}
``as a deeper conception of language, that it exists not for the sake
of communication but for the sake of thought itself, and that man does
not think except in and with language. What is true in this is that no
human being develops in isolation, but in a community of
representations with others prepared through the existing language,
from which it follows that the word, the mark of community, is
inseparably bound up with the shared representations, accompanies them
as already given, and preserves them precisely in \emph{this} form and
articulation; whereby it must indeed be noted that even for the living
representation the familiar word can sometimes fail. The new
representations which the individual forms and sets apart he will then
strive to give the same validity and place in the series, and yet this
striving first emerges, in the attempt at (oral or written)
communication, as a search for a new word \dots\ Even the writing down
of a new word as a reminder of the representation for oneself \emph{at
another time} has entirely the form of communication. Language loses
nothing of its essential character as a means of communication because
it also becomes a means for a surer and easier mastery of the
representations made common and objectified through communication. The
general
% --- p. 22 ---
\opage{22}feeling that accompanies thinking, that representations can
be uttered in language and one's own representations precisely in
one's own language, is by no means a thinking in and by the word; for
the word is sounding. (All speaking with oneself in actual words is a
consequence of the drive to communicate, or a kind of test of one's
own sureness in the series of representations, or of the latter's
capacity to be uttered in the words one possesses)''.

Two concessions are to be made here at once: 1) ``that no human being
develops in isolation'', and 2) that language is a ``means of
communication''; but the assertion that language is not to be ``for
the sake of thought itself'', that ``the general feeling that
accompanies thinking, that representations can be uttered in
language'', is by no means ``a thinking in and by the word'', is to
such a degree arbitrary and at variance with sound psychology that we
do not see how so distinguished a man as the author of the program
cited could have advanced such an assertion. When the question is how
far thinking does or does not take place in and by the word, one must
here expressly look to articulated thinking, the thinking that
develops itself in distinct thoughts. How, then, does one go about
apprehending a thought distinctly and determinately without
apprehending it in and by the word? One need only reflect on the first
necessary separation between image, representation and thought! An
image (the image of an animal, a human being, a cloud etc.)
consciousness can hold fast without the help of the word, precisely
because the image is an immediate imitation of the object; but image
is one thing, representation another. The images in the circle of
consciousness can be distinguished and recognized in the same
immediate way as the objects in the surrounding world, but the
representations cannot. Precisely as the pictorial is dissolved, the
representations must, in order not to run together, be made
recognizable by signs, non-pictorial signs, that is to say by words.
``Thus the word is a \emph{sign} for a representation that
\emph{could} have received another sign, just as this sound
\emph{could} also have
% --- p. 23 ---
\opage{23}become a sign for another representation''.\footnote{Anførte Program 1842. p.~11.}
That the word as sign is itself non-pictorial, that sound-symbolism
receives its significance precisely from the fact that the word is and
is to be a ``sign without imitation'', on this the author of the
program has also expressed himself clearly enough, in agreement with
Hegel.

What now follows from this? Indeed, from this it follows
incontestably that the mutual separation and connection of the
representations---the non-pictorial representations that pass over
into thoughts---is so originally conditioned by the word that the one
who represents and thinks needs the word just as surely in order to
understand himself as in order to make himself intelligible to others.
It is thus the making distinct of representations and thoughts within
the circle of consciousness, not their communication, on which the
emphasis falls. The hint with which the author of the program
accompanies his assertion that thinking stands in a contingent
relation to language, since ``the feeling that representations can be
uttered in language and one's own representations precisely in one's
own language'' is by no means ``a thinking in and by the word'', is
most singular; for it runs to this, that ``the word is sounding''. How
loudly, then, must the word sound in order to be a word? If it is the
sensuous in the sound that makes the word a word, then the word
uttered in a stentorian voice must thus be a word in a higher sense
than the word that is whispered. With such a conception it does indeed
agree that the loud-sounding word is a more powerful ``means of
communication'' than the faint-sounding one; but where now is the
boundary? If the word that is whispered is still a word, though a
faint word, how faint may a whisper then be without the word
disappearing? ``All speaking with oneself'' must, according to the
view put forward, either be audible, which among educated people is
readily taken for a sign that the speaker is growing somewhat
childish, or else a very soft whisper. To speak with oneself is thus
to whisper to oneself:
% --- p. 24 ---
\opage{24}where there is no whispering there are no words, where there
are no words there is no language, where there is no language there is
no speaking.

But what are we now to say of the one who sits quite silent and reads
in a book? Does such a reader perhaps sit and whisper, or, if he does
not whisper, does not use his voice but only his eyes, what then is it
properly that he reads? Words it naturally cannot be; if, be it noted,
it is to be understood literally that ``the word is sounding'', that
it is the sensuous sound that makes the word a word.

The view set forth in the program is the more singular in that the
Programmatarius himself in the preface (p.~4) has laid especial weight
on ``Hegel's remarks on language, weighty in all their brevity
(Encyclopædie, 2te Ausgabe, § 459),'' and has accorded them an express
recognition. With this recognition, however, there must be something
peculiar; either the recognition cannot apply to the whole of § 459,
or else the philologist cannot quite have understood the philosopher.
For so far is Hegel, in the paragraph cited, from having entered into
the view that the word, in order to be a word, must be sounding, that
on the contrary he will have the sound regarded as something
inessential to the word. Indeed, he goes so far in this respect that
in the very paragraph so commended he even sets written language
higher than spoken language. ``Die Buchstabenschrift'', says Hegel,
``ist an und für sich die intelligentere; in ihr ist das \emph{Wort},
die der Intelligenz eigenthümliche würdigste Art der Aeußerung ihrer
Vorstellungen, zum Bewußtseyn gebracht, zum Gegenstande der Reflexion
gemacht'' \dots\ What Hegel is driving at with this utterance is not
difficult to discover: he wishes, namely, to make the outward-directed
use of language as a ``means of communication'' the inessential, and
the inward-directed application, by which thought is articulated, the
essential. ``Die erlangte Gewohnheit'', he continues, ``tilgt auch
später die Eigenthümlichkeit der Buchstabenschrift, im Interesse des
% --- p. 25 ---
\opage{25}Sehens als ein Umweg durch die Hörbarkeit zu den
Vorstellungen zu erscheinen, und macht sie für uns zur
Hieroglyphenschrift, so daß wir beim Gebrauche derselben die
Vermittlung der Töne nicht im Bewußtseyn vor uns zu haben bedürfen;
Leute dagegen, die eine geringe Gewohnheit des Lesens haben, sprechen
das Gelesene laut vor, um es in seinem Tönen zu verstehen''.

All this now belongs, as was said, to ``Hegel's remarks on language,
weighty in all their brevity (Encyclopædie § 459)''; but that these
remarks contain the most decisive protest against the point of view
assigned by the Programmatarius to the significance of language as a
mere ``means of communication'' is certain. When, then, right at the
beginning of the program it is said: ``It is sometimes put forward as
a deeper conception of language, that it exists not for the sake of
communication but for the sake of thought itself, and that man does
not think except in and with language,'' the attentive reader asks:
what scientific phenomena may these be that are here aimed at? To
judge by the immediate impression, one would suppose that they were
only some obscure minds, fermenting mystics of little rank, to whose
view no further regard could on the whole be paid. To such a
conception one is led partly by the cursoriness of the expression:
``it is sometimes put forward'', partly by the circumstance that the
whole view is dismissed with a couple of remarks without further
investigation or grounding. Yet one sees (p.~6) that it is not merely
visionaries like (p.~8) a Schmidthenner or a K. F. Becker, but even a
Wilhelm Humboldt, who holds to the original unity of word and thought.
Even if it must now be granted that Humboldt, despite his great
services to linguistic research, can in his ``care for the dignity of
language'' here and there be so unclear that for that reason he does
not properly offer firm points of purchase for a motivated criticism,
this can nevertheless, as far as the fundamental question is
concerned, in no way hold of Hegel.

% --- p. 26 ---
\opage{26}Among the thinkers who expressly conceive language in such a
way ``that it exists not for the sake of communication but for the
sake of thought itself,'' Hegel, whose ``weighty remarks on language''
the Programmatarius himself has acknowledged in his preface, stands in
the first rank. The words become, according to Hegel,\footnote{Encyclopædie. Berlin 1845, § 462, p.~349.}
``zu einem vom Gedanken belebten Daseyn. Dies Daseyn ist unseren
Gedanken absolut nothwendig. Wir wissen von unseren Gedanken nur dann,
--- haben nur dann bestimmte, wirkliche Gedanken, wenn wir ihnen die
Form der \emph{Gegenständlichkeit}, des \emph{Unterschiedenseyns} von
unserer \emph{Innerlichkeit}, --- also die Gestalt der
\emph{Aeußerlichkeit} geben, --- und zwar einer \emph{solchen}
Aeußerlichkeit, die zugleich das Gepräge der höchsten
\emph{Innerlichkeit} trägt. Ein so innerliches Aeußerliches ist allein
der \emph{articulirte Ton}, das \emph{Wort}. Ohne Worte denken zu
wollen, --- wie \emph{Mesmer} einmal versucht hat, --- erscheint daher
als eine Unvernunft, die jenen Mann, seiner Versicherung nach, beinahe
zum Wahnsinn geführt hätte. Es ist aber auch lächerlich, das
Gebundenseyn des Gedankens an das Wort für einen Mangel des Ersteren
und für ein Unglück anzusehen; denn, obgleich man gewöhnlich meint,
das \emph{Unaussprechliche} sey gerade das Vortrefflichste, so hat
diese von der Eitelkeit gehegte Meinung doch gar keinen Grund, da das
Unaussprechliche in Wahrheit nur etwas Trübes, Gährendes ist, das
erst, wenn es zu Worte zu kommen vermag, Klarheit gewinnt''.\footnote{In
his work \textit{De l'origine du langage. Paris 1858. Préface.} p.~56
E. Renan has defended the original significance of the word for the
determination of thought against H. Ritter. \textit{M. Ritter
n'attribue au langage qu'un seul rôle, celui de communiquer la pensée;
il méconnaît une autre fonction non moins importante de la parole, qui
est de servir de formule et de limite à la pensée.}}

To explain the origin of language is in other words to explain the
origin of self-consciousness, of human
% --- p. 27 ---
\opage{27}subjectivity. The fundamental function is an ideal
objectification. Only through the objectification of another can
consciousness objectify itself for itself, can it raise itself from
sense-consciousness to self-consciousness. That the animal's
consciousness is only a sense-consciousness presupposes that animal
subjectivity is unable to raise itself above a merely sensuous
objectification, an objectification in which the ideal itself bears
the stamp of sensuous immediacy. This is seen from the fact that the
animal can indeed express itself by a sound corresponding to the
sense-impression, but that a whole system of animal modifications of
sound would nevertheless not constitute a language, since the animally
sensuous sound cannot become a word. In order that the sound may
become a word, it must ``\emph{be torn loose from that}'' ---
sense-consciousness's --- ``\emph{immediacy and be set free by the
human being who is not animally engrossed, but representing}'';\footnote{Madvig.
Indbydelsesskrift til Kjøbenhavns Universitetsfest 1842.
p.~10. The liberation of the sign from the first, so to speak
immediate, immediacy is emphasized in this work in the most emphatic
way, e.g. p.~11, where it is said: ``Many shrink from recognizing in
the word a mere \emph{sign} without a real impress of the essence and
character of the representation; they will not, as they say, let
language sink down into this bare and abstract shape from sensuous,
picturesque fullness and color''. To which the author remarks, ``that
the specific articulations of sound that constitute the word have no
necessary relation to the particular constitution of the
representation''.}
thus: in order that the sound may become a word, the representation of
the object must be separated from the mere sense-image, the
representation of the sensation from the mere sensation, the inner
objectification from the outer. That the inner, ideal objectification
is conditioned by a reciprocal action among those who speak, and that
language to that extent forms itself as a means of communication,
shows only that psychological subjectivity is from the beginning
sensuously turned outward, and is bent inward toward self-reflection
and ideal objectification only by influences coming from without, and
notably from other subjectivities. But from this it
% --- p. 28 ---
\opage{28}by no means follows that the language forming itself in the
course of communication should have significance merely as a means of
communication; no, from this it follows precisely that communication,
essentially regarded, is the condition and presupposition for the
development of the inward-turned consciousness, a development in the
course of which the means of communication is to be apprehended as a
means of objectification.

If communication, and not the inner, ideal objectification, were the
main thing in the formation of language, then the speaker would have
to be supposed capable of clarifying his representations and thoughts
in another way and by other means for himself than for others, namely
for himself without language, for others by the help of language; but
this is demonstrably not the case. Of him who communicates his
thoughts to another in an unclear and bungled language one of two
things holds: either the thoughts in his consciousness have an unclear
and bungled form, or, if they are clear and determinate, he thinks
them in one language and utters them in another. The language in which
the speaker thinks the thoughts is then for him the essential and
superior one; that in which he communicates them, the inessential and
subordinate one. The speaker cannot objectify his representations for
others without at the same time having to objectify them for himself;
but he can well objectify them for himself without objectifying them
for others: language can thus be a means of objectification without
being a means of communication, but not conversely.

The manner in which speech is determined coincides in principle with
the manner in which thought is determined. One speaks in propositions
because one thinks in propositions; the sounding word in speech and
the soundless word in thought are essentially the same word. If the
word ``lightning'', when it sounds, can count as a spoken proposition
``it lightens'', then the same word must also, without sounding, be
able to count as a thought-proposition. The subjectivity that speaks
to another about the phenomenon in question in the word
``lightning'' will, when no one is present, speak to it%
% --- p. 29 ---
\opage{29}self, objectify for itself the same phenomenon by the same
word; only with this difference, that whereas thought in the sounding
word is a speaking, speech in the soundless word is a thinking.
Restricted to the soundless word \emph{lightning} alone, the
thought-proposition is indeed without articulation, the subject goes
into one with the predicate, the substantive with the verb;
nonetheless the difference between subject and predicate, substantive
and verb is in an unconscious way understood. The same holds when the
word sounds as a spoken proposition; without an involuntary separation
of something about which one speaks and that which is said of it, a
separation which the word does not indeed designate but occasions,
there would here be no proposition,\footnote{``That language begins with the
proposition is entirely true in the sense that this is the task set,
toward which everything aims and strives: for man wants to speak, not
to name isolated representations; it is further true that single words
in the beginning involved propositions, in that it was presupposed
that the speaker wanted to say something, not purposelessly to throw
out a name. But when in the most recent time this has been presented
as if language had really begun its formation with the \emph{whole}
proposition, there lies in this the strangest confusion \dots\ The
proposition is whole only when it contains its constituents posited in
their singularity and thereupon comprehended together. If one will now
let language begin with such a proposition, then one assumes that at
the first attempt two separate words came forward, were understood and
held fast, and of these the one as a sign for a concrete
subject-representation (object), and that at the next attempt, the
next proposition, this repeats itself. From here one then turns to the
notion of propositions in which the constituents lay enveloped but
were afterward separated out, as from a root-tuber, that is, to
non-propositions, a monstrosity that could neither be uttered nor
understood nor be thought. For all speech and all understanding
proceed from the fact that the sound which one cannot resolve into
separately significant constituents is a sign for a single
representation, and out of this undivided mass no process can bring
forth separate, entirely different words, or out of many such
propositions always the same word for the same representation; if, on
the contrary, there are separately significant sounds in it, then they
are there as words prior to the proposition, or we fall back into the
absurdity indicated above.'' Madvig. Anførte Program. p.~27--28.
Renan, however, does not yet seem to have become attentive to this
certainly very important point, but lets the matter rest with rather
general reflections, such as: \textit{Dans nos langues modernes, le
sujet, le verbe, ainsi que plusieurs des relations de temps, de modes
et de voix, sont exprimés par des mots isolés et indépendants. Dans
les langues anciennes, au contraire, ces idées sont le plus souvent
accumulées dans un mot unique et exprimées par une flexion. De
l'origine du langage.} p.~154.}
and if there were here no
proposition, then no speech either.

% --- p. 30 ---
\opage{30}As the unity-in-difference of speech and thinking the
proposition is a product of ideal objectification, and the ideality of
the objectification is recognizable by the fact that what is
objectified undergoes a transformation. In propositions such as: the
horse runs, the tree grows green, it is not the sensuous objects, the
horse and the tree, by themselves, nor yet the universal categories,
running and growing green, by themselves, that are objectified; in the
objectification determined by the proposition the horse is the running
thing, and the running thing a horse. The complete proposition has its
separate constituents, but consists itself in their being comprehended
together into unity. The fundamental function of objectification is
thus conditioned by two separate functions: determination of the
members of the proposition and indication of the manner in which the
members are connected. In the determination of the members of the
proposition what essentially matters is word-formation; in the
indication of the mode of connection, form-formation. The
language-forming functions we \emph{therefore} apprehend in this
respect from two points of view, namely as \emph{word-forming} and
\emph{form-forming}, yet in such a way that the one of these
activities incessantly presupposes the other.

\emph{Word-forming Functions}. In order to explain the origin of the
multiplicity of differences of sound that a language requires for the
designation of the manifold phenomena of existence%
% --- p. 31 ---
\opage{31}, one has laid weight especially upon the natural man's keen
sense for sensuous differences and his great gifts for imitation. The
word-forming function thus came, on this way of viewing the matter,
originally to rest upon an imitation, especially of sound-phenomena,
so that the onomatopoetic words had to become the proper root-words
(\textit{τὰ γὰρ ὀνόματα μιμητικά ἐστι.} Aristotle). But however great a
role involuntary imitation may have played in the incipient formation
of words, we must nevertheless grant those philologists to be right
who reduce both the onomatopoetic words and the interjections to
subordinate moments.\footnote{``In any case the word here too first
becomes a word by laying aside this immediacy, precisely that of
designating what sounds so. From this only that determinate and
limited exception is to be made which rests upon the relation of the
linguistic means itself to a peculiar side of sensuous nature, the
sounding in it. The apprehension of the sound-representation itself in
its individual shapes (not as a general category) uttered itself
naturally through the sounds that were the content and object of the
representation, and thus arose the \emph{onomatopoetic} words. They are
so far from belonging to the kernel and foundation of language that,
because they remain standing in that immediate relation, they are in a
manner spurious words, often subjected, even in the otherwise
developed and consolidated language, to playful alterations and
formations''. Madvig. Anførte Program p.~16. But the interjections and
the onomatopoetic imitations would indeed be excellent words, if
language were a mere means of communication. Among the more recent
philosophers Renan in particular has spoken for sound-imitation.
\textit{La langue des premiers hommes ne fut donc, en quelque sorte,
que l'écho de la nature dans la conscience humaine \ldots\ La rupture,
par exemple, pouvait-elle s'exprimer d'une manière plus
pittoresque que par la racine ῥαγ (ῥήγνυμι, ῥήσσω, ῥώξ); sanscrit:
\emph{rug}; celto-breton: \emph{rogan}; ou par sa forme latine \emph{frac};
allemand: \emph{brechen}? \emph{Frem}, \emph{strep}, \emph{strid}, ne sont-ils pas
également la peinture naturelle du bruit dans ses diverses nuances?}
Anf. Skr. p.~136--37.} It is not imi%
% --- p. 32 ---
\opage{32}tation, it is on the contrary that designation by which the
sense-image becomes representation, a designation by which every
imitation falls away, that makes the sound into a word; instead of
being merely \emph{imitative} the word-forming function is actually
\emph{determining}.

That the word-forming function is actually determining must not be
understood as though the self-active subjectivity added something of
its own in the sense that it immediately brought forth ideal
designations for the ideal, peculiar words for relations, names for
pure categories and the like. Such an assumption is directly at
variance with the nature of the matter. It is true that it is not the
sensuous in the sound that makes the word a word; but for its first
immediate determinateness the word cannot dispense with an element of
sensuous punctuality. ``That all formation of language proceeds from
the designation of sensuous representations is not merely the result
to which all investigations into the connection of the significations
of words lead, even in the most developed languages, nor yet a mere
conjecture grounded upon a presupposition about the development of the
first language-forming men and the limits of their circles of
representation, but a proposition that follows directly from the
nature of language and its necessary manner of coming to be. The
confirmation of the community of representations by the help of their
sensuously present object was the starting-point of that whole
fundamental agreement in which language has its root''.\footnote{%
Madvig. Anførte Program. p.~18. How consistently the author holds fast
to and carries through this way of viewing the matter is seen among
other things from the following statement (Program 1856 p.~11). ``When
Humboldt somewhere, struggling with himself, declares that he
nevertheless does not see why there cannot be original words for
relations just as well as for representations, he forgets that the
presupposition of the word is the representation (and originally the
representation of what is sensuously demonstrable) \ldots\ originally,
and in accordance with the first presupposition by which sound is made
into word, everything that emerges in language becomes, by its
emergence, a norm of representation or a pronominal pointing-out of
representation.''} The word-forming function is thus, in designations
for the general, not an
% --- p. 33 ---
\opage{33}independently creating but an explaining, idealizing
activity.

The fact that every word in language has originally had a sensuous
signification has naturally been misunderstood by the sensualists to
such a degree that they have taken subjectivity's a priori
self-activity for an empty fancy. But subjectivity's self-activity in
the formation of language is recognizable precisely by the manner in
which the outward is converted into its corresponding inward, the
sensuous into its non-sensuous, the individual into its universal.
Subjectivity does not come to itself abstractly by itself; it must
begin with its other, but in the apprehension of its other, the
objective, it shows precisely its originally subjective essence. Thus
neither can thinking begin with an abstract determining of itself; it
must begin with its other, its real negation, the sensuous. So long as
the words still have only sensuous signification, thinking is lost in
sense-perception; but the very fact that the sound is stamped into
word and thereby into category proves that thinking has thought with
it even in its other. Only in the degree that the apprehension of the
relations in the outward is articulated by language does a
corresponding articulation of the inward take place. Why the speaker
does not originally choose words of his own for the designation of the
inward and the ideal is, insofar as language is apprehended as a means
of communication, clear from what has preceded; but why does the
thinker not now choose peculiar signs, independent of everything
sensuous, abstractly ideal private signs, in order to designate the
thoughts for himself? Answer: because it is impossible. What indeed
should such abstractly ideal private signs be? Representations of his
own! but then these representations, as was shown
% --- p. 34 ---
\opage{34}above, would themselves have to have their signs. Intuitions
of his own! but what is a peculiar intuition without a peculiar image,
without the semblance of something sensuous. It is thus not merely for
the sake of communication, but for the sake of thinking and of the
thoughts themselves, that words for the sensuous receive a transferred
signification and are employed for the designation of the
non-sensuous.\footnote{%
\textit{Ceux qui ont tiré le langage exclusivement de la sensation se
sont trompés, aussi bien que ceux qui ont assigné aux idées une
origine purement matérielle \ldots\ Le \emph{transport} ou la métaphore
a été de la sorte le grand procédé de la formation du langage.
Une analogie en a entraîné une autre, et ainsi le sens des
mots a voyagé de la manière en apparence la plus capricieuse;
souvent même la signification primitive a disparu, et n'a laissé
subsister que les acceptions dérivées, .. Le verbe \emph{être}, dont M.
Cousin disait hardiment en 1829: ``Je ne connais aucune langue
où le mot français \emph{être} soit exprimé par un correspondant qui
représente une idée sensible''; le verbe \emph{être}, dis je, dans
presque toutes les langues, se tire d'une idée sensible. L'opinion des
philologues qui assignent pour sens premier au verbe hébreu \emph{haia}
ou \emph{hawa} (\emph{être}), celui de \emph{respirer}, et cherchent dans ce mot des
traces d'onomatopée, n'est pas dénuée de vraisemblance. En arabe
et en éthiopien, le verbe \emph{kâna}, qui joue le même rôle, signifie
primitivement \emph{se tenir debout} (\emph{exstare}). \emph{Koum} (\emph{stare})
en hébreu passe aussi dans ses dérivés au sens \emph{d'être}
(\emph{substantia}). Quant aux langues indo-européennes, elles ont
composé leur verbe substantif avec trois verbes différents: 1) \emph{as}
(sanscr. \emph{asmi}, ἐμμί, εἰμί, sum); 2) bhû (φύω, \emph{fui}, allem. \emph{bin},
persan \emph{bouden}); 3) \emph{sthâ} (\emph{stare}, persan \emph{hestem}), devenu
partie du verbe \emph{être}, au moins comme auxiliaire, dans les langues
modernes de l'Inde et dans les langues romanes (stato, été). De
ces trois verbes, le troisième est notoirement un verbe physique
et signifie se \emph{tenir debout}. Le deuxième a eu
très-vraisemblablement le sens primitif de \emph{souffler}. Quant au premier, il
parait se rattacher au pronom de la troisième personne; mais ce
pronom lui-même, quelque abstrait qu'il soit, semble se rapporter à
un sens primitivement concret.} Renan. Anf. Skr. p.~122--130.}

% --- p. 35 ---
\opage{35}\emph{Form-forming Functions.} How intimately the
determination of matter and of form belong together in language is
seen from the formation of words. Already the first articulation of
the sound, by which it becomes the designation of a determinate
corresponding representation, presupposes form-activity. Since it can
moreover be a question how far the separate sound-elements out of
which a proposition is formed are, in their merely elementary
difference, to be apprehended as original ``words'' or rather as
``roots,'' out of which words in the proper sense can be formed, it is
seen from this that the words belonging to determinate classes in the
Indo-European languages must have undergone a twofold
forming.\footnote{%
Die Sätze der sanskritischen Sprachen bestehen aus Wörtern; in
diesen sind die Wurzeln aufgehoben oder enthalten, etwa in der
Weise, wie ein chemisches Element (z. B. Sauerstoff) in einem
chemisch zusammengesetzten Körper (z. B. Wasser), und die Wörter
sind die Glieder des Satzes. Nirgends, auch im Englischen nicht,
treten Wurzeln, als solche, in der Rede auf. Denn es liegt im
Wesen und im Begriff der Wurzel, etwas Vereinzeltes zu sein; in
der Rede aber hat alles Zusammenhang. Die Wurzel ist also
allemal nur ein abstractes Product der Analyse. Sie kann zwar
unter Umständen lautlich unverändert und ohne Zusatz Glied der
Rede werden, und im Chinesischen tritt allerdings die Wurzel in
ihrer vollen Nacktheit in die Rede, dann genügt es aber schon, daß
sie mit andern Gliedern des Satzes, wenn diese auch selbst wiederum
bloß nackte Wurzeln sind, zusammengesprochen werden, um ihr
abstractes Wesen als Wurzel abzulegen und lebendiges Element
der Rede zu werden. Durch das Zusammenfassen zweier oder
mehrerer Wurzeln in einem bestimmten Verhältnisse hört die Wurzel
auf eine solche zu sein. Es ist also ungenau, zu sagen, der
chinesische Satz bestehe aus Wurzeln; denn in den Satz aufgenommen
verschwindet das Wesen der Wurzel. In Bezug auf letztere
unterscheiden sich also das Chinesische und die sanskritischen Sprachen
dadurch von einander, daß in diesen die Wurzel im Worte, in
jenem aber erst im Satze aufgeht. H. Steinthal. Charakteristik
der hauptsächlichsten Typen des Sprachbaues. Berlin 1860. p.~112.}

% --- p. 36 ---
\opage{36}Straightforwardly and directly the form-forming function has
not begun by distinguishing substantives, adjectives, verbs, etc.; but
already in the formation of the proposition, in which each of the
roots received its function, such a distinction was originally
presupposed. ``When it has repeatedly been disputed whether the
substantive or the verb was the older in language, there usually
showed itself in the whole dispute a certain lack of clarity in the
representations. In and for itself the verbal root was just as much a
substantive as the name for any other representation, and so likewise
the original adjectival roots; but since in incipient language one did
not hold fast the activity or the quality as abstract representations
and speak of them, the roots did not come to be employed as
substantives, but became at once verbs or adjectives, in that they
were used assertively; for this is the verb, the name of activity set
as predicate \ldots\ In considering the setting of the root as verb to
a subject-substantive (when this had been formed) one must first come
to clarity about the so-called \textit{copula} between subject and
predicate \ldots\ This \textit{copula} lies, as the act by which the
representations are referred to one another, behind language and is
the presupposition of speech; in language itself it neither is nor can
be directly expressed otherwise than by the connection of speech.
Every sign would itself require a new sign in order that it should be
referred; least of all is the commonest predicate, \emph{is}, a sign
for the \textit{copula}''.\footnote{%
Madvig. Anf. Program 1842. p.~29--30.}

Straightforwardly thought and word can never become commensurable;
their unity is a unity-in-difference. The determinate thought,
indwelling in its determinate expression, is always reflected in a
thinking that hovers above it, and only in this reflection can it
first properly be grasped. Let the signification of a word be ever so
determinate through usage, it will nevertheless always be altered
according to the connection in which the word occurs, and the same
holds of the proposition. What then cannot be attained directly must
be attained
% --- p. 37 ---
\opage{37}indirectly; but in order that the attainment may take place,
there must occur articulations of the indirect by the direct, and
conversely; it is by these articulations that the form-giving
activities are especially to be known. Everywhere that the
thought-forms did not crystallize in a system of unfinished
designations, the satisfaction of the first elementary demands had, on
account of the disproportion, the far too great gulf between the
direct and the indirect, involuntarily to call forth new demands, the
satisfaction of these again new ones, and so forth. ``The task of
language, to present for apprehension and reproduction those
connections of single representations into intuitions in which the
thinking being moves, did not find its full solution in that which
constitutes the first result of the language-forming activity, the
words that designate the representations according to their content or
point them out according to their relation (the naming and the
pronominal words), since the intuiting consciousness contains more
than is expressed by these words and by their apprehension according
to that fundamental presupposition of speech, requiring no further
indication, of the predicative connection into an assertion, and
connects more representations in one intuition than could, on that
presupposition alone, be apprehended in their connection. To effect
the apprehension of that which is thus more in the intuiting
consciousness, the grammatical designations serve (word-order,
auxiliary words, inflection of the naming words)''.\footnote{%
Madvig. Om de grammatikalske Betegnelsers Tilbliven og Væsen,
1ste Stykke. Universitetsprogram. 6te October 1856. p.~7.}

But precisely by reflecting on the relation between the direct and the
indirect in ``the grammatical designations'' and noting ``the
motives'', as the acute philologist himself accurately states
them,\footnote{Anf. Program p.~7--8.} we obtain point for point the
proof that language is inseparable from thinking, and that the
development of language must for that very reason --- not for the sake
of communication --- coincide with the self-development of thinking
% --- p. 38 ---
\opage{38}subjectivity. ``The \emph{motives} which either require or
can occasion grammatical designations may, according to the nature of
the form of intuition, be referred to determinate principal classes,
in that they either (A) concern \emph{the extent of the single
representation and the form} under which it is, in a determinate
function (as substantive, verb, adjective), taken up into the
intuition in relation to the immediate signification of the word (such
as designations for the number of the substantives, for their
apprehension as individually determinate or indeterminate, for the
degrees of the adjectival representation, for the use of the verb in
the passive or in a similarly modified signification), or (B) \emph{the
relation of the representations in connection} into a whole in the
intuition (\emph{relation in the proposition}), which relation is
partly (a) predicative or attributive connection (designations at the
predicate or the attribute according to its belonging to the subject,
\emph{predicative or attributive designations}), partly (b) a relation
of determination of the substantival representations to and in the
assertion, either as predicated about or as entering in the form of
spatial intuition (\emph{special designations of relation}. Cases,
prepositions), or finally (C) \emph{the whole position of the
intuition for consciousness} with respect to actuality, presentness
and connection with other intuitions into a composite intuition
(\emph{the relations of the proposition}; designations for mood and
tense, conjunctions)''.\footnote{%
Madvig. Anf. Program 1856. p.~8.}

It is again the fundamental relation, the psychological-ontological
relation between image, representation and thought, upon which it
essentially depends when the question is what can be expressed
directly and what indirectly. Just as the representation is the image
idealized in the logical sense, so thought in turn is the idealized
representation. The less ideality a representation has, the more
direct the presentation can be; conversely, the more abstractly ideal
the representation is, the more indirectly is it from the beginning
contained in the direct presentation of the concrete. ``He who speaks
of \emph{men} has
% --- p. 39 ---
\opage{39}not a special abstract representation of plurality, but he
has, and demands to have apprehended, the representation \emph{man}
under the form of plurality; he who will speak of a horse's running as
past has not before consciousness an abstract representation of
pastness, nor would he, if he had it, speak of it as he speaks of the
horse and the running; by \emph{naming} pastness he would confuse the
whole meaning; he who says: ``The man lies in the house,'' has not a
representation of \emph{being} in brought into prominence in
consciousness, nor will he set it beside the representation of man and
of lying as a third; but both have the feeling of a beholding of the
named representations under a peculiar form, the form which, when it
is held fast by abstraction and made an object of discussion, will
require a name such as \emph{pastness} or \emph{being} in, and
likewise a demand that the hearer shall reproduce and behold the named
representations (in the first case: ``The horse ran'', the whole
assertion) under the same form''.\footnote{%
Madvig. Anf. Program 1856 p.~11.} --- He who speaks of \emph{men} thus
makes, be it consciously or unconsciously, a difference between what
is to be represented and what is to be thought. The shape \emph{man}
is to be represented, \emph{the plurality} is to be thought: in that
the representation \emph{man} is apprehended ``under the form of
\emph{plurality}'', the direct becomes determined by the indirect,
i.e. the representation determines the thought and the thought the
representation; the reciprocal determination of both is given in and
by language, and cannot be given by any other means. Should anyone
attempt to specify another means, the result is one of two: either the
means is inadequate, and then it has failed, or adequate, and then it
is a language. What is expressed indirectly in speech is just as much
indirect for the speaker himself as for the hearer.

If one attempts to place the separation between thought and word in
this, that the one speaker can have a more developed consciousness
than the other, although they both say the same thing,
% --- p. 40 ---
\opage{40}and thus, as regards the communication, stand on the same
level: then this attempt too will be wrecked on the cliffs of
philology. Let two communicate in the same proposition: ``The horse
ran''; let the one have ``an abstract representation of pastness,''
the other not: is it hereby made good that language as a means of
communication stands in a contingent relation to thought? By no means.
In the first place we have the philologist's express declaration that
he who, in uttering the proposition ``The horse ran,'' had an abstract
representation of pastness would nevertheless not mix it into the
speech, in order not to confuse the meaning. But if he confuses the
meaning by mixing in the irrelevant representation, then he confuses
it not merely for the hearer but just as much for himself. Next, and
this must be emphatically brought out, no one can have an abstract
representation of pastness and hold it fast in an abstraction without
his having, precisely for the sake of the abstraction, a corresponding
word. Pastness is an ideal representation; its essence is negative; it
cannot be intuited, it can only be thought, and the thought is
determined by the imageless sign, by the word.

What we are to see in this connection is that language originally
``exists for the sake of thought itself,'' and ``that man does not
think except in and with language''. Whichever of the motives adduced
for grammatical designations we consider, we nevertheless arrive at
the result that thought is determined by the language which it itself
determines, that thought is, by means of language, a function of a
function of itself.

If we look to designations for ``the extent of the single
representation and its form,'' thus to motive (A), why then is there
not in language, beside the name for the single representation itself,
a name for its extent and a separate name for its form? The
determinations extent and form are, it is said, not to be represented
directly but indirectly, that is to say, they are not to be
represented but thought. Nevertheless such determinations cannot
% --- p. 41 ---
\opage{41}begin in the language-forming consciousness as clear
thoughts; they must on the contrary begin as obscure intimations.
How, then, are the obscure intimations raised up into the light of
consciousness? By this, that as fleeting representations (in this case
accessory representations of a determinate extent, a determinate form)
they attach themselves to a representation of an object; this
attachment is possible only by the word. It is not man as individual
but man as word to which an article can be added; it is not the
individual man but the word man that by a vowel-change passes over
into the plural; but once the individual is exchanged for the word,
the obscure intimation of \emph{determinate} or \emph{indeterminate},
of \emph{single} or \emph{several} is involuntarily there in the
accessory representations. Now in that these accessory representations
accompanying the word of the principal representation are noticed as
essential, they must as indirect representations be determined either
in a direct or in an indirect manner. Only when the designation by
language has secured the indirect representations do they become an
object for abstraction and, through abstraction, for the determining
thinking. Without language there would be no consciousness of a
difference between a determinate and an indeterminate individual,
between singular and plural, between active and passive; without
language there would for thought be neither determinateness nor
indeterminateness, since thought for its determinateness just as
surely presupposes the word as the word for its determinateness
presupposes thought.

If we cast a glance at motive (B), surely no one will deny that the
difference between the subject and its predicate, between the
substantive and its attribute, is a concept-difference, thus a
thought-difference. Without such a difference the corresponding
thought would not be, but without a corresponding language the
difference itself would not be. The thing, antilogical as thing, is
neither subject nor predicate, neither substantive nor attribute; only
in language is there subject and predicate, only in language is there
substantive and attribute, only in language can there be thinking.

As regards finally motive (C), where is it then possible
% --- p. 42 ---
\opage{42}to think ``a position of the intuition for consciousness,''
an intuition composed of intuitions, thus an intuition of intuitions,
without words, without language? A position of the intuition for
consciousness is indeed something quite other than a position of the
object for consciousness. An intuition of an intuition, a
representation of a representation, is not like an image of an image,
a sight of a sight, in sensuous immediacy; here there is a doubling
that rests upon the transformation of accessory representations into
negative determinations of form, determinations of the image-like by
the imageless, of the imageless by the image-like,
reflection-determinations of word by proposition and of proposition by
word, of principal clause by subordinate clause, etc. ``In the power
with which the single members of the thought press forward
simultaneously there lies an incitement to extend the connection of
propositions by attaching several subordinate clauses to the same
principal clause, and subordinate clauses to subordinate clauses, an
incitement that becomes the stronger, the richer the thought-content,
the livelier and more comprehensive the communication, the greater the
need becomes for a sharp delimitation of the single thought by
accessory determinations''.\footnote{%
Madvig. Program 1856. p.~39.} --- He who recognizes this and still
assumes that determinate thoughts can be thought without words,
without propositions, consequently without any regard to the
difference between principal clauses and subordinate clauses, must
really, if his assumption is to count for more than an empty
assertion, undertake to show us philologically how a subjectivity sets
about thinking either a thought without a proposition, or a
proposition without words, or a proposition in words without language.

\textit{ββ)} \emph{Thought-forming Functions.} In that thought works
upon its means of objectification, it essentially works upon itself;
the thought-forming functions cannot then be kept in a rubric of their
own over against the language-forming ones. But thought is different
from the word as that which is designated is different from the
designation; the task of developing a
% --- p. 43 ---
\opage{43}system of designations is thus different from the task of
forming out and clarifying throughout the content reflected in the
designations: with respect to the first task the functions are
language-forming, with respect to the latter thought-forming. It is of
great significance for psychology and logic that philology in our time
has reached the turning-point at which it is able, with full
consciousness, to treat the language-forming functions by themselves,
without the admixture of irrelevant speculations.

``The question \emph{what} and \emph{how much} of thought and
representation lies in the grammatical designations of language'',
says Madvig,\footnote{%
Program 1856. p.~9.} ``and how it lies therein, by what path it is
attached thereto, forces itself forward as so important for the
scientific presentation of grammatical concepts, and its decision is
so essential a presupposition for the validity of considerations and
judgments upon the grammatical excellences and defects of languages,
upon the relation of their grammatical structure to the life of
spirit, that one should expect that for the sake of both these regards
eager work had been done to clear it up, that at the least one had
sought to state it and to answer it quite determinately''. That the
task really admits of being stated and answered determinately the
author of the Program of 1856 has shown by specifying and determining
the motives. But since language ``exists for the sake of thought
itself'', since man ``does not think except in and with language'',
the language-forming functions must just as surely reflect themselves
in the thought-forming ones as these in the former. Here there can be
no doubt \emph{that} it happens; but the question is \emph{how} it
happens.

It is, or has at least formerly been, a fairly widespread opinion that
the ancient classical languages, Greek and Latin, especially Latin,
are distinguished by strict logic, and that the strict logic displays
itself in their grammar. But how much lack of clarity can mix itself
into such assessments philology itself can best of all make
good.\footnote{%
``Sometimes one is referred to the fine nuances and turns of phrase of
which every language has its own, and which, as has been
developed above, testify more to naive receptivity than to consistent
logic; but more commonly appeal is made to the strict consistency with
which the parts of the proposition are joined together in exact
agreement with and dependence upon one another, the predicate
conforming to the subject in person, number, gender, the substantives
ordered in their determinate cases, etc.; in other words, it seems to
be apprehended as especially logical that language observes its own
rules and employs its means of designation each according to its
signification, as if the simpler languages did not do quite the same
within the domain of their system'' \ldots\ ``To this representation of
the greater logicality of the ancient languages there attached itself
next an obscure transference to the languages of other qualities
which one, each according to his own judgment, found in the ancient
peoples and in their literatures, partly because one wished to prove
from the languages' own grammatical constitution what can be proved
only from the great historical conditions of culture, namely that a
knowledge of the ancient languages is of the greatest and most
widespread importance for education, or because one sought
recommendations generally for the study of language that one pursued.
While Latin retained the reputation for a peculiar logical
determinateness, other predicates were sought for Greek. It is
possible that a comparison of the two classical languages with one
another in a purely grammatical respect --- which, moreover, when the
moments are to be rightly and unmixedly weighed, is extremely
difficult --- will show in Greek some preponderance in receptivity to
the influence of accessory motives (perturbations) upon the principal
analogy (see the preface to the German edition of my Greek Syntax.
p.~XI); very great it will hardly be, and an inference from it to the
whole people's spiritual disposition will be the more uncertain, in
that not all the particular influences lying in the conditions of
development can be taken into account''. Madvig. Om de
grammatikalske Betegnelsers Tilblivelse og Væsen. Andet og sidste Stykke.
Universitetsprogram 1857. p.~83--86.}
% --- p. 44 ---
\opage{44}``Language'', says Madvig, ``has not merely a grammatical
\emph{form}, but it has also within the limits of this form a
standpoint, a \emph{condition}, a raw unworked existence and a higher
development, an acquired stylistic capacity, and herein lies both a
\emph{testimony} that there has been more
% --- p. 45 ---
\opage{45}or less working in and with the language, that there has
been spiritual movement among the people, and an \emph{aptitude} for
future work. Here, not in the fundamental form, not even in the
directions and the position that have been given to single members of
the fundamental form in the course of the development, but in its full
and flexible development, flexible not merely in a single direction,
lies the grammatical excellence of a language \ldots\ But whereas it
can be discerned from Greek grammar, especially the syntax, that the
Greeks had and cultivated a detailed communication in speech and
writing, that the members of the people, at the least many of them,
were able to move in large and composite groups of representations,
and whereas it can be seen from metrics, as a part of grammar, that
they had a poetry clothed in sensuously beautiful and manifold
rhythmical forms, there can be drawn from the grammatical structure of
the language no further inference whatever either to the fullness of
content of the communication cultivated in the language or to that
direction of thinking (e.g. toward the ideal or the material, etc.)
which has made itself felt therein. Here there enters, as the far
richer, far more determinate linguistic testimony as to what has taken
place among the people, if one will, as to its view of the world (that
is to say, as to the concepts and representations that have
\emph{at all} been in movement and have come forth among the people),
the stock of words and the stock of significations in the words. It is
here that one notices that not merely discussions such as those
Demosthenes conducted, but also such as those at the head of which
Plato and Aristotle stood, have taken place among the Greeks, that all
the regions of the life of spirit have been traversed by them. The
lexical, the stuff-like, the material in languages, as it is sometimes
called, contains the only image, determinate to a certain degree, of
the people's \emph{direction of spirit}''.\footnote{%
Madvig. Program 1857. p.~88--89.}
% --- p. 46 ---
\opage{46}But however clear, definite and striking this statement may
seem at first glance, there is nevertheless here one point, and that a
principal point, which awakens misgivings. It shall be willingly
conceded that ``from the grammatical structure of a language no further
inference whatever can be drawn either to the fullness of content of the
communication carried on in the language or to that direction of
thinking which has made itself felt therein''; but --- we must ask ---
is it after all admissible to make such a separation of the grammatical
and the lexical, when what is in question is the relation of the
development of language to the life of spirit? It is surely just as
impossible to read off the definite image of a people's spiritual
direction from the lexicon of a language as from its grammar. In both,
the lexical as well as the grammatical, when each of these principal
sides is considered by itself, one sees only a multitude of conditions,
but not their application. Only in the application, in the definite
presentation of a spiritual content, does it show itself what a language
is capable of in a spiritual respect. In the presentation the
grammatical works point for point together with the lexical; it is in
this articulating co-operation that the reciprocal action of
language-forming and thought-forming functions is properly to be
recognized.

That it is not only the lexical, but also the grammatical, which in
every language acts upon the movement of thought, is especially striking
in those cases where one thinks in one language and communicates in
another; but it can moreover be shown literally by experiments with
translations, particularly of logical writings. Let one but try to
translate Hegel's Logic into French, and one will convince oneself that
changes in the lexical also bring with them changes in the grammatical,
that the paraphrases thereby occasioned cause alterations in the
thought, inasmuch as they cause changes in the mode of presentation, and
that the alterations of thought answering to the altered presentation
contain logical deviations which for the sake of the principal thought
must incessantly be corrected. One will fur%
% --- p. 47 ---
\opage{47}ther convince oneself that the difference which, in respect of
the train of thought, is a consequence of the diversity of languages is
of a quite different nature from that which is a consequence of what is
peculiar to authors who think in the same language. One will finally
convince oneself that different grammatical turns of phrase can give the
same meaning without on that account giving, in the logical sense,
exactly the same thought: the German \emph{ist gewesen} and the Danish
\emph{har været} give the same meaning, but not the same thought. ``On
the whole'', says Heiberg,\footnote{\emph{Perseus.} Journal for den
speculative Idee. Nr.~2. 1838. p.~44--45.} ``one will perceive that the
language in which one philosophizes cannot be without influence upon the
philosophizing itself. For language consists of the, so to speak,
embodied categories, so that grammar must be regarded as a collection of
illustrative figures to logic.
Hence Hegel too has categories at which he would not have arrived by
thinking in a foreign language, and of which some have no counterpart in
Danish either''.

There is in the formation and development of languages a difference
which can by no means be derived from definite stages of culture, as
little as from the contingency of a first more or less fortunate grasp
of the matter, but which points to what is peculiar in the nature of a
people. Let one imagine the Hebrew language cultivated as much as one
will, it will nevertheless never acquire the logical stamp that the
Greek has. Hebrew thinking is essentially psychological. The Hebrew's
thinking is in a special sense a speaking with himself; he thinks with
the \emph{heart}.\footnote{\textit{Prenons pour exemple l'hébreu, qui nous
représente un état fort ancien du langage. On sent que le phénomène qui a
servi d'occasion à la création des radicaux de cette langue, et en général
des langues sémitiques, a été presque toujours physique. ``Je conviens, dit
Herder, que le penseur abstrait ne doit pas
trouver la langue hébraique très parfaite; mais sa forme agissante en fait
l'instrument le plus favorable au poëte. Tout en elle nous crie: Je vis, je
me meus, j'agis! je n'ai pas été créée par le penseur abstrait, par le
philosophe profond, mais par les sens, par les passions!} Renan. Anførte
Skr. p.~124.}
Not only the Jewish people's cul%
% --- p. 48 ---
\opage{48}tural history, but the original character of the language too
bears the stamp of the fact that this people's highest thinking was, and
was to be, not a philosophical, but a religious thinking. It is
otherwise with the Greek language; for the
Greek knows that he thinks with the \emph{head}.\footnote{The Greek
knows too what it means that language exists ``for the sake of thought
itself, and that man does not think except in and with language'':
Οὐκοῦν διάνοια μὲν καὶ λόγος ταὐτόν· πλὴν ὁ μὲν ἐντὸς τῆς ψυχῆς πρὸς αὑτὴν
διάλογος ἄνευ φωνῆς γιγνόμενος τοῦτ᾽ αὐτὸ ἡμῖν ἐπωνομάσθη, διάνοια; . . Τὸ
δέ γ᾽ ἀπ᾽ ἐκείνης ῥεῦμα διὰ τοῦ στόματος ἰὸν μετὰ φθόγγου κέκληται λόγος.
Plato. Sophist. 263~E. Cf. Theæt. 189.~E.}
Long, long before there could here be any question of a philosophizing
thinking that was to work upon language, language itself was peculiarly
prepared for philosophical development.

A language's logical endowment is then not to be recognized simply from
the fact that in all immediacy it furnishes means for a development of
thought, for that every language does to a certain degree; but expressly
from the fact that in the course of its development it conditions and
furthers a thinking about thinking, a judging about the judgment, that
it, in other words, is fitted to express pure forms of thought by
definite shadings in definite transitions. Here, then, is an essential
difference between psychologically naive thinking and logical thinking
that is conscious of itself. In both there is a distinguishing and a
discerning, in both a separating and a connecting of representations;
but psychological thinking, which discerns and distinguishes
immediately, involuntarily mixes the contingent with the necessary, the
special with the universal, the image with the concept; logical
thinking, which is attentive to itself, is on the other hand sure of the
abstraction and exercises a constant self-criticism.

% --- p. 49 ---
\opage{49}The purest understanding of its own formal activity that
thought can attain it obtains through a consideration of the forms of
language in which it becomes transparent to itself. When formal logic
makes a difference between ``judging'' and ``non-judging'' propositions,
this distinction can in a way be acknowledged, insofar as the emphasis
falls exclusively upon the form. But now the obscurity shows itself
precisely in this, that logic, before one rightly knows of it, mixes the
logically universal with the logically individual, the abstractly
objective with the abstractly subjective, the form of the judgment in
language with the act of judgment in the judging subjectivity: this
yields a great confusion.

Propositions such as: ``it is raining'', ``he is sneezing'', ``someone
is ringing'', are, so it is said, ``not properly judging;'' ``they are
merely narrating, reporting, stating, without there being any question
of an assuming of one thing about another.'' If one looks at the purely
formal, then this must be conceded, for the form is immediate and gives
no clear mirroring of the logical function; but when it is then added
that such propositions ``will be able to transform themselves into
properly judging ones as soon as a question arises about the correctness
of the utterance or of the statement'', then the confusion is in full
swing. Insofar as the proposition: someone is ringing, gives occasion
for doubt, investigation and discussion, what is looked at here is
solely the content, not the form. One demands of the proposition a
guarantee that the speaker really does think something in what he says;
but such a guarantee the proposition, as logical judgment, cannot
possibly give so long as the judging subjectivity is seen in
psychological arbitrariness. A proposition such as: God is love, can be
uttered by a parrot just as fleetingly and thoughtlessly as the
proposition: someone is ringing. If one now, as regards the latter
proposition, demands a closer determination: \emph{it is so, it really
is so}, that someone is ringing, then one must, as regards the former
proposition, likewise add the assurance: \emph{it is so, it really is
so}, that God is
% --- p. 50 ---
\opage{50}love. But such assurances can in their turn be made in a
thoughtless manner; one must thus assure oneself of the thought in the
given assurances by new assurances: it is in truth so, that it really is
so, that it is so, etc.

The confusion which formal logic brings about in that it, not content
with the objective logical form of the judgment, strays over into the
abstractly subjective and would have in the utterance itself a separate
confirmation that the speaker really does think something in what he
says, that same confusion has then also, though in a reversed manner and
on an opposite ground, forced its way into speculative logic. The
judgments in Hegel are, namely, as we have seen, so objective that out
of sheer objectivity they are on the point of judging themselves and
making themselves subjective; it cannot then be otherwise than that the
judgments which are objective in ``Subjectivity'' must develop and
advance together with the concept that is objective in its subjectivity,
whose content is one with the form. Since Hegel now, as regards the
judgments and their division, has retained the forms of formal logic and
set himself the task of saturating the empty forms with immanent
content, we meet in him too the misunderstanding that the form, instead
of having its logical content in a reflection of pure determinations of
form, must fill itself with a content captured from a richer development
of categories already traversed;\footnote{--- ``Wenn Beispiele von
Prädikaten der Reflexions-Urtheile gegeben werden sollen, so müssen sie von
anderer Art seyn, als für Urtheile des Daseyns. Im Reflexions-Urtheil ist
eigentlich erst ein \emph{bestimmter Inhalt}, d.~h. ein Inhalt überhaupt
vorhanden; denn er ist die in die Identität reflektirte Formbestimmung, als
von der Form, insofern sie unterschiedene Bestimmtheit ist, --- wie sie es
noch als Urtheil ist, unterschieden. Im Urtheil des Daseyns ist der Inhalt
nur ein unmittelbarer, oder abstracter, unbestimmter. --- Als Beispiele von
Reflexions-Urtheilen können daher dienen: Der Mensch ist \emph{sterblich},
die Dinge sind \emph{vergänglich},
dieß Ding ist \emph{nützlich}, \emph{schädlich}; \emph{Härte},
\emph{Elasticität} der Körper, die \emph{Glückseligkeit} u.~s.~f. sind
solche eigenthümliche Prädikate.'' Hegel. Wissenschaft
der Logik. Zweiter Theil. p.~92. --- Further, p.~101--102, where the
discussion is of ``Das Urtheil der Nothwendigkeit'', it reads:
\emph{Das kategorische Urtheil} hat nun eine solche Allgemeinheit zum
Prädikate, an dem das Subjekt seine \emph{immanente} Natur hat. Es ist aber
selbst das erste oder \emph{unmittelbare} Urtheil der Nothwendigkeit; daher
die Bestimmtheit des Subjekts, wodurch es gegen die Gattung oder Art ein
Besonderes oder Einzelnes ist, insofern der Unmittelbarkeit äußerlicher
Existenz angehört. \dots\ Was daran \emph{nothwendig} ist, ist die
\emph{substantielle Identität} des Subjekts und Prädikats, gegen welche das
Eigene, wodurch sich jenes von diesem unterscheidet, nur als ein
unwesentliches Gesetzseyn, --- oder auch nur ein Namen ist; das Subjekt ist
in seinem Prädikate in sein An- und Fürsichseyn reflektirt. --- Ein solches
Prädikat sollte mit den Prädikaten der bisherigen Urtheile nicht
zusammengestellt werden, wenn z.~B. die Urtheile:
\begin{center}\begin{tabular}{r@{\ }l}
 & die Rose ist roth,\\
 & die Rose ist eine Pflanze,\\
oder: & dieser Ring ist gelb,\\
 & er ist Gold,
\end{tabular}\end{center}
in Eine Klasse zusammengeworfen, und eine so äußerliche Eigenschaft, wie die
Farbe einer Blume als ein gleiches Prädikat mit ihrer vegetabilischen Natur
genommen wird, so wird ein Unterschied übersehen, der dem gemeinsten
Auffassen auffallen muß.'' --- But that the logical form of the judgment,
when what is at
stake is the pure immanent development of form, has nothing to do with
the otherwise varied content of cognition, is easy to see. A judgment
such as: man is mortal, is in a formal respect just as much an
assertoric, categorical and positive judgment as it is a
``judgment of reflection''; indeed, which judgment is not, according to
its form, a judgment of reflection? The positive judgment is positive
only in reflection of the negative, the singular only in reflection of
the particular, the categorical only in reflection of the hypothetical,
etc.}
as though it did not follow of itself that all forms
% --- p. 51 ---
\opage{51}of logical judgments must at once be capable of application in
the development of the system's very first and poorest categories.

The confusion that has thus crept into logic's
% --- p. 52 ---
\opage{52}determination of the peculiarity of the several judgments in
respect of form brings with it in its turn another and corresponding
confusion when the question is of the divisions of the judgments.

In its ``Survey of the modes of division of judgments and judgings''
formal logic proceeds roughly as though what were at stake were an
arrangement of furniture, a distribution of kindling wood, a
classification of shoes and boots, and the like. Formal logic leaps from
one ground of division to another, enumerating now more, now fewer,
without the one being in any way developed out of the other.

Thus the judgments can be divided with respect to their matter into
mathematical, physical, astronomical, astrological, necromantic,
chiromantic, etc. etc. judgments.

Next the judgments can be divided with respect to the real nexus between
subject and predicate into ``not properly judging'' and ``properly
judging propositions''. The ``not properly judging propositions'' can in
their turn be divided into the ``merely propounding'' and the
``tautological''. The ``merely propounding propositions'' can in their
turn be divided into the ``merely narrating,'' the ``merely
describing'' and the ``merely concept-propounding.'' But the ``merely
propounding propositions'' can in their turn be divided into those in
which there lies no ``aggregate of judgings,'' and those in which there
lies an ``aggregate of judgings without its coming forward as
judging''. As regards the ``tautological propositions'', these can
according to circumstances be divided into the ``wholly vacuous'' and
the ``not wholly vacuous''. The latter can in their turn be divided into
those that serve ``to close the matter'', and those that are ``merely
concept-propounding'', etc. Further, the ``tautological propositions''
can be divided into the ``directly and the indirectly tautological'',
and these again be divided, etc. If now the ``not properly judging
propositions'' can bear so many divisions, subdivisions and
subdivisions of subdivisions, what must the ``properly judging'' ones
not be able to bear?

% --- p. 53 ---
\opage{53}Furthermore the judgments can be divided with respect to the
logical nexus between subject and predicate, and under this respect
again with respect to quality, quantity, etc. Finally these divisions
can ``in general be divided'' into those which at bottom do not concern
logic, namely the material ones, and those which do concern logic. The
divisions that concern logic can in their turn be divided into ``the
logical'' ones and then into ``certain divisions which are a consequence
of the fact that two or more judgments can stand in different relations
to one another,'' divisions which in the last resort then also belong
among the logical ones. Etc. etc. --- ``Thus'', says Heiberg, ``one is
astonished to see twenty to thirty species of judgments set up, each
distinguished by its own baroque name. Since furthermore the division of
these species proceeds according to different principles of division,
every species of a given division can combine successively with all the
other species of the other divisions, whereby there arises a host of
distinctions which are wholly vacuous, inasmuch as in the judgment there
lies no principle for this multiplicity.''\footnote{Den speculative
Logik §~144 Anm.~3.}

That formal logic, with its hotchpotch of sections, rubrics and
divisions, has had a repellent effect upon so clear and logical a head
as Heiberg is intelligible; it awakens more attention when we see him
turn even against Hegel's own attempt at an immanent deduction.
``Hegel, who gave scholastic logic its death-blow, has'' --- says
Heiberg ``by wishing to deduce the different species of judgments as
different stages of development through which the judgment becomes a
syllogism, nonetheless bestowed upon them an attention and allotted to
them an importance which they are far from deserving.''

But here Heiberg cuts the knot instead of untying it. In order to be rid
of formalism he demands that
% --- p. 54 ---
\opage{54}logic shall not enter at all upon any closer determination of
the mutual difference of the judgments, but go from the fixing of the
concept of the judgment straight over to the syllogism, so that the
further development of the judgment coincides with a development of the
syllogism. Hereby he overlooks that the logical development of the
different forms within the form of the judgment is an essential
development of form; were one in logic to skip over essential
developments of form in order to avoid formalism, it would presumably
end with one's skipping over the whole of logic at last.

``Every judgment'' --- it reads\footnote{Den speculative Logik §~144.
Anm.~3.} --- ``is positive, for whether the universal under which the
singular is subsumed be determined positively or negatively, the
subsumption itself is always a positive act.'' This must be conceded;
but in the concession there is then also seen at once the demand for an
essential development of form. For if reflection be made here upon the
act of judgment, hence upon what is subjective in the function, it shows
itself that the positive act of subsumption is essentially three
positings: the positing of the subject, the positing of the predicate,
the positing of the unity of subject and predicate. We may now readily
begin with the presupposition that here there is essentially only one
single logical judgment; but since the form of thought always alters
with the form of language, and no single form is absolutely adequate, it
is seen that the one judgment must necessarily develop into a plurality
of judgments, inasmuch as the one form develops into a plurality of
forms. If we reflect upon the positing of the unity of subject and
predicate, hence upon the \emph{copula}, then it may indeed be conceded
that the difference between subject and predicate is ``a difference of
reflection in which the one is only what the other is''; but the
difference of reflection is for all that none the less an essential
difference: the one is nevertheless not what the other is. To be sure
the judgment says that the subject is what the predicate is, that the
subject
% --- p. 55 ---
\opage{55}is its predicate; but the judgment also says that the subject
is not predicate, the predicate not subject. If now the being of subject
and of predicate is posited immediately as one and the same being, the
judgment is positive, and the positive judgment determined by the form:
S is P. If conversely the being of subject and of predicate is posited
as a being at variance with itself, differentiated in itself, then the
subject is not what the predicate is; the judgment is thus negative, and
the negative judgment determined by the form: S is not P. Immediately,
and according to the content of cognition, the positive judgment can
indeed be set up independently over against the negative: to that extent
there are here two different judgments. It is then upon this difference
too that logic reflects when it divides the judgments according to their
\emph{quality}: into \emph{positive}, \emph{negative} and
\emph{infinite}. But, essentially regarded, there is here only one
judgment. The positive judgment acquires logical significance only in
reflection of the negative, just as conversely the negative judgment
must become infinite and in its infinity meaningless if it cannot
reflect itself in its corresponding positive one. How the
language-forming functions prepare the thought-forming ones is then seen
in this connection from the fact that there is a linguistic expression
for a subject's being or not-being of its predicate.\footnote{That the
copula ``is'' is not found in the Hebrew and Aramaic languages has been
especially brought forward in the history of the religious controversies
with respect to the sacramental words: \cjRL{hU' dAmiy}, \cjRL{hU' gUpiy}
[this is my blood, this is my body].
% Hebrew set with cjhebrew, as in transcription.tex; the bracketed gloss is
% editorial. See that file's note at this site for the cjhebrew input scheme
% and for the standing uncertainty about the pointing.
 A language can have the abstract copula without therefore having an
abstractum answering to the category \emph{being} (thus, e.g., the
French \textit{être} is not being, but substance), but a language that
does not have the abstract copula will hardly furnish adequate
expressions for categories such as being and not-being.}

If we reflect in particular upon the \emph{positing of the subject},
then it may indeed be conceded that the positive in the act of judgment
is essentially determined by the negative, and conversely; but the unity
of the two-sided determinateness first acquires its punc%
% --- p. 56 ---
\opage{56}tual significance when it is expressly posited in the form of
the judgment. This S is P; this judgment is, immediately regarded,
positive; but the negative transforms itself into a limit and shines
through when weight is laid upon the fact that for the time being it is
only \emph{this, and no more} S of which P is asserted. Just as now the
positive judgment: this S (at least) is P, rests upon the possibility
that there are S which are not P, so conversely the negative judgment,
that there are S which are not P, rests upon the possibility that some,
indeed several, and ever more S may be P. Insofar as it is still
uncertain how many S are P, it is just as possible that they all are so
as that none is so. It is then evident from this that formal logic's
division of the judgments according to \emph{quantity}: into
\emph{singular}: this S is P, \emph{particular}: some S are P, and
\emph{general}: all S are P, has essential validity in respect of the
form, but becomes misleading when weight is laid upon the immediate
independence of the several judgments. It is the same, essentially the
selfsame logical judgment which in quantity \emph{distributes} the
negative and the positive, and which in quality lets the positive
exclude the negative, and conversely.

If we reflect in particular upon the \emph{positing of the predicate},
then it may readily be conceded that the unity of the two-sided
determinateness, of the negative by the positive, is decisive for the
judgment in every form; but then it also essentially depends upon how
far the form itself indicates the act of thinking by which the unity is
posited. That the unity of subject and predicate is a unity-in-difference
means then, in other words, that the predicate as predicate
\emph{belongs to} the subject insofar as it both is in the subject and
\emph{comes to} the subject. That the predicate belongs to the subject
means, then, that the predicate-determination is the subject's own
original determination; here, therefore, nothing is posited in the
predicate that is not already in the subject; the predicate-determination
is a subject-determination transparent in the form: \emph{the judgment
is analytic}. But that the predicate essentially belongs to the sub%
% --- p. 57 ---
\opage{57}ject presupposes then also that in the predication it is
really brought together with the subject; by the act of predication,
therefore, a new determination is brought to the subject itself:
\emph{the judgment is}, in other words, \emph{synthetic}. From the
standpoint of finitude it does indeed look as though the analytic
judgment were complete by itself over against the synthetic, and
conversely. But this is in the logical sense a pure illusion, and formal
logic must bear the blame for the fact that this illusion, so pernicious
for science, constantly receives new nourishment. How should it be
possible, when we look expressly at the concept, to find a judgment so
analytic that there were not to be found in it a moment of synthesis.
Let the predicate be in the subject itself as much as one will; the
subject cannot possibly be subject except in formal opposition to its
predicate: by the form of the judgment the formal independence of the
predicate is irrefutably asserted. Conversely the subject may, according
to its immediate reality, have whatever independence it will --- in
opposition to its predicate it has, in virtue of the form determined by
the act of predication, only formal independence. If now on the one side
it is not possible to conceive a judgment so analytic that it should
contain nothing synthetic, so on the other side it is just as little
possible to conceive a judgment so synthetic that it were not at the
same time analytic. A judgment so synthetic that the analytic were
wholly excluded would have to be a judgment in which subject and
predicate came forward each by itself with such immediate independence
that their union were to be regarded as wholly contingent; but such a
judgment is in its essence meaningless. The subject in the judgment is
subject only by reflecting its predicate, just as the predicate is
predicate only by reflecting its subject. What the predicate asserts,
the subject must be, were it only for a single moment; but what the
subject as subject itself is cannot possibly be indifferent to its
being: the synthetic judgment is essentially
% --- p. 58 ---
\opage{58}analytic, because the analytic judgment is essentially
synthetic.\footnote{It was by a true flash of genius that Kant became
aware that a priori judgments too can be synthetic; but he stopped
half-way. That synthetic judgments are \textit{a priori} possible rests
precisely upon this, that what is synthetic in the judgment is
\textit{a priori} one with what is analytic.}

The positive judgment is thus in its concept analytic-synthetic just as
well as the negative, the singular analytic-synthetic just as well as
the general. In the logical sense it is so far from being the case that
the analytic and the synthetic should here be something by themselves
over against the positive, the negative, the singular, the particular
and the general, that on the contrary they ground the
unity-in-difference of the positive and the negative, the singular and
the particular, etc. Herein lies then also the ground why the difference
between the analytic and the synthetic judgment has not, as the
difference between the positive and the negative, the particular and the
general judgment has, any direct expression in language.

Since it nonetheless makes an essential difference whether it is the
synthetic or the analytic side of the judgment that hovers before the
one who judges in the moment of thought, there is already hereby
awakened an expectation that language must surely have, if not a direct,
then at any rate an indirect expression for the distinction between the
merely synthetic, the analytic and the analytic-synthetic.

How this expression must be constituted can, although formal logic does
not let us surmise the least of it, hardly be difficult to find out. It
is then in the first place manifest that the immediately synthetic is
directly determined by the immediately given, S is really P (the wall is
really red; his eyes are really blue): \emph{the judgment is
assertoric}. But how now are the analysis and the consciousness of the
abstractly
% --- p. 59 ---
\opage{59}analytic to be expressed? Naturally in a negative and indirect
manner. When the predicate, immediately regarded, belongs to the subject
only because in actuality it is found united with it, there lies in this
that the opposite might be the case, namely that the union did not take
place; the subject then hovers between being and not being the
determinate predicate: S is perhaps P, perhaps not: \emph{the judgment
is problematic}. But when the judgment is problematic, the question
involuntarily forces itself forward: which predicate must the subject in
question then, according to its concept, incontestably have? In order to
decide what in this is necessary and what contingent, the analysis must
be undertaken. And with what then does the analysis end? Not with simply
discarding the contingent and retaining the abstractly necessary (for
such an attempt betrays but little power of judgment), but with taking
up the contingent as a moment in the concretely necessary: S is
necessarily P; here is the unity of the synthetic and the analytic:
\emph{the judgment is apodictic}.

Formal logic divides the judgments with respect to \emph{the real nexus}
between subject and predicate into \emph{analytic} and
\emph{synthetic}, and with respect to \emph{modality} into
\emph{assertoric}, \emph{problematic} and \emph{apodictic}; but it lets
each of the divisions stand as a rubric by itself, without any
suggestion that the one should in any way concern the other. It is this
parcelling out, this half thoughtless but nevertheless ``thoughtful''
rubrication, that does such irreparable harm. By overlooking the
connection between the two modes of division one naturally overlooks the
connection between the formal function and the form of the judgment. The
same judgment which for the one presents itself as assertoric presents
itself then for the other as problematic, for the third as apodictic,
etc. If on the other hand one perceives the connection between the
formal function and the form, then one perceives also the connection
between the different modes of division. The assertoric judgment is in
its concept assertoric precisely because it
% --- p. 60 ---
\opage{60}as immediately positive is immediately synthetic. In the
analytic judgment the predicate is taken up into its subject, in the
problematic it loosens itself from its subject; but the one of these
extremes presupposes the other. If the analysis is to begin, immediate
certainty must be at an end: the problematic judgment answers to a
judging consciousness which by the immediately synthetic is driven
towards the analytic. Finally, the apodictic judgment is not apodictic
because it occurs to a thoughtless man to use the form: S is necessarily
P, where he had better use the form: S is perhaps P; the form: S is
necessarily P, expressly presupposes that the judgment rests upon an
exhaustive analysis answering to the synthesis. In the apodictic
judgment the assertoric in the synthesis has become problematic, and
this problematic, by reason of the analysis, again assertoric, and
thereby more than assertoric.

If we reflect finally upon \emph{the posited unity of all three
positings}, hence upon the equilibrium between the logical independence
of subject and of predicate, then we come to the most concrete
completion of form that thought, by the help of language, is able to
give to the logical judgment.

The positive judgment is in its full determinateness \emph{categorical};
the categorically positive judgment cannot be immediately opposed to the
negative, for the positive has taken up the negative into itself as a
moment. The categorical judgment expresses the unity of the singular and
the general, which among other things comes to light from the fact that
the plural characteristic of the general judgment (All men are mortal)
is transformed into an essential singular, significative for the
universal judgment (Man is, as such, mortal). That the categorical
judgment is also determined as analytic-synthetic is seen from the
developed conceptual relation of subject and predicate; for the subject
is here posited in unity with its predicate, as the singular (the
existing) in unity with the universal (the concept). The categorical
judgment cannot, then, as assertoric have the problematic outside
% --- p. 61 ---
\opage{61}itself; the existent cannot, after all, fall out of its
category: the categorical judgment is essentially apodictic.

But just because the categorical judgment has all the judgment's
fundamental determinations in it, it can have no special
language-form: in its expression it is straightforwardly the
judgment as judgment. Its peculiarity first develops, and becomes
recognizable in language, when its own inner oppositions are asked
after. S is (categorically) P; the singular (S) is subsumed under
its universal (P); one and the same concept (P) comprehends a
multiplicity of particular S, namely S\textsubscript{1}, S\textsubscript{2}, S\textsubscript{3}, \dots.\ Thus:
S\textsubscript{1} is P, S\textsubscript{2} is P, S\textsubscript{3} is P. etc. But categorically
the converse holds as well, namely that many universal determinations
are found united in one and the same existing singular,
that, in other words, the predicate inheres in its subject. This
then gives us the following series: S is P\textsubscript{1}, S is P\textsubscript{2}, S is P\textsubscript{3} \dots\
From this it is seen that the predicate in the categorical judgment
can be partly one-sided, as a property of the thing, partly all-sided, as
the totality-concept that exists in the subject, the species that is
realized in the individual.

If the fundamental relation between subject and predicate, that
which exists and the category, has been brought to consciousness, then
it has at the same time been brought to consciousness that the
categorical judgment cannot possibly be seen in its significance
without a series of categorically connected judgments forming itself.
If now, e.g., S is (categorically) P, what then? Well, then s must
(categorically) be p; or still more abstractly and more formally:
if S is subject, then P must be predicate; thus
we have the \emph{hypothetical} judgment. In the hypothetical
judgment the categorical is doubled; at first
glance it looks as though there were two independent judgments here, one in the antecedent:
if S is P, and one in the consequent: then
s is p; but the unity consists categorically in the dependence%
% --- p. 62 ---
\opage{62}relation given between the two propositions, the condition and the
conditioned.\footnote{%
For the distinction of the categorical and the hypothetical judgment
Kant brings out that the former is determined by the categories substance
and inherence, the latter by the categories cause and effect. To this
Trendelenburg remarks (Logische Untersuchungen. Zweiter Band. Berlin
1840. p.~178.) that these concepts form no opposition.
``Die Substanz ist in der Eigenschaft causal. Die Eigenschaft ist
die an das Ding gebundene Thätigkeit. Die strengste Form der
Inhärenz ist das Verhältniß der im Ganzen inwohnenden Theile.
Da aber das Ganze die Theile trägt und zu dem macht, was sie
sind, so schlägt auch hier die Inhärenz in die Causalität über.'' But
what does this prove? Nothing other than that with respect to
the same object one can apply now the categories substance and
inherence, now the categories cause and effect. What was to be
proved here is that, when the categories cause and effect form the
ground, one applies the categorical form just as fittingly as the
hypothetical; but from proving that one will surely refrain. In
defense of Kant's conception Apelt (Die Theorie der Induction.
Leipzig 1854. p.~9.) declares it a downright impossibility to transform
the one of the two kinds of judgment into the other. ``Was man
von der \emph{Verwandelung} kategorischer und hypothetischer Urtheile
sowie kategorischer und hypothetischer Schlüsse in einander gesagt
hat, ist eine logische Fabel. Denn jede solche angebliche Verwandlung
besteht nur in einer Veränderung des \emph{grammatischen
Ausdrucks}, aber nie in einer wirklichen Veränderung \emph{der logischen
Form}.'' This too is a great misunderstanding;
for with the change in the grammatical there enters precisely an
alteration in the logical. Trendelenburg has then also (Anf.\ Skr.\ S.
% The Danish prints "Anfr. Skr."; read as "Anf. Skr." (as at p. 65).
179.) shown by examples how a categorical judgment allows itself to
be transformed into a hypothetical: ``Wenn ein hypothetisches Urtheil verschiedene
Gegenstände zusammenbringt, so sind sie immer unter
einem höhern Begriff eins, da sie als Bedingung und Bedingtes,
als Grund und Folge eins gesetzt werden. Vermöge dieser in
% The Danish prints "Wermöge"; given here as "Vermöge".
einer Thätigkeit oder einer Eigenschaft ruhenden Einheit kann es
geschehen, daß auch in diesem Falle das hypothetische Urtheil ein fast
gleichbedeutendes kategorisches neben sich hat. Man vergleiche z.~B.
die Urtheile: Wenn Bernstein gerieben wird, so entwickelt sich Electricität.
Der geriebene Bernstein entwickelt Electricität. Ferner:
Wenn ein Kegel durch eine Ebene geschnitten wird, so entstehen
regelmäßige Curven. Der Schnitt durch einen Kegel erzeugt regelmäßige
Curven''.

Thus: from the one side one seeks to level out the difference
between the categorical and the hypothetical judgment, because one sees
that different categories can be applied to the same object, and
does not notice that the judgment's form alters with the altered
categories; from the other side one wants to raise a partition wall between
the categorical and the hypothetical judgment, because there admittedly is
a difference between the logical and the grammatical, and does not notice
that the logical, in accordance with the reciprocal determination of the
language-forming and thought-forming functions, is altered by the grammatical. A
third standpoint is represented by Ulrici (Compendium der Logik.
Leipzig 1860. p.~177--178.): ``Hypothetische Urtheile giebt es in
Wahrheit gar nicht. Denn daß sie auf dem Verhältniß der Ursache und
Wirkung beruhen, ist (z.~B. von dem Urtheil: wenn dieß Mineral
ein Metall ist, so ist es schmelzbar) nicht nachweisbar und würde
keinen Unterschied begründen. Und ihre anderweitige Verschiedenheit
% The Danish prints "begrüuden"; given here as "begründen".
von den kategorischen ist entweder nur eine rein äußerliche des sprachlichen
Ausdrucks, oder beruht bloß auf einer rein subjectiven Beziehung
des Urtheilenden zu dem Gegenstande (Inhalt) des Urtheils,
--- ist also ohne alle logische Bedeutung.'' The judging consciousness's
subjective relation to the judgment and to the judgment's object
is thus --- ``ohne alle logische Bedeutung?''

To extricate oneself from these quackeries there is only one way out:
to hold fast the judgment's form in its unity with the abstract
logical content and to take note of the unity-in-difference of word and thought.}
That the hypothetical judgment can designate something uncertain,
something still undecided, in likeness with the problematic, must
% --- p. 63 ---
\opage{63}be granted; but the emphasis by no means falls on the uncertainty,
on the contrary! it falls categorically on the decisive determinateness
with which the antecedent entails the consequent: if
S is P, then (categorically) s is p, or if s is not
p, then (categorically) S is not P.

Inasmuch as every existing thing expresses its category, things exist
in categorical connection. The categorical connec%
% --- p. 64 ---
\opage{64}tion is contingent-necessary, contingent with respect to the
immediacy of the existent members, necessary with respect to
their mutual dependence. The contingent-necessary connection of
things is expressly brought to consciousness in the hypothetical
judgment; but the hypothetical judgment is in this respect
only a one-sided development of the categorical. The hypothetical
judgment has its point of support in the presupposition that the singular
is subsumed without the universal,\footnote{%
If we disregard the deduction and mix in representations of one
real content or another, totality-predicates can naturally
also come forward as subjects: if the plant is a rose,
then the fruit etc.}
% The Danish reads "subsumeres uden det Almene", as printed (cf. "under sit Almene", p. 61); translated as printed.
that it is thus the singular
things in their relations and reciprocal conditionedness on which,
in the totality-development, everything essentially depends: the totality-unity
is as yet only a connection-unity. But with this
the categorical judgment is by no means exhausted.

According to the categorical judgment the thing can first fully
correspond to its category, and the category to the thing, when both,
each in its own form, express the same totality of categorical
determinations. But when the predicate has the same totality
as the subject, only with this difference, that the predicate-determinations
are essential and logical, the subject-determinations existential and
antilogical, then the relation can of course be reversed: the subject, that
which exists, then becomes predicate, and the predicate, the concept of that
which exists, subject. The concept, the essence totalized in itself,
is actualized in a multiplicity of existences; these existences
now constitute the members of the predicate, while the unity of the concept
furnishes the subject. Thus there arises, when we look to
the essence of the subject, the conjunctive judgment: S is (in its essence)
both P\textsubscript{1} and P\textsubscript{2} and P\textsubscript{3}, and, when we look to the existing
members of the predicate, the \emph{categorically-disjunctive} judgment\footnote{%
That the disjunctive judgment can moreover also have that which exists
as subject and predicate in one-sided categories (the thing is either
black or white or red) is in this connection, where what is at issue
is the deduction, without significance.}: S is
% --- p. 65 ---
\opage{65}(in its existence) either P\textsubscript{1} or P\textsubscript{2} or P\textsubscript{3} \dots.\ Since now each
P is in its totality a realized S, each existing member is not only
a relative totality for itself, but at the same time a moment
in the comprehensive totality that penetrates all the members.
The penetrating totality-unity is a fundamental unity;
by the reciprocal determination of the hypothetical, the conjunctive and
the disjunctive judgment the form is brought to consciousness, the
categorical form, which serves to express that the connection of the members,
which exclude yet condition one another, essentially rests
upon a fundamental unity.\footnote{%
Alle disjunctiven Urtheile \dots.\ (die auf dem Verhältniß der Wechselwirkung
oder des Ganzen und Theils beruhen sollen, was wiederum
nicht einmal richtig ist) sind ebenfalls \emph{nur} sprachlich von den
s.~g. positiven (kategorischen) unterschieden, und lassen sich daher
unbeschadet ihres Sinnes in kategorische umwandeln: Das Urtheil
z.~B.: Cajus ist entweder gesund oder krank, ist nur sprachliche Abkürzung
für: Es ist unmöglich, daß Cajus zugleich gesund und
krank sey, oder: Es ist nothwendig, daß er entweder krank oder gesund
ist. Und das Urtheil: Die Kegelschnitte sind entweder Kreise
oder Ellipsen oder Parabeln oder Hyperbeln, läßt sich auch so ausdrücken:
die Kegelschnitte sind theils Kreise, theils Ellipsen oder:
Die Kegelschnitte können nur die Form von Kreisen, Ellipsen etc.
haben, nur diese Formen sind möglich. Ulrici. Anf.\ Skr.\ p.~178.
What is true in this reduction is that the judgment's particular forms
are modifications of one and the same fundamental form; but what is untrue
is that the particular forms should on that ground have no
logical significance. How arbitrary the reduction is strikes the eye
when the judgment: Cajus is either healthy or sick, is supposed to be a ``sprachliche
Abkürzung'' for: it is impossible that Cajus can at one and the same time be
both healthy and sick. It would really be interesting to see a
little proof that the judgment in the latter form is more logical
than in the former. Why can Cajus not at one and the same time be
both healthy and sick? Because a decisive either-or is presupposed here:
either healthy or sick; if doubt is raised about this presupposition,
doubt is at the same time raised about the judgment:
it is impossible that Cajus can at one and the same time be both healthy and
sick. The logical in the latter judgment can thus just as well be said to
rest upon the logic in the former as conversely.}

% --- p. 66 ---
\opage{66}When formal logic inculcates the necessity of
going through the formal element in thinking in a thorough
and penetrating manner, we must grant that it is right. The doctrine of
the judgment's modes of division ought in that respect just as
little to deter us as the doctrine of the judgments' ``subalternation'',
``equipollence'', ``opposition'', ``conversion'', ``contraposition''
etc., because it is not ``spirited and attractive'',
if only it were scientifically fruitful. But here is precisely
the great misunderstanding of which the defenders of formal
logic make themselves guilty. It is not the thoroughness and the
acuteness, but the barren, unfruitful element in the treatment of the
logical forms, that is here opposed. Instead of pointing out
a principle that itself sets thought in movement, and in the
movement of thought lets the one determination of form, in a given connection,
develop out of the other, school-logic on the contrary
aims at depositing and setting up results of thought that merely are,
with disregard of every consideration of principle and movement.
The frozen waves in a brook still always bear traces of movement,
namely of an arrested movement; but formal
logic's distinctions and rubrics bear no trace of movement;
its propositions are like finished figures, strokes and
crosses, whose number and order must be impressed on the memory.

Herein lies a great misdirection. It does indeed seem as though
the many abstractions, arrangements of concepts, exact definitions
etc. must serve precisely to develop the understanding,
to exercise acuteness and, as Herbart expresses it, to ``accustom
us to sustained attention to the natural course of
thinking''; but these advantages are of a very ambiguous kind. What
help is it that the understanding is developed, when it is developed in a
% --- p. 67 ---
\opage{67}skewed direction, in that it is developed so as to separate what is thought from
thinking, the logical forms from their corresponding functions.

It is not a schema of the judgment's divisions finished once and
for all, but a development and motivation of
the difference of form under which the judgment comes forward by turns,
that opens to us an insight into the logical judgment-functions.
Least of all can scholastic formalism ``accustom
us to sustained attention to the natural course of
thinking''; it is on the contrary the unnatural course of thinking upon which
formalism fixes attention. For what can be
more unnatural than to conjure the living forms into a crystallized
being? what indeed can be more unnatural and more pedantic
than the manner in which formal logic by its distinctions
sets abstractum against abstractum, without an inkling that
living thought at every moment demands the sublation of the abstraction
just as necessarily as the abstraction itself? If it
really were thought's natural, not its unnatural course
to which scholastic logic with its barbarous mnemonic verses
directed attention, then in scientific
presentations, especially in every philosophical treatise,
one would surely see the fruits of logic's teachings; but the author
who should seriously set about letting logic's
teachings shine through a logical treatment of a concrete
subject --- in him thought would come
to run about as natural a course as in a Magister Stygotius
or an Else David Skolemesters\footnote{%
One can become pedantic just as well by chasing after the abstractly
logical with a petty, unfruitful acuteness, and arbitrarily
fixing a schema for what is ``in his own imagining'' philosophically
stringent, as by visibly overexerting oneself to
express everything quite ceremonially, figuratively, ingeniously and ornately. If it is
comic when a Stygotius as a suitor corrects his expression ``draws
my heart's iron'' to ``thrusts my heart's iron'', then it is also
comic when a Stygotius as a logician grows shy of a categorical
judgment and imagines that a tautology lurks in it; expectorations
such as: ``Gold is --- if I may permit myself an implicitly tautological
proposition --- for when we look at the manner in which the concept
metal has formed itself, and at what is in general understood by
metal, we shall not be able to call the proposition `Gold is a
metal' anything but an implicitly tautological proposition, since it at \emph{bottom}
or, what would I say, \emph{in general}, says neither anything other nor
anything more than that gold is such a thing as gold, silver, copper
etc.\ --- a metal'', would be style for a Stygotius.
% The Danish prints "o s. v." without the point after "o"; rendered "etc."

Editors of formal logic (for a long straw that is threshed over again
only because it has been threshed so often can hardly be used
for productive authorship) not seldom make themselves guilty of the comic
contradiction that, whereas in the so-called principle of identity, where the tautology
is conspicuous, they cannot discover the tautological, but
for the sake of ``certainty'' and ``truth'', ``determination'',
``statuation'', ``whatness'', ``unanimity'' and ``concordance''
alone rejoice over ``the deeper logical element'' that lies in
logic's highest proposition, that A is what A is, they on the other hand
get scruples when they are to recognize it as a judgment that gold is a
metal, for it might after all possibly be a tautology. It
naturally does not occur to that pusillanimous acuteness to think
the matter through and let the decision of the case turn upon a dialectic;
for then it would soon become clear. If the following judgments: gold is
a metal, silver is a metal, copper is a metal, are implicitly tautological,
what follows from this? Then it follows from this that gold is
silver, and that silver is copper, a consequence whose absurdity is
just as evident to the natural understanding as its necessity.}.

% --- p. 68 ---
\opage{68}To thought's natural course it belongs that the language-forming and
the thought-forming functions enter into living reciprocal action;
only thus can the representation pass over into concept, and the
pictorial be illuminated by the non-pictorial. Logic's significance is
that thought catches sight of itself, that the abstraction is carried through.
First a judgment is passed on the object, and the judgment's
form is deposited in language; the thought-forming function is
% --- p. 69 ---
\opage{69}latent in the language-forming.\footnote{%
The language-forming functions are recognized, namely, so far as style is concerned,
not immediately by the fact that either new words or new grammatical
turns are formed here, but by the fact that the given means of language are used
in a peculiar manner. To every peculiar
turn of language there corresponds a peculiar turn of thought. He who
says: ``all metal is fusible'' speaks a different language and thinks
the thought otherwise than he who says: ``if this is a metal,
then it is fusible''; and the latter in turn speaks a different language and thinks
the thought otherwise than he who says: ``if this mass does not
allow itself to be melted, it cannot possibly be metal''.} In the judgment's linguistic form
thought then sees its own trace, an imprint of its invisible essence,
and now begins with a judgment upon the judgment: the language-forming
function introduces the thought-forming. Once thought has, by
its own turns of language, become attentive to the abstract and
has won certainty in abstraction, then it is so far from being the case
that fuller, more figurative expressions should weaken the light-effect of
ideality, that reflection, reflection carried through, on the contrary
has a satisfaction in striking into the way of intuition
and transforming the pictorial into a clear form, transparent
in its color, for the non-pictorial.\footnote{%
It is wonderful how intuition and thought, the pictorial
and the non-pictorial, can by the word's sorcery be interwoven with
one another, so that thought's clarity and the image's boldness stand in
direct proportion. The reciprocal action of language-forming and thought-forming functions
is then aesthetic. --- ``My sorrow, that is my knight's castle, which like
an eagle's nest lies high up on the mountains' peak among the clouds;
no one can storm it. From it I fly down into actuality and
seize my prey; but I do not remain down there, my prey I bring
home, and this prey is an image that I weave into the tapestries
of my castle. Then I live as one deceased. Everything that has been
lived through I dip down into the baptism of forgetfulness unto the eternity of recollection.
Everything finite and contingent is forgotten and effaced. Then I sit
like an old, gray-haired man, deep in thought, and explain the images
in a soft voice, almost whispering, and at my side sits
a child and listens, although it remembers everything before I tell
it.'' Victor Eremita. Enten-Eller. Kjøbenhavn 1849.
p.~21. --- What are we now to judge, from logic's standpoint,
of such a presentation? If there are here only images without thoughts,
then the whole is presumably an insipid fabrication, and we must, with Master
Mannegrimm, ask for ``one more lie''. If, on the contrary, there are here just as
many, indeed still more thoughts than images, if the colorful
image-propositions are essentially clear thought-propositions, should the thoughts
then perhaps gain by being stripped of the rich garb and presented
in their logical nakedness? But if this last is neither
desirable nor in a stricter sense possible, then we have here
an example of the aesthetic reciprocal action of language-forming and thought-forming
functions.}

% --- p. 70 ---
\opage{70}The logical expression too can have its essential beauty.
A presentation that in each and every respect should satisfy
logic's aesthetic demand is certainly a great rarity,
since it presupposes a peculiar, happy union of reflection and
imagination; but that logic, strict logic, makes aesthetic
demands is certain: the aesthetic is here grounded in the
living thinking's appropriation of everything that stirs in
the life of consciousness; it is this aesthetic that inspires thought,
gives the shadings their play of color and makes productivity recognizable
by the attractive harmony of language-forming and thought-forming
functions.

\textit{γγ)} \emph{Idea-forming Functions.} Between subjectivity
and the actual world of objects the chasm is so deep
that it opens anew every time reflection dreams
of having got it filled. \emph{The Logic of Functions}
ended with a concretion of the teleological unity of the ideal
system and the real system, namely as a syllogism-unity of nature and
spirit, a development in which the dualism of the objective and the
subjective provisionally found its explanation. From several sides
it has then also been shown in advance how the antilogical,
by the mediation of sense-perception, allowed itself to be transposed into the logical;
how the same content could assume a different form,
the same form a different content etc. But beneath
% --- p. 71 ---
\opage{71}all these turns, transpositions and explanations there is
understood a something which, just because it was in immediate
application, could not at once become the object of special reflection,
namely the word.

What then is the word? That the word is not without further ado the
same as language is clear from our foregoing; for there
are indeed languages without free-standing words, and the language that has
specially formed words has, after all, many and different words. The category:
\emph{word} therefore stands in this respect partly lower,
partly higher than the category: \emph{language}. Lower! for just as letters
are elements in syllables, syllables elements in
words, so the single word is in turn an element in the proposition,
and the proposition an element in language. But as an element of the proposition
the word is not yet essentially word; a language proper
cannot do without propositions, but a proposition can transform
roots into words: that there can be language without words is then
true insofar as propositions can be formed without words. Higher!
for the word is the expression of thought. Insofar as the proposition yields
a meaning, and the meaning a thought, the proposition is essentially
a word. ``There is a word, an old word, which says'' etc.
That the proposition is, according to its meaning and by virtue of the thought, word,
means then, in other words, that the word as word in a deeper sense
is higher than the proposition. From this standpoint there can
then be no language without words; it would have to be
a language without propositions, a pure animal language, or a language
of propositions without meaning, a meaningless language.

If now the word in a deeper sense stands higher than the proposition,
then it also stands higher than language. If we take
words as elements of propositions, that is, finitely and in the plural,
then one may certainly say that it is language which,
by a peculiar development, makes its elements into words; but
if we take the word as word in the absolute sense, that is, in the singular,
then we must say that it is the word that is the original,
language the derived; that it is the word which makes
% --- p. 72 ---
\opage{72}language into language. How unreasonable it is to want to make
language the absolutely original --- which is nevertheless what those do
who with F.\ Schlegel trace the origin of languages back to a
supernatural act, a miracle --- can already be felt when one,
instead of the Johannine: ``In the beginning was the Word, and
the Word was with God, and the Word was God'', attempts to put:
In the beginning was Language, and Language was with God, and
Language was God.

If it is now not language as finite language, but the word
in the absolute sense, that is the original, then it is seen from this
that what in the finite and relative sense is accomplished by language
must in the absolute and original sense be ascribed to the word.

What then does speaking subjectivity properly owe to language?
The answer sounds from the foregoing: an objectification-doubling,
whereby the opposition between the inner and the outer
world is essentially articulated! Certainly the shapes of the
surrounding world would be able to mirror themselves in consciousness
even without the mediation of language --- we see it, after all, in the animals; --- but this
mirroring would then have to proceed: partly in such a way that every image,
even if only for a moment, remained standing in antilogical outline,
relating itself to the object as the copy to its original;
partly in such a way that the images dissolved into representations as
clouds into mists, the flow of the representations proceeding on the one side
with spiritually determined natural necessity, and on the
other side losing itself and vanishing in empty, hovering reflections.
Only by the help of language does it become possible for consciousness to
articulate the objective content of existence in subjective form and,
conversely, to give the content of subjectivity objective form; only
by the help of language can the same content enter into a double
form, and the same form into a double content; only through
language does it become possible to form the idea, the absolute
content.

What in the relative sense is accomplished by language must,
as has been said, in the absolute sense be referred back to the word: the relation
% --- p. 73 ---
\opage{73}between the content and language presupposes a fundamental relation
between the idea and the word. Since the idea's absolute content is
at once both subjective and objective, the idea-forming
functions are to be determined specifically by the unity-in-difference of the
subjective and the objective. Insofar as this unity shows
itself as a \emph{magical} unity of idea and word, the idea-forming
functions are \emph{mythical}; insofar as the word's unity with
the idea has a \emph{fantastic} stamp, they are \emph{gnostic}; if
the unity is posited in the \emph{essence of the Godhead}, they become \emph{theosophical}.

\emph{Mythical Functions.} It is the unity-in-difference of
the objective and the subjective in the word on which the emphasis
falls. Essentially, with this unity-in-difference, everything depends
on the content. If the content is finite, and the separation between
subjective and objective unclear, then we have before us a consciousness
fermenting and groping among sensations and obscure representations,
a stage of consciousness corresponding to that of the child in its early years:
this standpoint cannot come into consideration here. If
the content is finite, and the separation between subjective and objective
clear, then we have the critically understanding stage of consciousness.
At this stage the word is one thing and the thing another, the representation
one thing and actuality another, and language then mediates in a
finite manner between the factors in this parceling out: this
standpoint likewise cannot come into consideration. What is decisive for
the idea-forming functions is that the content of consciousness
must be absolute.

Since now an absolute content cannot possibly move through
subjectivity without reflecting itself in relative oppositions,
it is further seen that the consciousness which is to objectify
the idea for itself in and by the word, and therein to cognize the absolute,
must already beforehand have reached the point at which it can in the finite
sense distinguish the things of the surrounding world, and on the whole
make a difference between actuality and sense-deception, between what
is outward and what is merely inward. The psychological
subjectivity which, after having attained, within narrow
% --- p. 74 ---
\opage{74}limits, a clear understanding of existence in the finite,
is seized and moved by the idea in so immediate a manner that
finitude's oppositions of dream and actuality, of
inward and outward disappear, in order to repeat themselves in higher
forms corresponding, by virtue of the word, to the absolute, is \emph{mythical}.

The distinguishing mark of the mythical is, in principle, by no means
obscurity about the outwardly finite; it is not sensuousness's
misty images, but the idea's lightning-flash, not the understanding's
night, but reason's daybreak, that lends the content of consciousness
its wonderful mythical colors. How peculiar the blending of colors
is strikes the eye when one really tries to characterize
mythology and the mythical standpoint. ``For mythology'',
says Heiberg, ``there is no other, all-embracing point of view
than the poetic. Mythology is poetry, and purely original
folk-poetry''. Later on this folk-poetry is then
at the same time apprehended as a kind of speculative philosophy. ``Every
mythology'', it is said, ``begins with cosmogony;
for philosophizing humanity's first endeavor is to
explain the wonder of existence. The first question is:
how has the world come into being? What was there before all
this came into being? Or what was there when as yet there was nothing?
The understanding can so little answer these questions that it must even
declare them foolish and self-contradictory.
But imagination hovers boldly out over the beginning of things,
without fearing to lose itself in nothing. It peoples the void,
it lets light arise out of darkness, existence out of non-existence''\footnote{%
J.\ L.\ Heibergs nordiske Mythologie. Paa Dansk udgivet af Chr.\
Winther. Kbhvn.\ 1864. p.~4 og p.~24.} etc.

That mythology is at once both poetry and philosophy
means then, in other words, that it is at bottom neither
poetry nor philosophy; not poetry, because it is philosophy;
not philosophy, because it is poetry. The closer determination is seen
% --- p. 75 ---
\opage{75}from the functions. The poetic is here not apprehended as
poetry. Poetry presupposes a separation of illusion and actuality
bearing upon the absolute content, a separation that
expressly sets illusion above actuality, because it
sets the idea above finitude; but with so reflected a separation
the myth as myth cannot possibly subsist: the mythical
figures maintain themselves, by the idea's magic, at the stage of a higher
actuality; here at this stage there is not even an inkling of illusion.
Just as little is the philosophical in myth-formation
apprehended as philosophy. What kind of
philosophy would that be, which began with an express discord between
understanding and imagination, and which on account of this discord
gave up the understanding and let imagination rule? Instead of
beginning with a dualism of imagination and understanding, myth-formation
on the contrary begins with an imagination impregnated by
understanding and reason; one may say with full right of the mythical idea-movement
that thought fantasizes because
\emph{imagination thinks}.

When, then, the imagination impregnated by understanding and reason
``hovers boldly out over the beginning of things,
without fearing to lose itself in nothing'', what steers its
flight, what is it that preserves it from plunging into the void's
abyss? Essentially it is \emph{the idea}; immediately seen, it is
\emph{the word}. The abyss of the beginning is equally deep, whether we begin
with \emph{Chaos} or with \emph{Ginungagap}; it is in both infinitely
deep; but in the infinite deep the idea is wedded to the
word.

If we begin with Chaos, then it is naturally not the
empty appellation on which weight is laid here, for the appellations
can vary. But the synonymous appellations all bear
upon the same fundamental word. That Chaos as a fundamental word
is a revelation of the idea, posited and moved by the idea's power,
can easily be shown. Chaos is in the mythical sense neither an
abstract logical concept nor an expression for a state of physical
% --- p. 76 ---
\opage{76}confusion, but a totality-unity held fast in the word, a
subjective-objective First which in pictorial and yet non-pictorial
form renders intuitable what cannot be intuited, namely the origin of
things, the idea of all actual beginning, or the actual beginning's
idea.

That chaos as concept is not in the form of the concept, but in the
form of a non-pictorial image, follows expressly from intuition's
unity with the word. What is there in chaos? Nothing, and yet
everything! This self-contradiction of thought, with the consequences
springing from it, is not apprehended in the form of a thought, but is
seen in an image, a mighty thought-image formed by the word. In that
thought, by help of the word, lets the imagination play over nothing
and everything, the imagination thinks essentially of everything and
yet intuits in chaos nothing, the empty; it thinks essentially of
nothing, and yet intuits in chaos everything, namely the seeds of
everything (\textit{non bene junctarum discordia semina rerum}). The
content, the content void of content in chaos, is essentially
subjective by means of the word: a product of imagining thinking, of
thinking imagination; but it is at the same time intuited objectively,
strikes the intuiting consciousness with the power of objectivity in
the word, and takes on the shape of physically immediate reality.

But that chaos in an objective sense nevertheless does not have
physically immediate reality is precisely a consequence of the
non-pictorial in the image, of intuition's mythical unity with the
word. In an objective, physical significance chaos would have to be
the actual dissolution in which all-existence would have to find
itself when all forms were disordered and the swarm of atoms in
lawless confusion. For an understanding, critical consciousness of
such a chaos the word has lost its magical force. The word chaos is in
this case only a name; chaos, on the other hand, as an outward actual
state is the thing, and the thing, not the word, is what matters.
Separated from the imagination that rests upon the word, the
understanding will then on its own account investigate the actual
chaos, the outward actual beginning of things. But this enter%
% --- p. 77 ---
\opage{77}prise of the understanding will totally miscarry. As an
actual state chaos is utterly unactual; there is actually no existence
that can begin with an actual chaos; what the understanding is to get
out of such a chaos it must in actuality itself smuggle into it: it is
an idle play, the understanding plays with itself and with the actual.

In the mythical chaos objectivity is subjective; the objective content
is by the word transposed into subjective form. How different the
mythical chaos is from the outward physical one is then also
recognizable from the figures that proceed out of chaos. ``First is
chaos, next the wide earth, dark Tartarus --- and Eros, the fairest
among the immortal gods.'' He who sees the old gods step forth out of
chaos cannot be in doubt about what sort of being chaos is. In that
one here begins with chaos, a run-up is here in a manner taken toward
an objective presentation of the world's coming to be, of the
formation of things, of the finite realities; but the view turns
about. Cosmogony swings over and becomes theogony; chaos becomes a
divinity, the divinity from whose womb the particular gods are born.
This regression of the cosmogonic into the theogonic is in a high
degree characteristic of myth-formation. In order that the world may
come to be, the gods must first come to be; in order that the finite
realities may in actuality come to be, the real ideas must first come
to be: the finite becoming of the real ideas reflects in itself the
eternal being of the fundamental ideas.

Chaos is thus in its ideality eternal, in its immediate reality
temporal, divinity in being, nature-ground in becoming, totality in
original unity, the ground of totalities in original plurality; in
chaos the relative is absolute, the absolute relative; in the
chaos-word motion is rest and rest again motion, fullness emptiness,
and emptiness fullness, the inward an outward, and this outward again
an inward: in chaos the idea is word, and the word idea.

% --- p. 78 ---
\opage{78}It is then not a contingency, but an essential necessity,
that the oldest gods issue out of chaos. In logical form the necessity
is recognizable by this, that plurality proceeds out of unity, since
unity itself is originally plurality, that the particular in the
totality proceeds out of the universal, since the universal is
originally a totality of particulars. But in the logical form the
content lacks immediate reality; reality, the mythical reality, comes
with the word, for the word determined by the idea determines the
idea.

Such figures as Gaea, Tartarus, Erebus, Nyx and Eros are in their
originality more than empty personifications. ``Figures of the
figureless'', ``elementary forms of the first One, in themselves
formless'' --- such famished abstractions fall just as far below the
idea as below actuality. No, by the magic of the idea in the mythical
word the oldest gods are, despite what in form is unformed, typical
figures full of reality, each divinity a totality in itself, a being
of its own, a determinate physiognomy. It is the idea that in
particular totality-unities penetrates, moves and determines the word;
but it is in the immediate sense the word, the particular fundamental
word, that in its unity takes up, determines and gathers the manifold
of the ideal and real, non-pictorial and pictorial moments that make
up the idea's absolute content.

The goddess Gaea, who receives her name after the earth, takes up by
this name all the earth's reality into herself, becomes by this name a
divinity with immeasurable capacities, mighty forces and great
possibilities; but that the goddess has received her name after the
actual earth says at the same time that it is the actual earth that
has received its name after the goddess. The earth-goddess and the
actual earth have not merely one name, but they are one being,
designated by one word: they are one in the word, since they are one
in the idea. That the actual earth is a goddess shows us in a
strikingly forcible manner what it means that the idea's objective,
real content is --- by means of the word ---
% --- p. 79 ---
\opage{79}transposed into subjective, ideal form; that the goddess,
with all the traits of elementary subjectivity, is the earth itself,
the actual earth, renders intuitable on the other side what the word
is capable of where it is a question of translating a subjective
content --- for the goddess's being is in its ideality a subjective
being --- into immediate, objective form.

% Errata: p. 79 — printed "udvikes", corrected to "udvikles"; translated as corrected.
Not merely a positive magnitude like the earth, but even such negative
forms of sense as darkness and night, such essentialities as Tartarus,
Erebus and Nyx, absorb the idea's reality through the word, are
individualized in the word and become living figures by means of the
word. Breathed through and moved by the idea, the word has not merely
the peculiar power that, as the highest and deepest of language, it is
able to gather the relativity-determinations, which in the
propositions of speech are piecemeal presented, described and
developed, into pregnant unity-expressions; but as a member in the
proposition, as an element in language, the word has at the same time
that suppleness, flexibility, mobility and relativity, that in the
most diverse connections it can be now for sense, now for thought, now
figurative, now literal, can unite what it has itself separated,
separate what it has itself united, can, member by member and point by
point, illuminate the real by the ideal, the pictorial by the
non-pictorial, the immediate by the reflected. Once the word has
become a god, and the god a word, then it stands in the power of
subjectivity to let this god of the word appear in the most
adventurous situations, undergo the most remarkable transformations,
show itself from the most diverse sides, thinking, desiring, acting,
suffering.

The mythical word, which transforms the particular ideas into
god-figures, at the same time transforms the ideal movements into
actual events; this transformation is decisively characteristic with
respect to the mythical essence of the idea-forming functions. In that
the ideas' strife with material actuality, and through actuality's
elementary
% --- p. 80 ---
\opage{80}oppositions their mutual strife with one another, is
transformed into a history of the gods' battles, the magical unity of
idea and actuality is upon the ways of the word brought to that
turning point at which the mythical identity of the ideal and the real
swings over into dualism, and the illusion shows itself.

The colossal audacity with which the thinking imagination,
impregnated by the idea, defies probability at the sight of working
gods is an eloquent testimony to the word's power in the objectified
events. The age of the gods, the distant past, is the golden age of
the ideas, of the ideal causes; the present with its facts is the age
of actuality, of the conditions and of the material causes: can the
two ages indeed become one age? In the mirror of the word, i.e.\ of
narration, of recollection, the past does indeed reflect itself in the
present as the great fundamental ideas in the actual results of
development; but actuality with its thing-causes reacts with the power
of the present against those ideal presuppositions; the consequence is
a discord between the thing and the word, a discord between the
objective and the subjective, a discord between content and form,
imagination and thought, a discord that awakens criticism, sublates
the standpoint and sets the idea-forming functions new and higher
tasks.

\emph{Gnostic Functions}. In principle mythology is different from
poetry; for the myth, which has magical objectivity, is without
illusion, whereas poetry as poetry essentially presupposes
consciousness of illusion. Since consciousness, however, can advance
continuously from dim intimation to clear cognition, can as in a
waking dream be at one and the same time in the power of illusion and
have a knowledge of the illusion: it is intelligible enough that
mythology has been able to pass over just as imperceptibly into folk
poetry as folk poetry in turn into art poetry. But for that reason the
separation of illusion and actuality nonetheless remains
% --- p. 81 ---
\opage{81}just as decisive for the stage of consciousness with which
poetry begins.

In proportion as the mythical content is transfigured in the forms of
poetry, the understanding, which sees the illusion, will tear itself
loose from the imagination, make front against the actual and strive
to orient itself in the world of things. But the world of things is
sensuous, finite, material and to that extent idea-less. Separated
from the imagination, the understanding is at the same time separated
from the idea. But the idea-forsaken understanding has an essential
need of ideality; the understanding is in principle rational: here
philosophy begins. Poetry with its idea-filled illusions and
philosophy with its formal concepts now form extremes with the
sensuously concrete world between them. In order to vindicate the
idea, the ideal powers, poetry and philosophy, work each from its own
side at dissolving, transforming and transfiguring the outwardly
actual; but actuality is strong, and the battle hard. For
reinforcement in the hard battle, reflection, enfeebled among
universal concepts, then feels a need to fill the empty forms with
intuitable content by renewing the old covenant between intuition and
thinking, reflection and imagination; on these terms the idea-forming
functions become \emph{gnostic}.

It is gnosticism's essence, and not a determinate historical figure of
gnosticism, that we here have in view; only insofar can regard be had
to the historical phenomenon as this phenomenon is in a peculiar sense
fitted to make the essence manifest. The re-established union of
understanding and imagination is recognizable by the word, the word
recognizable by the categories that bear the intuition. Here one does
not essentially begin with chaos. Chaos is platonically reduced to
pure emptiness, an expression for the essential nothingness of the
sensuous; thus chaos is nothingness, negative being, \emph{Kenoma}.
That the material and sensuous in all-existence is reduced to Kenoma
is in gnosticism an indication that the beginning, the real beginning
of development, cannot be sought in the physically objective, in
% --- p. 82 ---
\opage{82}sensuous reality, but must be sought in the essentially
subjective, in overreaching ideality. In categorical opposition to
\emph{Kenoma}, the objectively sensuous empty, there stands then
\emph{Pleroma}, the realm of ideas, the region of power and of
fullness.

Insofar as gnostic thinking gives itself over to intuitions of the
imagination, Pleroma is an objective thought-image in likeness to the
mythical chaos; but a closer consideration of this thought-image soon
shows us that in Pleroma we have to do with a realm of subjectivity,
of spirituality and inwardness. Pleroma's original unity is in
\emph{the Depth} and \emph{the Silence} (βυθὸς and σιγὴ). In and for
itself the depth can now indeed designate something objective just as
well as something subjective; but already \emph{the Silence}, not the
stillness, lets us surmise that the depth is in subjectivity. The
surmise becomes certainty when the opposition's first abstract form is
more closely determined by a following, more concrete opposition, the
opposition between \emph{the Primal Father} and \emph{the Thought}
(προπάτωρ and ἔννοια). If we cast besides a glance at the ever
increasing concretions in which the original opposition repeats itself
by degrees, the opposition of \emph{Reason} and \emph{Truth} (νοῦς and
ἀλήθεια), \emph{the Word} and \emph{Life}, etc., then we recognize in
Pleroma a realm of subjectivities and of ontologically subjective
content.

How peculiar the fusion of image and concept is in the gnostic Pleroma
strikes the eye when we note the immediacy in which the word-meaning
of the categories set up, indeed even the gender of the words, is
apprehended and employed. That the categories: the Depth and the
Silence, the Primal Father and the Thought, Reason and Truth, etc.,
are hypostatized essentialities, logical personifications, can be seen
from a translation too; but the linguistic contingencies which the
imagining thinking has employed in these personifications are
discovered only by a glance at the original text. The designations of
the pleromatic hypostases are namely chosen in such a way that
% --- p. 83 ---
\opage{83}the one is always a masculine word, the other
correspondingly a feminine word: βυθὸς, προπάτωρ, νοῦς,
λόγος \dots\ are masculine words, whereas σιγὴ, ἔννοια, ἀλήθεια,
ζωὴ \dots\ are feminine words.

By this division, a fundamental division which in its kind corresponds
to Hegel's objective judgment (das Urtheil, derived from urtheilen), a
certain equilibrium between reflection and imagination is brought
about. Reflection demands that Pleroma's fundamental division shall in
all repetitions yield oppositions of totality-concepts. This is in a
certain manner attained by the categories indicated; λόγος, νοῦς,
ἔννοια, ἀλήθεια are true totality-categories. Likewise there is here,
by the categorical oppositions, designated a progression from the
abstract toward the concrete, from the universal through the
particular toward the individual; for it must be conceded that βυθὸς
and σιγὴ are more abstract than νοῦς and ἀλήθεια, these in turn more
abstract than ἄνθρωπος and ἐκκλησία. Indeed, so ingeniously is the
series laid out that reflection can even, if it be demanded, give a
logical account of the progress. The more abstract doublings form,
namely, the foundation for the more concrete; these are, to speak in
Hegel's language, mediations of the former. Thus λόγος and ζωὴ are
evidently mediations of νοῦς and ἀλήθεια, since νοῦς by means of
ἀλήθεια becomes λόγος, ἀλήθεια must by means of νοῦς naturally become
ζωὴ, etc.

But if reflection is satisfied in a spirited manner, the imagination
is no less so: ingenuity and spiritedness are in general essential
marks of all gnosticism, the more recent as well as the older. That
the imagination has had its share in the spirited choice of categories
of opposition is seen from the manner in which the chosen masculine
and feminine words support the personification. According to the
personification, every categorical pair is a married pair, husband and
wife: deified forms, and yet not mythical gods!

The difference between the mythical and the gnostic is now striking;
% --- p. 84 ---
\opage{84}what a distance between a pleromatic hypostasis with its
abstract-metaphysical habitus, in its ideality ingeniously measured
out by reflection, and one of mythology's gods or goddesses, which
individualizes itself in the finite, plays with the contingent and
teems with reality! In gnosticism the imagination no longer has
thought in its power, and just as little can thought be said to have
the imagination in its power. The alliance in which both, imagination
and thought, were properly to strengthen each other is conditioned by
an accommodation whereby they mutually weaken each other: the
imagination becomes thoughtful, and thought fantastical.

That gnostic reflection aims at idea and ideality with a quite other
consciousness than the mythical does is a matter of course. The
ingenious and spirited is then also seen in the manner in which the
particular ideas are represented by the hypostases in Pleroma. In this
respect too the word makes its immediacy felt by the name. Every one
of the pleromatic hypostases bears, as is well known, the name of aeon
(αἰὼν). What does aeon mean? A world-age and an eternity! thus an
eternity enclosed by time, set in the ring of time, an individualized
absolute. As an individualized absolute every one of Pleroma's aeons
is then a particular totality, a relative Pleroma in itself, in other
words: an idea. That the one aeon has the other for its opposition
then means that the one idea has the other for its opposition: the
particular aeons have each its place in Pleroma, the particular ideas
have each its place in the system of the ideas.

It is a distinguishing mark of the gnostic functions that they cannot
reconcile the unbounded with the boundary or master concretion. As
regards the unbounded and the boundary, gnosticism hovers incessantly
between the oriental and the Greek way of viewing. It will in Greek
spirit individualize the absolute, and one should then to that extent
suppose that the most concrete aeons would have to occupy the highest
place in Pleroma; but on the contrary! Gnosticism, with its oriental
respect for
% --- p. 85 ---
\opage{85}the unbounded, the infinite, has a decided predilection for
the abstract; it is precisely the abstract aeons (βυθὸς and σιγὴ) that
occupy the highest or, if one will, the deepest place, the central
place, in the realm of ``the fullness''. The strife between the
bounded and the unbounded, a strife that originally smoulders in the
whole of Pleroma, then also comes at last to open breakthrough in the
catastrophe with Sophia.

It is an ingenious, profound trait in the Valentinian gnosis that
Sophia, the deeply grounding aeon of the aeons, this idea
comprehending the ideas, is seized by an irresistible longing to
behold the Propator. Instead of an absorption in the concrete,
wisdom's feminine, receptive aeon thus chooses a return to the
abstract, as to the originally productive; instead of cognizing the
indeterminate as negative determinateness in the concrete and thereby
determinate idea, Sophia, the bounded, will immediately embrace the
boundless one: the guardian of the boundary (ὅρος) is not able to hold
her back. Sophia's fall and restoration; matter, which in Kenoma's
darkness is formed out of her anguish, her tears and laughter; the
birth of the Demiurge (the world-aeon) from the fallen Sophia; the
creation of the world by the imprisonment of light in the darkness;
the redemption by the return of the soul-light to Pleroma, etc., all
this together presents a manifold of ingenious traits, which all show
how the imagining thinking has transformed the movements in the realm
of ideas into illusory events.

With the old gnostic systems scientific criticism has long since
finished; but gnosticism's marks of recognition: an ambiguous unity of
thoughtful imagining and fantastical thinking, ingenious turns and
spirited expressions instead of clear concepts and categorically
decisive determinations; direction to deeper beholding and dramatizing
depiction of supernatural events instead of penetrating investigation
of eternal fundamental relations: these and similar peculiarities
continue
% --- p. 86 ---
\opage{86}still\footnote{As e.g.\ in speculative dogmatics.} to draw
the logician's attention to the gnostic functions.

\emph{Theosophical Functions}. In the mythical word it is the
imagination that unconditionally has thought in its power; in the
gnostic word there is a unity of imagination and thought reflected in
approximation; in the theosophical word it is thought that immediately
masters the imagination, yet in such a way that, where logical
determinations are concerned, it prefers to make use of the
imagination. That theosophical thought has the imagination in its
power by a magic of the idea appears at once from this, that whereas
the mythically absolute is involuntarily parceled out into relative
figures, the finitizing relativity is consumed by the theosophically
absolute. But when thinking divests the absolute of all relative
determinations, the form of course vanishes: the form of the absolute,
the absolute form, is formless.

The theosophical is in this respect more nearly akin to the mystical
than to the mythical. The mythical too has indeed a strife between the
form and the formless; but it is the imagination, not thought, whom
the strife most nearly concerns. For theosophy it is expressly not the
imagination, but reflection, that is essentially interested in giving
form to the formless. The formless has in thought's absolute no
immediate, elementary determinate being, such as it has, so far as
mythology is concerned, in chaos and the oldest gods. The formless is
for the theosophist on the contrary reflected in the negation that has
removed the finite. And yet the theosophical consciousness cannot be
sated with empty infinity; theosophy's gaze is fixed upon the
in-itself absolute, upon God as God, thus upon an affirmative
determinateness.

Theosophical reflection is naive; the naivety is recognized by the
restless unity of the negative and the affirmative: this
% --- p. 87 ---
\opage{87}reflection of the form in formlessness recalls
mysticism\footnote{Mysticism makes the word itself negative. ``Gott ist
ein lauter Nichts, Ihn rührt kein Nun noch Hier, Je mehr Du nach ihn
greifst, Je mehr entwird er Dir.'' Angelus Silesius.}. But theosophy
is nevertheless different from pure mysticism in this, that it lays
weight upon determinateness and wills an absolute form for the
absolute. Thought's absolute is the idea; language's absolute is the
word; the task is to find an immediate expression for the absolute
ideality, or --- in naive reflection --- to find the name by which God
names himself, the word in which God reveals himself, the language
that the Godhead speaks.

But is it now also certain that God himself speaks a language? Is the
Godhead not the infinite depth (βυθὸς)? Must not the depth be grounded
in eternal silence (σιγὴ)? To this the theosophist must answer: ``A
Godhead without language, a God without a word is a dreaming Godhead,
a thoughtless God; for language is thought's revelation, all thinking
is essentially a speaking.'' --- But is it possible for a human being
to form for himself any concept of the Godhead's language, to think
the divine speaking? See, human language has in sound, in the
trembling of the air, a sensuous element, an immediate material, which
thought can naturally form into words and propositions. But now there
is surely no one who can represent to himself that the Godhead's mouth
should set either the air or the water or the ether in vibrations, in
order in such vibrations to have a sensuous element, an immediate
material, in which the divine thought could articulate its first
revelation? To this the theosophist must answer: ``The divine
language-element has its immediacy in the ideal power, but this
immediacy is not outwardly sensuous. How far the language-element
original in the Godhead's speaking ought most nearly to be compared
with something sensuous, with light-phenomena or with sound-phenomena
or with both, is not decisive for the ontology of language;
% --- p. 88 ---
\opage{88}but this is decisive, that the language-forming functions
must condition the thought-forming --- just as originally in the
ontological subjectivity as in the psychological; this is decisive,
that the language-element in its divine immediacy must bend itself for
thought itself into words and propositions and judgments.'' But how is
it possible that the divine language can, in likeness to the human,
consist in a system of propositions and judgments, when the Godhead is
therein to utter its essence's inwardness, an original unity of
knowledge and power? If the word is with God and itself God from the
beginning, then this word must surely, in divine unity of knowledge
and power, be absolutely substantial; but the human word, the impotent
word, is without substantiality. To this the theosophist must answer:
``The Godhead-word and the divine language are in the idea the same
and yet different. In language the manifold is uttered, the relative,
the many visions, representations, reflections and judgments; from
this side of relativity and manifoldness human language has in formal
respect an undeniable likeness to the divine. But what in language is
a relative is in the word an absolute. The Godhead-word is the
absolute itself. Here is seen the deep chasm between the human word
and the divine Logos. In the human word, the content-expression of
human thoughts developed by language, the idea can indeed mirror
itself; but the mirroring is only formal: the word developed in human
language cannot become one substance with the idea. It is otherwise
with the almighty Godhead-word: it develops its content in wisdom's
language, it unfolds its riches in systems of systems of divine
judgments; but while it during the development determines all
judgments, moves all propositions, penetrates, fulfills and completes
the infinite totality of concepts and form-determinations, it is
itself alone the totality's absolute and yet immediate form-unity.
With this absolute but immediate form-unity the word is the idea's
first revelation,
% --- p. 89 ---
\opage{89}is in absolute substantiality one with the idea, the idea
itself: the word is God.''

In the above speeches we have let the theosophist utter his guiding
thought according to the sense, for the sake of understanding, but not
in his own language, thus not in theosophical form. For the
elucidation of the theosophical form we cast a glance at Philo.

1) \emph{The Word and God in Ambiguous Unity}. The Philonic principal
category is, as is well known, the category Logos; but whether Philo
for the rest apprehends Logos as identical with God or as a being
different from God, whether this being is to be assumed to be personal
or impersonal, about this the interpreters have hitherto not been able
to agree\footnote{Cf.\ G.\ Steenberg. Om Synspunktet for Opfattelsen af
Philos Gudserkjendelse. Kjøbenhavn 1849. S.\ 84--125.}. The
theosophical reflection chooses in all naivety different points of
departure, without heeding the contradictions that arise when
heterogeneous ways of viewing are to be resolved into one another. On
the one side Logos is then presented as a being identical with God, as
God's thinking (λογισμὸς τοῦ θεοῦ), as God's word (ῥῆμα), God's oath
(ὅρκος), God's law (νόμος), God's covenant (διαθήκη). But just as
immediately Logos is set up on the other side as a being for itself,
as the bond (δεσμὸς) that binds together all things (περιέχων τὰ
πάντα), which, ``in order to hinder everything from going astray, is
thought of as stretched out from the one end of the world to the
other''; Logos is ``the divine wisdom that streams through the world;
the world is this wisdom's garment; wisdom itself is ὁ τῆς φύσεως
ὀρθὸς λόγος.'' Although Logos in both these apprehensions is
impersonal, yet other points of view again present themselves, from
which Logos is expressly seen as a personal being, as a god
subordinated to God himself (δεύτερος θεὸς, πρωτόγονος υἱὸς), etc.

% --- p. 90 ---
\opage{90}Theosophical reflection is everywhere recognized by this,
that thought, that which fundamentally moves, pushes itself in under
abstract intuitions, and lets the totality-moments remain standing in
immediate independence. Theosophical reflection is in its naivety
uncritical; if the naivety now vanishes, while the uncritical use of
intellectual intuitions is forced up into a higher science of
revelation resting upon faith, then we have dogmatic speculation. In
dogmatic speculation Logos is then also, as ``the heart of God the
Father, at the same time the eternal world-heart through which the
divine life streams into the creation. He is ground and source of all
reason in the creation'', etc.\ etc.

2) \emph{The Word and God in Ambiguous Opposition}. The theosophical
reflection, which posits the unity of God and Logos immediately,
posits at the same time the opposition immediately, and sublates again
the posited opposition just as immediately; hence that swarm of
contradictions which criticism so easily discovers. If these
contradictions are to be resolved, then the indeterminate unity must
be developed into opposition, in order that the opposition, the
totality-division, the idea-filled judgment, may again be resolved
into developed unity. It is, be it noted, developments of the concept
and not new revelations upon which in this respect what essentially
matters depends.

Speculative dogmatics assures us that on the ground of its new
revelation it finds itself highly exalted above the old ``Alexandrian
philosophy of religion''. ``One has often'', it runs, ``compared the
Logos-doctrine of revelation with the Alexandrian philosophy of
religion; indeed the assertion has often been uttered that John must
have borrowed his Logos-doctrine from the Jew Philo. But here the most
determinate opposition takes place. For Philo history transforms
itself into a world of seeming, the historical facts into an allegory
for universal thoughts and ideas. For Philo there is a yawning abyss
fixed between `word' and `flesh' (λόγος and σὰρξ); for sensuousness
and nature are to him only a prison of the spirit, and
% --- p. 91 ---
\opage{91}only by a mystical rapture away from actuality, only by the
mortification of sensuousness can the thinker be united with the divine
reason.''

If now in science it were merely a matter of bringing ``flesh'' into the
Word and ``word'' into the flesh; if the knot were untied as soon as
speculation, with pathos, merely went on repeating its λόγος σὰρξ
ἐγένετο: then dogmatics would undeniably celebrate a great triumph over
the Philonic ``philosophy of religion''. But if development of the
concept is really required here, it will be easy to show that dogmatics,
despite its many triumphant assurances, has not come a hair's breadth
further than Philo.

If ``sensuousness'' and ``nature'' are to be able to be taken up as
integrating moments in the Logos, the concept of the Logos must
originally be so determined that the real is present in the ideal, that,
in other words, immediacy breaks through from the beginning. This is now
not the case in Philo. The defect in Philo is just this, that the Logos
is one-sidedly apprehended as thought (λόγος νοητὸς), not at the same
time as the Word, that Word which in the identity of immediacy is one
with thought only insofar as it is thought's other, opposed to thought.
But the same defect repeats itself in speculative dogmatics; here too
there is no hint of the ideal reality of the language-element, or any
trace of the primitive immediacy of the Word; immediacy comes first when
we come to ``the flesh''. The defect in Philo is further this, that the
ideal world, the original content of the Logos (κόσμος νοητὸς), lacks
essential articulation, because here all consciousness is lacking of the
method valid for such an articulation\footnote{Cf.\ Grundideernes Logik
I. p.~178--220.}. But the same, the very same lack repeats itself in
speculative dogmatics. Here too we peer, in the midst of ``playing
possibilities'', ``plastic prototypes'' and ``magical actualities'',
% --- p. 92 ---
\opage{92}in vain for an articulating principle; for a reference to
``the free, artistic activity of spirit'' is naturally idle, so long as
one is not, by means of something resembling a concept, in a position to
state what is really meant here by the plastic art and the magical arts
of spirit. In Philo wisdom (σοφία), the Logos and the powers (δυνάμεις)
are separated and mingled with every possible arbitrariness: now Sophia
is the source from which the Logos springs, now again the Logos is the
source of Sophia. The same, the very same up-there-and-down-here of the
representations we find again in speculative dogmatics, when ``the
magical visions that play before God's countenance in the inner
self-revelation must be thought to determine themselves for God as
decrees for the creation'' etc. Philo lacks consciousness of the fact
that the logical in the ideal system (κόσμος νοητὸς) lets itself be
developed only in the reflection of the antilogical moments of the real
system\footnote{Grundideernes Logik I. p.~224--229.}. The weak point in
Philo is therefore the obscurity with which he hypostatizes the Logos
now as a divine person (δεύτερος θεὸς), now again as world-essence, the
cosmic principle, (ὁ τῆς φύσεως ὀρθὸς λόγος), the reason indwelling in
the world, which sows the germs of life (λόγος σπερματικὸς). But the
same, the very same obscurity reigns in speculative dogmatics, and for
the same, the very same reason. Without theosophical naïveté dogmatic
speculation knows how, in a naïveté of its own, to operate both with
``the Logos of the Father'' and with ``the world-Logos'', and with ``the
world-germs'' (λόγος σπερματικὸς) and with ``the second God'' (δεύτερος
θεὸς) and with ``the originating ideas'' and with other things of the
sort.

3) \emph{The Word and God in Developed Unity-in-Opposition.}
Logic is not theosophy; but in the theosophical functions there works
something rational which logic must strive to clarify. When the form of
the judgment is to contain the idea, the absolute content, there shows
itself a fundamental division, whereby the one
% --- p. 93 ---
\opage{93}and the same totality of totalities divides itself into two
originally opposed forms, the form of selfhood or of subjectivity and
the form of otherness or of objectivity. With the form of selfhood the
absolute content is reflected into unity, essentiality, inwardness and
ground; in this inwardness the depth is the silent thought that breaks
the silence. As the infinite subjectivity which in the deep silence
lifts the silence and utters its content, the Godhead is not the Word
itself, but \emph{the origin of the Word}. With the form of otherness
the absolute content is articulated in manifoldness, in difference, thus
a totality distinctly marked in its members; the logical element of
speech, that is to say, of the revelation of thought, a primitive
immediacy in which the judgment separates what is united and unites what
is separated, is as totalized unity \emph{the Word}. The Word in its
immediacy is not the Godhead, but its revelation; not subjectivity
itself, but its objectified content. The absolute totality has by the
two opposed forms become two opposed totalities: the Godhead and the
Word.

But when the form of the judgment is to contain the idea, the absolute
content, the fundamental division, whereby the subjective totality
resting in selfhood is set over against the objective totality in the
immediate-mediate unfolding of otherness, must necessarily have a
connecting link, a middle, a unity-in-opposition, in which the extremes
meet. By virtue of the totalized unity-in-opposition the Word in its
opposition to God is one with God: the Word is God. In its unity with
God the Word is the objectified God (ὁ δεύτερος θεὸς).

As the self objectifying itself God is the origin of the Word; as the
objectified self God is identical with the Word. But as the God
objectified in himself the Word is not merely one with God, but by
reason of the objectification at the same time opposed to God: in its
objectified opposition to God the Word is not God, but a world in God,
the ideal world (κόσμος νοητὸς).

% --- p. 94 ---
\opage{94}But how does it really come about that the same absolute
content, this same totality objectified in the Godhead, can become both
God himself (πρωτόγονος υἱὸς, δεύτερος θεὸς) and a world in God
different from God? It lies in the idea of the judgment! Here there is
in the objectified totality a decisive tension between selfhood and
otherness, reflection and immediacy, the unity of the content and the
content of the unity. If the emphasis falls upon selfhood, and upon the
content-unity reflected in selfhood, then we have in the objectified
totality the objectified Godhead, God himself in his own objectification.
If, on the other hand, the emphasis falls upon otherness and upon the
development by which the content of the unity is totalized in the
immediate-mediate relational determinations of the parts, then we have
the objectified totality as a connection-unity articulated in
manifoldness and difference, thus a world in ideal determinate being, a
world in God, different from God, an ideal world.

But the ideal world, the ideal world-totality, could not possibly form
an essential opposition to the Godhead objectified therein, if this
ideal world were not originally in totalized opposition to its real
world. The ideal and the real world have in reality the same content,
but in ideality, by the difference of form corresponding to the idea of
the judgment, to the fundamental division, they have each its own
content: in its ideal world the Word is a word of knowledge, but in its
real world it is a word of power. Only by unity-in-opposition with the
word of power does the word of knowledge, and with it the ideal world,
acquire ideal independence.

It has been developed from several sides in the doctrine of
subjectivity\footnote{Grundideernes Logik I. p.~210 ff.} that ``the acts
of thought'' in the ideal system cannot possibly coincide with ``the
acts of nature'' in the real system. But in the first treatment of the
problem it could not yet be seen
% --- p. 95 ---
\opage{95}in what manner the idea-movement in ontological subjectivity
could free itself from the real movement in the outer world of actual
objects; now, on the contrary, the problem is solved. If the Word in the
articulation of propositions and judgments is a divine language; if the
thoughts of the Godhead have inner, ideal existence in the ideal
language in which they have their first revelation; if the word in which
the infinite wealth of uttered thoughts has totalized unity is the word
of the divine knowledge: then it follows from this that the ideal world
must have inner independence in a divine thinking and speaking.

The ideal and the real world are, however, not merely in opposition, but
--- what indeed also lies in the idea of the judgment --- at the same
time in unity-in-opposition. If the judgment corresponding to the idea
demands an opposition of the Word in knowledge and the Word in power,
then the idea of the judgment must necessarily also demand a totalized
unity of the word of knowledge and the word of power.

\begin{center}{\bfseries β)\quad Judgment and Object: Functions of the Word of Power.}\end{center}
\addcontentsline{toc}{paragraph}{β) Judgment and Object: Functions of the Word of Power}

It has been shown from several sides in ``the Logic of
Subjectivity''\footnote{Grundideernes Logik I. p.~145--166; 201--204;
255--263.} how the same content can, in the form of knowledge and in
that of power, in the logical form of the concept and the antilogical
form of existence, transform itself and become a double content. This
doubling was then reflected through by means of the word, without its
yet being possible to reflect upon the word. In the present connection,
where reflection is expressly directed upon the word as word, the task
returns with heightened emphasis; the task is to show how the
concept-content in the word of knowledge and the existence-content in
the word of power presuppose and determine each other to infinity.
Inasmuch as the concept acquires ideal independence in the word of
knowledge, while existence steps forth as an object for itself
determined by the word of power, they at first glance fall so far apart
from
% --- p. 96 ---
\opage{96}one another that it seems as though they were wholly
indifferent toward one another. On the one side it indeed looks as
though the concept, after it has received ideal determinate being in the
word, could be formed, developed and held fast either for itself or in
connection with other concepts, without regard to whether there were in
existence any object corresponding to this concept or not. On the other
side it appears as though the existing objects, subsisting in the word
of power, were wholly indifferent, not indeed toward being objectified,
but yet toward the peculiar development of the concept to which their
objectification may give subjectivity occasion. How, then, do we
determine the relation between concept and thing, judgment and object?

Formal logic treats the task in its own way, in that, for orientation in
the domain of thinking, it makes a distinction between our concept of
the object and the object's concept, but not, on the other hand, between
our judgment about the object and the object's judgment; it
distinguishes between the object's concept and the object itself, but
not between the object's judgment and the object; thus between the
concept in the subjective and the concept in the objective sense, but
not between the judgment in the subjective and the judgment in the
objective sense. We shall now see what value can be attached to these
distinctions.

The difference between the concept in the subjective and the concept in
the objective sense cannot possibly, when the concept itself is looked
at, be any logical difference. Were there any logical difference here,
it would have to be conditioned by our subjective concept of the thing
containing either more or fewer or at any rate other determinations than
the thing's objective concept; but in that case the \emph{subjective}
concept would become the same as a \emph{false} concept, which cannot be
the meaning. But if the subjective concept has the same logical content
as the objective, then the difference falls away: the subjective concept
is in its truth objective; our true concept of the thing is the thing's
concept.

% --- p. 97 ---
\opage{97}Just as untenable is the distinction set up between the
object's concept and the object itself. The object's concept, it is
said, is the ``intellectual'' and the ``underlying'' in the object, that
which constitutes the object's ``nature and essence.'' But how is it
possible to separate an object's nature and essence from the object
itself? What remains of an object when its nature and essence are
thought away? If that in the object which does not constitute its nature
and essence does not belong to its concept, then it must be something so
contingent that it really does not concern the object at all; if, on the
other hand, there must necessarily be here, besides that which
constitutes the object's nature and essence, yet something else,
something which, without concerning its ``nature and essence,''
essentially and necessarily belongs to its existence, then it lies in
the object's nature and essence to have something which does not concern
its ``nature and essence:'' what then may that something be?

It is said that the object itself is ``the real outward'', the sensuous
exterior, whereas the concept is ``the inner, ideal and underlying''.
But when ``the real outward'' constitutes the object itself, the actual
object, then ``the inner, ideal and underlying'' must in actuality fall
outside the object\footnote{``A. Here, then, in one and the same
existence there are two existents, namely the existing object and the
concept existing in the object? C. No, it cannot be taken that way; it
is not the object itself, but its material, that is opposed to the
concept. A. Then presumably we get three existents: the existing
material, the existing concept, and that existence of the object in
which both exist? C. No, neither can one take it that way. A. How then
is one to take it? C. The material by itself is not a separate existent,
any more than the concept is, but the material and the concept are two
sides, namely two existence-sides of the same existing One, the object.
A. The existing One is, then, in opposition to and over against the two
existence-sides, that which exists in and for itself? C. No, neither can
one take it that way, for if we let abstraction take away the two
existence-sides, the existing unity will itself become an abstractum, an
existence-side that does not exist. A. But how then is one to take it?
C. In such a way that the existing unity is seen to exist in the two
existence-sides, which then precisely thereby come to exist in
inseparable unity \dots. A. But what is material in the object is surely
something sensuous, which as sensuous cannot be thought, and what is
intellectual in the object something non-sensuous, something thinkable,
which cannot be sensed? C. Perfectly correct. A. In the existing unity,
then, the unthinkable is nevertheless a thinkable, the non-sensuous a
sensuous. C. That I do not understand!'' A. But you do, on the other
hand, no doubt understand that the ideal, the inner and the underlying,
which constitutes the object's ``nature and essence'', is not the real,
the outer, the existing, which precisely constitutes the nature and
essence of the object, of the actual object! Cf.\ Forelæsninger over
phil.\ Propæd. Universitetsaar 1860--61. p.~26--27.}. Nor can
% --- p. 98 ---
\opage{98}it help that one would make the inward ``the constitutive,''
the outward ``the consecutive''; for in the object as object, and here
there is question only of the actual object, the inward is object just
as much as the outward: here there is not an object that goes around a
concept, not a concept that sits inside an object; here there is, in the
object, within and without, only object.

How obscure the distinction between the object and the concept existing
in the object is, is best seen when the relation between concept and
judgment is to be more closely determined. One cannot --- so we are
assured in formal logic --- use the word judgment in an objective sense,
as one uses the word concept: ``one cannot say that in the same manner
in which a concept can come to meet us in an object, a judgment too can
come to meet us.'' And why not? Because one has once for all got it into
one's head that inside the objective object there sits an objectively
existing concept. But as soon as one
% --- p. 99 ---
\opage{99}by carrying reflection through has got the better of this
fixed idea, and has seen that the concept which ``can come to meet us in
an object'' is precisely our own concept corresponding to the object,
one will also at once see that the judgment, the true judgment
corresponding to the object, has the same objectivity as the true
concept. That no judgment can come to meet us in the object is indeed
certain, insofar as there cannot be given a judgment about the object
objectively existing in the object; but the same holds of the concept.
If the object cannot come forward with a judgment about itself, with a
judgment objectively judging itself, then neither can it come forward
with a concept of itself, with a concept objectively comprehending
itself. If the judgment about the object falls within the judging
subjectivity, then the concept of the object, the object's objective
concept, assuredly falls within the comprehending subjectivity as well;
if the concept does not lose its objective validity because it must
absolutely be comprehended by subjectivity, then neither does the
judgment lose its objectivity because it is a subjectivity that judges.

The relation between concept and existence, between judgment and object,
is in principle a fundamental relation between word and word, between
the word of knowledge and the word of power. The content is the same;
but the form is different: in the word of knowledge the content is
concept, and the concept is uttered in the judgment; in the word of
power the content is existence, and the existent is the object of the
judgment. Seen from the side of the concept it is nonetheless the form
that becomes the same, while the difference falls in the content; and
the same holds when the matter is seen from the standpoint of the
object. To be sure, in the concept of the object, if one will, in the
object's concept, a difference must be made between what in the act of
comprehending belongs to the concept and what on account of existence
belongs to the object; but so long as the separation of conceptual and
object-determinations falls under the concept, the concept of the
object-determinations must anew resolve itself into conceptual
determinations. The form for both
% --- p. 100 ---
\opage{100}kinds of determination thus becomes the same; the difference
falls upon the content: in the totality of the conceptual determinations
the content is logical; in the concept of the totality of the
object-determinations the content is antilogical.

``But the antilogical content must surely have a corresponding
antilogical form; the object-determinations are antilogical in form as
well as in content;'' so runs the judgment about the object from the
standpoint of the object. Since now the judgment about the object cannot
be separated from the judgment about the judgment without the separation
becoming idle, the judgment runs to the same effect from the standpoint
of subjectivity as well. The object-determinations expressed in the
judgment will then, for the judging consciousness, set themselves
sharply against the universal, will in the object's name individualize
themselves with antilogical punctuality, and the form will thus be
antilogical. The form turned toward the logical then now stands for
consciousness in a logical mirroring of an antilogical and yet logical
content. The doubling of the content shows itself again: in itself, as
object-content, the content is antilogical like the form, but in the
\emph{logical} mirroring of the \emph{antilogical} form the content is
logical.

To determine the relation between judgment and object is by no means the
same as to show the identity of judgment and content. The latter task is
logical in the sense that it can restrict itself to a development of
pure conceptual determinations; the former task, on the other hand, is
logical in the sense that the antilogical itself demands a logical
explanation. It is the logic of the antilogical forms, one of the most
difficult portions in the whole of philosophy, that matters in this
connection. If we had to do only with logical forms and corresponding
content-determinations, the functions would all without exception be
functions of the word of knowledge; the highest would be attained when
the form of the judgment was so far developed that the fundamental
division corresponded to the idea. The antilogical sting, which, so to
speak, palpably affects the
% --- p. 101 ---
\opage{101}content and brings it about that the object of the judgment
sets itself sharply against the judgment, without the identity of
judgment and content being allowed to be given up, awakens on the
contrary everywhere a knowledge of opposition between power and
knowledge, and makes the word of power the foundation of the functions.
The functions are logical; the word of power is antilogical; the
infinite reciprocal determination of the logical and the antilogical
which hereby arises requires an illumination of judgment and object from
the standpoint of knowledge, from that of power, and, by
unity-in-opposition of knowledge and power, from that of the idea. We
designate these standpoints in what follows by: \emph{Knowledge of
Judgment and Object}, \emph{Power in Judgment and Object}, \emph{The
Idea in Judgment and Object}.

\textit{αα}) \emph{Knowledge of Judgment and Object: Theoretical
Functions.} In psychological dualism the surrounding world with its
objects stands outwardly over against the contemplating, thinking,
judging consciousness. That the object in its immediate independence
falls outside the judgment; that the judgment concerns only the
consciousness of the object, while the object itself is indifferent
toward consciousness: this one-sided way of regarding the matter is to
such a degree rooted in so-called sound common sense that philosophy
must for that very reason find itself called upon from all sides to work
against its one-sidedness. No doubt in the immediate sense the actual
object falls outside the judging consciousness; but what essentially
makes the separation between the object and consciousness is precisely
the judgment about consciousness and object. That with respect to the
object's independence it depends just as essentially upon the judgment
as upon the object, is something that escapes the immediate sense. To
the immediate sense the matter appears as though the object came to meet
consciousness by itself from without, while yet it is just as originally
consciousness that comes to meet the object from within; ``sound
understanding'' overlooks that the object is only outside consciousness
because it is judged and by the judgment
% --- p. 102 ---
\opage{102}set to stand over against consciousness. Take away the
judgment, that judgment of instinct upon which all separation between
the subjective and the objective, the outward and the inward,
essentially rests, and the object will instantly coincide with
consciousness in such a way that here there is neither object nor
consciousness.

If the judging subjectivity posits the separation between itself and the
object in the judgment, then it is also seen from this that the
distinction between judgment and object must likewise rest upon the
judgment. But that the distinction between judgment and object rests
upon the judgment presupposes in turn that the identity of judgment and
object rests upon the judgment. That there is here in the judgment a
knowledge of the object in its express opposition to the judgment is
grounded, then, in there being here in the judgment a knowledge of the
object's identity with the judgment: the double knowledge is in the
judgment itself one and the same knowledge. In the reflection of this
two-sided knowledge, hovering above judgment and object, the functions
are \emph{theoretical}.

It is a distinguishing mark of all theory that the object of knowledge
determines knowledge, and not conversely; but here there is now a double
knowledge, a knowledge of the judgment in its relation to the object,
and a knowledge of the object in its relation to the judgment;
knowledge has, in other words, a double object: an \emph{ideal} one,
namely the judgment as judgment, and a \emph{real,} the object as
object. Theory will not arbitrarily alter anything either in the
judgment or in the object, but take them in their reciprocity, as they
are in themselves. The problem is seen from the following questions: Is
it the judgment that essentially determines the object, or the object
that determines the judgment, or is there here a reciprocal
determination? Does the truth of cognition rest upon the mere likeness
between the determinations in the judgment and the object-determinations,
or must a transformation take place here in order that the
object-determinations
% --- p. 103 ---
\opage{103}may become identical with the knowledge-determinations
contained in the judgment?

If theory starts from the presupposition that the determinations in
judgment and object are heterogeneous, and that between heterogeneous
determinations there can be given only a parallelism on the ground of
likeness, not a transition conditioned by transformation, then agreement
as mere and bare parallelism steps into the place of identity, and
psychological dualism must then seek its principle of cognition in a
presupposition standing above consciousness. Thus:

1) \emph{Parallelism of Judgment and Object: Supposing Functions.} The
existing object is in its outward immediacy something sensuous, the
content of the judgment something thinkable; the sensuous is
heterogeneous with the thinkable: it lies near at hand, then, to
apprehend the sense-phenomena and the thought-determinations as two
series which may indeed yield a parallelism, but no identity.

Thus it is, according to Plato, not the sensuous in things, but the
ideas in which the things of sense have come to partake, that let
themselves be thought. That judgment and object can agree in true
thinking is conditioned by the same ideas which reveal themselves
subjectively in thinking revealing themselves objectively in things.
Whence, now, does the judging subjectivity know that the ideas reveal
themselves objectively in things, or, what is the same, that things
partake of the objective ideas? The assumption rests upon a
presupposition immediately slipped under by thought itself. How
immediate and objective this slipping-under is, is seen from the myth of
pre-existence with the corresponding doctrine of recollection
(ἀνάμνησις).

What is characteristic of the supposition is the naïveté with which
subjective reflection comes objectively by its own assumptions. First it
is assumed that thought can think only the ideas, since the ideas alone
are beings of thought. By holding fast to this assumption thinking would
obviously not
% --- p. 104 ---
\opage{104}get beyond its own subjectivity; the judgment would in other
words have only a logical content, not an antilogical object. Since,
however, the object in its outer independence obtrudes itself
involuntarily, thought is, without rightly knowing of it, moved to lay
its subjective content objectively into things and thus to make these
partake of the ideas. The ideas are now in things themselves objective,
not subjective, since in thinking they are subjective, not objective.
The original agreement of subjective and objective, upon which thought
builds as upon an objective given, it has subjectively given itself.

The slipped-under foundation of the supposing functions comes forward
still more strongly in Cartesius and, after him, in Malebranche and
Leibnitz. The presupposition is a separation of the intellectual and the
material, determined by a gulf between thinking and extension
insurmountable for the relative substances. A direct reciprocal action
of the representations in the circle of consciousness with the things in
the surrounding world is from this standpoint wholly impossible. As
little as a representation or a concept can become an efficient cause of
a change in the world of things, just as little can a physical
impression bring about the formation either of a corresponding
representation or of a concept. That the representations in
consciousness and the things in the surrounding world, without
thoroughgoing reciprocal action, merely form parallels, is equally
decisive for theory, whether with Malebranche we appeal to an
omnipotence working unceasingly through occasional causes, or with
Leibnitz to a \textit{harmonia præstabilita.}\footnote{``Material
objects'', says Malebranche, ``cannot possibly be apprehended by
themselves, they are apprehended only by the ideas corresponding to them
in the soul. The opinion of the Peripatetics, that the outward objects
should send us their images (\textit{des espèces, qui leur resemblent}),
and that these should then through the outer senses be brought to the
common sense (\textit{au sens commun}), is unreasonable; these
impression-images (\textit{espèces impresses}) would then themselves
have to be small bodies, and these bodies, with which space would be
filled, would then, by reason of the impenetrability of matter, have to
crush and jostle one another, some going to one side, others to another,
and so could not make the objects visible.''\par
The opinion of those who assume that the ideas of objects are created
with us, so that God does not need at every moment to awaken them in the
soul when we have need of them, Malebranche refutes with the following
consideration: ``We have, even as regards simple figures, an infinite
multitude of ideas: ellipses, triangles, polygons and the like can with
respect to proportions of magnitude be imagined in countless ways. It is
then highly improbable that God should have created such a swarm of
ideas in and with the human spirit. But suppose that we also had such a
store of co-created ideas (\textit{un magazin de toutes les idées}), it
still remains wholly inexplicable how we come to apply the right idea
when the corresponding object shows itself. We see, e.g., the sun. The
image which the sun impresses in the brain has no likeness to the idea
we form of the sun; nor does the soul notice the very least of the
impression which the sun makes either upon the bottom of the eye or in
the brain: it is thus impossible to explain how among the countless
ideas we come precisely to choose the idea of the sun when we see the
sun'' \dots\ By these and similar reasonings Malebranche then arrives at
the result that spirit can see what it is that is in God which presents
``the created things'' etc. Cf.\ Om Theologiens Naturbegreb med særligt
Hensyn til Malebranche: \textit{De la recherche de la vérité.} Af R.\
Nielsen. I Anledning af Reformationsfesten 1855. p.~13--15.\par
What Malebranche refers to God immediately, Leibnitz refers to God
mediately. The presupposition is that the monad of monads has so
constituted each single monad that it can as a living mirror
(\textit{miroir vivant}) mirror the whole universe. But in this
mirroring one must not think of any immediate communication between the
monad and the surrounding world. The physical mirror receives from the
object itself the light by whose reflection the image is formed; but the
metaphysical mirror forms its images out of its inner constitution
without immediate influence from without: between these images and the
objects of the surrounding world there can therefore be only an
agreement secured by parallelism, not a causal relation.}

% --- p. 105 ---
\opage{105}Although Cartesius must admit to himself that the essence of
a triangle, e.g., is so perfectly determined that the many
% --- p. 106 ---
\opage{106}qualities which can be proved of it depend upon thought in such a
way that the thinker must cognize them whether he will or not, he still dare
not, when the question is of the truth of the mathematical propositions, rely
straightforwardly upon the testimony of thinking, but builds upon the
assurance that \emph{God is and does not deceive.} In its relying upon God,
thus in being certain, not of itself, but of an other than itself, thought is
immediately certain of itself: this is the contradiction. That thought has
itself slipped in underneath the Absolute to which it appeals in the last
instance is, however, a fact about which the supposing functions can
themselves give no information.

If thinking must seek truth's guarantee outside itself when it occupies itself
with its own logical content, how much more then when the talk is of the
things existing in the corporeal world. That corporeal things --- so Cartesius
reasons --- can exist follows from this, ``that God can create everything that
I am thus able to apprehend. But now there is in me a certain passive capacity
for sense-perception, or for receiving and cognizing the ideas of sense-things.
For this capacity I should have no use if there did not also exist a certain
active power, either in myself or in others, to bring forth these ideas. This
power cannot, however, be in myself, \dots\ since those ideas are brought forth
without my co-operation, often against my will. It must thus be in another
substance, in which all that reality is contained objectively
(\textit{formaliter}) or in a higher degree (\textit{eminenter}) which is
subjectively in the ideas brought forth by it. This substance must either be
God, or
% --- p. 107 ---
\opage{107}a created being nobler than the corporeal, in which these ideas are
contained in a higher degree, or the corporeal itself, in which they are
contained precisely objectively. Since now there is given me no capacity
whatever for cognizing that the ideas should be brought to me by God himself,
either immediately or through another created being, but on the contrary a
great inclination is given me to believe that they are sent out by the
corporeal things, I cannot see otherwise than that God would have to be a
deceiver if they were not really sent out by the corporeal things: the
corporeal things therefore exist.''\footnote{Cf. \textit{Meditatio Sexta.}
p.~35; 39--40. (\textit{Opera philos.})}

On the basis of the supposing functions thing and thought, judgment and object
do indeed stand outwardly against each other; the mechanical and physical is
heterogeneous with the intellectual and metaphysical, the material with the
spiritual: but what is nevertheless the point of the problem, the opposition
between the concept and its existence, between the same content in logical and
in antilogical form, between the same form with logical and antilogical
content, cannot on this standpoint come to clear consciousness all the same.
The antithesis of the concept and the object is blunted when the category idea
is applied just as immediately to the object as to its concept. The question
how the idea is able to come forward in two such dissimilar shapes as a
corporeal thing and a representation evaporates when both shapes are after all
assumed to resemble each other in this, that they are both ideas. To this
comes the fact that, in consequence of the ground slipped in underneath, one
has at every difficulty omnipotence, thus a word of power, a divine power-language, to refer to in all arbitrariness.

In the supposing functions thought has originally gone out from itself, and
yet it supposes itself to be at home with itself; this contradiction must
become manifest.

2) \emph{The object determines the judgment: abstracting
% --- p. 108 ---
\opage{108}functions.} If thought goes literally out from itself in order to
begin with an other, then the first and nearest thing it encounters is not the
Absolute but the relative, not a God but a sensible thing. Insofar as one here
begins with the sense-impressions proceeding from the object, one here begins
with the existential extreme, the antilogical. The a priori concepts, ``the
innate ideas,'' have here vanished; only the immediate forms of existence, the
sensible object-determinations, remain. We have then again the sensualist
doctrine of experience. That Locke lets the soul begin as a \textit{tabula
rasa} acquires in the present connection the significance that he makes the
attempt to let all conceptual determinations arise through the transformation
of individual object-determinations. It is then not the judgment which by
presuppositions \textit{a priori} is to determine the object; it is on the
contrary the object which by existence's word of power is to determine the
judgment. If we consider the theory somewhat more closely, as it is laid out by
Locke and developed by Hume, it is chiefly two factors that play the principal
role, namely: \emph{the impression} and \emph{the abstraction}.

By the impression the existing object is supposed to imprint its individual
form in the soul in a manner similar to that in which the seal imprints its
characters in the wax. If it now remained at mere impressions, there would
surely be found as many individual forms in the soul as there are objects the
soul has apprehended. The form of the object's image in the soul and the
object itself in existence would in this way be the same; but the content
would be different, insofar as the object-image in the soul only was and
remained an image, a copy of the actual, while actuality itself, the copy's
original, would have to be sought in the existing object. But however passive
one may wish to make the soul that is subjected to the impressions, the soul
is nevertheless essentially
% --- p. 109 ---
\opage{109}subjectivity, and subjectivity self-activity.\footnote{Cf.
Grundideernes Logik. I. p.~25.} Beneath the psychical passivity self-activity
comes to view in the transformation the object-form undergoes as it is taken
up into the soul: in the object itself the form of existence has the
individual stamp of singularity; apprehended by the soul it assumes the stamp
of generality; it is a thoroughgoing transformation, and this thoroughgoing
transformation is due to abstraction.

Since abstraction does not add anything but on the contrary takes something
away, it seems at first glance as though it were the form of existence itself
that was preserved pure and unfalsified in the form of consciousness. What is
it, then, that falls away through abstraction? The contingent, the merely
individual, the changeable, the inessential! And what is it that is expressly
held fast through abstraction? The common, the universal, the constant, the
essential! But the object-form is in its objective determinateness
antilogical, the form of consciousness in subjective mobility essentially
logical; through abstraction the miracle must then have taken place that an
antilogical form has become logical; thus an \emph{other} form, and yet
\emph{the same} form --- how is this to be understood?

Abstraction, this receptive, negative function, which was supposed to aim
solely at holding fast the object-form as a pure, unadulterated form of
existence freed from the contingent, has evidently done something to the form,
has, so to speak, on its own account reworked it, has transformed it from an
antilogical individual form into a logical universal determination, from a
form of existence into a form of the concept. Abstraction is then not so
abstractly receptive that it cannot also be productive. Abstraction is a
function of a priori selfhood; but the thinking that is lost in abstracting
functions does not discover it.

The result of the theory is that both the form of the concept and the form of
existence are dissolved, the judgment is suspended and the object is
% --- p. 110 ---
\opage{110}lost. The concept must necessarily vanish when the universal, that
which is necessary in itself, vanishes. Necessity and universal validity are,
namely, distinguishing marks of all true conceptual determinations; but the
necessary and universally valid does not allow itself to be abstracted from the
immediate singularities of facts. Hume does indeed maintain that categories
such as likeness, degree, magnitude, etc., yield immediate certainty through
abstraction and the corresponding reproduction; ``for when an object is like
another, the likeness will at once affect the eye, or rather the mind, and
seldom does it require a second examination; in the determination of magnitudes
or numbers we can proceed in approximately the same way,'' etc. With the
categories cause and effect, on the other hand, matters are supposed to stand
in quite another way. This is now evidently an error; for if causality does not
allow itself to be abstracted from facts, then the same holds of likeness; and
if likeness does allow itself to be abstracted, then the same holds of
causality.

Why does the causal relation not allow itself to be abstracted from facts?
Because this relation as relation cannot become an object of immediate
experience, because in other words the relation itself is not a sensible
object. But now likeness! is likeness perhaps a sensible object? Hume assumes
that the likeness in what is alike immediately affects the eye; that is a
childish way of representing the matter. The categories likeness and unlikeness
arise only through comparison; but things do not compare themselves with one
another: comparison is a function of the judging, abstracting subjectivity.
Comparison is a reflection that finds likeness only by presupposing likeness;
in order to discover the likeness in the thing, the reflecting subjectivity
must itself bring the likeness with it, must proceed from the category of
likeness. The same, the selfsame, holds of causality. The difference is only
this, that the causality-reflection is involved, while the likeness-reflection
falls out as if of itself.

% --- p. 111 ---
\opage{111}What does it now mean that likeness is abstracted from like
objects, unlikeness from unlike? But no! it is in this connection more
important to determine what it does not mean. For it by no means means that
the sensible thing, besides its other qualities, should also have the quality
likeness, and that abstraction could then, by removing the other qualities,
retain the likeness. Let two things be alike in color, let them e.g. both be
red; there is nevertheless no one who assumes that in the like color there are
two sensible qualities, the color and the likeness; the likeness in the color
is not sensible in the way the color is. The sensation of color rests upon a
sense-impression, but the likeness rests upon a judgment. Abstraction is
determined by the judgment; in a judgment such as ``this is a color'' the
color is sensible in the subject, non-sensible in the predicate, antilogical
in the subject, logical in the predicate; only thus can the category in the
predicate be an abstraction from the object in the subject.

This Hume and the sensualists have now quite overlooked. From the color on A
and the color on B they will without more ado abstract the likeness, just as
though abstraction consisted in scraping the likeness off the color, as one
can in certain cases scrape the color off the object. Only when abstraction is
reflected in the judgment does the illusion disappear: A and B are alike in
color; here the determination ``color'' points back to immediate
sense-perception, but the predicate ``alike'' belongs to reflection.

It is, logically regarded, the axiom of sensualism that the object is to
determine the judgment. In this lies the tacit presupposition that the existing
object has in and for itself a sensible objective independence, independent of
every judgment. But this presupposition too dissolves for the abstracting
functions. Already Locke was embarrassed by substance; his incipient doubt was
developed in Hume: for Hume there is no substance. ``If,'' he says, ``the
concept substance is conveyed to us through the senses, then I
% --- p. 112 ---
\opage{112}ask, which sense gives us this concept, and how does it give it to
us? Were it received by the eyes, then it would have to be a color; by the
ears, then a sound; by the palate, then a taste, --- and so with all the
remaining senses.''

Hume comes consistently to the result that by the concept substance nothing
else can be understood than a collection of single qualities (\textit{a
collection of particular qualities}), a gathering together of simple concepts
which are united by the imagination and designated by a name of their own.
``Thus we first discover in gold the yellow color, the weight, the
malleability, the fusibility; afterwards we discover that it allows itself to
be dissolved in aqua regia; this quality we then connect with the others, and
reckon it just as much to the substance as if the concept of it had from the
beginning constituted a part of the composite whole.'' What then is gold? As
object-unity gold is a name, an aggregate of heterogeneous qualities which not
actuality, not existence itself, but ``imagination'' has united. By isolating
the qualities and apprehending each single quality by itself, abstraction is
led to deprive the object of substantiality; but deprived of its
substantiality the object is at the same time deprived of its objective
independence.

Substance is an imaginary magnitude which, according to Hume, ``arises from
this, that one ascribes to the particular qualities an unknown \emph{something}
in which they are assumed to inhere.'' The word of power that was to secure for
the object a factually determinate, substantial objectivity is thus through
abstraction transformed into a word of impotence, a name; for substance, this
\textit{caput mortuum} of abstraction which remains when all qualities fall
away, is only a name: if the object is to determine the judgment in this way,
then there is here neither judgment nor object.

3) \emph{The judgment determines the object: determining functions.} By going
out from itself and seeking its determinations of content in the existent
object,
% --- p. 113 ---
\opage{113}thought has from the ground up disturbed both its own essence and
that of the object. In order to remedy this fundamental disturbance, thought
turns back from the dissolution of the object to its a priori point of
departure. The beginning is made from the subjective extreme, in that the
judgment determines the object.

That the judgment determines the object must naturally not be understood as
though subjectivity could now judge arbitrarily according to its own fancy, and
the object were then without more ado to conform itself to the judgment. The
truth is on the contrary that the essence-determinations decisive for the
cognition of the object must, in order to be object-determinations, absolutely
be conceptual determinations; it is in consequence of this theory that the
functions of the judgment show themselves as \emph{determining}.

How impossible it is to assume object-forms that are not originally forms of
subjectivity, Kant has sufficiently demonstrated. The one-sided element in Kant
lies in this, that by subjectivity he can think only a psychological
subjectivity, whereas the task is conversely, in the psychological
subjectivity, to think only of the subjectively universally valid, of that
which is essential to all subjectivity. When Kant declares the intuitions of
space and time to be subjective forms, we grant him his right; but when he
thereupon at the same time assumes that these forms cannot be objective because
they are subjective; that they are only forms for \emph{our} subjectivity,
without its following therefrom that these forms also have validity for a
subjectivity of a higher order, if such there be: then we cannot grant him his
right.\footnote{Grundideernes Logik. I. p.~77--80.}

The acuteness with which Kant in his ``transcendental analytic'' derives
categories of the understanding from the judgments has awakened admiration and
deserves admiration. It is hereby demonstrated once and for all that no one can
judge about the thing's objective being without his, by reflecting upon the
judgment, at last
% --- p. 114 ---
\opage{114}having to become conscious that objective being is posited by a
category lying in the understanding itself. The being object \emph{is}, insofar
as subjectivity adjudges being to it. If anyone will hold that the object's
objective being is in reality independent of the subjective being in the
judgment, let him go back to Hume and take what follows upon it. But when Kant
then sees in the distinction of being and non-being decisive for the judgment's
quality --- ``reality and negation'' --- only a presentation of
psychological-subjective forms, forms of the understanding which are so
subjective that they cannot possibly fit the thing in itself, then he is in
error.

When the thing is judged about as phenomenon: this thing is hard, red, etc.,
then Kant knows very well that the subjective being in the judgment's copula
does no damage to actual existence, but on the contrary determines the
phenomenon's objective being. For matters do not stand thus, that the
phenomenon, the object of cognition, has a determinate being independent of
the judgment and therefore indifferent toward the judgment, and that this
object-being is then, when it so happens, apprehended and designated by the
judgment's copula. Far from it! The relation is precisely the reverse, that
the phenomenon is subjectively judged to be objective: the phenomenon has its
objectivity from the judging subjectivity.

Otherwise, when the thing is spoken of, not as an object of cognition, but as
``the thing in itself.'' Here Kant now absolutely insists upon a being so
objective that it shall in no way be able to allow itself to be apprehended by
any subjective category of the understanding, and thus not by the being
corresponding to the judgment's copula either. By this assumption the author of
the critique of reason comes into infinite contradiction with himself. When
Kant says of the thing in itself that it \emph{is}, he passes a judgment. He
may now introduce as many distinctions as he will between reality and being,
between objective and subjective, incomprehensible and comprehensible being: it
still stands fast that
% --- p. 115 ---
\opage{115}the distinctions by which the thing's being is to be placed outside
the category of the understanding rest upon nothing but categories of the
understanding; further it stands fast that he cannot declare it an
impossibility to judge about what the thing in itself is except precisely
thereby, that he judges about what the thing in itself is. If he says: the
thing in itself is an X, then by this judgment he has just as surely let a
category of the understanding be determining for the thing in itself as, by
saying: the thing is hard, he has let it be determining for the thing as
phenomenon.

What has here been said of the categories which, according to Kant, lie at the
ground of the judgments according to their quality holds also of those which
are derived from the division of the judgments according to quantity, namely:
\emph{allness}, \emph{plurality}, \emph{unity}; from the division according to
relation: \emph{substance} and \emph{inherence}, \emph{cause}, \emph{effect}
and \emph{reciprocal action}; and finally according to modality:
\emph{possibility}, \emph{actuality} and \emph{necessity}.

If it is once proved that the category of the understanding \emph{being} must
be determining for the thing in itself, then the same holds of all the
remaining categories. If the thing in itself belongs to a sphere in which the
category \emph{being} has validity, then it belongs precisely to a sphere in
which the categories \emph{non-being}, plurality, \emph{unity}, \emph{cause},
\emph{effect}, etc., have validity. Conversely: if the categories
\emph{non-being}, \emph{plurality}, \emph{unity}, \emph{cause},
\emph{effect}, etc., have no validity when the talk is of the thing in itself,
then the category being has no validity either. One can then just as little say
of the thing in itself that it is as that it is not; in setting itself up
against the judgment, the thing in itself ceases to be an object.

The Kantian ``thing in itself'' can be regarded as a peculiar representative of
the antilogical, insofar as the antilogical stands against thought and for that
reason does not allow itself to be thought. But in Kant the apprehension of the
antilogical is nevertheless in a high degree mistaken. In the first place
``the thing in itself'' is an%
% --- p. 116 ---
\opage{116}tilogical in the sense that it does not allow itself to be
apprehended in any thought-determination whatever, and therefore maintains
itself as an X inaccessible to human knowledge. This is now, for the
determination of the antilogical, both too much and too little. Too much; for
the antilogical does indeed stand against thought, but in such a way that
precisely through the opposition it both determines thought and allows itself
to be determined by thought. Too little; for ``the thing in itself'' falls to
thought with its determining functions in two quite different ways. For here
it is presupposed that this which is unknown to us, unthinkable for us, can
possibly be thought by a divine thinking: the thing in itself is thus not in
itself absolutely unthinkable. Here it is further presupposed that the thing in
itself can in a certain manner coincide with our forms of cognition and
subjective categories of the understanding. For what is it, then, that fills
out the forms of cognition and gives experience its determinate content?
Evidently ``the thing in itself.'' But then the phenomenon must, as has been
shown in an earlier connection, precisely be an identity of the thing in itself
and the forms of cognition; or in the phenomenon the concept and the thing must
be identical. --- Next, and this is our second principal objection, Kant has
missed the antilogical in this, that, with all the fuss he otherwise makes
about the separation between thing and concept, he nevertheless in complete
indeterminateness lets the thing in experience and the concept in thinking have
the same form, only that the concept's form still lacks the content of
experience. Thereby he has overlooked the real problem. If the actual contains
in the concept nothing more than the merely possible; if there is here in the
concept's form nothing, nothing at all, that points to a difference between
possibility and actuality, then we ask: by what form and in what manner is the
determination of the determinate actual to be distinguished from the
determination of the many-sidedly determinate possible? Here there must
necessarily be a distinguishing mark in the form itself; for to add empirical
intuition to the pure concept is surely about
% --- p. 117 ---
\opage{117}as difficult as to subtract the sky from the blue sky.

The peculiar manner in which Kant in his critique of reason treats the relation
between the formal and the material conditions of experience is well suited to
incite metaphysical reflection in a one-sided direction. Everywhere it looks as
though there were in the object of experience, besides the form of cognition,
an indeterminable something that set barriers to cognition, a foreign element
that thought could not get the better of; and yet it is not anything --- the
resistance yields everywhere before the thought that works its way forward in
analysis. Thus attention was now in Fichte directed exclusively upon thought's
overreaching power: thought came to itself and saw in everything only itself.
Once reflection had dissolved the thing in itself and, despite actuality's
protest, forced back every representation of the antilogical character of
existence, one could no longer repel subjective idealism either by appealing to
that power of immediacy with which the outwardly actual comes forward; the
immediate power of the actual was then cognized in its essence not as an
antilogical factor, but as a logical power, a pure thought-power: the objective
returned, not as the judgment's object, but as the judgment's content, indeed,
in Hegel, as the judgment itself.

\textit{ββ)} \emph{Power in judgment and object: practical
functions.} If the Hegelian development of the concept's power were
many-sidedly satisfactory, the opposition of judgment and object would dissolve
into an illusion corresponding to the psychological standpoint. Instead of
antilogical power-determinations we should then get an existential immediacy;
this immediacy, although in the phenomenon it seems to be opposed to the
concept, would then essentially be only a moment in the concept itself, since
the concept's reality, its logically overreaching power, consists precisely in
positing and presupposing the elements through whose mediation it gives its
content the form of infinity and reveals the
% --- p. 118 ---
\opage{118}absolute idea. What is one-sided and untenable in this theory has
already been pointed out from several sides at earlier stages of the Logic of
the Fundamental Ideas.\footnote{Grundideernes Logik. I. p.~148--155 and
p.~162--168 etc.} Nonetheless the theory returns at every new stage with
renewed emphasis, and quite naturally, so surely as the Hegelian philosophy is,
in its kind, a consistent conclusion of all preceding systems.

If there is none among the thinkers of the past who can rightly be said to have
raised the problem of the antilogical; if Kant is among them all the one who
has come nearest to the matter without yet having grasped the matter; if
everything that earlier systems have set up under various designations as
reality's mark of distinction over against thought is mistaken, partly in such
a way that in its absolute unthinkability it must become self-contradictory and
meaningless, partly in such a way that as an ambiguous possibility it can
belong only to a subordinate sphere and thus be destined to disappear: what
wonder then that a logical carrying through of reflection ends by transforming
existence into a concept.

Only in the logic of the antilogical forms can panlogism find its essential
corrective. Earlier the knowledge-concept has been set over against the
power-concept under various lights; the opposition is now heightened by this,
that each of the concepts receives its corresponding development in its
corresponding word, the concept of knowledge in the word of knowledge, the
power-concept in the word of power. If the content of the power-concept could
be developed in the word of knowledge, then the opposition of existence and
concept would have to dissolve into a seeming; only in the light of
representation, thus in the light of untruth, would the existence-determinations
appear as antilogical object-determinations; seen in the light of truth they
would on the contrary show themselves to be pure conceptual determinations;
panlogism would then in principle be completely right.

If panlogism is to be overcome, the power-concept must be developed in the word
of power. But where do we now come to? As
% --- p. 119 ---
\opage{119}power is, so must the word be; power is antilogical, the word
likewise. As the word is, so must language be; the word is antilogical,
language likewise. As language is, so must its propositions, judgments and
concepts be; language is antilogical, propositions, judgments and concepts
likewise. Antilogical concepts, antilogical propositions, antilogical
judgments! and all this is to be treated in a logic, the antilogical thus
treated logically: is there not here a web of insoluble contradictions? ---
Before we enter upon the reality of the matter, we will consider the
difficulties from the formal side.

Formal logic makes, as is well known, a finite separation between thought and
what is thought, between the concept and what is comprehended, a separation
that finds its confirmation in turns of speech such as: to think \emph{of}
something, to think \emph{over} something, to form for oneself a concept
\emph{about} something and the like. In displeasure at the relation between
thought and thing thus made finite, objective logic has in turn from its side
held fast the identity of thinking and being as so immanent that what is, by
being thought, became thought, and the existent, by being comprehended, became
concept. Objective logic does not think of something; it thinks something, and
this something is thought itself (ὥστε ταὐτόν νοῦς καὶ νοητόν. Aristoteles.).
Objective logic does not form for itself a concept about something, --- that
would be far beneath its dignity, --- but it comprehends something, and the
something it comprehends is the concept itself.

Both ways of regarding the matter are, when they are held fast each by itself,
equally one-sided and by their one-sidedness equally harmful to true
development of the concept. Formal logic's separation of thought and its object
is one-sided insofar as it overlooks the identity; objective logic's identity
of thought and what is thought is one-sided insofar as the opposition is
overlooked. In order to defend the paradox that thought's object can only be
thought itself, objective logic appeals to immanent negation. By its immanent
negation thought is able
% --- p. 120 ---
\opage{120}to raise itself above itself and make thought into object; by its
immanent negation thought is conversely able to transform the thought object
into thought. The object is, in other words, thought's other; but this other is
through the immanent negation thought itself: here there is thought of thought
of thought in a thinking of a thinking of a thinking, and further one
fundamentally does not get.

In order to get beyond this boundless formalism, we must necessarily
distinguish between mere formal negation and real negation. So long as
thought's other is only thought itself, and all form-determinations to that
extent logical, the negation is formal; only when thought's other can no longer
be transformed into thought, thus when the antilogical enters, is the negation
real. If one will now, in virtue of real negation, declare that other of thought
which is not itself thought to be something in itself unthinkable, then one
will also easily be able to defend the paradox that thought in actuality thinks
the unthinkable. The matter is that this so-called unthinkable is in reality
thinkable: thinkable in that it is apprehended by the thinker, unthinkable in
that it is not thought and cannot think itself. Upon real negation, upon a
thinking of that which is not thought, thus upon a logical determination of the
antilogical, rests the reality of all thinking.

That the reality of thinking must coincide with the reality of the matter goes
without saying. In the reality of the matter there must here, then --- this
lies in real negation --- be an actual opposition of the concept in the word of
knowledge and existence in the word of power; the concept in the word of
knowledge is the judgment, existence in the word of power is the object: the
actual opposition of concept and existence is an actual opposition of judgment
and object. If the actual opposition of judgment and object is now not to sink
down into the parceling out of finitude and be fixed in psychological dualism,
then to the actual opposition of the word of knowledge and the word of power
there must necessarily correspond an actual unity;
% --- p. 121 ---
\opage{121}such an actual unity cannot be merely theoretical; the actual
unity is here only insofar as it is actually brought about; it is brought
about only by an actual transition, a transformation: the bringing-about
is essentially practical. In order that logic may now be able to give a
mirroring, determined by the word of knowledge, of the praxis answering
to the word of power, it is necessary to begin with the relation between
thought and power.

Theory, pure theory, must content itself with thought-movements; its
functions are only thought-acts. By pure thought-acts it is, according to
the real negation, impossible to unite knowledge and power. The actual
unity which in the infinite has been brought about and is to be brought
about presupposes an essential unity which neither can nor is to be
brought about, because it eternally is; it is in subjectivity itself, the
absolute principle reflected into itself. According to the finite
opposition between knowledge and power the thought-act is powerless and
the power-act in its immediacy thoughtless; but in virtue of the original
unity of power and knowledge the thought-act is a transformed power-act,
and thereby itself powerful; the power-act, conversely, a transformed
thought-act, and itself thought-fast. Inasmuch as the power-act is
converted into thought-act, and the word of knowledge assimilates the
power-determinations, the movement goes from the antilogical to the
logical, the object-forms become concept-forms, and the judgment logical.
Inasmuch as the thought-act conversely passes over into the power-act,
and the word of power absorbs the ideality of the conceptual
determinations, the movement occurs from the logical to the antilogical,
and the judgment itself becomes antilogical. The practical functions are
simultaneous in both movements, insofar as the movements form an infinite
circuit; the reciprocal determinations of logical and antilogical moments
answering to this circuit are explained in a comprehending reflection,
which makes the judgment at once logical and yet antilogical.

1) \emph{In the logical judgment the object is preformed: antici%
% --- p. 122 ---
\opage{122}pating functions.} If the judgment is really to be able to
determine the object, knowledge must become one with power; this unity is
brought about by a system of anticipating functions. In teleology there
is talk often enough of anticipations and preformations, without its
becoming clear just on that account how the preforming activity actually
proceeds, and how far anticipation is originally a theoretical or a
practical function. What relation, then, obtains here between
preformation and anticipation? what must anticipation necessarily
presuppose? how does it take place?

Since preformation, as the word runs, aims at giving the object an inner
existence preceding the outer, and therewith at once an inner shape
preceding the outer, it follows of itself that all preformation must rest
upon anticipation; for what is anticipation but a first forming of a form
that presupposes itself, a becoming form without actual existence. In
preformation it is impossible to point out an immediately first form: for
preformation says in its concept that every existing form must result
from a preceding form. Could there in actuality be given a first existing
form, then the one existing form could go before the other, prepare the
other and be relieved by the other without any proper preforming. In
preformation, by contrast, the preparing form is so essentially
determined by the resulting one that it is at once both its copy and its
model, and can be the latter only by being the former. The actual shape
is always, if the word may be used, a this-shape, just as the existent is
always a this-being. If now it is not a shape in indeterminate
generality, but this, just this shape, that is to be preformed, then it
must necessarily be presupposed in order to be posited: it must be
anticipated. That a shape which is not yet, and consequently has not the
% --- p. 123 ---
\opage{123}determinateness of the existent, should act determiningly upon
the initial form by which it is determined --- this presupposing, on the
part of preformation, of that which is to be posited, is precisely
anticipation.

From the infinity reflected in anticipation it shines forth that its
being is itself a reflecting, and that this reflecting must necessarily
presuppose a reflecting subjectivity. All anticipating functions are
subjectivity-functions. But just as selfhood with necessity presupposes
otherness, so must subjectivity's anticipating activity necessarily
presuppose objectivity. That anticipation is determined by consideration
of a coming objectivity is evident; think away what follows, and what
goes before ceases to be an antecedent. But not only in a formal, also in
a real sense must anticipation presuppose a given existence. What is
posited subjectively in anticipation is posited precisely with the
express determination that it shall one day become objective. But nothing
can become objective save that which in its subjective beginning already
has conditions for objectivity. Conditions for objectivity can only that
have which is originally determined by objective considerations; but
objective considerations, consideration of the objective, presuppose a
world of already existing objects. Divided between the subjective and the
objective, or rather in its indivisible unity determined in part by both,
anticipation thus falls under the relative concepts: an absolute
anticipation is just as impossible as an absolute beginning.

From the relativity of the anticipation-concept it follows that actuality
in infinite extent cannot possibly be anticipated at once. Every relative
totality, every actual member, must indeed proceed from its possibility;
but for that reason the whole of actuality cannot simply have its
beginning in a possibility going before all actual things: possibility
has actuality not only before it, but also behind it; the actu%
% --- p. 124 ---
\opage{124}al possibility rests upon actuality. So too with anticipation:
it prepares an existent by means of an already existing one. From the
relation between the actual and the possible, the objective and the
subjective in anticipation, the relation between the theoretical and the
practical is seen as well.

Insofar as anticipation proceeds from the objective, the anticipating
functions are essentially theoretical. The theoretical is that the
judgment is determined by the object, thought by the given, that
actuality, the objective, rules. With its gaze upon the objectively
actual the anticipating subjectivity recognizes its actual limit, its
barrier. Over against the actual objective the abstractly subjective is
powerless. Or what should pure subjectivity be in a position to alter in
objective actuality? The actual objects would have to be no objects, but
mere figments, pure dream-images, if they could without more ado be
dissolved and transformed by a purely subjective determining. What holds
of the objective members holds of course also of the relations in which
the members essentially bring their objectivity to light. Every
alteration of the objective must occur upon objective conditions, in
consequence of objective presuppositions. The objective must indeed be
objectified; but it itself determines the content of its objectification,
and thereby proves its reality.

But if subjectivity is really in every respect powerless over against the
objective, then anticipation is an impossibility: a unity of knowledge
and power is required, a unity which anticipation itself both
presupposes and brings about. The anticipating activity is preforming;
preformation presupposes precisely in subjectivity an overreaching power
over the objective. Without such an overreaching power the subjective
forming would not be determinative for the objective form. As the
subjective activity which prepares and shapes a coming actuality,
% --- p. 125 ---
% Errata: p. 125 --- printed "en", corrected to "er"; translated as corrected.
\opage{125}anticipation must thus be an ideal praxis; but from which side,
then, can the objective fall liable to determination by the subjective?
From that of possibility, of real possibility! For real possibility is at
once both objective and subjective.

Insofar as possibility is objective, the objects can in countless ways be
reshaped and altered by one another. As existing object, as outer
actuality, the object has in itself the possibility, on different
conditions, of playing a highly different role. In its immediate
actuality the object is of determinate peculiarity, sets itself against
subjectivity and asserts its independence; but according to possibility
the objective is powerless, according to its possibility the object is
stuff-like and indeterminate; to know possibility is to know the side
from which objectivity lets itself be attacked subjectively: by its
knowledge of objective possibilities the anticipating subjectivity
acquires an ideal power over the objects.

How stuff-like objectivity is, is seen from the objective possibilities.
Never are so many individual shapes formed that there could not, by
reason of the possibilities, be formed more; actuality is never shaped in
types so particular that it could not, by reason of the possibilities, be
shaped in other types: here there everywhere remains in the objective
formations an infinite surplus of possibilities. It is this surplus of
possibilities that always stands at subjectivity's disposal; subjectivity
judges, and the judging subjectivity knows how to handle possibility, for
subjectivity \emph{liberates} possibility. In its objective form
possibility is bound, and the bound possibility is abandoned to
contingency; only in the subjective form does possibility become free:
the liberated possibility has its expression in the judgment.

The actual can as actuality be expressed both in the object and in the
judgment about the object; but the possibility latent in the object can
be expressed in the form of possi%
% --- p. 126 ---
\opage{126}bility only insofar as it is expressed in the judgment. To be
sure, objective possibility always rests in a corresponding actuality;
but in actuality possibility is not; were possibility actually there, it
would be precisely an actuality, no possibility. In outer actuality, in
the object itself, possibility must content itself with a being without
determinate being; in the judging subjectivity, by contrast, possibility
acquires ideal existence, its ideal existence is the judgment itself, the
judgment in which it is uttered. Here now is seen anticipation's original
unity with the judgment.

In the knowing subjectivity many images may form, also images of things
which are not yet; but no image can in its immediacy as image pass for an
anticipation. Anticipation expressly demands a knowledge that the given
image designates an object not yet existing, an object whose existence is
precisely introduced and prepared by the image; but such a knowledge is
inseparable from a judgment; for only in a judgment can the relation
between the image and the object be expressly determined. That the image
of anticipation does not remain standing in pictorial immediacy, but
resolves itself into a representation and passes over into the judgment,
is seen from this, that the demand indicated by the image, to be
actualized as object, is in the judgment expressly posited as a task. It
is not the image, but the task, the task determined by the judgment about
the image, about which the anticipating functions essentially move.
Transformed by the judgment into a task, the anticipation-image is more
closely determined by the category's pregnant unity-expression, thus by
the word, the practical word of knowledge, which carries the whole task.

With the task in the practical word of knowledge the anticipating
functions now work at a closer determination of the possibility of the
object in question. In the determination of this possibility the
subjective is articulated point by point by objective considerations.
The articulation is expressed in the judgment,
% --- p. 127 ---
\opage{127}inasmuch as the judgment in turn expresses a peculiar
reciprocal relation of freedom and necessity, arbitrariness and terms. If
the task, categorically determined, is to construct, e.g., a clock, then
the choice by which it became just this and not another task is doubtless
arbitrary; but the terms announce themselves at once when one looks to
the solution of the task. The fundamental demand is to bring about a
series of isochronous events and to indicate these in numbers. From the
task thus posed it is seen that its solution rests upon the terms of
actuality, upon the conditions objectivity offers. Nor, then, does the
image of the clock stand fast; it is as yet only a possibility, and
possibility varies in the alternation. ``A clock is (can be) both a
sand-clock and a water-clock and a mechanical clock, etc.'' The
conjunctive judgment: \textit{A} is (can be) both \textit{B} and
\textit{C} and \textit{D}, sets several possibilities before a choice; the
disjunctive judgment: \textit{A} is (must be) either \textit{B} or
\textit{C} or \textit{D}, expressly demands a choice among several
possibilities.

Inasmuch as the anticipating subjectivity in the judgment surveys the
possibilities offered by objectivity and bethinks itself of a choice, it
must likewise see in the judgment how it may best adapt its choice to the
conditions of actuality. The judgment now assumes the hypothetical form:
if the clock is to be a sand-clock, then \dots\ if \textit{A} is
\textit{B}, what then? Then, answers formal logic, \textit{A} is neither
\textit{C} nor \textit{D}. This answer is in the present connection idle;
the idleness of theoretical formalism is most palpable when the question
concerns the solution of a determinate task; the determinate task is
essentially practical. Thus: if \textit{A} is (is to be) \textit{B}, what
then? then a system of conditions of its own holds, the system
\textit{b}; if \textit{A} is (is to be) \textit{C}, then by contrast
another system of conditions holds, the system \textit{c}, and so forth.

To the judgment, the decisive judgment which goes through a system of
judgments, there answers the anticipated possibility, which
% --- p. 128 ---
\opage{128}with its determinateness consists in a system of
possibilities. In order to determine possibility the anticipating
subjectivity works with nothing but imaginary elements, but in every one
of these imaginary elements a reflex of the real is seen, of the reality
in the objective conditions. The objective conditions in their reality
are nothing but power-determinations; the subjectivity immersed in
anticipation's deliberations has thus, through these deliberations,
appropriated to itself a system of reflected power-determinations; with
this system of reflected power-determinations possibility is articulated
into an inner, imaginary actuality; with the complete development of the
inner, imaginary actuality preformation, and with it anticipation, is
carried through.

2) \emph{With the antilogical judgment the object is realized: executive
functions.} In the task determined by anticipation actuality is posited
in advance as concept; by the solution of the task the concept is to be
posited as actuality. How a concept is converted into actuality is seen
wherever a work of subjectivity is objectively accomplished, wherever a
purpose is realized. What is essential in the actualization of the
concept is however misunderstood from two equally one-sided standpoints,
by everyday consciousness and by the philosophy of the immanent method.

For everyday consciousness it is human manufactures, not natural
products, that must serve as examples of the actualization of concepts in
material objects, houses, household utensils, machines and the like. What
is orienting in these examples is that they point expressly to an
anticipating subjectivity, an actual surrounding world in which the means
answering to the purpose must be sought, and an active individuality
which by voluntary use of its limbs procures itself tools and by
voluntary use of the tool works its material in a manner answering to the
task, etc.

% --- p. 129 ---
\opage{129}By such exemplifications the task now becomes parcelled out to
such a degree that one can, so to speak, anatomically separate the
individual moments. First the master-workman has the whole concept of the
work in question anticipated in the purpose; next, by means of thought's
unity with the will, the execution of the plan lying in the purpose
becomes actual resolve; in the will's immediate action upon the organs,
in the voluntary use of arm and hand, there is finally seen the identity
of the ideal with the real, the subjective with the objective, the
concept with actuality. Inasmuch as the will is determined by the concept
by means of thought, the motion in hand and fingers is in turn determined
by the will and, by means of the will, by the concept; every
finger-motion now answers to a thought-movement; the concept's ideal
power has transformed itself into real power: the transition from
concept-form to object-form, from the logical to the antilogical, has
begun.

But however illuminating such exemplifications may be in a propaedeutic
sense, they are nevertheless insufficient for the logic of the ideas. So
long as attention dwells one-sidedly upon the individual master-workman
and his work, there will always be a multitude of collateral
representations here which can work disturbingly. It constantly looks as
though the plan and subjectivity stood in a merely contingent relation;
the plan must indeed be in the master-workman's consciousness, but the
master-workman's consciousness may very well exist without such a plan.
Just as arbitrary and just for that reason contingent is the relation
between anticipation and actualization, between plan and execution; the
master-workman cannot indeed proceed to execution without a plan, but he
may go about with the plan without proceeding to execution. Finally, and
this is in the main the most misleading thing, by such examples the point
of coincidence of the subjective with the objective in the will's use of
the limbs is grasped in so restricted and particular a manner that the
fundamental relation between concept and actuality, between pur%
% --- p. 130 ---
\opage{130}pose, means and the attainment of the purpose, the fundamental
relation which runs through the whole of existence, is not brought to
light: here there is and remains a wide gulf between \emph{manufactures}
and \emph{natural products}.

% Errata: p. 130 --- printed "des", corrected to "das"; translated as corrected.

In sharp opposition to the popular exemplifications, Hegel, in his
doctrine of the concept's actualization by itself, notably in his
doctrine of the purpose, has to such a degree removed all considerations
of finitude that the opposition between the subjective and the objective
purpose, between the concept in itself and the actual world of objects,
falls under the concept's general self-movement. The mystification which
has already been pointed out in ``\emph{The Logic of
Subjectivity}''\footnote{Grundideernes Logik I. S. 226.} can in this
connection be illuminated from a new side. --- ``Die Bewegung des
Zwecks'' --- says Hegel\footnote{Wissenschaft der Logik. Zweiter Theil.
S. 220.} --- ``kann \dots\ so ausgedrückt werden, daß sie darauf gehe,
seine \emph{Voraussetzung} aufzuheben, das ist die Unmittelbarkeit des
Objekts, und es zu \emph{setzen} als durch den Begriff bestimmt. Dieses
negative Verhalten gegen das Objekt ist ebenso sehr ein negatives gegen
sich selbst, ein Aufheben der Subjektivität des Zwecks. Positiv ist es
die Realisation des Zwecks, nämlich die Vereinigung des objektiven Seyns
mit demselben, so daß dasselbe, welches als Moment des Zwecks unmittelbar
die mit ihm identische Bestimmtheit ist, als \emph{äußerliche} sey, und
umgekehrt das Objective als \emph{Voraussetzung} vielmehr als durch
Begriff bestimmt, \emph{gesetzt} werde.''

According to the Hegelian conception all difference then vanishes between
the anticipating subjectivity and the plan, the subjective purpose
determined by anticipation. Insofar as the purpose presupposes
anticipation, and anticipation subjectivity, the purpose-concept in Hegel
is this subjectivity. The concept has, as ``the determinate concept,''
all the requisites it
% --- p. 131 ---
\opage{131}essentially needs. Here there is to that extent a smuggled-in
anticipation, but no anticipation-problem. ``The movement of the purpose
aims at sublating its presupposition, namely the immediacy of the object,
so that the object may expressly become determined by the concept;'' what
does this now mean? The object is as object already concept; inasmuch as
the purpose as concept has the immediacy of the object for presupposition,
this means, in other words, that the subjective concept reflected into
itself has the objective, immediate concept for presupposition. If the
emphasis is now laid upon this, that the presupposition of the reflected
concept is itself concept, then the whole movement by which the immediacy
resulting from the concept is again sublated in its reflected concept is
only a repetition of the worn-out story of reflection's conversion into
immediacy, and immediacy's into reflection; here, then, there can be no
question of the purpose realizing itself, for here there is no reality at
all, and of course no realization either. If, by contrast, the emphasis
is to fall upon the object's objective immediacy, thus upon reality, then
it is asked whence the purpose as concept acquires the independence to
step forth in free ideality against a given reality, to intervene in the
connection of things and to master the outward objective. The concept as
purpose is a subjective thought which abstraction has torn loose from the
thinking subjectivity and set up as an essence for itself. The essence
thus set up, ``the subjective purpose,'' is a phantom, a powerless
shadow, which will surely refrain from realizing itself.

The Hegelian concept, for all its divinity, evidently lacks the essential
opposition of ideal existence in the judgment and real existence in the
object. The concept's self-realization is a feigned movement, due to
abstraction and to abstraction's stepwise sublation. That ``the
subjective purpose turns itself negatively against the object,'' that it
therewith ``turns itself negatively against itself,'' that it, in other
words, ``sublates its subjectivity and posits itself negatively as
realized
% --- p. 132 ---
\opage{132}purpose,'' means only that the thinking spectator is, by help
of an explicative dialectic, to fix his gaze now upon the one of these
moments, now upon the other, now upon both together.

If Hegel with his whole system is not in a position to state the
difference between an object's concept and the object itself, between the
concept in knowledge and the concept in power, how should he then be in a
position to point out a transition from concept to object? Formally there
are transitions enough here; but here there is no concept of a real
transition. A real transition namely presupposes a real distance; where
there is no real distance, neither is there, in the logical sense, any
real transition. The transition-illusion, an illusion not unlike that by
which one who sails sees the coast move, is maintained chiefly by this,
that the concept in the \emph{general} now steps into opposition to, now
again melts together with, the concept in the \emph{particular}; but
concept is always concept after all: the dialecticizing reflection posits
and sublates the separation with equal ease. Once a separation between
the objective and the subjective has here been posited in consequence of
the concept in the particular, as is now the case with the objects and
``the subjective purpose,'' such a separation is again sublated by the
concept in all generality. Of the concept in all generality it may
doubtless hold that it turns itself negatively against itself in turning
itself negatively against the object; for the concept of the objective
belongs undeniably just as essentially under the concept in all
generality as does the concept of the subjective; but how should such a
thing be able to solve a realization-problem, how should it be able to
explain to us the transition from ``the subjective'' to ``the realized
purpose''?

What we need to have explained in this connection is not a tautology such
as this, that the subjective purpose is in its concept subjective and
thereby opposed to the object and to the objective; no, what ought to be
explained to us is the problem
% --- p. 133 ---
\opage{133}how it actually comes about that the concept, not the
subjective-objective concept in logical ambiguity, but the determinate
concept, the concept of this particular, determinate purpose, can effect
a transition from ideal existence to actual existence, where by the
transition there is then thought no logical paraphrase of the subjective
by the objective, but an event, a production, in which the concept is
seen to determine actuality and itself to become actual. For the
answering of questions belonging here the Hegelian logic contains no
instruction, not even a hint.

In Hegel the opposition between concept and object, when it is driven to
its utmost, cannot reach further than to an opposition of essence and
phenomenon, only with the express addition that the totality is
completely developed from both sides, indeed has in the course of the
development even become two totalities. Since now the totality-unity of
the two totalities is posited as the judgment, the logically objective
judgment (the fundamental division, ``das Urtheil''), one should expect
that Hegel would hold fast in the judgment to the opposition between the
concept in the form of essence and the concept in the form of phenomenon,
and would thus arrive at a cognition of the object; but no! in the
judgment reflection is directed upon the filling-out of the copula, thus
upon the totalization of the identity in which the separate totalities
are sublated. ``So ist die Form des Urtheils untergegangen, erstens, weil
Subjekt und Prädikat \emph{an sich} derselbe Inhalt sind; aber zweitens,
weil das Subjekt durch seine Bestimmtheit über sich hinausweist, und sich
auf das Prädikat bezieht, aber ebenso drittens ist \emph{dieß Beziehen}
in das Prädikat übergegangen, macht nur dessen Inhalt aus, und ist so die
\emph{gesetzte} Beziehung oder das Urtheil selbst. --- So ist die
konkrete Identität des Begriffs \dots\ im \emph{Ganzen} hergestellt''
\dots\footnote{Wissenschaft der Logik. Zweiter Theil. S. 117.}

With the filled-out copula the judgment now passes over into the
% --- p. 134 ---
\opage{134}syllogism.\footnote{Cf.\ Grundideernes Logik. I. S. 214.} If
the judgment ended in its immanent development with the two totalities,
the subject and the predicate, going together in the totalized unity of
the concept, then one would now have to expect that the content-rich
totality-unity which, in the course of the syllogism's self-development,
is divided and united in three totalities should lead to insight into the
opposition between the concept as concept and the concept as object; but
no! the syllogism ends thus: ``Die ganze Formbestimmung des Begriffs ist
in ihrem bestimmten Unterschied und zugleich in der einfachen Identität
des Begriffes gesetzt.'' And now is added: ``Dadurch hat sich nun
\emph{der Formalismus des Schließens}, hiermit die Subjektivität des
Schlusses und des Begriffes überhaupt aufgehoben''.\footnote{Wissenschaft
der Logik. Zweiter Theil. S. 169 \& 71.}

Should anyone hold that here at last a thoroughgoing totality-opposition
must show itself, since it is expressly said that the whole
form-determination of the concept is posited, not merely in simple
identity, but also in determinate difference, then let him note how the
chapter on the syllogism ends: ``Das Resultat ist \dots\ eine
\emph{Unmittelbarkeit}, die durch \emph{Aufheben der Vermittelung}
hervorgegangen, ein Seyn, das ebenso sehr identisch mit der Vermittelung
und der Begriff ist, der aus und in seinem Andersseyn sich selbst
hergestellt hat. Dieß Seyn ist daher eine \emph{Sache, die an und für
sich ist} --- die \emph{Objektivität}''. And what, then, is gained by
this? A concept which is itself object, because it is Object, but an
Object which is nevertheless not object, because it is concept; thus a
concept which is neither concept nor object; an object which is neither
object nor concept; a dialectical story which can go on as long as one
pleases.

If the object, seen from the side of logic, is to assert its independence
against the concept; if the concept, seen from the side of existence, is
to assert its independence against the object:
% --- p. 135 ---
\opage{135}then they must both have each its own existence, the concept an
ideal existence in a comprehending subjectivity, the object a real
existence in actual objectivity. The concept's existence in subjectivity
is determined by a word of knowledge; the object's objective existence is
determined by a word of power. The concept's realization, the actual
transition by which the determinate concept comes to existence as
determinate object, occurs through an intellectual-real praxis whose
peculiar character must rest upon the executive functions.

The executive functions are themselves determined by a difference-unity
of antilogical judgments and antilogical propositions. The antilogical
judgment, which has the concept in its subject and existence in its
predicate\footnote{Grundideernes Logik. I. S. 160--62.}, first acquires
its full significance in thought's transition into action, theory's
transition into praxis.

From the standpoint of theory the antilogical judgment will look only
like an inverted presentation of the logical. For when someone, instead
of saying: this is an axle, this is a wheel, this a roller, says: an axle
is this, a wheel is this, etc., it certainly sounds inverted enough. It
is otherwise when the task is to construct a wheel, a roller and the
like. The construction presupposes the concept: the concept is the
original, the object the derived. The concept is had in the subject of
the judgment; the one who judges knows what a wheel is in its concept, he
does not then wish to have the concept determined over again as concept
in the predicate, so that the predicate should give him a logically exact
definition. Such theoretical formalities are here superfluous. The
constructing master-workman will indeed define his concept, but the
definition is to be practical, not theoretical: the practical definition
is one with a corresponding production, a construction.

It is said in logic that it is not the singular but the particular, not
the object itself but its concept, that
% --- p. 136 ---
\opage{136}admits of definition (\textit{species sola definitur, non
individuum}); this is, in the theoretical sense, quite correct; the
logical judgment cannot define the object. But in practice the matter
stands the other way round; here the endeavor aims precisely at
defining the concept's reality by producing the object corresponding to
the concept: \emph{the produced object} is \emph{the concept's practical
definition}. To such questions as: what is a triangle? a square? a
circle? one can answer practically by constructing a triangle, a square,
a circle. How far such practical definitions will satisfy theory shall
not be investigated here; only this shall be emphasized, that the
practical definition is guided and determined by an antilogical judgment.

He who answers the question: what is a circle? by constructing a circle,
translates the concept of the circle straightway into existence. ``A
circle is a figure like \emph{this}, is \emph{this figure}''; with such a
judgment the construction begins, but such a judgment is antilogical. The
logical judgment coincides immediately with its corresponding
proposition; in the word of knowledge the proposition is judgment insofar
as it is reflected in thought, and the judgment proposition insofar as it
is reflected in language: both reflections are theoretically the same
reflection. It is otherwise when the question concerns the relation
between judgment and proposition in the antilogical; here the identity is
practical, and therefore there must here also be a practical difference.

The practical difference is a difference between the thought that guides
the action and the action itself. The constructing, and hence acting,
subjectivity must \emph{think}; but the thought is practical, its
functions executive. The executive functions all aim, from the standpoint
of the concept, at showing the possibility of the practical definitions.
Here the antilogical judgment divides itself into a multiplicity of
particular judgments. The constructing, and hence thinking, subjectivity
must \emph{act}.
% --- p. 137 ---
\opage{137}The actual acting, determined by thought, is not itself a
thinking, but on the contrary is so opposed to thought that precisely by
the opposition it can correspond to thought. The actual action
corresponding to thought cannot possibly, since it is not itself thought,
coincide with the judgment as judgment; it must necessarily have its
expression in an antilogical proposition determined by the judgment. As
the antilogical judgment branches out into the multiplicity of particular
judgments, so the antilogical proposition branches out, in the execution
of the work, into a multiplicity of particular propositions; in the
articulated reciprocal determination of antilogical judgments with
practical-antilogical propositions is seen the full significance of the
logic of the executive functions.

That the object produced in and by the judgment, by the antilogical
judgment in its system of corresponding propositions, is now also
actually opposed to the judgment, can be seen from the executive
functions. In every executive function the judgment decisive for the
proposition is an expression of the identity of power and knowledge in
the form of knowledge: the power in the judgment is ideal. But in every
executive function the proposition determined by the judgment, the action
governed by thought, is likewise an expression of the identity of
knowledge and power in the form of power; the power in the action is
real: as many actions, so many power-determinations. Since now the unity
of the proposition through the judgment gives, point for point, a unity
of the power-determination with its corresponding conceptual
determination, the totality of the power-determinations in the produced
object must in the execution itself necessarily correspond to the
totality of the knowledge-determinations in the producing subjectivity.
The object in its system of antilogical propositions stands as a word of
power over against the system of judgments which, in the practical word
of knowledge, has given the concept existence.

If the unity-in-opposition of thinking and acting, of judgment and
proposition, is a unity-in-opposition of knowledge and power, then it is
seen from this that the producing subjectivity must be
% --- p. 138 ---
\opage{138}at one and the same time both a capable and a knowing
subjectivity. If the power is restricted but the knowledge extended, this
will be seen in the production; if conversely the power is extended but
the knowledge restricted, this will likewise be seen in the production.
Here is a turning-point at which the relation between \emph{manufactures}
and \emph{natural products} can be more closely illuminated.

All real productions, even the apparently meaningless ones, are
determined by concepts and realize concepts. To assume that there could
in actuality be productions so confused that they fell outside every
thought, every conceptual determination, is just as unreasonable as to
assume that in sand, water or air there could be eddies that fell outside
all laws. If it is now true that every production rests upon a
unity-in-opposition of concept and actuality, and if actuality is the
objective and the concept the subjective side of the production, then it
is seen from this that production is a two-sided category, and that it
makes a great difference whether this category is regarded from a
subjective or from an objective standpoint.

``The manufactures, the human productions, are'', it is said,
``productions that take place with consciousness, but the natural
products are unconscious productions''. Such a separation of manufactures
and natural products is misleading, insofar as it gives the appearance
that there were here two productions different in principle, subjective
and objective, whereas the truth is that every production is in the last
instance both subjective and objective.

That the human, conscious productions are subjective is evident; but with
all the arbitrariness that distinguishes them they are nonetheless
objectively conditioned. Whoever will abstract from the master builder's
subjectivity and during the work fix his attention upon the hand, the
tool and the material, will convince himself that what is subjective in
the production is determined, link for link, objectively. That
% --- p. 139 ---
\opage{139}natural productions are objective is evident; but whoever will
attend to what observers say about ``natural forms'', ``natural
concepts'', ``natural thoughts'', ``natural purposes'' etc., will soon
see that no thinker can consider the objective side of natural
productions without looking about him for the producing subjectivity.

``But'', it is said, ``natural productions surely take place with
objective necessity; how then can a production that occurs so objectively
rest upon subjectivity?'' We answer: Not upon a finite psychological
subjectivity which, with a limited knowledge and a still more limited
power, stands in finite opposition to the world of objects; but indeed
upon a subjectivity which, with an original unity of knowledge and power,
stands in an infinitely articulated fundamental relation to the
objectivity in everything objective. If natural productions were to lose
their reality because they are produced ideally and in a subjective
manner in the concept, then natural objects too would have to lose their
reality, because they can be objects only insofar as they are objectified
by an infinite subjectivity. Since it is now so far from being the case
that objects become less objective by being objectified, that on the
contrary they must have their objectivity in the objectification itself:
then it will also be possible to understand that natural productions do
not lose their objective validity by being produced by a subjectivity
that is indwelling in them and yet hovers above them; it is precisely the
subjective in the production that grounds the
objective.\footnote{Grundideernes Logik. I. p.~444--456.}

% Errata: p. 139 --- printed "I den", corrected to "3) I den"; translated as corrected.
3) \emph{In the logical-antilogical Judgment the Object is explained:
explicative Functions.} It is a peculiar matter with observations of
nature; everything that the ingenious observer puts into nature he gets
out again, neither more nor less. The well-known dispute between the
Hallerian and the Goethean view concerning ``the inner of nature''
science has accordingly
% --- p. 140 ---
\opage{140}not been able to settle hitherto either. Indeed, the natural
products in their immediate objectivity have no liberated inwardness, as
if there could here be an outer that, like a shell, surrounded and
concealed the inner; but the immediacy with which all activity of nature
takes place nonetheless gives the production so objective a stamp that
the subjective is at one and the same time both missed and presupposed.
When the one mass attracts the other, the one stone presses the other,
the one substance combines with the other: then there is here, point for
point, a determining of other by other, of objective by objective upon
objective possibilities, and the same holds when the seed germinates and
the organism thrives; --- what then is the inner of nature? In its
immediate corporeality the organism is an object just as much as the
stone; that there is here in the object an opposition of inner and outer
must be conceded; but the object's inner is, as we have already seen
above, just as much object as the outer. That in nature's outer an inner
should everywhere come to meet us is a figure of speech that can be just
as ambiguous as when it is asserted, in clumsy immediacy, that ``in the
object a concept comes to meet us''.

It is with the inner of nature as with the natural thoughts; when we look
at the outer, then there is an inner; but if we would see the inner
itself, then it has vanished. The inner has essentially a reflected, not
an immediate being; but nature does not reflect, and in its immediate
naturalness cannot have an essential inner. And yet nature has an inner;
when Goethe says: ``die Natur hat weder Kern noch Schale, Alles ist sie
mit einem Male'', it sounds striking, but is nonetheless veiled; the
poet's clear words stand in great need of an explanation.

That nature is without an inner, without free inwardness, means that it
is without subjectivity; that nature has its inwardness, its essential
inner, means that it nevertheless has subjectivity; that the inner of
nature is obscure, its
% --- p. 141 ---
\opage{141}inwardness a riddle, means that the subjectivity which nature
has is a subjectivity that has nature, and hence a subjectivity doubled
within itself; a subjectivity that is of course in nature itself, and a
subjectivity hovering above nature, a subjectivity in which the producing
actions are reflected as determining thoughts.

To the doubled subjectivity of natural productions there corresponds the
difference-unity of antilogical judgments and practical-antilogical
propositions. Without a system of antilogical judgments the concepts
decisive for the productions could not possibly be. The concepts are only
in a comprehending, the comprehending only in a comprehending
subjectivity. The antilogical judgments are determined in the
anticipation by the logical ones; the reflection of the logical and the
antilogical judgments is in the knowing subjectivity that hovers above
the productions and reflects itself in everything and everything in
itself, in the divine selfhood whose knowledge is power and whose power
is knowledge. But without a system of antilogical propositions, of
actual power-actions expressly determined by the anticipated concepts of
practical knowledge, natural productions could not possibly have the real
character of otherness, of actuality, of real existence. Where, then, do
we find the original unity of the antilogical propositions with the
corresponding judgments? In the essence-unity of the hovering
subjectivity with the powerful subjectivity indwelling in the
productions; determined still more closely, in this subjectivity's
reflected unity with the immediately producing cause, the principle of
nature that is present in objective working.

The actual relation between the principle of nature and the infinite
subjectivity must in logic be determined by \emph{explicative functions}.

Every natural object is a produced thing. To cognize the produced is to
cognize the production; to cognize the production is to repeat it in
cognition;\footnote{Grundideernes Logik I. p.~169--70 ff.} to
% --- p. 142 ---
\opage{142}repeat it in cognition is to reflect through the system of
antilogical judgments that have guided the producing actions, and to
imitate these actions in particular imaginations answering to the demands
of the concept. To the system of the actions, of the antilogical
propositions, there corresponds the principle of nature. The principle
--- this holds both of the principle in general and of the particular
principles --- works but does not think, does the deeds of the concept
but does not itself comprehend it; the principle is, as principle of
nature, selfhood in the objective form, a disposition of subjectivity, a
selfhood working selflessly, the productive unity of the antilogical
propositions. To the system of the judgments, of the logical-antilogical
judgments, there corresponds selfhood in the form of the self, the
reflecting, knowing subjectivity. The knowing subjectivity --- this holds
both of the powerful subjectivity, infinite in its ontological knowledge,
and of the powerless subjectivity, restricted in its psychological
knowledge --- cannot in its pure knowledge act, for that lies in the
nature of knowledge: knowledge is not acting; the knower can in his
knowledge only think and judge.

To the pervasive difference between thinking and acting, knowledge and
power, theory and practice, there corresponds, then, a pervasive
opposition between the comprehending subjectivity, which thinks the
action but does not act, and the working principle, which does the deeds
of thought but does not think. A dualism of subjectivity and principle of
nature cannot deter the logical inquirer; for such a dualism serves
precisely to ground the identity. The identity rests in the power.
Insofar as the power is, according to its ideality, originally one with
knowledge, it belongs to subjectivity: the ontological subjectivity is,
in infinite knowledge, infinitely powerful; insofar as the power is, in
its immediate reality, originally opposed to knowledge, it founds
otherness, posits objectivity and is identical with the principle.

In subjectivity the objective power is subjective on
% --- p. 143 ---
\opage{143}account of its ideality, its unity with knowledge; in the
principle the subjective power is objective on account of its otherness,
its reality, its unity with working. Here, then, there are not two
powers, an ideal and subjective one alongside a real and objective one;
it is one and the same power that in powerful self-opposition is ideal in
thinking and judging, real in working and producing.

In its system of antilogical judgments subjectivity thus has ideal power
over the principle with all the activities proceeding from it. By virtue
of the ideal power every natural action is determined, link for link, by
the corresponding judgment. The principle does, in all expressions of
power, the deed of thought. In its system of antilogical propositions the
principle has the conceptual determinations reflected in the judgment as
real motives of power, so that the production satisfies the concept
precisely by the object's opposition to the concept.

Every natural production, the organic just as much as the chemical and
the mechanical, thus takes place objectively in accordance with the
principle that moves in the concept, but is determined subjectively by
the concept reflected in subjectivity. The objective in the production is
thus not without further ado subjective, but receives its significance
through the subjective; the subjective in the production is not itself
objective, but receives its significance through the objective. That the
objective is punctually explained by the subjective, and this by that,
that the objective is objective only by signifying the subjective, and
conversely, is a theoretical moment in the practical fundamental relation
that gives the explicative functions their peculiar stamp.

The explicative, i.e. the explaining, appraising, interpreting functions
presuppose precisely a dualism of subjective and objective, of inner and
outer, concept and object. The object is, as object, not concept, but
signifies the concept, and is therefore to be interpreted and understood
by means of the concept. The concept is indwelling in the production
(immanent) and yet hovering above it
% --- p. 144 ---
\opage{144}(transcendent), indwelling as a word of power, hovering above
as a word of knowledge; upon the true relation between the indwelling and
the hovering concept rests the truth of the interpretation.

In all activity of nature the concept is immanent; but for that reason it
still cannot possibly, as concept, coincide with the activity itself. If
the concept were so absorbed into the activity that in what is active we
got only the activity itself, then the activity would become
concept-less, and hence meaningless. In order that the concept indwelling
in the activity of nature may be concept, it must necessarily hover
above it; in order to move in the real system as a word of power, it must
be reflected in the ideal system as a word of knowledge.

There is an essential difference between \emph{being} something and
\emph{signifying} something; an object's significance can, insofar as it
is arbitrarily attributed to the object, be wholly indifferent to its
existence; but an object's significance is, when it corresponds precisely
to ground and essence, decisive for the object's being.

The significance that is ascribed to the natural object by the
explicative functions is one with the object's own essence. In such
judgments as: this is a stone, a plant, an animal, the one judging
declares that he recognizes the concept in each of the objects he
encounters: the judgments are logical. According to the logical judgment
the object is its concept, its significance is stated immediately in the
judgment as judgment; here, to that extent, there can be no talk of any
difference between signifying and being: what the object is, it
signifies, and what it signifies, that it is. It is otherwise when the
question concerns the origin of the natural object, when the produced is
to be understood from the production. Here the relation reverses itself:
instead of beginning with the object and grasping it theoretically in the
concept, the judging subjectivity here starts out from the concept and
wants to understand practically how a given concept has been realized in
the determinate object. ``The stone is a
% --- p. 145 ---
\opage{145}natural product like this''; ``the plant is a natural product
like this''; such antilogical judgments presuppose that the one judging
is originally certain of the concept's objective validity, and the
judgment now decides in each individual case how the concept in question
has realized itself in the given object. From this standpoint there can
here be talk of the object's representing its concept, signifying the
concept, receiving its significance and therewith its interpretation from
the concept. ``Is this formation a leaf, or is it perhaps a branch?''
here is in botany a matter that must be investigated more closely. What
now does the botanist do? He sets the concept branch, and the manner in
which this concept realizes itself morphologically, sharply over against
the concept leaf and the morphological appearance of the leaf-concept,
and the judgment then falls that this, e.g., ``must be a branch!'' Here
it is seen expressly that the object signifies its concept, and receives
its interpretation and significance from the concept.

The reciprocal determination of judgment and object has shown itself to
be both logical and antilogical. In the logical judgment the
determination is anticipating, in the antilogical it is realizing, in the
logical-antilogical it is explicating. ``But'' --- it will be asked ---
``how can a judgment at one and the same time be both logical and
antilogical? The logical judgment has its form, the antilogical has its;
but has logic any form for a judgment that is both logical and
antilogical?'' We answer: a special form for the logical-antilogical
judgment logic assuredly does not have, insofar as by special form one
will think exclusively of a special proposition. But no weight is laid
here upon the special proposition; the whole weight must fall upon the
special thought. The special thought utters itself in all judgments that
expressly indicate that the one judging recognizes, in the produced
object, the concept of the production. Such judgments announce themselves
under very different turns of phrase: ``this must be a branch''; ``so
must a building be''; ``this is a
% --- p. 146 ---
\opage{146}poor specimen, an excellent specimen of'' ``an abnormity of''
etc.

If reflection is now directed particularly upon the concept immanent in
the production, then other is determined point for point by other,
objective by objective upon objective possibilities; through the
objective possibilities and their conditioned realization the observer
sees a reflected gleam of the thoughts that guide the activity of the
principle; the inner is absorbed into the outer, and we must, in the
objective sense, affirm the poet's words: ``die Natur hat weder Kern noch
Schale''. If, on the other hand, reflection is directed upon the
difference between the concept and the activity, then human knowledge
meets a twofold barrier. In the first place it must acknowledge it to be
impossible to pursue the reflection of the immanent natural concept in
the transitional forms of transcendence through all intermediate
determinations into the infinite depth of the ideal system. In the second
place it must acknowledge it to be impossible to reproduce productions
that presuppose a creative unity of knowledge and power, of antilogical
judgments and practical-antilogical propositions.

How the concept of a triangle guides the construction of a triangle can
become clear to us, for we can both think the concept and construct the
triangle; but how the concept of the organism can move the organic
principle, and the principle moved by the concept the material parts,
will, under the present conditions of human cognition, never become
transparent: we can see \emph{that}, but not \emph{how}. Whether, then,
we consider the concept's immanence or its transcendence, nature has in
both respects its mystery, its inner, an inner that loses itself in the
unfathomable; with this explanation it sounds, not like a philistine
remark, but like a voice from the deep: ``Ins Inn're der Natur dringt
kein erschaffner Geist''.

But although the interpretation of the concept's relation to the object
thereby produced must remain imperfect, insofar as the unity of thinking
and real producing is
% --- p. 147 ---
\opage{147}imperfect, it nevertheless stands just as firm that the
produced object can be explained only by a subjectivity that is able, in
a logical system of antilogical judgments, to reproduce the production.
The explicative functions have significance not merely for human
knowledge, they have absolute significance for the divine: the knowing
power that produces the object is eternally one with the powerful
knowledge that explains the production and in the produced sees the
concept satisfied.\footnote{It is said in the language of religion: ``And
God saw everything that he had made, and behold, it was all very good''.
That the explication must in the logical sense take up determinations of
such categories as: appraisal, acknowledgment, recognition, etc., goes
without saying.}

\textit{γγ}) \emph{The Idea in Judgment and Object: ideal
Power-Functions.} The real world is a world of objects, the ideal a world
of thoughts; thought has as thought its expression in the judgment, and
therefore the relation between judgment and object is, from the
standpoint of totality, decisive for the determination of the relation
between the ideal and the real world. The ideal world has inner existence
in the word of knowledge, the real outer existence in the word of power:
by the power of the word the dualism is carried through; by the power of
the word the dualism must again be sublated. In the sublation of the
dualism the Absolute, the idea, is seen again; but the idea, the
Absolute, cannot with its deep self-opposition possibly come to
revelation in any merely finite and relative form; the Absolute demands
an infinite and absolute form; but an infinite and yet absolute form can
be given only in the word.

Through the word, the word of omnipotence, the Absolute moves in a double
circuit, a subjective and an objective. In the subjective circuit the
power is ideal-real; the real world proceeds from the ideal and returns
in infinity to
% --- p. 148 ---
\opage{148}it: the Absolute is posited in subjective form, and this
Absolute, the Absolute in subjective form, corresponds to the ideality of
the power-word. In the objective circuit the power is real-ideal; the
ideal world mirrors the real, as the latter encircles itself through
subjectivity: the Absolute is posited in objective form, and this
Absolute, the Absolute in objective form, corresponds to the reality of
the power-word. Such is the self-opposition of the Absolute. The
self-opposition of the Absolute then shows itself, according to its
actual unity, objectified in a power-language corresponding to the
principle of doubling. Thus:

1) \emph{The Ideality of the Power-Word: imperative Functions.} In the
course of the development of the relation between judgment and object the
power-word has already presented both its ideal and its real side: the
development has proceeded in such a way that it runs back into itself
from anticipating through executive and explicative functions: it begins
and ends with the anticipation. Seen cursorily, it might seem as though
the explicative functions placed themselves in direct opposition both to
the executive and to the anticipating ones, and in this way formed a
final conclusion; but this is not how the matter stands when it is seen
in the principle: the explicative functions form, on the contrary, an
essential presupposition for the anticipating ones.

That the anticipation is determined by objective considerations, a look
at the actual, has been shown above; it is the objective possibilities
that are at issue. But the objective possibilities are discovered only by
the interpretation of the actual, and hence only by explicative
functions. That the explication cannot satisfy the infinite subjectivity
by finite results is then also easy to see. The result satisfying in the
finite sense would have to be a determinate object, a completed work, in
which the judging subjectivity, explicating the work's significance,
could see the concept perfectly realized; but the object, the Object,
which here represents the satisfied concept, is in the quality of an
actual object
% --- p. 149 ---
\opage{149}subject to change; through the change the object adequate to
the concept becomes precisely inadequate: the actuality determined by the
concept is, so to speak, pregnant with possibilities. If the realization
of the new possibilities is now also to satisfy the concept, the
comprehending subjectivity must by anticipating functions again determine
these possibilities and preform the actuality in which the concept is to
be realized. The new actuality is then once more an object adequate to
the concept, which by being changed becomes inadequate; the inadequate
object then conditions a new anticipation, and so on to infinity.

If the eternal circuit of the anticipations is now intuited in the light
of the Absolute, then the whole of actuality, with all the elements
belonging to its infinite totality, becomes shone through by ideality, a
medium for the idea's self-revelation; every actuality-element is by the
anticipation a moment of the idea from the beginning: every link in the
whole of actuality, every actual shape has, according to its
objective-subjective possibility as idea-moment, originally been
anticipated and objectively preformed in the infinite subjectivity. In
the eternal circuit of the anticipations the real world is seen to
proceed unceasingly from the ideal; the real system is preformed in the
ideal system: the Absolute is in its subjective circuit.

The actual productions find themselves, then, in a double system of
relations; from the objective side they are subjected to the conditions
of finitude and forfeited to relativity; from the subjective side they
are announcements of the unconditioned, utterances of the Absolute.
Insofar now as a production directly utters the Absolute in the word of
power, it shows itself as a primitive natural production; insofar as it
belongs, in the circle of its conditions, to relativity and to finite
causality, it ranks in a class with ordinary natural products. The
primitive natural production is accordingly also called, insofar as
emphasis is here laid in a special sense upon the ideality of the
power-word,
% --- p. 150 ---
\opage{150}a \emph{creation}. Since now religion has in its language
impressed upon the category creation a quite peculiar stamp, the question
arises in this connection how the religious significance of the category
stands to the scientific one, and what dogmatics, especially the
speculative, has done to clarify this relation.

How ``revealed religion'' expresses itself in the narrative of the
creation is well known. Whether we translate the first verse of Genesis
by: In the beginning God \emph{created}, or by: In the beginning God
\emph{had} \emph{created} the heaven and the earth, the creation is in
both cases conceived as an event, for which reason it is also aptly
presented in the form of a narrative: to the category event there
corresponds the category narrative. If the religious narrative is now to
be transformed into a link in a doctrinal edifice, then the
representation uttered in the narrative must, if not exactly changed into
concept, yet be fixed by a conceptual determination, in order to be able
to be set up as dogma. Thus the doctrine of creation passes from religion
through theology over into dogmatics.

If we choose as an example one of the older orthodox definitions of
creation: \textit{Actio Dei triuni externa qua is res omnes visibiles ex
nihilo, sex dierum spatio, solo liberrimæ voluntatis suæ imperio
omnipotenter et sapienter produxit, in nominis sui laudem et hominum
utilitatem} (Quenstedt), it is seen from this that dogmatics, as
ecclesiastical doctrinal concept, conceives the creation as an event at
the beginning of time; to this there then corresponds the creation out of
nothing, a creation that is referred exclusively to God's free,
omnipotent will.

Between the fixed creation-dogma of the doctrine of the Church and the
rational creation-category of which the doctrine of the idea can avail
itself, when it sees in primitive natural productions a determinate
utterance of the Absolute, a preformation of the infinite subjectivity,
there is a yawning chasm, an engulfing abyss.
% --- p. 151 ---
\opage{151}This abyss, this chasm, speculative dogmatics has nevertheless
pretended to be able --- of course, to a certain degree --- to fill up
and level out; we shall now see what ``filling'' it uses for the filling,
and how it sets about levelling the chasm.

According to the doctrine of the Church the creation is a miraculous
event, not a fundamental relation; every accessory representation of an
eternal creation-relation, a creation from eternity, must be rejected as
an intolerable heresy. According to the doctrine of the idea the creation
is a fundamental relation, not a miraculous event; every accessory
representation of miracle and miraculous event must be rejected as an
imagination.

It is now in characteristic agreement with the ambiguous standpoint of
speculative dogmatics --- for the standpoint is from first to last an
ambiguity --- that in all principal points it is moved by heterogeneous
considerations, motives that work against one another and end by blunting
and neutralizing one another mutually. This is seen especially in the
treatment of the creation-dogma. Dogmatic speculation will not break with
orthodoxy, will at any price avoid the appearance of heresy; here is the
first objective consideration, the ecclesiastical motive. By virtue of the
ecclesiastical motive speculation makes a show of being in earnest about
the creation as a miraculous event, about the creation out of nothing,
about the creation out of unconditioned freedom. ``The Mosaic
creation-history expresses the fundamental thought of the creation in
this, that the world came to be by God's almighty word. He spoke: ``Let
there be light''! and there was light. Each of the six days of creation
arises only by virtue of the almighty creator-word''. To a certain degree
the creation out of nothing is accordingly emphasized in agreement with
this as well: ``In creating, God calls the non-being into existence. This
is the significance of the old doctrine that God created the world out of
nothing''.

But, of course, these emphases naturally hold good
% --- p. 152 ---
\opage{152}only to a certain degree. Here there is, besides the objective
consideration for the Church, yet another objective consideration, the
objective consideration for science. Moved by the scientific motive,
speculation turns gently, for the uninitiated almost imperceptibly, away
from the ecclesiastical. The creation is now not merely a miraculous
event, but a fundamental relation, the relation in which God must
originally stand to the world. By virtue of the scientific motive ``each
of the six days of creation'' then becomes ``each new epoch in the
world-system''. The nothing out of which the world is created is not = 0,
against which representation there holds: \textit{ex nihilo nihil fit};
the nothing out of which the world is created is then neither the orthodox
nothing, that \textit{nihil negativum, negatio omnis entitatis}, according
to which a Hollazius can meet Aristotle with a proposition such as this:
\textit{ex nihilo nihil per vires naturæ generatur; at per vim Dei
infinitam omnia creata sunt ex nihilo}. Nor can the scientific motive
agree well with the orthodox when the talk is of the unconditioned freedom
(\textit{solum liberrimæ voluntatis imperium}) with which God in the
ecclesiastical sense has created the world. To be sure God has on the one
side created the world with full freedom; but on the other side the world
has nevertheless come to be with indefeasible necessity. To be sure it is
a love-necessity; but necessity is nevertheless always necessity. ``Just
because God is love, he \emph{cannot} close himself off as the God ``of
the ideas'', but \emph{must} determine himself, etc. Yes, he must
determine himself, for it is ``freedom's exceeding riches, which cannot do
otherwise than will'' \dots\ See, now the result is just as ambiguous as
the point of departure.

When the Church now asks after freedom, necessity goes up into
``freedom's exceeding riches''; and when science demands necessity, why,
then ``freedom's exceeding riches, which cannot do otherwise than will''
\dots, goes under in necessity. ``From this standpoint it is evident
% --- p. 153 ---
\opage{153}in what significance we deny, and in what significance we
affirm, the proposition: Without the world God is not God''.

In what significance do we then deny the proposition? In the
ecclesiastical orthodox significance, naturally! for necessity goes under
in freedom. In what significance do we affirm the proposition? In the
scientific, that is to say: in the speculative scientific significance,
naturally! for freedom goes under in necessity. Do we then really deny the
proposition? Yes, for the ecclesiastical emphasis falls on freedom. But is
it now also certain that we really deny the proposition? No, for the
scientific emphasis falls nevertheless on necessity. Thus Yes and No
become at last a good theology.

``But'', one might object, ``Yes and No are in this connection
dialectical determinations of one and the same thing; freedom is a
dialectical determination of necessity, necessity a dialectical
determination in freedom: why should dogmatics not have the same
scientific right to avail itself of dialectic as philosophy has''? To this
it must be answered that philosophy, as far as dialectic is concerned, is
bound to make a distinction between the two-sided and the ambiguous. This
distinction can, in the present connection, where everything turns upon
determinations of the ideality in the word of power, be made clear by an
elucidation of the power-word's \emph{imperative} functions.

In the executive functions there showed itself a practical difference
between the determining thought and the actual action determined by
thought; the unity corresponding to the difference is to be sought in the
practical self-determining, in the will. But the will's ideal acting upon
the bodily real must immediately rest upon a bidding identical with
acting. One can paraphrase pure willing in many ways, but the determining
that proceeds from the will is essentially commanding and bidding. Insofar
as the actual production is conceived as a realized concept, it stands at
the same time as objective expression for a
% --- p. 154 ---
\opage{154}Become of subjectivity, a sign of the power of ideality. The
fact is a factum; every factum presupposes a \textit{fieri}, every
\textit{fieri} a \textit{fiat}: in the Absolute's subjective circuit the
executive functions go up into the imperative ones.

In consequence of the imperative functions every factum is thus led back
to a \textit{fiat}, the objective in the production back to the
subjective, the power-word in the object back to the word of knowledge in
the judgment; but to refer factum to \textit{fiat}, the objective with its
conditions to the subjectivity determining in the concept, is, in other
words, to subordinate necessity to freedom, to lead the activity of nature
back to creation. The validity of the imperative functions must in this
respect be acknowledged by the doctrine of the idea just as much as by the
doctrine of faith; but when the acknowledgment is more closely determined,
the chasm opens anew, the deep chasm that sets a separation between faith
and knowledge.

In the doctrine of faith the divine Become is taken absolutely; the act
of creation is unconditionedly free, the divine power and will knows
nothing, will know nothing, of any impelling necessity whatever. God is
God from eternity; but the world is not from eternity; God is God from
eternity --- without the world. The motive for the creation, possibility
for the creation, plan and purpose for the creation, all is gathered in
one single moment and concentrated with infinity's energy in one word, one
single word, the almighty word: \emph{Become}! With this word time begins,
prior to this word is eternity; with this word the world begins, prior to
this word is God alone, God without the world.

This absolute conception of the absolute Become on the part of faith the
doctrine of the idea is not able to appropriate. That God is God from
eternity means, when the doctrine of the idea is to translate it into the
concept, that the divine essence, the subjectivity free in its inner
infinity, is with necessity. No almighty Become can bid the Godhead itself
become something;
% --- p. 155 ---
\opage{155}no almighty Become can bid the Godhead itself become nothing.
But in the concept as concept selfhood could not possibly be if otherness,
the self's opposite, were not originally; seen in the concept as concept
subjectivity, the subjectivity free in its inner infinity, could not be if
objectivity, the real objectivity, the object's outwardly actual world,
were not actually. The almighty Become that was to call the real
objectivity from non-being to being, this same almighty Become would then
also have to call the infinite subjectivity from non-being to being; the
almighty Become that was to call the outwardly actual world from nothing
to something would at the same time have to call the inner, divine essence
itself from nothing to something. This almighty Become, unconditioned in
its absoluteness, is to science an incomprehensible.

But when science cannot comprehend a coming-to-be out of nothing, how can
it then acknowledge the category of creation? Here weight must again be
laid upon the imperative functions. What is rational in the creation
rests, like what is rational in the ideality of the power-word, the
almighty Become, upon the fundamental relation of the subjective and the
objective: in this fundamental relation selfhood, the infinite
subjectivity that through another relates itself to itself, is eternally
the overpowering, that which overreaches in everything. By the power that
overreaches into the infinite the subjective determining, the divine
Become, acquires original significance for all natural productions; out of
this significance there shines forth what is rational in the category of
creation. The doctrine of the idea can from this standpoint cognize that
the world is created, that it has come to be out of nothing in original
opposition to God; but the coming-to-be out of nothing of which the
doctrine of the idea speaks is in categorical conflict with the creation
out of nothing upon which the doctrine of faith builds.

It has been treated in several places in our foregoing how philosophy has
striven with itself in order to bring matter's reality into accord with
reason's demand;
% --- p. 156 ---
\opage{156}it has been shown that the solution of the riddle must be
sought in the concept of power: the masses with their atoms are facts,
punctual manifestations of power's immediacy. Since now power's real
immediacy is a moment in power's own reality, and the real power in
subjectivity itself identical with the ideal, it is understood from this
that the subjectivity which anticipates everything actual and preforms all
existent shapes cannot have either eternally independent atoms or material
or chaos outside its power, which overreaches in all things. All
determinations of material, all material elements are real negations in
the Godhead's ideal power, antilogical moments in the all-capable
ideality; everything is formed and produced by power itself, the divine
power: outside the divine power there is nothing, nothing at all, out of
which something can be formed.

But with such a nothing the doctrine of faith cannot be content. For
faith the creation out of nothing means that the whole world at
omnipotence's bidding proceeds forth out of nothing; herewith it is
presupposed that the world, before the Creator uttered his Become, was
nothing, nothing at all, that there goes, as shown above, an eternity
before the final creation, and that in that eternity there was no actual
world. This religious fundamental thought no doctrine of the idea will be
able to raise to scientific concept. Were God to be thought as being an
eternity before the world, then the ideal power would indeed have to be an
eternity before the real. With the real power, with the real manifestation
of the divine power, the world is there instantly; --- for what is the
world, seen from the standpoint of creation, other than the real
manifestation of divine power? --- If, then, a whole eternity elapses
before the creating Become lets the world come forth, then a whole
eternity also elapses before the creating Become gives the ideal power
reality. But the ideal power by itself, without the real power, is only
half-power, not omnipotence; only the identity of ideal and real power is
omnipotence.

Seen from this standpoint the miracle of creation now becomes
% --- p. 157 ---
\opage{157}two miracles, the one, if possible, more wondrous than the
other. First it is a miracle that the Creator's absolute Become posits a
world in nothing's place, a world where formerly nothing was; this miracle
presupposes omnipotence. Next, the coming-to-be of omnipotence is itself a
miracle. In order to work the miracle with the world, the Creator must
work a miracle with himself, in that he must call the real power from
non-being to being; in order to work omnipotence's miracle, the Creator
must first by a miracle make himself almighty.

And what then does this prove? it does not prove that religion is in the
wrong against science, for religion does not undertake to make the miracle
of creation comprehensible; but it proves that every scientific attempt to
explain the almighty Become by ascribing abstractly unconditioned validity
to the imperative functions must end in the meaningless.

2) \emph{The Reality of the Power-Word}: \emph{indicative Functions.}
In the system of the imperative functions subjectivity has disposal over
the produced at all stages of the production. The production begins with
the explication's transition or rather return to anticipation; from both
sides the beginning is determined by the determining subjectivity. The
beginning production is subjected to the executive functions; the actual
actions, the practically antilogical propositions proceed from antilogical
judgments, both in their unity from the thinking, acting subjectivity. The
result of the completed production, the produced, which according to the
appraisal acquires significance in this, that it both satisfies and yet
does not satisfy the subjectivity explicating its production, then becomes
a moment for new anticipations, new productions, etc., thus constantly
subjected to subjectivity. Thus the Absolute has taken up all the elements
and conditions of the production as moments in its own subjective circuit:
subjectivity's superiority culminates in an all-capable Become.

% --- p. 158 ---
\opage{158}But that the objective elements of the production are immanent
in subjectivity presupposes at the same time that the conceptual
determinations, the subjective moments of the production, are immanent in
objectivity. The whole circular movement must thus, in order to be seen in
its reality, be considered from the objective side. In the anticipating as
well as in the explicating functions it is just as much the object that
determines the judgment as conversely the judgment that determines the
object; but that the object determines the judgment means, in other words,
that the objective determines the subjective. The infinite reciprocal
determination of objective by subjective, and conversely, must then in the
objective sense end with the objective determining itself, through the
subjective, point for point, from the production's beginning to its
completion. Herewith objectivity is expressly posited as the Absolute's
own form, as the idea's immanent form: the Absolute is in its objective
circuit, and this is the reality of the power-word.

By the ideality of the power-word the functions become imperative, by the
reality of the power-word they become indicative; the task is to show how
the indicative relates to the imperative. Everywhere that the objective
determines the subjective, the judgment aims at a straightforward
indication of what the objective is; in this indication the indicative is
seen. The judgments treated in formal logic are all indicative. That the
determination of the indicative nevertheless cannot come forward in the
usual logic has its ground in this, that it lacks its opposition, namely
the imperative: we shall now see how the imperative fades away and comes
to rest in the indicative.

The objective must be objectified, the objective productions preformed;
but however great a part subjectivity may have in preformation and
execution, its co-operation can nevertheless never go further than to
determine the objective rightly objectively. To determine the objective
objectively is not to
% --- p. 159 ---
\opage{159}determine it immediately at any point, but to let it determine
itself at all points. The objective in the production now comes forward
from all sides in such a way that the producing subjectivity seems to
vanish altogether. The only thing subjectivity still has to do is to
objectify for itself the objective production; but this objectifying is so
theoretical, so contemplative, that nothing is altered in what objectively
takes place.

The difference between the finite and the infinite subjectivity is
decisive on this point too. When the finite, the psychological
subjectivity is to act determiningly upon objective productions, one
notices at once the separation between the subjectively arbitrary and the
objectively necessary; just because the objective sets such immediate
barriers to the impotent subjectivity, the latter must in turn just as
immediately break through these barriers at single points; here there
arises a finite conflict between what the subject wills and what lies
immediately in the objective; therefore in all manufactures and
art-products the influence of subjectivity is so evident.

It is otherwise with the infinite subjectivity; its acting upon the
objective does not take place in a finite manner and is for that reason
not immediately recognizable. Although everything is anticipated,
preformed and subjectively determined by the mighty Become, although there
is not an element in the objective whose being has not been prepared
subjectively, yet in actuality it looks as though the objective went by
itself. Yes, the objective really does go by itself objectively; here in
the appearance there is no delusion, no deception: in all productions the
actual determination is a determination of objective by objective upon
objective possibilities; subjectivity has only to look on, its functions
are indicative.

The indicative in the functions has nevertheless not displaced, but
transfigured the imperative: the indicative shows that
% --- p. 160 ---
\opage{160}the development of the objective is objective; but the
imperative shows that the objective is objective only by means of
subjectivity. The indicative shows that the actions of nature are guided
by principles and immanent concepts; but the imperative shows that the
principles are selfhood-factors and the immanent concepts
logical-antilogical utterances of a comprehending subjectivity. The
indicative shows that every creation is really a natural production; but
the imperative shows that every natural production is originally a
creation. Here is identity of nature and creation.

But speculative dogmatics too wills identity of nature and creation; let
us then again listen a little to speculative dogmatics. ``As the Johannine
view contains that the world's existence has its ground in a creative
\emph{producing}, so it contains that the world is in consequence of a
\emph{transition} from non-being to being, by a \emph{becoming}, a
coming-to-be, a birth, a \textit{fieri}, a \textit{γίγνεσθαι}. The world
thus has a double beginning, a cosmogonic and a creative, a natural and a
supernatural. The cosmogonic \dots\ is the relative, the finite, which
therefore is also split up into \emph{sporadic} multiplicity \dots\ But
the innumerably many natural beginnings are all grounded in the one
creative, supernatural beginning, in the divine Logos-will, which encloses
within itself the wellspring of life and of light, and out of its
spiritual depth releases the whole multiplicity of life-forces to free
self-movement. And only because this supernatural beginning, only because
the creating will continues to stir in the many finite beginnings and in
free omnipresence illuminates and works through the natural development,
can the latter thoroughly get beyond the chaotic ferment, can the sporadic
oppositions be united and harmonized into an organic and systematic whole.
The world must therefore in every moment of its existence be considered as
\textit{natura} or as organic self-development, and as \textit{creatura}
or as progressive revelation of the
% --- p. 161 ---
\opage{161}divine \emph{will}, and it is the one only \emph{because} it is
the other''.

Against this utterance two objections hold: first, that it is adapted to
the rational concept of creation, but stands in conflict with the
ecclesiastical dogma; next, that it is kept in a merely descriptive form,
and is for that reason conceptless.

As regards the first point, the creation out of nothing (\textit{creatio
ex nihilo}) is in the ecclesiastical sense an expression for the divine
will's absolute, unconditioned power over everything that we in the sense
of actuality call otherness, objectivity, thing-causes, etc. So far is it
from being the case that the divine self should essentially need such an
otherness, that on the contrary it can abolish it whenever it will; so far
is it from being the case that the divine subjectivity should for its own
actuality need an actually objective, that it can, at whatever moment it
will, abolish objectivity. That the things and the thing-causes are
created out of nothing means that they can again, at whatever moment it
pleases omnipotence, become nothing; to \textit{creatio rerum} there
corresponds in the old dogmatics --- and that with true consistency ---
\textit{annihilatio rerum}.

``There is in the history of theology no dogmatician who has held fast
the concept of God's omnipotence with greater rigor than Malebranche, none
who has expressed with greater scientific consistency what lies in the
creation-dogma than precisely Malebranche.\,\dots\ To be sure Malebranche
has not been able in a speculative manner to see into the Father's primal
ground, to discover world-germs in the Son's Pleroma and to mediate the
process of creation in the plastic force of the Spirit; but the dogma he
has, the pure dogma of the creation out of nothing. Of this dogma he has
not merely declared that it must be thought (for because it is now and
then said of a dogma that it must be thought, it is not thereby said that
it really is thought); no, he has in earnest attempted to think it, to
hold the thought fast and to apply it to a scientific understanding of the
created.
% --- p. 162 ---
\opage{162}To apply the thought of creation to an understanding of the
created means: to see the whole world under a double point of view:
\emph{nothing and yet really something}. The concept in which this
contradiction is resolved is the concept omnipotence; the created is in
itself nothing, but at omnipotence's bidding it has become something. He
who now looks upon what has come to be so one-sidedly that he sees only
its nothingness, for him the world is semblance and shadow. Whether
theology will call him Pantheist or Acosmist is a matter of indifference;
the point is that he denies the creation in denying omnipotence. He who
looks upon what has come to be so one-sidedly that he sees only its
reality, forgets its nothingness and ascribes essentiality to it, for him
the coming-to-be is a natural process and the world a realm of nature.
Whether dogmatics must for the rest reckon him among Deists, Naturalists
or Atheists is in this connection a matter of indifference; the point is
that he denies the creation in denying omnipotence.\,\dots\ It has indeed
been supposed that one could go a step further and improve the doctrine of
creation by ascribing to the created some independence after all, some
disposition to inner self-development, some capacity for activity \dots\
It is however a dubious improvement \dots\ For what is a relative
independence? A degree of determinate being outside omnipotence, a
self-acting which, however slight it may for the rest be, is nevertheless
not omnipotence's acting, thus: an actual limit to omnipotence. If this
actual limit is now meant in earnest, then omnipotence is an actually
bounded power and no longer omnipotence. If on the other hand the Almighty
still has, as before, the limit in his power, then the actual limit is
still, as before, a nothing in itself, and the miraculous relation between
Creator and creature must continue. A relation that has begun by a miracle
must also be continued by a miracle: not merely for some moments, not
merely, if I may so put it, for a week or a fortnight after the creation,
but for centuries, for millennia, in the eternity of eternities must the
created things continue to be in themselves what they
% --- p. 163 ---
\opage{163}originally are and have been, a nothing, a pure nothing for the
Almighty. It is this consistency that Malebranche holds fast, and
therefore he cannot think self-activity in any creature, be it angel or
man; for essential self-activity is a divine matter, and God himself
cannot create gods: \textit{il n'en peut faire de véritables causes, il
n'en peut faire des Dieux.}''\footnote{Om Theologiens Naturbegreb af R.\ Nielsen. I Anledning af Reformationsfesten 1855. p.~22 ff.}

Are the causes of nature, the actual thing-causes, created, namely in the
dogmatic-ecclesiastical meaning, or are they uncreated? This is the
question, the inescapable fundamental question. If the question is
answered with Yes: the causes of nature are created! then the world quite
rightly has ``a supernatural,'' a ``creative'' beginning; but then it is
nonsense to speak of its \emph{also} having a ``cosmogonic'', a
``natural'' beginning. If the question is answered with No: the causes of
nature are uncreated, then the world quite rightly has a ``natural'' and
``cosmogonic'' beginning; but then it is nonsense to speak of the world's
\emph{also} having a ``supernatural'' and ``creative'' beginning.

Only on the condition that the category of creation is conceived
rationally can the identity of nature and creation be comprehended. The
contradiction that the created should be both something and yet in itself
nothing is then resolved in this, that power is self-power only by having
power over another, and can have power over another only when this other
itself stands in force. Separated from the infinite self-power the
subsisting other is a semblance, a nothing; but in unity with self-power
the existing other is substantial; its substantiality is, by opposition to
self-power, the latter's real moment.\footnote{Grundideernes Logik. I. p.~112--113.}
The contradiction that there should here be an actual limit to omnipotence
is resolved
% --- p. 164 ---
\opage{164}in this, that the actual limit, power's own real negation,
precisely expresses power's reality: without real negation power would be
only ideal, not real; but the abstractly ideal power is only half-power,
not omnipotence. Instead of making omnipotence into an actually bounded,
thus restricted power, the actual limit is on the contrary a categorical
expression for the original doubling of power, for the power-word's
eternal reality.

When speculative dogmatics here and there makes a run at explaining the
created's ``given independence'' from this, that God renounces his power,
or that he releases ``the magical possibilities,'' then this is an
unmistakable testimony that ``the dogmatic talent'' has renounced thought
and released all manner of fantastic notions. Let one but try to think
something by a notion such as this, that the Almighty in the creation has
really --- of course, to a certain degree --- renounced his power. What
concept is one now to form of such a renunciation; it must surely have
happened either by power or by impotence? If it has happened by impotence,
thus from lack of power, if God during the creation has, on account of the
``given independence'' of the created, fallen into a certain voluntary
impotence, then the miracle of creation becomes not a miracle of
omnipotence but a miracle of impotence. If on the other hand it has
happened by omnipotence; if God during the creation has renounced his
power in such a way that this his renunciation has precisely become a
proof of his infinite power, unrestricted in all respects, then he has
also, by renouncing his power, retained his power, and consequently not
renounced his power at all.

To speculate about the creation without insight into the logical
functions belonging here is a dubious business; only by reducing the
relation between Creator and creature to a fundamental relation between
subjectivity and the objective actuality, and thereby to a logical
relation between imperative and
% --- p. 165 ---
\opage{165}indicative functions, will it be possible to come to clarity
about categories so rich in content as \emph{nature} and \emph{creation}.

In consequence of the imperative functions the natural objects, the
actual objects, are moments for the producing subjectivity, the infinitely
overreaching self-power; in consequence of the indicative functions the
natural objects, the existing objects, have essential independence over
against subjectivity; in consequence of the imperative functions the
thinking, judging subjectivity is \emph{acting}, in consequence of the
indicative functions the acting subjectivity is thinking, \emph{judging}.
The fundamental relation between the indicative functions and the actual
world of the natural objects then lets itself be reduced to the idea in
judgment and object.

If the actual object could really determine the judgment without itself
being determined by the judgment, then subjectivity would have to relate
itself, in its system of indicative functions, to the natural objects as a
mirror to the mirrored objects. But even in the purest contemplation the
subjective mirroring is nevertheless more than a mirroring; the subjective
mirroring is in its ideality an incessant thinking and judging. Outside
the idea, thus on the level of finitude, the object about which judgment
is passed stands in an outward and indifferent relation to the judgment
and the judging subjectivity. The object determines the judgment, but the
judgment's indicative function does not make it recognizable that the
judgment at the same time determines the object. Where this moment is
psychologically emphasized, the distance between subjectivity and ``the
thing in itself'' becomes, as we have seen above, only the greater.

It is otherwise when the relation between judgment and object is seen in
the light of the idea. Here it becomes clear that the object can determine
the judgment only insofar as the judgment determines the object. The most
decisive expression of the fact that the judgment essentially determines
the object is had in the imperative functions. On the basis of the
imperative functions the
% --- p. 166 ---
\opage{166}divine subjectivity consider the objects of nature in their
objective immediacy as productions of the word of power. The objective
element in these productions, which is characteristic of the
power-word's reality, is determinative for the judgment's indicative
functions.

The judging subjectivity, which contemplatively lets the object
determine the judgment, rests with infinite superiority in the
contemplation of the objective development of the objective, because in
its self-absorption it eternally knows that the object which determines
the judgment is originally itself determined by the judgment; the
contemplative subjectivity, which in the system of indicative functions
conceals itself beneath the productions of nature, then knows itself
identical with the practical subjectivity, which breathes spirit through
the principles and in every natural production recognizes its own
creature.

3) \emph{The categorical Power-Language: interpreting Functions.} What
it signifies that there are thoughts in things, concepts in objects,
must now be clear from the foregoing. But if there are thoughts in
things, then there must surely also be \emph{meaning} in things; is it
then certain that there is, in an objective sense, meaning in things? In
relation to the pure concept, opining stands far beneath knowledge; but
with respect to the relation between concept and existence an insight
into the meaning of things can stand above an objective knowledge of the
essence of things. For the cognition of the objective essence of things
it is sufficient to lay finite causality at the ground and to explore
the laws of things; for the apprehension of the meaning of things there
is required a glance at the task which existence has solved, an insight
into the purpose that was to be attained by means of things, a purpose
which alone can be said to give them, not an immediate being, but an
essential significance; that there are \emph{thoughts} in things,
physics teaches; but what it properly means that there is
\emph{meaning} in things must be cleared up in logic.

In the logically explicative a distinction has been made between
% --- p. 167 ---
\opage{167}a comprehending of the object as a whole and a comprehending
of the object's parts. In this respect there is said to be ``a
\emph{twofold kind of comprehending}.'' --- ``As regards, namely, the
given object in the \emph{whole}, there is discerned,'' it is said,
``\emph{how} the connection therein is constituted, or \emph{how} the
manifold therein is bound into unity. As regards, on the other hand,
the \emph{singularities} and the single relations, there is discerned,
a) \emph{why they are} and b) \emph{why they are} so constituted
\emph{as they are}, in that there is discerned how the whole object's
mode of connection \emph{carries or can carry them along with it in
this way}; so that they, with the constitution they have, are seen as
\emph{co-requisite} (\emph{integrating}) and at the same time
\emph{co-engaging} or at all events \emph{co-belonging}, or naturally
\emph{entailed}, members in the whole. When therefore, as regards the
object in the whole, merely the \emph{inner constitution} or its
\emph{mode of existence} is comprehended, there are on the other hand
discerned, as regards the singularities and the relations, at the same
time their \emph{grounds of existence} (\textit{rationes essendi}), or
at all events their grounds of possibility, and moreover their
\emph{existence-significance} or their \emph{real significance} in the
given whole. The whole is here only \emph{presented} or
\emph{described} according to its inner constitution; the
singularities, on the other hand, are at the same time
\emph{interpreted} and \emph{explained}. The first kind of
comprehending we will'' --- it is said further --- ``call a
\emph{connection-comprehending}, perhaps also a \emph{formal
comprehending}, and the concept a \emph{connection-concept}, perhaps
also a \emph{formal concept}; the latter, on the other hand, with
respect to the use of the word \emph{reality}, a \emph{real
comprehending}, and the concepts \emph{real concepts}. A comprehending
of the former kind lies here at the ground of the comprehension of the
latter kind; for it is the connection-concept concerning the whole
object's how, its inner how namely (which must not be confused with its
outer constitution), from which the insight into the why of all the
\emph{singular},
% --- p. 168 ---
\opage{168}from which therefore the real comprehending in respect of
all this singular, proceeds.''

That the distinction thus set up between connection-concept and real
concept is superficial, unhelpful and one-sided must, with a little
reflection, easily strike the eye. The whole cannot possibly determine
the parts without the parts conversely having to determine the whole;
if the why of the parts is to be derived from the how of the whole,
then the why of the whole must also be derived from the how of the
parts. Everywhere that there can be talk of a comprehending, there must
be a grounding, and thus a why. Without why no ground, without ground
no grounding, without grounding no comprehending, without comprehending
no concept. A connection-concept concerning the whole object's inner
how without a why cannot be laid at the ground of any insight into the
why of the singular, and that for the simple reason that such a
connection-concept cannot be given. That there can be, e.g., two kinds
of tailors: ladies' tailors, who do not sew for gentlemen, and
gentlemen's tailors, who do not occupy themselves with sewing for
ladies, is not contrary to sound human understanding and is not to be
censured in logic; but that there should now also be two such kinds of
concepts as: pure how-concepts, which do not concern themselves with
any why, and peculiar why-concepts, which float on top of ``the inner
how'', is contrary to sound human understanding, and must be censured
in logic.

The connection-concept concerning the whole object's inner how is
nothing else, and indeed cannot be anything else, than the concept of
the mutual connection of the \emph{single parts}; but the concept of
the parts' mutual connection is nothing else and cannot, regarded
essentially, be anything else than the concept of their constitution as
determined in the connection: as the connection is, so is the
constitution, as the constitution is, so is the connection. If it is
now impossible to form for oneself a connection-concept of the whole
% --- p. 169 ---
\opage{169}without comprehending the connection of the parts, then it
is also impossible to comprehend the connection of the parts without
comprehending \emph{why} each single part has just this and no other
constitution. ``The insight into the why of all the singular'' can only
be laid at the ground of ``a connection-comprehending'' of the object's
inner how insofar as this connection-comprehending is itself laid at
the ground of an insight into the why of all the singular. --- But if
we regard the distinction set up between ``connection-comprehending''
and ``real comprehending'' as a groping indication of the difference,
decisive for the cognition of existence, between the objective concept
and the meaning uttered in the concept, it does after a fashion admit
of application.

``The grounding reflection'' has already taught us that the sufficient
ground is objective as ground of existence only insofar as it is at the
same time subjective as motive-ground. To the two-sided ground there
corresponds here a two-sided why. The objective why reflects itself in
the objective concept, the connection-concept which expresses the
object's inner how; from this how there then shines forth the why of
the whole as well as of the parts: this reflection of why and how gives
us the concept, but not its meaning. One may think of natural products
such as feldspar, granite and the like. It will yield no meaning to lay
the concept of feldspar's inner how at the ground of an insight into
the why of its single parts, when this why is not to concern the ground
of existence but the motive. As it goes with the feldspar, so with the
whole earth, so long as the comprehension keeps to the purely
objective. If, on the other hand, we ask after the meaning of the
feldspar and of the earth-globe's existence, it is at once noticed that
what is aimed at here is not the objective ground, the ground of
existence, but the subjective ground of that which exists, the motive,
the motive-ground holding for the production of the feldspar and of the
earth-globe.

If the object is not a natural product but a manufactured article,
% --- p. 170 ---
\opage{170}e.g.\ a clock, then it certainly lies near at hand to lay
the connection-concept concerning the whole's inner how at the ground
of the cognition of the single parts' why, for the whole's how
involuntarily takes on the stamp of an arbitrarily chosen task. In the
clock the master-workman has solved a determinate task in a determinate
way; therefore the clock has become such a one, just such a whole, with
such parts in such a composition. The whole is, on the master-workman's
side, purpose and the parts means: it is thus a subjective, not an
abstractly objective why that is here reflected upon. Yet even from the
teleological standpoint it holds that it is not merely the why of the
means that must be sought in the how of the purpose, but also
conversely the why of the purpose that must be sought in the how of the
means: in the anticipation it is the purpose that motivates the choice
of the means, insofar as the determining considerations are subjective;
but it is conversely the means offering themselves that motivate the
choice of the purpose, insofar as the determining considerations are
objective.

Everywhere that the relation between a whole and its parts admits of
being apprehended teleologically, the work in question will speak to
the observer in a language that interprets for him the master-workman's
meaning. So surely as there is meaning in the thing, it must be just as
possible for the productive subjectivity to express itself in facts as
to express itself in words; but speech in facts is different from
speech in words: the language of facts is a \emph{categorical
power-language}.

A swinging pendulum represents an objective concept, but this concept
is as yet without meaning; if, on the other hand, the swinging pendulum
is set by the escapement in reciprocal action with a clock's wheelwork,
then the union of pendulum and escapement in the actual composition is
just as eloquent as the union of subject and predicate in the logical
proposition: the pendulum's swing then yields a determinate meaning.
That a barrel, wound about with a cord whose free end bears a weight,
is set in
% --- p. 171 ---
\opage{171}motion by gravity presents to the reflecting person an
objective concept; but the concept for itself is as yet without
meaning. If, on the other hand, the barrel is connected with the
clockwork's barrel-wheel, and in such a way that the latter follows the
barrel only when this is turned by the weight's heaviness, but not
when, during the winding up, it is turned in the opposite direction:
then we have in the mechanical connection of such members just as clear
a meaning as in a grammatical connection of the words of language. Just
as the single mechanical member for itself can yield no meaning, so the
single word for itself can of course yield no meaning either. The
single word for itself can, by the mere naming, designate a concept;
the single member for itself can, by its individual shape, its mere
appearance, represent a concept. The propositions of language can,
according to the meaning, be partly principal clauses, partly
subordinate clauses, partly protases, partly apodoses, etc.; the same
holds in corresponding fashion of the compositions, the
practical-antilogical propositions in the machine that solves its task.
Every composition, every practical-antilogical proposition, is
determined by a logical-antilogical judgment. What the master-workman
has through the system of such judgments said to himself during the
production of his work, that he now says to the thinking observer
through the system of real propositions in which the completed work
yields meaning.

To the language of facts, the categorical power-language, there
corresponds a system of \emph{interpreting} functions. The interpreting
functions have at first glance a striking similarity to the explicative
ones. In both it is a matter of a recognition of the concept uttered in
the production, a judging of what is produced, a development of its
significance; and yet the difference between explicative and
interpreting functions can be brought out with categorical emphasis.

By the explicative function the produced object is measured chiefly by
the production's indwelling
% --- p. 172 ---
con\opage{172}cept: ``this wheel is badly made''; ``this wheel is well
made''; ``thus shall a wheel be''. The explication motivates all
judgments by the concept as concept; it keeps to the concept of
production, not to the producing subjectivity. By the interpreting
function the matter goes the other way round. The concept of
production could not, as shown above, be indwelling in the production
if it were not at the same time hovering above it. The hovering
concept, subjectivity's proper meaning-concept, is decisive, and that
in a twofold sense, for the interpretation.

In the first place the producing subjectivity must necessarily
anticipate the solution of its task in objective as well as in
subjective respects; it then lies in the nature of the case that the
conditions in the objective can in many ways alter and restrict the
execution of the original plan. Here there shows itself a difference
between what is attained and what was intended, between the
production's indwelling and hovering concept, between the thought
uttered in the work and the master-workman's proper meaning. From this
standpoint the productions of nature in their immediate finitude may
even be said to stand in a certain disproportion to the divine meaning,
a disproportion which the all-capable subjectivity nevertheless
sublates again by divine interpretation.

If there is now on the master-workman's side a finite difference
between the completed work and what was originally intended, then ---
and this is the next principal point --- this difference must surely
acquire a heightened significance for the outsider who sets about
understanding the factual language of the production and reading the
master-workman's meaning out of the work: from this standpoint the
difference between the merely explicative and the interpreting
functions is unmistakable. Through the explicative functions the
contemplating subjectivity has directly to do with the relation between
the work and its concept, the regard to the master-workman is a
secondary regard; by the interpreting
% --- p. 173 ---
func\opage{173}tions the observer will, on the other hand, interpret in
the work itself the master-workman's thoughts, grasp his original
meaning; in this way he then has through the work chiefly to do with
the master-workman's subjectivity. --- Here is a turning point where
the significance of the interpreting functions for teleology must
properly strike the eye. Why is it that scientific consciousness has so
much difficulty in coming to agreement with itself concerning the
purpose of nature, the wisdom revealing itself in nature? Evidently
because two strong considerations here come into conflict with each
other: the regard to the objective, for nature is from first to last
objective, and the regard to the subjective, for teleology is
essentially subjective: means and purposes exist only for a judging
subjectivity. Those who fix their gaze exclusively upon the objective
must then consistently reject teleology's significance for the
cognition of nature; those who one-sidedly fix their attention upon the
subjective expose themselves to slipping their own individual meaning
under the productions of nature, a meaning which then in turn comes
into conflict with the objective concept indwelling in the productions.
Seen more closely, the conflict is grounded in teleology's own essence;
for teleology will on no account let go of the objective; it will only
lay it under the overreaching subjectivity. So long as the distinction
between the objective and the subjective is critically fixed: either a
subjective or an objective, teleology's problem will never be able to
be solved. Kant felt this when, instead of making teleology's principle
a constitutive principle, he made it a regulative one: the interpreting
functions shine through, but the interpretation is as yet only
psychologically subjective.

Sciences such as botany and zoology cannot refrain from applying
categories of typically teleological origin; but so long as the
ontology of subjectivity is not thoroughly cleared up, it will be
impossible to decide how far the teleological principle is
``constitutive'' or merely ``regulative''. One demands an
% --- p. 174 ---
\opage{174}objective knowledge of the subjective: this is the
contradiction; one applies indicative functions where one absolutely
must use the interpreting ones: this is the obscurity.

``If we hold fast'' --- to adduce an example from natural science ---
``to the presupposition that, with respect to the elementary as well as
to the organic wholes, a distinction must be made between the rational
and the meaningless: then the consequences fall out as of themselves.
Just as an animal's organization would necessarily become meaningless
if its organs stood in mutual conflict with one another, if it had,
e.g., teeth like a beast of prey but limbs like a hoofed animal, so
would a totality such as the solar system also become meaningless if
the planets, instead of moving in a closed line about the sun, came to
describe curves in which they removed themselves to infinity from it
\dots\ Indeed, if one --- granted, what cannot without further ado rest
upon accidental circumstances either, that the orbit absolutely had to
be a conic section --- had possibilities only for ellipse and parabola,
one might even so let chance rule; for, it goes without saying, while
the probability of a parabola is infinitely small, since among the
infinitely many possible cases only a single one can be favorable to
the parabola, the probability of the ellipse, which has an infinite
multitude of favorable cases, is so decidedly preponderant that even
the accidental would in this connection have to look like the in itself
rational. But now there is here, as is well known, not merely a choice
between ellipse and parabola, but also, what especially comes into
consideration, a choice between ellipse and hyperbola; and a hyperbola
makes the system just as meaningless as a parabola. If we therefore go
back to the beginning, the alternatives show themselves: either an
ellipse or a hyperbola, that is to say: either a rational or a
meaningless totality; here is an essential difference between the
dominion of chance and that of reason''.\footnote{Forelæsn.\ over
``Phil.\ Propæd.'' Universitetsaaret 1861--62. p.~96--101.}

% --- p. 175 ---
\opage{175}Here in the interpretation of the phenomenon it is expressly
aimed at giving the proof of the teleological as objective a character
as possible, and yet the objective in the conduct of the proof is
precisely what is insufficient. The proof is essentially subjective,
built namely upon an \emph{interpretation}. Were the proof to be
completely objective, then it would have had to be precisely the
parabola, not the ellipse, that corresponded to the in itself rational.
The proof would then come to run thus: ``The parabola is, insofar as
the orbit depends on accidental circumstances, the most improbable of
all curves; since now the parabola is nevertheless the only rational
one, it is seen from this that the rational cannot be explained by
blind natural causes, but must be referred to a divine wisdom
intervening in the course of things''.

What then does one demand, in order that the proof may be objective?
The meaningless! First a discord is to be set between the natural
causes and the rational; next the divine subjectivity is to use the
discord in order to intervene in finite fashion and bring its wisdom to
light; one demands, in other words, a discord between the indicative
and the imperative functions in the Godhead itself: for the revelation
of wisdom one demands wisdom's self-contradiction.

The scientific foundation of all natural teleology is the cognition
that the objective productions are productions of a divine
subjectivity; that the divine subjectivity in these productions utters
eternal thoughts with the stamp of divine meaning. Every object is a
word of power, every member a syllable, an element, a tittle in the
system of practical-antilogical propositions that make up the
categorical power-language. How far a human interpretation is able to
hit upon the original meaning in this categorical power-language of
actuality becomes a question for itself --- here there can in any case
only be talk of an approximation; but this is the absolute, that there
is in things a meaning, a divine meaning, that the divine
% --- p. 176 ---
\opage{176}meaning corresponds punctually to a divine interpretation,
that the interpreting functions have ontological validity for the idea
itself.

\begin{center}{\bfseries γ)\quad The Validity of the Judgment: the Modality of the Functions.}\end{center}
\addcontentsline{toc}{paragraph}{γ) The Validity of the Judgment: the Modality of the Functions}

Throughout the whole preceding treatment of the judgment, with respect
to its form as well as to its object, it has of course been constantly
presupposed that the judgment as judgment must have validity. Only the
valid judgment can logic acknowledge; the invalid it must reject. When
here, notwithstanding, there is in what follows a special question
about the judgment's validity, validity does not stand immediately
against invalidity; but validity of one kind stands against validity of
another. Insofar as one judges with full consciousness, there will in a
judgment about the object always hover before the one judging a
judgment about the judgment's validity. This judgment will, in relation
to the constitution of the object as well as of the cognition, be
different in essentially different cases; to investigate this
difference is, in other words, to investigate \emph{the Modality of the
Functions}.

That, for an all-sided illumination of thought's activity, a
distinction is made between functions of different kinds must not be
understood as if the functions of different kinds had a special
independence and an existence of their own in consciousness alongside
one another; on the contrary! they are different forms under which the
activity of thought comes forward, manners in which it brings itself to
light. That thinking must become different by the different manners of
operation is a matter of course; it is indeed, as one says, the manner
on which in all things everything depends. But that the thinking which
is of different kinds by its different manners of operation is
notwithstanding one and the same thinking shows at the same time that
the manner is not the matter itself, that essentially everything
depends on the thinking, not on the manner. The contradiction that
% --- p. 177 ---
\opage{177}the manner is at once both essential and yet unessential is
a striking proof of thinking's substantiality; for only in thinking's
substantiality does the contradiction find its solution.

In the imperative, indicative and interpreting functions, e.g.,
thinking is the same and yet not the same; that is to say: thinking
becomes in the different functions essentially different from itself.
The difference is so thoroughgoing that it looks as if each
function-system must have its own principle; and yet there is here only
one principle: thinking's essence is in principle one and the same
essence. Through an essential self-distinction thought is able to
develop and hold fast its essential unity; thinking's ideal power over
its oppositions gives the unity substantiality: the deeper the
oppositions are, the stronger is the substantial unity. Since now
thinking with its logical substantiality does not hold itself back
behind, but passes wholly over into its functions, the functions must,
despite their mutual essential difference, at the same time pass wholly
over into thinking's unity.

In relation to one another, to keep to the example, the imperative, the
indicative and the interpreting functions are different in ground and
in operation. What is accomplished by the one system cannot be
accomplished by the other; the relation in which the all-capable
subjectivity sets itself to the objects by virtue of the imperative
functions is in actuality quite another than that which is determined
by the system of the indicative ones.

But with respect to thinking's substantial unity it is so far from
being the case that the functions of the different systems should
exclude one another or stand in conflict with one another, that they on
the contrary are contained in one another, pass over into one another
in an immanent manner, and harmonize. Thus the anticipating functions
passed through the executive over into the explicative, and these again
over into the former; thus the imperative functions came to rest in the
indicative, in that the meaning conditioned by
% --- p. 178 ---
\opage{178}the unity-in-difference of both systems found its
clarification in the interpreting functions.

From the functions' reciprocal relation it is seen that their relative
independence is one with their modality: to thinking's substantiality
there corresponds the functions' modality; to the functions' modality
there corresponds in turn the judgment's validity.

In formal logic the judgment's validity is treated under ``the division
of judgments according to their modality;'' but modality itself is not
investigated. In Kant modality has a merely subjective significance; in
Hegel it is treated partly as measure, partly as mode of the
Absolute\footnote{Cf. Die objektive Logik. Erstes Buch. Dritter
Abschnitt. p.~395: \emph{Das Maaß}. Zweites Buch. Dritter Abschnitt.
p.~191: \emph{Der Modus des Absoluten}. In the 2nd Book's 3rd Division,
p.~192, Hegel expresses himself among other things thus: ``Der Modus,
die \emph{Aeußerlichkeit} des Absoluten, ist \dots\ die als
Aeußerlichkeit \emph{gesetzte} Aeußerlichkeit, eine bloße \emph{Art}
und \emph{Weise}; somit der Schein als Schein, oder die \emph{Reflexion
der Form in sich}; somit die \emph{Identität mit sich, welche das
Absolute ist}. In der That ist also erst im Modus das Absolute als
absolute Identität gesetzt; es ist nur, was es \emph{ist}, nämlich
Identität mit sich, als sich auf sich beziehende Negativität, als
\emph{Scheinen}, das \emph{als Scheinen} gesetzt ist''.} exclusively
from the objective side; it shall now show itself that modality's
significance for the judgment and its functions is both objective and
subjective.

Modality's subjective significance with respect to the judgment's
validity shines forth immediately from the distinction between
assertoric, problematic and apodictic judgments, insofar as the
distinction is determined exclusively by the corresponding certitude of
the judging subjectivity.

So long as one now keeps one-sidedly to the subjective, the assertoric,
the problematic and the apodictic can in manifold ways run into one
another: the one may for his
% --- p. 179 ---
\opage{179}part have ground to express himself problematically about
that concerning which the other straightforwardly judges assertorically,
and conversely; indeed, so obscure can subjectivity be that the one
judging imagines himself to know with apodictic certainty what he
nevertheless possesses only in the form of assertions and immediate
assurances.

It is otherwise when the matter is seen from the objective side.
Hereby the judgment's matter, now in the form of object, now in the
form of content, acquires a peculiar significance. That about which one
judges can namely in part be of such a constitution that it one and
alone is suited to an assertoric judgment, while another, according to
its nature, expressly answers to a problematic one, and another again
to an apodictic one. If it is a fault to use the assertoric judgment in
a matter that demands an apodictic one, then it is also a fault to
demand an apodictic judgment where there only can and shall be an
assertoric one.

When the ``dogmatic talent'' assures us that the truth of faith is to
him ``the most certain of all'', he utters his conviction in an
assertoric judgment, and against this nothing can in and for itself be
objected; but when the same ``talent'', from an inward urge toward
speculation, declares that the doctrine of faith seriously stands in
need of scientific grounding, that its content, in other words, is
suited to the apodictic judgment: then everything goes round for the
``talent''. For of what kind is the certitude now? If faith as faith
gives the full certitude, then the testimony is the highest, and the
believer then expresses himself in assertoric judgments. ``Verily,
verily I say unto you'', such an opening announces precisely the
assertoric. If, as concerns the certitude, faith as faith cannot be
enough for the dogmatician, although the truth of faith is to him ``the
most certain of all'', then by scientific grounding he must thus be
seeking a still higher certitude, a certitude by which the truth
cognized in faith can become to him still more certain than the most
certain of all, a seeking that is incontrovertibly meaningless.

% --- p. 180 ---
\opage{180}The meaninglessness in the alleged seeking becomes still
more conspicuous when we consider that the way from faith's assertoric
to science's apodictic judgment goes through the problematic.

When Anselm with his \textit{credo, ut intelligam} builds upon the
assertoric, Cartesius with his \textit{de omnibus dubitandum est}
expressly points to the problematic. If the one of these directions is
true, then the other must, as concerns the cognition of faith,
evidently be false. If Anselmus is right when he says: \textit{fides
præcedit intellectum}, then Cartesius must necessarily be wrong when he
says: \textit{dubitatio præcedit intellectum}; unless speculative
dogmatics will apprehend the relation between faith and doubt in the
same profound manner in which it apprehends the relation between
creation and cosmogony, so that faith is essentially faith because it
is doubt, just as creation is essentially creation because it is
cosmogony, and doubt conversely doubt because it is faith, just as
cosmogony is cosmogony because it is creation.

Let it now be with speculation as it will; we shall in what follows
show modality's objective significance, in that we shall show that the
essential difference between assertoric, problematic and apodictic
knowledge depends originally on the difference between the immediately
determined, the undetermined and that determined in reflection. To the
assertoric there correspond functions of immediate determinateness, to
the problematic functions of sublated determinateness, to the apodictic
functions of reflected determinateness.

\textit{αα)} \emph{The assertoric Judgment: Functions of immediate
Determinateness.} That there is an immediate determinateness which does
not admit of being dissolved in reflection has in our foregoing been
emphasized every time we have had occasion, from one side or another,
to treat of the antilogical. The antilogical has its ground of
existence in itself, breaks
% --- p. 181 ---
\opage{181}through by power, is because it is, and to that extent does not
admit of being grounded. The sensible qualities in their existential
immediacy do not admit of being grounded: no rational person can ask why
the red is red, the hard hard, the heavy heavy, as little as anyone can
ask why the like is like, the height high and the depth deep. There are
in all reflection reflexes of immediate determinateness, and in all
immediate determinateness something that must be apprehended
straightforwardly as it shows itself: to every such apprehension there
answers an assertoric judgment. It is characteristic that the assertoric
judgment has in its form a direction to actuality, that this direction
coincides immediately with an assurance, and the assurance with immediate
certainty: S is really (I assure you) P; the actual determinateness is
namely not merely immediate, but in its immediacy antilogical, and no
antilogical determinateness can give any other assurance than the one
that is expressed in the assertoric judgment.

The assertoric judgment is logical-antilogical in a peculiar sense. Its
form is logical, its content both logical and antilogical; in actuality
it always gives us an antilogical content in logical form, but the
logical form gives the content in turn a logical stamp. ``The stone is
really gray''; this judgment seems at first glance to be logical in
content as well as in form, but if we look more closely, we easily
discover that the logical content is only a determinateness reflected in
the form: the actual content is in its actuality an existing object with
an existing property, i.e. with a property inhering in existence. The
predicate ``gray'' designates here not at all the general, but on the
contrary the singular. A general judgment such as: ``the stone is gray'',
can be an expression of the fact that the thing, the singular, is
subsumed under a property (the general); but the assertoric judgment:
``the stone is really gray'', does not fit such designations. He who
says: ``the singular is really the universal'', is thoughtless; for in
actuality the
% --- p. 182 ---
\opage{182}actual is never the universal, and that for the good reason
that the actual as actual is always an immediately determinate, and the
immediately determinate is not the universal. The assertoric judgment:
``the stone is really gray'', does not aim at gray in general, but at the
gray in this immediate determinateness, the actual gray, the gray stone.

Since now the content of the assertoric judgment must always be of a
two-sided character, \emph{antilogical} by reason of the object,
\emph{logical} by reason of the form, it follows of itself that the
development of the judgment's validity is laid out with particular regard
to each of the principal sides and to their mutual relation. Thus we come
to a separate treatment of the antilogical or of the objective side of
modality, of the logical or of the subjective side of modality, and
finally of the logical-antilogical essentiality of the content, whereby
the validity of modality is determined on all sides.

1) \emph{The antilogical Element in the assertoric Judgment: the
objective Side of Modality.} The assertoric judgment comprehends a sphere
of its own, a sphere of infinite extent, namely the whole of actuality
from the standpoint of the power-determinations, the antilogical
elements. Actuality's world has indeed its essentialities, its grounds,
its possibilities, laws, species, genera, etc., categories that in the
last instance withdraw from the immediacy of assertoric knowledge; but
here in this connection the express talk is not of essentialities in the
actual, but of actuality's actual facts in their given diversity and
infinite multiplicity. It must indeed be conceded that the various
determinations of actuality with their antilogical immediacy incessantly
alternate, are altered and transformed; the present vanishes, the
existent is altered, that which exists goes to ground, etc.; but in all
these alternations the one antilogical element is relieved point for
point by the other. The present that vanishes gives place only to another
% --- p. 183 ---
\opage{183}present; every present is now really \emph{thus}, not
otherwise; it determines the assertoric judgment and is determined in
turn by this. That which exists passes away, it goes to ground; but the
ground forms on that account no chasm between that which exists and
passes away and that which exists and arises; at every point in the
actual transition that which exists follows upon that which exists,
everywhere \emph{this}! and \emph{this}! and \emph{this}! Here is a
stream of antilogical elements that knows no interruption: there are in
the wall of actuality no cracks, no fissures.

If truth were immediately one with actuality, and actuality immediately
one with the totality of the antilogical determinations, then the
assertoric judgment would be truth's logically adequate expression; now,
on the contrary, since the determination of actuality is only a moment,
even if a real moment, in the actually true, the assertoric judgment can
indicate only one side of truth, not truth itself, the whole truth. It is
an error when the immediate determination of actuality is without further
ado made one with truth; but it is also an error when the determination
of actuality is volatilized in its immediacy and transformed into
semblance because it is not truth itself. The relation between
sense-perception and thinking will in this connection be able to be
illuminated from a new standpoint.

The Sophist Protagoras started, as is well known, from the principal
proposition that ``man is the measure of all things, both of those that
are, that they are, and of those that are not, that they are not''. It is
a peculiar relation between sense-perception and thinking, between
judgment and object, that thereby comes to light. That Protagoras starts
from the sensible and originally lays emphasis on sense-sensation is
certain; just as it is also unmistakable that he sophistically disparages
thinking and undermines all cognition of truth. But it is also certain
that Plato, however much he asserts the right of thinking,
% --- p. 184 ---
\opage{184}has yet not quite thought himself into that which here lies at
the ground of the sophistic way of viewing.

``Were man'', it is said in the Theaetetus, ``the measure of things
merely as a sense-perceiving individual, then an animal just as fully as
the wisest man, indeed as a god, could cognize the true (ὅτι πάντων
χρημάτων μέτρον ἐστὶν ὗς ἢ κυνοκέφαλος ἢ τι ἄλλο ἀτοπώτερον τῶν ἐχόντων
αἴσθησιν).''

In this objection two very essential points are overlooked. In the first
place the relation between the sensible and the actual, and therewith
between the sensible and the true, is overlooked. In that Plato misjudges
the sensible, he must, as was shown in the first
chapter\footnote{Grundideernes Logik. I. p.~182--183.}, necessarily
misjudge actuality, and, in that he misjudges actuality, misjudge truth.
In the second place the relation between judgment and object, and thereby
the relation between a human and an animal sense-perception, is
overlooked. In that Protagoras makes man, not the pig, not the werewolf,
not the frog, the measure of things, he nevertheless apprehends, be it
with clear or unclear consciousness, man expressly as judging
subjectivity. When the one says: ``the wind is really cold'', the other:
``the wind is really warm'', then it is two assertoric judgments that
bring forth the contradiction and thereby call forth the question about
the true. The assertoric judgment is in this way, by reason of the
sense-impression, set in no contingent but in an essential relation to
the question. By merely setting the pig against the frog, the bird
against the fish, the question of truth cannot arise at all: the question
arises expressly through the contradiction, the contradiction through the
assertoric judgment; but the pig, the frog, the bird, the fish are only
sense-perceiving, not judging subjectivities.

Only when the relation between actuality and the
% --- p. 185 ---
\opage{185}assertoric judgment is apprehended clearly can Protagoras's
sophistry be fully illuminated. It is the reciprocal determination of
object and judgment upon which the emphasis falls. In that a beginning is
here made with the sensible, a beginning is here made immediately with
the actual: ``the wind is cold, it is really cold''; in such a statement
the judging subjectivity does not think at all of prescribing to the wind
whether it is to be cold or warm, whether it is to \emph{be} at all or
not. The judging one will not act upon the wind, but cognizes himself as
affected by the wind; it is, in other words, not the judgment that
determines the object, but conversely the object that immediately
determines the judgment: the judgment is correctly assertoric. Wherein
then lies the sophistical? In a fraudulent use of the judgment, for the
judgment conceals a contradiction. The judging one is affected by the
actual object; but he is affected only in the manner that exactly answers
to his individuality: in that he will judge assertorically about what
affects him, he comes involuntarily to judge about the affection he
receives, thereby about his own subjective state, thus literally about
himself. It is just the characteristic thing about the assertoric
judgment that from being immediately objective it suddenly becomes
immediately subjective, that from being a judgment about the object it
transforms itself into a judgment about the judging one, the judging
one's judgment about himself.

The Protagorean sophistry has excellently laid the contradiction bare,
but not resolved it. Instead of cognizing the actual as a moment in the
true and seeking the true in the law, the Sophist remains standing at
actuality's antilogical punctuality, and now uses the contradiction, in
which individuals through assertoric judgments bring their contingent
subjectivity to light, for a game with truth.

From another side and in another manner the relation between actuality
and the assertoric judgment has been misjudged by the Epicureans. In
Epicurus's Canonic too the
% --- p. 186 ---
\opage{186}criterion of truth is sensation or sense-perception
(αἴσθησις). ``A motion is produced in us, which is sense-perception, and
where there is a motion, there must be something moving. Therefore all
sense-perceptions are true; for things really move us just as we perceive
them. The visions of the madman and of the dreamer are also true; for
such images are really conveyed to them. Orestes really saw the
Eumenides, but he erred in this, that he took them for palpable
objects''.\footnote{Efterladte Skrifter af Poul Møller. 4de Bd. 3die Udg.
Kjøbenhavn 1856. p.~256.}

But despite the apparent agreement between Epicurus and Protagoras as
regards the point of departure, Epicurus has nevertheless noticeably gone
a step further. The advance is marked by the fact that weight is here
laid upon the difference between effect and cause, between the objective
and the subjective in the appearance. Had the Eumenides that Orestes saw
really been palpable, then they would also have had actuality for others
than Orestes; they would, in other words, have had objective actuality.
Already from this it is discerned what conditions actuality itself offers
for the closer determination of the objectively valid. Instead of rushing
overhastily into syllogisms from what is perceived to what is not
perceived, the observer is called upon to let the one kind of
sense-perception support the other, until with the united testimony of
the different senses he has so many marks of the objectivity of what is
perceived that he may with full assurance dare to say: what is perceived
is really an existent object.

The manner --- for with modality everything depends upon the manner ---
in which heterogeneous sense-perceptions correspond in the apprehension
of one and the same objective unity speaks in actuality for the unity's
substantial being, its independence of the sense-perceiving subjectivity;
this
% --- p. 187 ---
\opage{187}testimony of actuality is confirmed still more closely by the
fact that several perceivers observe the same sensible under the same
conditions in the same manner. If now in addition the existent object's
relations to other existent objects, thus the mutual reciprocal action of
things, are made the object of immediate observation, then actuality has
furnished the contemplating subjectivity all the contributions to the
establishment of the objective and the true that it is in a position to
furnish through the totality of immediate impressions. But however these
contributions may for the rest be apprehended and combined, it is
nevertheless evident, when we lay due weight upon the antilogical
elements, that the logical ground on which the assurance rests is a
totality of assertoric judgments. The assertoric is recognized precisely
by this, that the one judgment is just as original as the other; here in
this respect there is no transition from assertoric to apodictic.

The striking immediacy with which actual impressions of actual objects
determine the assertoric judgment explains the certainty with which the
judging subjectivity, despite all Humean reasonings, orients itself in
the actual world. If we here let the emphasis fall upon the objectivity
of things, their objective differences and objective relations, then we
have the standpoint of empiricism. The empirical judgment is essentially
an impression-judgment, the impression-judgment essentially assertoric.
But the impression-judgments, bound to the antilogical and secured by the
facts, are at the same time hovered over by a reflection that frees
itself from the antilogical punctuality and in the logical sense demands
the universal: the dualism that thereby forces itself upon consciousness
lets us see Hume's misunderstanding from a new side.

All true concepts must, this is Hume's fundamental thought, be concepts
of actuality, thus original impression-concepts. But whence then come the
untrue concepts? From the distortion that the impression-concepts undergo
in the
% --- p. 188 ---
\opage{188}following reflection! First, then, is the impression, then
follows reflection; first the sense-impression, then the judgment; first
the assertoric judgment, then the universal-logical one. Here is an
evident confusion of what belongs to principle with what belongs to the
phenomenal. As regards the phenomenal, Hume may psychologically be right;
but in principle he is wrong. According to the phenomenon it must indeed
seem as if the sense-impression already stood clear before consciousness
before any judgment formed itself; but in principle it is, as we have
seen above in Protagoras, just as originally the judgment that determines
the impression as the impression that determines the judgment. According
to the phenomenon it should indeed seem as if the assertoric judgment,
the judgment about this singular, determinate object in this singular
determinate case, preceded the universal-logical judgment: the judgment
about what holds for all possible cases; but in principle it is just as
originally the universal-logical judgment, the judgment as judgment, that
must determine the assertoric one, as conversely. If, then, an error is
made in the judgment about this substance with these accidents, the error
is not, as Hume maintains, that the categories substance and accident are
laid at the ground in the judging of actual objects; this error does not
in any case concern the assertoric judgment. What here concerns the
assertoric judgment is solely and alone the application of the general
categories to the singular case. If it is untrue that A in a given case
is cause, B effect, then the error is not, as Hume maintains, that the
categories cause and effect are laid at the ground in the judging of
something actual, but the error is that the assertoric judgment has in
the given case made a false application of the categories of causality;
whether it has for the rest applied them to a relation in which they do
not fit, or has confused the actual cause with something else, e.g. with
conditions, attendant circumstances, etc.

According to Hume it should be the universal-logical judgment that had
the assertoric one for its ground; but it is conversely
% --- p. 189 ---
\opage{189}the assertoric judgment that has the universal-logical one for
its ground. According to Hume the categories of actuality should make
good their reality by excluding the category of cause; but the matter is
just the opposite: it must, even according to Hume's own method, be
precisely the category of cause that lies at the ground of the categories
of actuality. What then is the infallible mark of a true category of
actuality? That it mediately or immediately has its origin from a
corresponding impression! With this it then agrees that the assertoric
judgment is laid at the ground of the universal-logical one, for the
assertoric judgment with its antilogical element stamps precisely every
concept of actuality as an impression-concept. But hereby there is
presupposed precisely a relation between impression and concept which, so
far as the concept is concerned, must be recognized as a relation of
dependence as necessary as it is actual: the impression is related to the
concept as cause to effect. If we think the causal relation away, the
concept stands in a contingent relation to the impression; a determinate
concept of actuality can just as well be thought to follow from the one
impression as from the other; the impression that in one moment awakens
the representation yellow can in another moment awaken the representation
blue, and the impression that the one time gives a color can the other
time give a tone: where then is the standard that is to display the
difference between true and false concepts?

What is true in Hume's view is the moment characteristic of empiricism
and of the assertoric judgment, namely that every existence in every case
individualizes and precisely thereby modifies the concept, so that every
singular phenomenon in every moment must necessarily demand its immediate
apprehension, its separate judging, if the cognition of the actual is to
be quite adequate. The question of an adequate cognition of actuality
stands then in the very closest connection with the logical treatment of
the problem of omniscience. In relation to actuality's antilogical
punctuality omniscience must
% --- p. 190 ---
\opage{190}be a collective knowledge, a particular immediate knowledge
about all existent things in all the singularities of existence: ``Yea,
even all the hairs of your head are numbered''; these words logic must in
this connection acknowledge literally.

``But'', someone might object, ``an omniscience that counts the hairs of
the head is surely just as incomprehensible to us as an omnipotence that
by its mere word creates the world out of nothing; if science can
acknowledge the religious expression for omniscience, why should it not
just as well acknowledge the religious expression for omnipotence? If
revelation's word about omnipotence is to be rejected because it is
incomprehensible, then the same must surely hold when the talk is of
omniscience''. We answer: The Logic of the Fundamental Ideas everywhere
makes a distinction between what is for us unfathomable and what is in
the idea incomprehensible. An omnipotence that unconditionally creates
the world out of nothing is not merely unfathomable for us, but
incomprehensible in the idea; an omniscience, on the contrary, that
immediately holds of the singular is comprehensible in the idea, for its
concept is a necessary consequence of the articulated reciprocal
determination of knowledge and actuality.

Already at an earlier stage\footnote{Grundideernes Logik. I.
p.~257--263.} the relation between the category of omniscience and the
ideal of knowledge has been particularly illuminated with regard to the
ontology of sense-perception\footnote{Cf.\ Grundideernes Logik I.
p.~423--427.}. In the present connection the illumination is heightened
by the fact that the infinite multiplicity of sense-perceptions and
sense-determinations receives its ideal form in a corresponding
multiplicity of assertoric judgments. The assertoric functions are then
also, as regards the objective side of modality, originally akin to the
explicative and the indicative, and appear as modifications of both.

The assertoric function modifies the explicative in that, with emphasis
on the given, it lets the taut
% --- p. 191 ---
\opage{191}opposition of concept and existence, upon which explication
supports itself, disappear in immediate unity. If we let the judgment:
``this is really a branch'', receive its significance from the explicative
functions, then it begins, as we have seen, with the concept's furnishing
the standard for existence, and ends with the given existence's being
acknowledged as the representative of the realized concept: in this
existing thing the concept branch is realized. Such a mediacy, such an
opposing and mediating is not noticed on the assertoric standpoint; the
concept is indeed recognized in the existing thing, but the concept is
not weighed as concept, the reflection is like a glance: the recognition
is quite immediate.

The assertoric function modifies the indicative in that, with emphasis on
the given, it makes the logical identity of concept and existence
recognizable by intimating the opposition. From the indicative standpoint
the actual existences, seen in the light of reflection, answer to the
eternal categories: the satisfaction is that the logical clarifies the
antilogical and is enriched by antilogical reflexes. The complete
equilibrium between the logical and the antilogical, which is the result
and the presupposition of the indicative functions, would bring the
heterogeneity of concept and existence into forgetfulness if the
assertoric function did not, by its antilogical element, give the
immediately actual a preponderance.

2) \emph{The logical Moment in the assertoric Judgment: the subjective
Side of Modality.} Vibrating between the logical and the antilogical, the
assertoric judgment will constantly give the antilogical element the
preponderance so long as the object in question is outwardly objective.
If, on the contrary, the judgment's object is inward, and in its
inwardness thus of subjective nature, the relation will reverse itself:
the logical moment will become the preponderant one; but this
preponderance will then bring it about that the assertoric acquires an
ambiguous character. The inner is, as inner, with all its immediacy still
always a
% --- p. 192 ---
\opage{192}reflected. Pain is immediate and yet in reflection;
enjoyment, fear, hope, joy, etc. likewise. Here, with these
determinations of inwardness, very different combinations of immediacy
and reflection can be given, but in all of them immediacy is reflected,
reflection immediate. He who says: ``I am really glad'', utters himself
in an assertoric judgment; his state of mind is given, he has an
immediate assurance of it, and more than an immediate assurance of his
momentary state of mind no one can have.

The inner has its outer and reflects itself therein. As reflection,
reflection is always logical; but the outer stands over against
reflection with antilogical determinateness; in the clash it depends upon
which party in actuality gains the upper hand, either immediacy over
reflection or, conversely, the latter over the former.

If immediacy gains power over reflection, then the inner will stand
straightforwardly painted in the outer, will be recognizable by the
utterance. In a judgment such as: ``the man is really sorrowful'', the
judging one's immediate, though reflected, assurance answers to the
immediate, though reflected, appearance of sorrow in the sorrowful one's
outer. If reflection will weaken the assurance, then immediacy reacts at
once with heightened emphasis: ``it does not help, what you say, he is
really sorrowful''.

Under the tension between reflection and immediacy, a tension which is
here conditioned by actuality's own doubling into inner and outer, there
is seen a fusion of assertoric and interpreting functions. When the
judgment's object is not the judging one's own but another's inner, then
the judgment rests upon an interpretation of the utterances through which
the two subjectivities correspond. ``The man is really angry, I assure
you, I have indeed heard it from his own mouth; yes, he is thoroughly
angry, I can see it from his mien''. That such judgments expressly
% --- p. 193 ---
\opage{193}presuppose a determinate interpretation of the speaker's words
and mien is noticed involuntarily when someone comes forward with the
opposite assurance: ``the man is really not angry, he means nothing by
it''. Insofar as reflection is set in movement, it seems indeed as if the
judging ones must grasp at proofs and grounds; but on account of the
prevailing immediacy the grounding is a semblance; the grounding consists
in the one assertoric judgment's being exchanged for the other. ``The
man'', whose anger is disputed, ``has really meant what he said''; ``no,
he has really not meant what he said.'' ``His mien really betrays
anger''; ``no, his mien really betrays no anger''; ``it is really
earnest!'' ``it is really all a jest!'' That the assertoric function,
however, does not straightforwardly coincide with the interpreting one,
but becomes an alteration of it, a subjective modality, one of the
manners in which the interpreting function can utter itself, follows
straightforwardly from the fact that the grounding which interpretation
essentially demands essentially fails to appear so far as the assertoric
function is concerned. On the stage of immediate actuality no essential
grounding can here be given and none is to be given: the assertoric
interpretation is a straightforward translation of actuality's own
language, and actuality's language is a categorical power-language.

How striking, how unambiguous this categorical power-language is can be
seen everywhere where sound common sense is proved right against
contorted school-theories. According to the Cartesian doctrine of the
relation between soul and body the animal could, as is well known, have
no soul: every animal was to be regarded as an automaton. When one struck
a dog and it whined, its whining was not at all any utterance of pain; it
was mechanically a sound, as when one blows on a flute. It is thus,
according to this way of viewing, an absurdity to issue a prohibition
against cruelty to animals; one might just as well issue a prohibition
against cruelty to locomotives. And yet cruelty to animals is held
% --- p. 194 ---
\opage{194}to be something ignoble and punishable: has Cartesius then
really been refuted? Yes, but not theoretically (the theoretical
refutations do not namely come into consideration here), but practically,
that is to say: by a power-language; not by an arbitrary one, no, by a
categorical one, actuality's categorical power-language. In that the
language of facts is categorical, the immediate understanding must be
assertoric. ``The dog really has a soul, it shines out of its eyes; the
poor animal can really feel pain''. These are only assertoric judgments,
but judgments that, so to speak, strike speculation to the ground with
doomsday-blows.

If, conversely, reflection gains the upper hand over immediacy, then the
subjective side of interpretation becomes predominant; the assertoric
interpretation gives, on this standpoint, what one understands in the
ordinary sense of the word by an \emph{opinion}.

% Errata: p. 194 — printed "høist", corrected to "høiest"; translated as corrected.

That the category opinion hovers in a peculiar manner between subjective
and objective, between reflection and immediacy, can easily be shown. The
immediate apprehension of the immediately objective is not yet an
opinion; for there is surely no one who will call the assumption that
water is fluid, or that the sun shines, an opinion. If someone were to
point at the sun when it stood highest in the heavens and say: ``see, it
is now my opinion that the sun is shining'', then one would take it
either for an utterance of caprice or for a sign of affectation. On the
other side, a thoroughgoing reflection of essential relations likewise
does not lead to opinion, but to that which is more than opinion: it
leads to knowledge. That semblance must presuppose essence, that the
ground must have its consequent, the cause its effect, is with the
judging subjectivity not an opinion but a knowledge. Where reflection, on
the contrary, is set in movement in such a way that it neither denies
immediacy nor unconditionally bows under it, there is the proper place
for opinion, there it has its proper scope.
% --- p. 195 ---
\opage{195}If we set the two assertoric judgments: ``the man is really
angry''; ``no, the man is really not angry'', over against each other,
then it is, properly speaking, not a contradiction between two concepts
but a strife between two opinions that we have before us.

It is the immediacy of an actuality reflected into itself that
constitutes the objective ground of opinions. Such actualities reflected
into themselves we have in state, art and religion. On the stage of
piety, naivety and simplicity the judgment about these actualities is
still immediately assertoric. ``What the authorities command is really
right''; ``what the master has there produced is really beautiful'';
``what the priest says is really true''. In these judgments the power of
the immediate is so preponderant that subjective reflection is held
bound: it is dominating prejudices, secure propositions of authority,
fixed dogmas that we here find before us; but they are not yet opinions
proper. That liberation of subjectivities, that mobility in reflection,
that facility in hovering considerations which favors the emergence of
opinions in society, just as the warmth of June favors the emergence of
flies, sets in only at a later stage.

% Errata: p. 195 — printed "belyste", corrected to "belyse"; translated as corrected.

It is the stage of culture, the stage of the subjective reflections, that
is properly to be apprehended as the stage of opinions. ``Now what the
authorities command is really not always right''; ``nor is everything
that an artist produces really beautiful''; ``one really cannot believe
everything the priest says either''. What then is right? what then is
beautiful? what then is one to believe? Yes, it is precisely about this
that there are so many different opinions. While the interpreting
functions, in virtue of the grounding, illuminate the objective meaning
of things, the assertoric ones on the contrary turn attention --- not
toward the objective, toward the meaning in things --- but toward the
subjective, toward individuals' opinions about things: we get a host of
political, aesthetic and religious opinions; every opinion is an
assertoric judgment, determined by subjective reflection.
% --- p. 196 ---
\opage{196}In opinion the objective and the subjective elements are in a
peculiar way set free and yet bound together. Subjectivity is emancipated;
for what is easier for the singular than to form an opinion for himself,
and what is the singular's spiritual property, if it is not his
``unpresuming opinion'': every individual, even a schoolboy, has at this
stage an unpresuming opinion. And yet the opinion points objectively
straight at the thing; it is the opinion about the thing that decides what
the thing is. The thing is the actual, and the opiner lays weight upon the
actual, upon the matter; it is of course not about his individual opinion,
but about the matter, the important matter, that he is concerned. Could
anyone convince him that his opinion was incorrect, ``he would truly, for
the sake of the matter's importance, gladly change it''; but that is
impossible, for, since it is now once for all his unpresuming opinion that
the matter really stands as he opines it to stand, it would indeed be
treachery against the good cause if he changed his opinion.

Opinion is an assertoric judgment; when the assertoric is strongly
stressed, opinion rises to assertion; in the assertion subjectivity will
then expressly vouch for the objective: the opiner becomes assertive. In
that the opiner becomes assertive, it looks undeniably as though he wanted
to force his contingent subjectivity upon others; but that is not the
meaning of it. The meaning is that the assertive one is zealous for the
matter, is the only one who sees the matter as it really is. The assertive
one only opines that his opinion ought unconditionally to be
\emph{everyone's} opinion, and that for the simple reason that it is the
only correct opinion. Opiners are liberal, for everyone surely has the
right to keep his opinion; but when the many opinions run against one
another, and each holds to his own because it is the one correct one, the
many unequal opinions must nevertheless fall at last into a serious
conflict: the task is then to find the opinion valid for all, the gene%
% --- p. 197 ---
\opage{197}ral opinion, more precisely the public opinion, \emph{Opinion}.

It is of logical interest to characterize the struggles of opinions, the
struggles for Opinion, because these struggles, seen at a distance, look
as though truth were at stake, and yet in principle they are indifferent
to truth. It is not truth, but actuality, not actuality as it utters
itself in the categorical power-language of facts, but actuality as it,
subjected to the interests, veiled in the cobwebs of reflections and
colored by the opinions, addresses subjectivities, that the strife is
about.

The manner in which the strife is conducted is in many respects
instructive. Here there is investigating, here deliberating, here grounds
and counter-grounds are adduced; the whole armory of logic, with
judgments, syllogisms, proofs, etc., is ransacked in order to procure
weapons for the combatants: one reasons, discusses, argues; one crosses
swords with the other. But, if we look more closely, the strife is
fleeting; reflection plays upon the surface of actuality, often
brilliantly, like sunbeams upon a rippled sheet of water, but the ground
is shallow, absorption is wanting. The characteristic thing is that
however much is said here about principles and questions of principle, it
nevertheless never comes, and in the sharply bounded sphere of opinions,
assertions, assurances, assertoric judgments never can come, to any real
investigation of principles. Reflection has a dexterity in moving as
though it went to the bottom of the thing and reached the ground; but this
bottom is no essential ground, it is a crust of crystallized opinions.
Reflection has the facility of stacking the assertoric judgments in such a
way that to the uninitiated it looks like batteries of striking arguments,
whole fortifications of irrefutable proofs. It is semblance, but a
semblance of reflection, which corresponds to the immediate in the essence
of actuality.

In every strife, let it move otherwise upon the precincts of politics,
% --- p. 198 ---
\opage{198}of art or of religion, it makes an essential difference whether
what is striven for is really a truth or an Opinion. This can never be
heard directly from the battle-cry, for the battle-cry of course always
sounds upon truth; but it can be seen from the method, it can be known by
the manner in which the strife is conducted. For a method that must be
rejected when truth is at stake may be precisely the one that recommends
itself as the only right one when Opinion is at stake. We can illustrate
this by a couple of examples.

If the strife is political, and one of the standard-bearers comes forward
with the declaration that he, as regards the opposing party's standpoint,
would make it a matter of conscience to weary the hearers' patience with
adducing proof for what all who are gifted with sound sense have long
since agreed upon, that it belongs among the spurious standpoints, and
will therefore confine himself to a demonstration of facts that are so
conspicuous, so eloquent, etc. --- then we hear at once that it is not
truth, but Opinion, that he is concerned about. Were the spokesman
concerned about truth, it would have to be of the greatest importance to
him that the difference of principle between the standpoints should become
properly clear to the hearers, for the facts appealed to will of course
look different according as the light falls in from the one side or from
the other. In a struggle for Opinion, on the contrary, the speaker's
procedure is practically---tactically perfectly correct. It is an
assertion, an assertoric judgment, that the opposing party's standpoint
belongs among the spurious standpoints; an assertion, an assertoric
judgment, which rests upon the assertion that all who have sound sense can
see it without further ado. This assertion, this assertoric judgment,
rests in turn, if anyone among the hearers should get scruples and opine
that one might after all also be of another opinion, upon the assertion,
the assertoric judgment, that he who cannot see that the opposing party's
standpoint is a spurious one really cannot be said to have
% --- p. 199 ---
\opage{199}sound sense. This assertion, this assertoric judgment, rests
then in the last resort upon the assertion that an individual who in
actuality does not have sound sense is not entitled to have a say about
political standpoints. With such a substructure of assertions, upon such a
timberwork of assertoric judgments, the speaker is sure of himself, sure
of his standpoint, sure of his strife: he strives for Opinion.

If the strife is aesthetic, the method must likewise vary according as the
aim of the strife is a truth or an Opinion. If the aim is a truth, the way
can at times be long and troublesome enough. For an assessment of the
works of art it is of course not enough to have natural taste and to have
visited a number of museums; what matters in the assessment, if it is to
have scientific worth and to serve truth, is above all that the judgment
be motivated. In order to be able to motivate his judgment, the one who
judges must before all else be sure of the relation between concept and
phenomenon, between the universal and the individual; but to that there
belongs a study of the concept in science as careful and as penetrating as
that of the phenomenon in art. It is otherwise when the strife concerns an
Opinion. Here a standard-bearer and trumpeter can spare himself the
trouble of science's analyses. ``What indeed comes of these dry
abstractions, these colorless philosophical concepts? No, phenomenon and
feeling and feeling for the phenomenon, that is and remains after all the
main thing.'' Since now the finer feeling is developed only through
contact with the manifold phenomena of art, while the eye of the
understanding is opened by reading descriptions and composing descriptions
of the many objects of art, kinds of art, artists and judges of art from
the various ages, the standard-bearing art critic must, for Opinion's
sake, so arrange his performance that the public can both have confidence
in his pure taste, his fine sense, and at the same time be impressed by
his great learning and detailed descriptions. In such an apparatus, an
apparatus which a quick head
% --- p. 200 ---
\opage{200}will be able to procure for itself, especially on a journey
abroad, in the course of a few months, Opinion's standard-bearer has
everything that he aesthetically needs for opinions, assertions,
assertoric judgments.

Yet here, for the method's and Opinion's sake, there is still one thing to
observe. An apparatus of opinions, assertions and assertoric judgments may
be as large and imposing as it will, so long as one keeps to the
phenomenon: it soon shows itself insufficient when one in earnest engages
with something scientific. The method here demands a certain caution; for
the method is that one, without directly engaging with science, preserves
a certain scientific tinge, the tinge that suits precisely the assertoric
judgment. When the assertoric judgment, instead of being set up in all
crudity as an assertion, acquires through a tincture of ``feeling'' and
learning a scientific tinge, it transforms itself and advances from simple
opinion to fine remark: let one work above all by fine remarks!

That a remark can have a scientific tinge, indeed can become almost of
full scientific weight, may then be seen from the esteem that Opinion
grants, not of course to the philosophical systems, but to the
philosophical remarks. Yes, Opinion has really a taste for philosophical
remarks. No wonder that, for Opinion's need, philosophies are delivered by
the ell, philosophies that are woven together out of nothing but remarks.
In such remark-philosophies there may then in many things be much to
become aware of and much to remark, for the becoming-aware goes upon
immediacy, and the remark runs precisely along the boundary, the
assertoric boundary, between opining and knowledge. If now
remark-philosophy notoriously enjoys great esteem in public opinion, why
should a remark-aesthetics, an art criticism of fine remarks, especially
when they are exhibition-remarks, not enjoy the same esteem?

If the strife is theological, then one cannot sufficiently praise,
% --- p. 201 ---
\opage{201}recommend and admire the Opinion-method. Here, e.g., is a
writing that is rightly regarded as one of the most remarkable products of
spirit in our century. It is a Christian philosophical writing,
charmingly written toward a reconciliation of faith and knowledge. As
eloquent proof of the writing's deep philosophical spirit there serves the
fact that, although it is rejected by those who really understand
philosophy, it has nevertheless, in spite of the severe criticism, won
immense circulation and scientific recognition among all those who do not
understand philosophy; for its philosophy is theological. ``We dare,''
says the reviewer, ``freely assert that this writing is in the hands of
all who unite with higher culture a deeper sense for true Christianity. We
can prove our assertion; for we assert that if there really are cultivated
persons in whose hands the writing is not, then these cultivated persons
have no sense for true Christianity. We can prove our assertion, for we
assert that if there really should be true Christians in whose hands the
writing is not, then these Christians are without higher culture.'' Such
assertions are reprehensible, indeed contemptible, when the strife is for
a truth; but as assertoric judgments with reinforced emphasis they produce
an excellent effect when the strife is for an Opinion.

3) \emph{The assertoric judgment's logical-antilogical essentiality:
double-sided modality.} To articulate a unity-in-difference of the
antilogical with the logical is by no means an exclusive distinguishing
mark of the assertoric functions; for the same can in a certain sense be
said to hold of all functions. No, what is peculiar to the assertoric
functions is precisely this relation between actuality and reflection,
that the actual in its antilogical immediacy must either govern thought by
a categorical power-language or, insofar as thought takes on the stamp of
subjective independence, then by goading opposition bring about such a
rising of reflection that the accentuating interplay of logical and
antilogical determinations of form, which conditions real knowledge,
% --- p. 202 ---
\opage{202}becomes visible from all sides; only thus can modality's
objective as well as subjective validity be fully illuminated.

In three ways will actuality be able to enter into ideal process with
reflection, according as it either lets itself be transformed into a
phenomenon immediately transparent for its individualized content, and
thereby into the idea's mirroring; or, as regards its character, indeed
retains its finite reality, but undergoes a transformation in that it is
apprehended not as something present but as something past, and thus damps
down the antilogical edges; or, with the stamp of a divine reality that is
raised above the finite difference between present and past, forces itself
upon subjectivity as an actuality of a higher order, a superactual
actuality, an actuality that makes a breach with the actual. In the first
case the assertoric function is \emph{aesthetic}, in the second case
\emph{historical}, in the third case \emph{religious}.

The \emph{aesthetic} unity of essence and phenomenon is assertoric; herein
lies this, that the difference to which the unity corresponds is likewise
assertoric. If we consider the phenomenon, the form is antilogically
determined as the actual itself: this Jason in marble has, e.g., just as
immediate an individuality as this workhouse inmate in flesh and blood;
but ``Jason is really beautiful,'' ``the poor wretched man is really not
beautiful'': the one assertoric judgment has here the same validity as the
other.

But upon what does the validity now rest? That both figures really are,
the one as statue, the other as existing individual, places them both, as
regards actuality's antilogical immediacy, in a certain sense upon an
equal footing: the assertoric judgment that apprehends the immediately
actual shows us modality's objective side. The difference between the
beautiful and the unbeautiful is thus an actuality-difference; it is the
actuality-form itself that turns the scale. The actuality-form is in its
objective immediacy independent of the
% --- p. 203 ---
\opage{203}considering subjectivity; the really beautiful is thus really
objective: beauty is in the beautiful one with the actual.

But unbeauty is in the unbeautiful also one with the actual. Thus: the
actual is one with the beautiful, and the actual is one with the
unbeautiful; thus: the actual is different from the unbeautiful, for it is
one with the beautiful; the actual is different from the beautiful, for it
is one with the unbeautiful; thus: the actual is, as actual, different
both from the beautiful and from the unbeautiful, can just as well be
unbeautiful as beautiful, stands in an external and indifferent relation
to beauty. What then becomes of beauty? It transforms itself into idea, is
taken up into imagination and reflects itself in subjectivity. It is thus
the beholder, not the beheld object, that itself decides what is beautiful
and unbeautiful. From modality's subjective side man is again ``the
measure of things,'' both of the beautiful ones, that they are beautiful,
and of the unbeautiful ones, that they are unbeautiful.

In the aesthetic domain the assertoric function's original kinship with
the explicative and the interpreting ones is quite conspicuous. From
reflection's, hence from subjectivity's logical standpoint the task is to
recognize the concept realized in the object; the function is to that
extent explicative. But in the aesthetic sense the recognition must be
immediate. In the immediate recognition the form goes into one with its
content; the concept vanishes as concept and transforms itself for
intuition into prototype: to recognize the concept in the object is to
measure the object by the prototype. If now the immediate appraisal says
that the object corresponds to the prototype, the judgment falls that it
is really beautiful; if the object deviates notably from the prototype,
the judgment falls that it is really unbeautiful; if the object stands in
striking conflict with the prototype, the judgment runs that it is ugly.
On account of the form's antilogical determinateness and the recogni%
% --- p. 204 ---
\opage{204}tion's immediacy the aesthetic judgments are originally
assertoric.

The assertoric character of the aesthetic judgments makes it explicable
that about the beauty-phenomena and about the beautiful there can prevail
an agreement as striking as the disagreement.

The assertoric judgment is immediately subjective; the aesthetic judgments
are judgments of taste: each of those who judge has his taste; what does
it help to want to dispute about the difference between a good and a bad
taste, when the judgment about taste is itself a judgment of taste? And
yet there is really agreement here; people of good taste join together,
understand one another and support one another in combating the bad: that
is a striking agreement. The striking agreement must be ascribed to a
categorical power-language; but this time it is not actuality in its
contingency, but actuality impregnated by the idea, the idea in
actuality's antilogical form, that speaks to subjectivity.

The assertoric judgment is immediately objective. The aesthetic judgments
are judgments of intuition; it is not the judgment that determines the
object, but the intuited object that determines the judgment. Intuition's
object is common to the many who intuit: ``see, this figure, how beautiful
it is! these features, these lines, how lovely they are!'' --- ``And now
this statue, how clumsy, how disproportioned!'' etc. But although the
objects are the same, although the aesthetic outbursts are exchanged,
among the sober-minded and the expert, for interpretations, whereby the
assertoric functions rightly prove their kinship, not only with the
explicative ones, insofar as the interpretation proceeds according to an
assumed prototype, but also with the interpreting ones, insofar as the one
subjectivity will address the other through the work of art, there is
nevertheless disagreement here, just as the interpretations get under way,
a striking disagreement; even the experts can now, with their best will,
not come to agreement.

% --- p. 205 ---
\opage{205}Has beauty then no objective validity? Yes, certainly! it is
precisely what is peculiar in its objective validity that makes the
judgment about it so immediately subjective. But then must not beauty's
validity rest upon subjectivity? Yes, certainly! it is precisely what is
peculiar in its subjective validity that the judgment about it is so
immediately objective.

The unity-in-difference of objective and subjective has its universal
validity in the idea; it is neither the objective for itself nor the
subjective for itself, but the idea itself, that is beauty's measure;
therefore Plato speaks of self-beauty, the invisible beauty by which all
visible beautiful things are beautiful. And yet Plato sees wrongly; for
the idea in its invisible universal essentiality is not beauty; only as an
idea of power in intuitable form does the beauty-idea have existence. As
beauty-idea the idea, the prototype of that which exists, is itself an
existing figure: the existing figure immediately determines the assertoric
judgment. Thus the accentuating interplay of subjective with objective in
logical-antilogical moments continues unceasingly; but it is explained by
the idea, for the aesthetic actuality is transparent in the idea.

The \emph{historical} unity of essence and phenomenon is assertoric; but
in another way than the aesthetic. For the idea the aesthetic actuality
must be perfectly transparent; this is now not so with the historical. The
historically actual bears the stamp of finitude and thereby of
contingency. The fact is indeed by its contingency in the idea's power,
but by no means the idea's adequate phenomenon; to that extent the
reflection repelled by antilogical immediacy cannot go directly to the
idea either in order to find the explanation. Actuality has, on the
contrary, in becoming historical, the peculiarity of being transposed into
pastness; by this transposition it occupies reflection without giving up
immediacy.

% --- p. 206 ---
\opage{206}That the past, although it does not have presence's immediacy,
nevertheless has an immediacy, is involuntarily discerned from such
judgments as: ``Caesar was really murdered,'' ``Canute the Great really
conquered England.'' However many proofs and grounds one may scrape
together for the confirmation of such judgments, the immediacy-element
will always be decisive. The vanished has indeed lost its immediacy
insofar as it has vanished; but it has preserved it insofar as it has not
vanished. In the present narrative the vanished is surely not vanished;
the vanished actuality has in the narrative a reflected presence and works
in this presence with reflected immediacy.

But the fabricated can in a poetically intuitable narrative just as well
take on the stamp of actuality as the vanished actual: upon what does the
separation between the fabricated and the past actual now rest? Evidently
upon a manifold of immediate and mediate actuality-impressions! If the
reader is not himself an investigator, he takes the narrative upon the
narrator's authority; but an impression of authority is, with all the
determinations of immediacy comprised in it, always an
actuality-impression. If the reader is an investigator, manifold
considerations --- of the source or the sources, of the event's character,
its relation to other events, etc. --- will come into account; but the
main thing is that the many particular impressions gather themselves at
last into a total impression, and that this total impression, which speaks
for the event, is a decisive actuality-impression.

Through the total impression the vanished works in a mediate and yet
immediate manner; it works immediately, for the monument stands there, the
report captivates attention, and the attentive hearer sees the event
before his eyes; it works mediately, for the monument sets at a distance,
the narrative reminds in every particular of the distance
% --- p. 207 ---
\opage{207}between the present and the past: reflection is in movement.

The accentuating interplay of reflection and immediacy is here of a
peculiar kind. Upon the aesthetic ground reflection can melt immediacy;
for the phenomenon is essentially semblance; that the assertoric function
nevertheless plays the principal role lies in this, that the sublated
immediacy is instantly restored to antilogical determinateness, since the
phenomenon is an existent figure. Upon the historical ground, by contrast,
reflection cannot melt the immediate, for a fact is still a fact, even
though it belongs to the past; immediacy's ideal process with reflection
must then be seen in the mediate. Between the beholder and the event the
relation is restless: constantly a distancing that alternates unceasingly
with a corresponding approach. Mediately the beholder approaches the
event, and the event out of the past comes to meet him; during the
approach immediacy grows: the actual event makes an actual impression. But
what is thus present is nevertheless something past, it removes itself
again; it must go back to its place when reflection begins, just as the
revenant must go back when morning dawns. It is this restlessness that
comes to rest in faith, as historical faith (\textit{fides historica}).

Historical faith is an assertoric function, peculiarly determined by
reflection's transposition of the mediate into the immediate. Historical
faith is determined by actuality-impressions, its ground is immediacy, its
judgment is assertoric. No reflection can give immediate assurance; only
the actual can speak for the actual; this is modality's objective side:
the assertoric judgment is immediately objective. But the
actuality-impression that founds historical faith is, on account of the
mutual distance of beholder and event, nevertheless only a mediate
impression, an impression of an actuality really absent and present only
in reflec%
% --- p. 208 ---
\opage{208}tion: this is modality's subjective side; the assertoric
judgment is immediately subjective.

The \emph{religious} unity of essence and phenomenon is assertoric in a
quite peculiar sense. Historical and religious faith resemble one another
in this, that both rest upon immediate impressions of something
immediately absent: faith is an assertoric assurance of a non-present
actuality. Only when the question arises what it really is that hinders
the immediate presence of the actual does the difference between
historical and religious faith come to light. What prevents the
historically actual from working with the whole immediacy of the present
is solely and alone pastness; could pastness be sublated, the past would
of course work with the same immediacy as the present, for in essence they
are of the same kind. It is otherwise with the religiously actual. The
religiously actual cannot, even though it be something now existing, come
forward as something immediately present. A miracle cannot be intuited as
immediately present; its presence proclaims itself, on the contrary, as a
breach with everything actually present; such a breach terrifies, and the
terror makes the present into something absent. What is thus absent is, in
its terrifying presence, the hidden, the mysterious, the inexplicable.

``What is it then?'' Contemporaneous with the miracle this question is
only an exclamation; sober deliberation cannot be contemporaneous with the
miracle, for the miracle's immediacy overwhelms reflection. When composure
sets in and deliberation begins, the miracle is over; deliberation
therefore meets, even in the contemporaries, never something present, but
always something past.

``What was it then?'' Here it shows itself that faith, the religious faith
too, is an assertoric assurance of a non-present actual. It is not a
grounding of faith, but faith's logical character as assertoric assurance,
% --- p. 209 ---
\opage{209}that is here treated of.\footnote{As regards the concept's religious validity, we refer to S.\ Kierkegaard's writings, in particular to ``Uvidenskabelig Efterskrift'' and ``Indøvelse i Christendom''. --- ``If there is to be any meaning in saying that the miracle proves who Christ is, then we must surely begin with our not knowing who he is, hence in the situation of contemporaneity with a single human being who is as other human beings are, in whom there is nothing \emph{directly} to be seen, a single human being who then works the miracle and says of himself that he works the miracle. What does this mean? It means that this single human being makes himself into more than a human being, makes himself very nearly into God: is this not offensive? You see something inexplicable, miraculous (nor anything more), he himself says that it is a miracle --- and you see before your eyes a single human being. The miracle can prove nothing; for if you do not believe that he is the one he says himself to be, then you deny the miracle. The miracle can draw attention --- now you are in the tension, and it depends upon what you choose, the offense or the faith; it is your heart that shall be made manifest.'' Indøvelse i Christendom. p.~106.}
What the fact is objectively must be decided immediately subjectively.
Between the fact and the judging subjectivity the distance is infinite, by
reason of the miracle. To sublate the infinite distance is here expressly
posited as subjectivity's infinite task; and yet the solution of this task
is immediately determined by an impression of the miracle, hence of
something actual.

Seen from this standpoint faith has not doubt but offense for its decisive
opposition. Offense is, logically seen, an assertoric judgment just as
much as faith is: ``This is really a miracle!'' ``This is really a
deception of the eyes!'' The believer and the offended one both appeal to
the fact, the actual in its immediacy, for the assertoric judgment is
immediately objective. ``This is really (I can surely feel it upon myself)
a work of God!'' --- ``This is really (sound sense teaches it) a natural
event!'' The believer and the offended one both appeal to an inner sense
for the difference between semblance and actuality,
% --- p. 210 ---
\opage{210}falsehood and truth; for the assertoric judgment is immediately
subjective.

In the strife between faith and offense it is not merely fact that stands
against fact, but in an infinite sense subjectivity stands against
subjectivity, explanation against explanation, understanding against
understanding: it is the assertoric functions that here furnish the ground
for the interpreting ones. So long as the ground is assertoric, faith has
a spiritual superiority over offense, insofar as the believing
subjectivity with its positive forces lives a higher life than the
offended one, who can employ forces and passions only negatively. But if
the ground is to be changed, if \textit{credo, ut intelligam} is to be
thrown together with \textit{de omnibus dubitandum est}, if the opposition
is to be not faith and offense but faith and doubt, then the way goes from
the assertoric to the problematic, and then offense has triumphed.

That offense triumphs with doubt goes without saying; for doubt, which is
faith's torment, is offense's relief. If the assertoric ground is given
up, the reflecting subjectivity must either remain in the problematic or
work its way through to the apodictic. The apodictic rests in rational
necessity, but rational necessity is only in science. Can faith-science
then attain the apodictic? Speculative dogmatics answers: to a certain
degree! but the apodictic to a certain degree is nonsense, an assertoric
nonsense that stands far below the problematic.

\textit{ββ)} \emph{The problematic judgment: functions of sublated
determinateness.} If we, after Kant's example, were merely to employ the
judgment of modality in order to assure ourselves of the a priori validity
of the corresponding categories of the understanding, a detailed treatment
of these judgments at science's present stage would doubtless be
altogether superfluous. But it is not the assurance about the categories'
validity, it is a many-sided development of the categories valid in
themselves, that is here
% --- p. 211 ---
\opage{211}treated. Hegel has indeed held that the categories actuality,
possibility and necessity had to receive their complete development in
``the objective Logic,'' a development which in his system precedes the
doctrine of the judgment. On this occasion we must remark that
actuality-reflection\footnote{Grundideernes Logik. I. p.~105--133.} with
all the categories belonging under it may very well be treated
independently of the doctrine of the judgment, only that it not, by a
misleading abstraction, be treated as independent of all thinking and
judging subjectivity; but that the development of the categories
actuality, possibility and necessity which corresponds to
actuality-reflection by no means renders superfluous a development
corresponding to the doctrine of the judgments, will already be seen from
what has been adduced under the assertoric judgment, and will press itself
forward still more strongly when the talk is of the problematic one.

From whatever standpoint we conceive the problem of possibility, be it as
a general essence-problem, as an objectification-problem, or in particular
as an anticipation-problem: from all sides it shows itself that
possibility and subjectivity stand in so close a connection that
possibility cannot possibly be developed except in reflection of
subjectivity. It is the subjectivity free in itself that frees possibility
and gives it existence; but only in the word, by the judgment, is
subjectivity free, in the logical sense; if now the doctrine of the
judgment can alone develop the freedom of subjectivity, and if
possibility's liberation to ideal determinate being rests upon free
subjectivity, then it is evident that the doctrine of possibility must
stand in intimate connection with the doctrine of the judgment.

According to the assertoric judgment, actuality stands antilogically over
against thinking: the antilogical immediacy can in many ways repel,
attract, move and determine reflection, but the moved reflection is not in
a position to
% --- p. 212 ---
\opage{212}dissolve actuality. According to the problematic judgment,
reflection gains complete power over actuality, nothing actual withstands
the dissolution: here abstraction is made from the antilogical. The
characteristic mark of the actual is determinateness, the immediate,
antilogical determinateness; the characteristic mark of the possible is
the dissolution of the actual's determinateness, hence the indeterminate.

Seen from the subjective side, the indeterminate is the same as the
uncertain. The standpoint of uncertainty is represented by skepticism; the
possibility corresponding to uncertainty we shall therefore call the
skeptical. From the objective side, indeterminateness is a categorical
expression for the actuality that is not, hence either vanished or not yet
come to be, actuality concealed in objective essentiality; this
possibility, conceived as objective essentiality, we call the dogmatic. By
determining indeterminateness, from the subjective as well as from the
objective side, we see that indeterminateness, and with it possibility,
must originally be determined to sublate itself; upon this insight rests
the ontology of possibility.

1) \emph{The skeptical Possibility: the subjective Side of Modality.} On
the assertoric foundation, immediacy's action upon reflection is partly
attracting, determining and binding, partly repelling, loosening and
liberating; the two-sidedness is seen already from this, that the
assertoric judgment is at once both immediately objective and immediately
subjective: through the immediately subjective the assertoric passes over
into the problematic. So long, however, as the immediate is in a position
from one side or another to master, determine and bind reflection, the
assertoric has decisive preponderance: the positive, the actual, is
secured by the fact that thinking goes in the leading-strings of fact.

If thinking is to be liberated, and the reflecting subjectivity to become
aware of its infinite relation to itself through another, then the actual
beginning can be made only by this,
% --- p. 213 ---
\opage{213}that actuality itself shows itself repellent toward reflection.
The problematic does not have its point of departure in itself; it springs
by immediate consequence out of the assertoric. One cannot properly say
that skepticism abstractly begins with doubt; on the contrary! it begins
immediately with immediate certainty; but since immediate certainty has in
itself the negative determination of turning over into uncertainty, doubt
naturally has at once an actual beginning in this coming-to-be of
uncertainty.

How immediately and distinctly reflection in ancient skepticism took aim
at actuality in its manifoldness of heterogeneous phenomena is seen
already at a fleeting glance at its ``Tropes.'' What is adduced in the
``Tropes'' about ``the difference of animals,'' ``the difference of men,''
``the difference of the senses,'' ``the various circumstances, distances,
relations'' etc., in which the thing is found, shows sufficiently what it
means that the problematic proceeds from the assertoric, that uncertainty
has its origin in an immediate certainty.

The task is to think the actual, that is to say: by considered thought to
assure oneself of the immediate certainty; in this task a contradiction
lies concealed. To think the actual is, in other words, to transform the
singular actuality-determinations into thought-determinations. But the
actuality-determinations in their immediate singularity are contingent,
the thought-determinations in their reflected generality necessary; the
thought-determinations are logical, the actuality-determinations
antilogical: here is the contradiction. This contradiction comes to light
on the assertoric foundation in that the objects in their existent
difference occasion the judging subjectivities, in their likewise existent
difference, to mutually divergent and conflicting judgments: ``The wind is
really cold''; ``The wind is really not cold''; that is the contra%
% --- p. 214 ---
\opage{214}diction, the old contradiction, which was already to be seen in
Protagorean sophistry.

Repelled by actuality and bent back into itself, reflection now seizes the
assertoric judgment from the subjective side with heightened energy. Not
enough that reflection transforms the judgment into opinion and opinion
into assertion; the subjective element is intensified to the utmost by
this, that opinion is set against opinion, assertion against assertion,
and in such a way that the one opinion, the one assertion, receives
precisely the same validity, hence the same invalidity, as the other. The
immediately determined actuality has in this way become something wholly
indeterminate, something for reflection indeterminable.

Skeptical reflection has not in an objective manner penetrated the essence
of actuality, hence has not in the objective sense dissolved actuality
either; on the contrary! what the actual is in itself, something or
nothing, essence or semblance, substance or accident, of this the skeptic
knows nothing, nothing at all. It is thus possible that the actual is in
itself something; but it is also possible that it is nothing; the one of
the two conflicting assumptions is in actuality just as possible as the
other; actuality with its immediate determinations is thus only in itself
an indeterminate, insofar as it is for subjectivity indeterminable: the
dissolution is subjective, the skeptical possibility uncertainty.

That every immediate determinateness, and therewith every
actuality-determination, in the moment it is touched by reflection
transforms itself into indeterminateness, that consequently no true
certainty can be had about actual things, is now the discovery to which
skeptical reflection has led. In what sense this reflection can be said to
signify an advance will be evident when we take everyday consciousness as
our measure. For everyday consciousness it holds that we have immediate
certainty about one thing while we are in uncertainty about another, so
that
% --- p. 215 ---
\opage{215}the uncertainty that obtains can under certain conditions be
raised to certainty; for the skeptical consciousness immediate certainty
has once and for all vanished, that is to say: it has been transformed
into semblance, its credit is undermined; it turns over into uncertainty,
it is a deception. For the unmasking of this deception the ancient
skeptics employed a multitude of examples, heaped together from all the
various nooks and corners of actuality. On these exemplifications, with
their repetitions, variations and alterations, no particular weight can
however be laid so far as the principle is concerned; they are a kind of
empirical filling-mortar, prolix enumerations calculated to make up for
the lack of progressive development and to fill out the emptiness.

``But,'' one might ask, ``when skeptical reflection, repelled by
actuality, was thrust back into itself, must we not then expect that it,
after having lost the certainty about the nature and constitution of
existing things, would with all the greater energy set about understanding
itself and coming to certainty about its own essence?'' To this it must be
answered that the manner in which reflection here gives up actuality is
precisely what makes it impossible for it to become transparent to itself.
Reflection is repelled by immediate actuality, hence by the phenomenon,
and has in turn thrust the phenomenon away from itself; just thereby it
has become a phenomenon to itself, and for that reason cannot conceive
itself as essence. Only a phenomenon can be thrust back by a phenomenon;
essence as essence the phenomenon cannot thrust back. --- The thought of
the Eleatics could nevertheless in a way hold fast to itself in the One;
but neither had it let itself be thrust back by the manifold; it was not
uncertain about the nature of the manifold, it was on the contrary quite
certain that the manifold-being was a non-being, and just for that reason
it could be certain of itself in the one-being. Fichtean reflection too
makes an end, in its own way, of the whole outward actuality, in that it
makes the thing into a barrier for the I; but this dissolution of the
actually objective
% --- p. 216 ---
\opage{216}signifies no uncertainty; on the contrary! thought in
subjective idealism gains the force to reduce the thing to a barrier
precisely by this, that in its other it is able to hold fast to itself: it
knows what semblance is, and therefore it also knows what essence is. It
is otherwise with skeptical reflection.

Skepticism gives up actuality and declares it to be semblance; but since
the declaration has only subjective significance, is only a designation
for uncertainty, it being just as possible that the actual is in itself
not semblance as that it is semblance: then skepticism concedes that it
does not know what semblance is. But a reflection that does not know what
semblance is does not know what essence is either, and cannot possibly see
through its own essence.

It is in this respect characteristic that the skeptics conceive the inward
and the intellectual, cognition, understanding etc., just as phenomenally,
outwardly and immediately as sensible things. Thus e.g. when they speak of
the possibility that man should be able to cognize himself;\footnote{``For
if man is to cognize himself, then either the whole man must be the
cognizing subject, or the whole man the object of cognition, or one part
of man the cognizing and another part the object of cognition. But if the
whole man is the cognizing subject, then the object of cognition is
sublated, and if the whole man is the object of cognition, then the
cognizing subject has vanished. If, on the other hand, only one part of
man is the cognizing and another part the object of cognition, then man,
since his parts are the body, the senses and the understanding, must
either with the body cognize the senses and the understanding, or with the
senses and the understanding cognize the body.'' E. Frisenberg Nielsen. Den
antike Skepticisme. Kjøbenh.\ 1862. p.~75.} the understanding and
cognition are considered in such investigations in the same immediate
manner as blue and yellow. For so immediate a conception the understanding
is only a phenomenon; of such a phenomenon one has only immediate
certainty, the fleeting certainty that is instantly uncertainty.

% --- p. 217 ---
\opage{217}Uncertainty-reflection can moreover be driven to a remarkable
degree of acuteness and subtlety without the reflecting subjectivity
thereby coming a hair's breadth nearer to essence-cognition. Thus e.g. in
the reflections on cause and effect. ``If anything is a cause, then either
a body must be the cause of another body, or an incorporeal of an
incorporeal, or a body of an incorporeal, or an incorporeal of a body.''
The facility in setting up and assuming various cases is conspicuous; it
is everywhere the phenomenal immediacy, with the corresponding immediate
assumption, that yields the point of departure. Only assume something:
every assumption shall at once give place to the opposite assumption; only
posit a possible case, and the opposite case shall show itself just as
possible; set up whatever determinate case you please as actual: the
actual case shall at once change into an absurd, hence into an unthinkable
case. ``If a body is to bring forth another body in and for itself, then
this its being in itself must, which is contradictory, bring forth a being
outside itself, and thus, which is absurd, act outside its own nature. If,
on the other hand, a body is to bring forth a third in combination with
another, then out of one there will come not merely two, or out of two
three, but an infinite number, which is impossible''\footnote{E. F.
Nielsen. Den antike Skepticisme. p.~117 \&\ 118.} etc.

The dialectical dexterity with which the ancient skeptics, for the most
part by straightforward use of what they found in their predecessors, the
great philosophers, know how to exhibit grounds for and against every
assumption, has led Hegel to an overvaluation of the whole
spirit-phenomenon. ``Only ancient skepticism,'' he says, ``is of a true
and deep nature, while the more recent is rather Epicureanism. For whereas
ancient skepticism is directed precisely against the thinking of the
understanding, which lets determinate difference hold as ultimate, as
being, the more recent skepticism is on the contrary directed against
thoughts, against
% --- p. 218 ---
\opage{218}the concept and the idea, hence against the higher philosophy.
The more recent skepticism, which in this way lets the reality of things
stand quite undoubted, is then not even a peasant philosophy; for even
peasants know that all earthly things are perishable, that their being is
thus just as good as their non-being.''

Of the misunderstanding contained herein the following statement gives
what seems to be an apt and quite satisfactory characterization. ``In
order to explain the Hegelian overvaluation,'' it is said\footnote{E. F.
Nielsen. Den antike Skepticisme. p.~190--192.}, ``we must expressly note
the difference there is between a negative use of the categories and a
higher consciousness of the significance of the negative itself. With
reference to the seventh trope, where it is said e.g.: ``Wine strengthens
when it is enjoyed in moderation, but weakens and dissolves the body when
it is enjoyed to excess,'' Hegel emphasizes that the skeptics have, if
only indirectly, intimated that quantity and composition are so far from
being a matter of indifference with respect to quality and dissolution,
that with the change in quantity there follows a change in quality as
well. But that Hegel in this connection reads his own logic out of the
skeptical reflections is conspicuous. It has assuredly as little occurred
to any of the ancient skeptics to lay a special logical weight on the
essential connection between the quantitative and the qualitative as it
is on the whole in keeping with their spiritual tendency to ascribe
essential validity to any category whatsoever. To adduce yet another
example of Hegel's soaring interpretation, we refer to the following
statement in him: ``The skeptics seek everywhere to exchange the
expression ``it is,'' which contains in itself something fixed, something
that maintains itself under all circumstances, for the expression ``it
seems'' or ``everything is movable.'' But if they thereby sublate sensible
likeness and identity, and consequently \emph{this}
% --- p. 219 ---
\opage{219}generality, then another enters; for generality, or that which
is, lies precisely in this, that one knows e.g. that to the jaundiced that
shows itself yellow which to another shows itself e.g. white, etc.
Against such a generality, which is indeed not true because it is
immediate, its non-generality is therefore demonstrated with full
justification within this generality itself by means of another
generality.'' ``But'' --- so runs the objection --- ``the whole context
nevertheless shows plainly enough that the skeptics by no means strive to
set the one generality against the other in order thereby to clarify the
consciousness of a truer and higher generality, but much rather, in
accordance with their all-dissolving tendency, to set the one assumption
against the other, and just thereby to make the general unrecognizable.
The ideality of skeptical negativity is plainly overshadowed by that
darkness of consciousness in which the skeptic moves. There is in the
dissolving acumen a certain unconscious --- half-conscious blindness,
which makes one easily take the skeptic's words literally when he declares
that he does not even know whether he knows anything.''

``Skepticism remains,'' according to Hegel\footnote{Cf. Vorlesungen über
die Geschichte der Philosophie. Zweiter Theil. Werke. Vierzehnter Band.
p.~475 \&\ 514.}, ``standing at the negative as result and thereby fails
to recognize that negation has a determinate affirmative content in
itself, whereas the speculative method determines negation as the negation
referring itself to itself, as the negation of negation, or more precisely
as the infinite affirmation. Skepticism exercises its dialectic according
to contingency, since it, just as the material presents itself to it,
shows in this material that it is in itself the negative, whereas the
speculative method exercises its dialectic in virtue of necessity, in that
it proves that the idea as speculative idea is neither a straightforwardly
positive nor a straightforwardly negative, but such a union of both that
the negative and the positive in this union everywhere reflect one
another.''

% --- p. 220 ---
\opage{220}Skeptical negation separates itself into a swarm of isolated
negations, just as skeptical contingency into a swarm of particular cases.
To the assumption of each single case corresponds the assertoric; to the
rejection of each assumed case, the problematic. The beginning is
everywhere, for here one begins everywhere with a case; the end is
everywhere, for the certainty about every case transforms itself
everywhere into uncertainty. It is the negatively pointed interplay of the
assertoric and problematic functions that characterizes skepticism: only
in this interplay can uncertainty incessantly come to be.

``It might well seem as if his\footnote{E. F. Nielsen. Den antike
Skepticisme. p.~3--5.} (the skeptic's) seeking led at least to the
determinate result that no truth is to be found. But to let the skeptic
fix this proposition is again to misunderstand skepticism and to confuse
it with a negative dogmatism. The thetic conception of the negative is
precisely characteristic of negative dogmatism, in that it is not merely
assured that something is uncertain, but also that it is certain that
something is uncertain. Negative dogmatism then knows what is certain and
what is uncertain, and is thus certain of its proposition and of its
opposite. The skeptic, on the other hand, sees with his double vision the
doubling everywhere, even in truth's negative expression. A statement such
as: everything is uncertain, is therefore, on his view, itself uncertain,
and this last again uncertain, and so on to infinity. It stands
correspondingly with the expressions: everything is false, nothing is
true, etc. \dots\ But even though the Aporetic indicates uncertainty never
so much, the indication will always take on the semblance of a fixing, a
positing, since it must after all remain a statement, a negative
expression being able to stand fast just as well as a positive one. He
must nevertheless, as it seems, concede and assert that he has a sure and
certain cognition of the inse%
% --- p. 221 ---
\opage{221}curity and uncertainty of cognition. As one who now would avoid
the reproach of sublating, by his mere statement, the equilibrium of
grounds and counter-grounds, \dots.\ the skeptic calls himself an
Ephectic, \dots.\ for only by making epoche, by holding back every assent,
every expression, i.e., in other words, by keeping silent, can he guard
himself against every misunderstanding and satisfy his ideal demand.
Silence is thus the highest dissolution into uncertainty. It is not a
silence in which one rests as in inwardness and self-absorption, for this
pure silence is not skeptical; but the silence in which the negative is
consummated, in that it shows that nothing is decided, the silence in
which there is undecidedness, uncertainty, in which there is
possibility.''

It is the skeptic's task to bring about an equilibrium between the
conflicting assumptions, the mutually contradictory opinions, an
equilibrium which means that the one is just as reasonable, and therefore
just as absurd, as the other. When this equilibrium has been brought about
for the single case, a new case must at once present itself, so that the
proceeding can begin afresh. Were there here but one single case,
attention would have to fasten continually upon the uncertainty attained
by the equilibrium: it would then have to be certain that uncertainty had
once and for all been attained; this standing uncertainty is insufferable.
The skeptic cannot straightforwardly deny uncertainty; for by a denial of
uncertainty he would confirm certainty. Neither can he straightforwardly
confirm uncertainty; for a confirmation of uncertainty would then itself
yield a certainty. All contradictions uncertainty can balance out;
\emph{the contradiction of certainty and uncertainty alone it cannot
balance out}: if certainty is posited, then uncertainty is excluded; if
uncertainty is posited, then certainty is excluded. Every time the skeptic
is to acquiesce in a last decisive uncertainty, he makes epoche; he keeps
silent, and the silence expresses that uncertainty is held fast, not in
the form of certainty, but in the form of uncertainty. To speak would in
such cases be a rashness; to
% --- p. 222 ---
\opage{222}declare oneself straightforwardly for uncertainty would be to
place it in the form of certainty, to assume uncertainty as the certain,
hence to contradict oneself.

Everything is uncertain; this is skepticism's principal proposition. But
no one can get him to fix the certainty that everything is uncertain. No,
he fixes nothing, he keeps silent; the silence means that he cannot at all
enter upon \emph{general} answers to \emph{general} questions. To enter
upon a general answer to a general question would for skepticism be the
same as to deliver a definitive declaration; but one single such
definitive declaration overturns his whole skepticism.

But does there not then lie in this proposition: everything is uncertain,
a definitive declaration? By no means! that is to say, in the skeptical
sense. Instead of engaging with the universal and saying: uncertainty is
universally valid, the skeptic goes at the singular, and consoles himself
that in every single case he will manage to get the uncertainty out, will
manage to hold it in continued becoming: only in continued becoming is
uncertainty infinite.

The skeptical possibility is infinite uncertainty; \emph{the infinite,
i.e. endless uncertainty, is in infinite becoming}. If becoming ends in
determinate being, if uncertainty is, so to speak, permitted to stand
still and settle down into an existence, then in reflection it is all over
with it: an uncertainty that is, a certainty that uncertainty is, is a
certainty negatively reflected into itself. With uncertainty's
transformation into negatively reflected certainty, skepticism is
transformed into negative dogmatism: the skeptical possibility has become
dogmatic.

2) \emph{The dogmatic Possibility: the objective Side of Modality.}
Indeterminateness must always have its foundation in some determinateness
or other, uncertainty in a certainty, the problematic functions in a
non-problematic. So long as the non-problematic coincides with actuality's
immediate determinateness and the certainty corresponding to it, the
foundation is
% --- p. 223 ---
\opage{223}assertoric. If the subjectivity reflecting on its own
reflections gives up assertoric immediacy, reflection must itself posit
and presuppose its reflected foundation.

The reflected foundation is now in the first place conditioned by the
removal of the immediate; this is the negative side of the matter. If the
immediate determinateness is taken away, then only indeterminateness is
left; if certainty is taken away, then here there is only uncertainty. How
then shall the negative now by itself be able to yield a positive? The
indeterminateness reflected into itself is itself negative determinateness;
the uncertainty determined in itself is negatively determined certainty;
the new foundation is thus a determinateness won within indeterminateness,
a certainty affirmed within uncertainty: instead of actuality's immediacy
we have obtained an immediacy of thought; this is the positive side of the
matter.

The reflected foundation is, on account of negation and ``the negation of
negation,'' infinitely affirmative, but not antilogically positive; it is
in actuality itself not an outer, but an inner; not a
phenomenon-determination, but an essentiality. Only in actual essentiality
can actuality-reflection come essentially to itself; actual essentiality,
the actuality anticipated in reflection, is dogmatic possibility.

The skeptical possibility was negative with respect to the objects as well
as to the reflecting subjectivity; the dogmatic possibility is affirmative
for the objects as well as for subjectivity. The affirmative possibility
is objective as the possibility of things, the possibility out of which
the actual proceeds, the possibility that only takes up the conditions for
becoming actual, and only becomes actual in that it at once reaches and
oversteps its own limit, for the limit of possibility is indeed actuality.
The affirmative possibility is subjective as the thinkability of things,
for only in the form of thinkability, of thought, of the concept can an
actuality with all its determinations be there as possibility. The affir%
% --- p. 224 ---
\opage{224}mative possibility receives its particular dogmatic stamp from
this, that the supposing functions pass over into the problematic ones.

Just as the skeptical possibility arose from the fact that the problematic
function sprang out of the assertoric, so the dogmatic possibility now
arises from the fact that the problematic springs out of supposing
functions. What the assertoric functions are in relation to immediate
actuality, the supposing ones are in relation to the reflected immediacy
of essentiality. The movement is this: here, by logical slipping-under, a
thought-determination is posited in reflected and yet immediate unity with
a corresponding being-determination: the immediate unity of thinking and
being is dogmatic. Once the being-determination in question has been set
up, the testing follows; reflection places itself over against its own
presupposition as before a logical object; the thought is represented, the
represented is thought; this is the unrest; in this unrest lies the
negativity: the negativity at once represented and yet thought awakens the
problematic.

For dogmatic reflection the possible is quite generally one with the
thinkable; the supposed unity of thinkability and possibility corresponds
straightforwardly to the supposed unity of thinking and being: as the
thinkable, possibility is subjective; as the possible, thinkability is
objective. Hence it comes that the test of the possible coincides with the
test of the thinkable, just as the test of the impossible coincides with
the test of the unthinkable. The mark of the thinkable is freedom from
contradiction; the same mark holds for the possible (\textit{possibile
est, quod non implicat contradictionem}). The mark of the unthinkable is
contradiction; the same holds for the impossible (\textit{impossibile est,
quod implicat contradictionem}). But this supposed unity of thinkable and
possible, of unthinkable and impossible, now calls forth doubt about the
objective essence of possibility.

% --- p. 225 ---
\opage{225}``Everything in which no contradiction lies is possible,
whether it exists or not. Indeed everything that exists has first been
possible before it came to be: otherwise it could never come to any
actuality. That which has existed is still possible, although it no longer
exists, and it remains possible without end \dots\ A man would never have
come to be had it not been possible that his parents could have brought
him into the world; and, after he is dead, one can say the same, he would
never have lived had his life not been possible. The possible is
consequently eternal, and every possibility is eternal, because no time
can be thought, either before or after, at which it should either begin or
cease again.'' These words of Eilschow, an independent thinker\footnote{Fr.
Chr. Eilschov. Philosophiske Breve. Kjøbenhavn 1748. p.~147.} from the
Leibnitz-Wolffian period, are well suited to characterize the dogmatic
possibility. Thus: possibility is eternal; it is a thought's agreement
with itself in the form of a being's agreement with itself; that which is
self-agreeing agrees with itself for all eternity, whether or not there is
anyone who thinks upon it.

First, then, the possible is posited as identical with the thinkable; next
the thinkable is, on account of thought's supposed unity with being,
posited in the form of the possible as so being, that the dogmatist falls
into doubt whether he shall let possibility rest eternally upon itself, or
make the eternally thinkable dependent upon an eternally thinking one. The
eternally thinking one is God; it lies in the immediacy with which the
supposing functions presuppose God, that the problem takes on a tinge as
much theological as logical. One might then be tempted to emancipate
possibility from God. ``Yet we shall,'' says our dogmatic philosopher,
``not make every possibility into a little divinity that has the ground in
itself for itself, which seems to follow from what I have said. Some
% --- p. 226 ---
\opage{226}old school-sages have fallen into the perverse thought that
the possibilities did not depend upon God or have their ground in his
essence, because one could think of them and prove them without
representing to oneself that they had their ground in some
eternally-necessary being, or in God. A triangle is possible because a
figure can be enclosed by three lines: one can say neither that these
three lines nor that their meeting in as many points is possible because
there is a God. It seems much rather that this possibility would be
undeniable even if he had never existed; and an atheist who denies God
does not for all that deny the triangle's possibility. What then are we
to lay down?''

Three preliminary possibilities now present themselves concerning the
eternal possibilities: either they are grounded in themselves, or in
God's will, or in the divine understanding. ``Are we to say with
Cartesius that they are grounded in God's will and good pleasure alone,
so that a thing is possible merely because God wills that it shall be
so, then it must also be conceded that the same thing can become
impossible as soon as he finds it good. A triangle became possible, then,
when God willed that three lines in as many points should be able to
enclose a space; but how could it now come about that the same could
become impossible, or how could God resolve that three lines should
henceforth no longer be able to compose a figure? If one concedes that
the latter resolve is impossible, as one cannot but concede, then the
first, which makes the triangle possible, must be eternal; afterwards at
the very least: that is to say, what God has once willed should be
possible must thereafter remain so through all eternities \dots\ But that
is not yet all that can be called to mind. Just as Cartesius must concede
that what has once become possible according to God's will always remains
so, so can we also maintain that what is now possible has been so always;
for as great a violence as one must do to the understanding in order to
represent to oneself that the time could come when a triangle should be
as little able to enclose a space
% --- p. 227 ---
\opage{227}as two straight lines now can do it, just so completely must
one also violate the understanding in order to represent to oneself that
there has been a time when God could have made a two-sided figure
possible and a triangle impossible. From this it follows that just as one
must concede the eternity of the possibilities afterwards, so must one
also concede it beforehand, or that all possibilities are, without
beginning and end, unchangeable, necessary,
eternal''.\footnote{Eilschov. Philosophiske Breve. p.~147--49.}

What is here said of the eternal possibilities holds dogmatically for all
the abstract a priori concepts, the so-called eternal truths
(\textit{veritates æternæ}). But although possibility is in its universal
concept eternal, and in this its eternity indifferent toward the
difference between actual being and non-being, possibility cannot
possibly, in relation to the thing, be indifferent toward the thing's
possible determinate being.
If the possibilities are made independent of God, and things are only
actualized possibilities, then things seem thereby to become independent
of God as well: how now is this possible? ``The possibilities of things
are dependent upon God, but not upon his bare will and good pleasure.
They are grounded in his infinite understanding and the same's eternal
representations of all that could be brought forth by his omnipotence.
Were God not to exist, then no thing would be possible, that is true
enough; but a possibility is for that reason as little created by God as
he can be said to have created his own
understanding''.\footnote{Anf. Skr. p.~150.} --- The dogmatic possibility
is indeed an eternal concept, but with all its eternity nonetheless,
remarkably enough, a restless concept. If the possibilities of things
cannot rest upon God's ``bare will and good pleasure'', since the eternal
concept is raised above all arbitrariness; if the possibilities cannot be
independent of God, since no thing would be possible if God did not
exist, then they must presumably be in the divine understanding. But with
this assumption too the
% --- p. 228 ---
\opage{228}dogmatic reflection strikes upon difficulties. Were the
possibilities of things to depend upon the divine understanding because
all possible things depend upon God, then God, who has created things,
must surely also have created the possibilities; but that is altogether
impossible. The understanding cannot then determine possibility; the
possibility-concept is objective: the understanding has to respect the
concept's unchangeable objectivity. ``I say it is enough that the
properties of things do not contradict or overturn one another; therefore
they are possible, and they need no other ground for their possibility
than just this non-contradiction. Or how am I to comprehend the
connection in this syllogism: this or that thing has been, is and remains
eternally possible; therefore there exists a being which eternally
represents its possibility to itself. Our forefathers did not think that
many of the things which the moderns have proved in the mechanical and
physical sciences were possible before art in our days discovered them;
but have they for that reason not been so? The ancients knew, for
example, only quite a few electrical bodies, and did not represent to
themselves that water could be electrified \dots.; yet nature was then,
as now, filled with electrical bodies. We know it; but they did not know
it. Their representation and ours alters just as little in nature as one
can say that a mirror in a room makes the things that are found in it, or
that there is a chair, table, benches because one can see all this in the
glass \dots. When we say that the infinite representation-force
represents infinite things to itself, ours by contrast only certain
things \dots. then we set down all the difference that there is between
God's understanding and ours. The former represents to itself at one
stroke all possible things, the latter only a few, and those by degrees;
the former is unrestricted, the latter is restricted, etc. But in all
this we still find no ground in God's understanding for the possibility
of things. So much the famous Wolf shows, that God's will without
antecedent necessity cannot have invented the ideas in which the
possibilities are contained; since they are unchange%
% --- p. 229 ---
\opage{229}able and therefore eternal; but it is by no means a
consequence that they are therefore grounded in his
understanding.\footnote{Eilschov. Philosophiske Breve. p.~154--55.}''

From these utterances it will be possible to see what the relation is
between dogmatism's slipping-under and its problematic functions. The
dogmatic thinking is hypostatizing. First it makes the possibilities into
pure thought-beings and gives them an objective determinate being;
alongside this it posits the Godhead as the absolute substance; and now
the question arises how far the possibilities are in God or outside God.

Thus the dogmatic reflection can, according to the foregoing, apprehend
the possible as the thinkable, stare at this in-itself thinkable as at a
being, without seeing the impossibility that the merely thinkable should
be able to have any determinate being other than that which it obtains
through thinking: thought's logical objectivity impresses the dogmatic
thinker. The possible that I think, says the dogmatist, does not become
possible because I think it; but I think it possible because it ---
independently of my thinking --- is possible. ``A triangle is possible
even before I learn to represent it to myself in thought. I can afterwards
represent it to myself; but the figure's possibility is older than my
representation, and the latter does not contain the ground of the former.
If I inscribe it upon a surface, then it has me to thank for the fact
that it exists; but by no means my preceding representation for the fact
that it is possible\footnote{Anf. Skr. p.~155.}''. Thinking and
representation clash with one another. In that thinking objectifies its
concepts to itself, these concepts take on a represented objectivity;
through the represented objectivity thinking's own products become
independent of thinking: the dogmatic thinking comes as a stranger to
what is its own and thereby becomes a stranger to itself.

% --- p. 230 ---
\opage{230}So long as representation clashes with thought, an altered
point of departure will constantly lead to an altered result: here, in
dogmatism as in skepticism, there is possibility for a multitude of
highly different points of departure and therefore also possibility for a
multitude of mutually conflicting results. The difference is only this,
that while the skeptical reflection finds itself repelled by every
result, because in every result it strikes upon the antilogical, the
dogmatic reflection will every time find self-confirmation in the result
it has won, since this result is nothing other than its own logical
essentiality, slipped under by supposing functions. What infinitely
interests the skeptics, that something problematic incessantly results
from every possible assumption, surprises and torments the dogmatist. The
dogmatic thinker supposes at every moment, once a demonstration is
concluded, that he has found something fixed, something whereby he can
set himself at rest, and does not heed that contradiction smoulders in
possibility, that possibility is like a snake in the grass.

If the dogmatist, e.g., starts from the consideration that the thing's
possibility can exist prior to the thing only in this, that it has
existence in a representation independent of the thing's actuality, then
he will soon see that it is the representation that makes possibility
independent of the thing. But the general representation cannot possibly
represent itself: every representation must presuppose a representing
subjectivity. Thus, in that the possibility independent of the thing is
dependent upon representation, and representation upon the subject's
representing, the independence of possibility from the thing must
precisely be grounded in its dependence upon subjectivity: the eternal
possibilities, independent of all possible things, must be grounded in an
eternal, thinking, willing substance independent of all possible things.
Now it is thus certain that the eternal possibilities are in God.

But in this certainty there is concealed once again a good deal of
uncertainty.
% --- p. 231 ---
\opage{231}For when the divine will is posited as dogmatic positive
reality over against or alongside a corresponding reality, namely the
divine understanding, then it is still uncertain how far the eternal
possibilities that are in God are to be assumed to rest either upon God's
will or upon the divine understanding. Everything here depends upon the
point of departure: one can dogmatically start from the understanding,
but one can also start from the will. Here, as we have seen, Wolf and
Cartesius stand sharply opposed to one another. Wolf assumes that all
possibilities are grounded in the divine understanding, Cartesius that
they rest upon God's will. The author of the ``Philosophiske Breve''
attaches himself most nearly to Wolf, but attempts a grounding of his
own. ``It is of no use,'' he says, ``to defy an Aristotelian opponent
under the customary philosophical armor, with which some believe they can
ward off every thrust. It is of no use, I would say, to come forward with
the general fundamental rule: \emph{Everything that is must have a
sufficient ground why it rather is than is not}, and thereupon to
syllogize thus: every possibility must have a ground why it is rather a
possibility than an impossibility: this it cannot contain in itself;
otherwise it were God, consequently in something else'' \dots.\ ``the
said principal ground concerns only the realm of actual, and not of
possible things.'' The author's agreement with Wolf is then, in
consequence of the corrective applied, grounded in a disagreement; the
result is the same, only the ways are different. But the dogmatic
separation of way and result leads precisely into new entanglements and
new contradictions. In the corrective applied there is concealed a
contradiction.

Dazzled by the pure concept's formal objectivity, the author of the
``Breve'' sets the eternal possibility self-subsistently against all
actualization; relying upon the axiom of contradiction he does not notice
the contradiction that a possibility which lacks a ground of existence,
and therefore, so long as the ground is lacking, can ne%
% --- p. 232 ---
\opage{232}ver become actual, is impossible. Between possibility and
actuality he interposes the ground, and does not notice that possibility
itself stands just as sorely in need of the ground in order to become
possible as ``the realm of actual things'' does in order to be actual.
``Everything that is, namely everything that is actual, everything that
exists, must have a ground why it exists, and rather in that manner than
in any other; but not so long as it remains in possibility. --- So long
as it is merely possible, it has ground enough in itself for itself'';
--- \emph{without God things are possible}. ``God necessarily represents
to himself all possibilities, such as they are, that is, as
possibilities, as things in whose properties and constitutions no
contradiction is found: further ground they do not need. If any of them
is to come into being outside God's representation, then God himself must
prove that it is possible, then he must create, then he must will that it
shall exist\footnote{Eilschov. Philosophiske Breve. p.~156 ff.}'':
\emph{without God things are impossible}. Thus: \emph{without God things
are possible; without God things are impossible}; this is the
contradiction.

After Eilschov has in this manner, by tacking among contradictions,
pursued by contradiction, brought movement into Wolf's congealing
syllogism, he then goes over to setting forth his own view of the matter,
and arrives at a result that is quite characteristic of the dogmatic
reflection; for in the result contradiction is again seen. He will, as he
expresses himself, show ``that all possibilities are grounded in God's
understanding, and that they consequently have their ground in God's
possibility alone, above which there is none.'' First he elucidates the
matter by an example from ``numbers and magnitudes'', for of these it
holds that ``the less is always enclosed in the greater down to the first
unity.'' ``The possibility of the one is grounded nowhere but in itself;
but it has in itself the ground of all other numbers' possibility;'' it
is therefore so apt ``that Pythagoras has called God
% --- p. 233 ---
\opage{233}by a name which is full of emphasis, \emph{Unity}''. Now
follows the explanation. God's possibility is solely and alone grounded
in itself, because it is determined by itself without regard to anything
else, just as unity, the original unity, is determined by itself without
regard to any finite number. Were now a thing's possibility to be thought
independent of God's possibility, then it would likewise have to be
thought to be determined by itself ``without the help of any other
possibility'': but this is, on account of the thing's finitude,
impossible. ``When I ask whether a thing is possible, I become
unreasonable if I do not first posit the constitution of that thing about
which I am asking.'' The finite is in its limits undetermined, and must
therefore be determined; only in being determined does it become
possible. It is indeed in general possible that a line can ``lie evenly
between its points'', but an infinite straight line is impossible. ``When
one then names an infinite straight line, one says a line which is to
have points, that is, bounds, for otherwise there is no line, and which
at the same time is to have none, for otherwise it is not infinite; that
is, \emph{one contradicts oneself}.'' --- By these considerations the
author of the ``Philosophiske Breve'' then arrives at the result that the
possibilities of things are contained in God's possibility, just as the
finite numbers are in unity. ``Were no infinite representation-force
possible, then no finite thing would be possible: because every finite
possibility is not of necessity determined; \emph{it must first be
determined by the infinite}. This determination cannot take place except
in representation, so long as the possible is not actual, or does not
exist. \emph{Consequently the possibilities of finite things have their
ground in the infinite representation-force and by no means in
themselves}\footnote{Eilschov. Philosophiske Breve. p.~160--61.}''.

But a syllogism which in this manner is to solve the problem will
precisely on account of the manner --- with the modality of the func%
% --- p. 234 ---
\opage{234}tions everything depends upon the manner --- be a new
beginning of the problematic. For if the possibility of things must be
determined by an infinite understanding, because the finite, that which
is in itself undetermined but nonetheless determinable, can be determined
in infinitely many manners, then there must, dogmatically understood, be
something arbitrary in the manner in which the infinite understanding
expressly determines each single possibility: if the single possibilities
are dependent upon the divine understanding, then they are thereby
dependent upon the divine will as well. In this manner --- and everything
depends upon the manner --- the strife between Cartesians and Wolfians
can never end.

\emph{Thesis. It is God's arbitrariness that has determined the finite
possibilities. Proof.} The possible is in its finite determinateness
opposed to the necessary; for the possible is determined \emph{thus} only
because it could also be determined \emph{otherwise}. What is determined
\emph{thus}, although it could be determined \emph{otherwise}, is
determined with arbitrariness; thus \dots

\emph{Antithesis. God's arbitrariness has not determined the finite
possibilities. Proof.} Necessary is that whose opposite is impossible or
contradictory. Now it is contradictory that the possible should be
impossible; consequently every possibility is necessary. The necessary
cannot be determined arbitrarily; thus \dots

It is characteristic that Kant, in the course of his attack upon
dogmatism, set up ``the antinomy''. The antinomy is dogmatism's
consequence; its secret is an interchange of essential points of
departure: in the antinomy the dogmatic possibility becomes impossibility;
the task is to get this impossibility transformed back into possibility.

3) \emph{The Ontology of the Possibilities: two-sided Modality.} It has
come to be regarded as a sign of the Hegelian dialectic's superiority over
``the finite understanding'' that it
% --- p. 235 ---
\opage{235}has been able to solve the Kantian antinomies. In a certain
manner this superiority is unmistakable. In two points Kant shows himself
entangled in dogmatic prejudices. He is transfixed by the immediacy in
the dogmatic points of departure. In his thesis one standpoint holds, in
his antithesis another and opposed standpoint; between the opposed
standpoints no essential transition is here given: one comes by a leap
from the one to the other. He is transfixed by the immediacy in the
dogmatic limit. The dogmatic limit is a fixed, existent limit, secured by
a distinction that supports itself upon the axiom of contradiction ---
and upon \textit{principium exclusi medii}.

In both respects Hegel's logic has illuminated the error and by this
illumination earned itself immortal merit. The Hegelian logic can, so far
as the form is concerned, with full right be called a logic of the
possibilities. Possibility is at bottom transition; he who cannot
understand the Hegelian logic cannot understand any essential transition
in the concept. There may be much to object to in many of this logic's
transitions in detail; but one thing is the transition determined by
determinate categories, another is the transition-concept: the general
transition-concept is possibility. Should anyone object that it is not
possibility but becoming that is in Hegel the proper and first
transition-category, then one must answer to this that becoming as
transition from being to nothing can acquire no significance at all if it
is not seen how it is possible that being can become nothing, and nothing
being. All who have raised objection to this transition have really
raised objection to its possibility.

The controversies about the logical becoming, the immanent transition,
the axiom of contradiction, \textit{principium exclusi medii} and the like
can all be brought under the one decisive standpoint, the standpoint of
possibility. It is an old joke that a dialectician can make black into
white. Insofar as an actual opposition is here thought of, he who
% --- p. 236 ---
\opage{236}makes black into white is plainly a sophist. In actuality
black is not white, white not black; in actuality \textit{principium
exclusi medii} holds without contradiction. But when the talk is not of
the actual but of the possible, is he then also to pass for a sophist who
should point out the possibility of a transition from black to white?
Should it be sophistical to point out how the opposed concepts pass over
into one another with inner necessity, so that the one, essentially
regarded, is the other's possibility: then it would have to be
sophistical as well to point out connection, ground, grounding, etc.

How inwardly the question of possibility hangs together with an altered
apprehension of the axiom of contradiction is seen by setting the
possibility-concept against its oppositions. The possible is opposed to
the impossible, the necessary and the actual. That the possible is the
necessary is of course a contradiction; but if the possible is set
abstractly against the impossible, then it becomes, according to Wolf's
own definition, precisely the necessary: necessary is that whose opposite
is contradictory; the impossibility of the possible is contradictory;
thus the possible is necessary. But the possible can be altered; the
unalterability of the alterable is contradictory, consequently the
possible cannot possibly be the necessary. If the possible is set against
the actual, then it is a contradiction that the possible should be the
actual; but if the possible cannot be actual, how can the actual then be
possible? The contradiction hidden in possibility the ancients have
illuminated from manifold sides; but only Hegel has seen that the
solution of the contradiction lay precisely in the transition, that in
other words the possibility-concept is a transition-concept.

The opposition between possible and impossible is relative; only between
necessary and impossible is the opposition absolute; the opposition
between possible and actual is again another than that between possible
and necessary. Had Hegel been atten%
% --- p. 237 ---
\opage{237}tive to the different strength and the differently
constituted character of the oppositions, he would with his
possibility-forms have found quite other transitions than those now to be
found in his logic; but that he has dissolved contradiction into
possibility and raised the possibility-form to immanent transition-form
must be appreciated. We shall now test the different oppositions.

The opposition between the possible and the actual is, in the newer as
well as in the older philosophy, botched in manifold ways. In that the
dogmatists without further ado identify the possible with the thinkable,
the actual becomes merely a modification of the possible: ``everything in
which no contradiction lies is possible, whether it then exists or not''.
In possibility the difference between actual and unactual is then in this
manner entirely effaced: in possibility the unactual is actual, for it is
indeed actually possible; and the actual unactual, for in possibility it
passes only for a possible. It is the subjective modality that is here
without further ado confounded with the objective. In subjective modality
the possible is indeed the thinkable, and the thinkable indifferent
toward the opposition of actual and unactual. To think the actual is to
transform it into a possible, that is, to apprehend it logically as a
concept; it is, in other words, to abstract from its actuality; for the
actual is in actuality not a concept, it is an existence; it is not
logical but antilogical: possibility is logical, actuality antilogical.

When the antilogical is posited, the logical is sublated; when the actual
is posited, the possible is sublated: the antilogical has been logical,
for existence has been concept, reality has been ideality, power has been
knowledge, the actual has been possible. But the actual as actual is not
what it has been before it became actual. He who says that the actual as
actual is still possible says, after all, that the actual still is what
it was before it
% --- p. 238 ---
\opage{238}became actual; he speaks ambiguously, and easily becomes a
sophist. The ambiguity is that the subjective modality is given out for
objective, the concept for its existence. The concept has not by an
alteration become existence, the logical not by an alteration become
antilogical; it is a manifestation of power, a real reversal, upon which
the transformation rests; it is a manifestation of power, a real
reversal, whereby the possible has become an actual: this is the
character of the objective modality. To reduce the possible and the
actual to indifference in the thinkable is an error which Hegel with all
his dialectic has in common with the dogmatists. The only one who has
been rightly attentive to the opposition between possible and actual is
Kant; but he has not been able to master it.

If the opposition of possible and actual is apprehended in all sharpness,
then the objection lies near ``that since there can be given no even
transition from possibility to actuality, the above commendation of the
Hegelian possibility-concept as an immanent transition-concept must
presumably be retracted''. This is however a misunderstanding. It is true
that the immanent in the transition cannot be apprehended so one-sidedly
that the difference between the subjective and the objective in modality
should vanish in pure objectivity; but the reciprocal determination of
subjective and objective, logical and antilogical, is in the transition
so infinite, the conversion so continuous, that the immanent is seen in
transformation: as transition-concept possibility is essentially a
transformation-concept.

Possibility is a word of knowledge, actuality a word of power; the
opposition of possible and actual is a decisive opposition of knowledge
and power. But then the transition from possibility to actuality is
thereby determined as a conversion of knowledge into power, an incessant
transformation of conceptual determinations into existence-determinations.
Since now, as we have seen from many sides, the same totality of
determinations can be both in the form of knowledge and in the form of
power,
% --- p. 239 ---
\opage{239}the opposition of possible and actual cannot straightway be
referred to a lack in possibility, that it e.g. has some conditions but
lacks others. By one-sidedly emphasizing this distinguishing mark one
could easily arrive at the Megarians' assertion that the possibility
which lacks conditions is, on account of the lacking conditions,
impossible. No, if the opposition of possibility and actuality is itself
to bear the stamp of actuality, then it must take up into itself another
opposition, namely an opposition between possible and possible, between
the one possibility and the other. Where an actuality is possible, its
opposite is also possible: possibility has alternating terms.

The significance of possibility's alternating terms has already shown
itself in anticipation; the further development is conditioned by the
alternating modality of the anticipating functions. The relation between
possibility, anticipation and preformation has been indicated in an
intuitive and popular manner by Leibnitz. In his Theodicy Leibnitz has,
as is well known, started from the presupposition that God, in order to
be able to create the best world, has made a choice among infinitely many
possible worlds. With regard to this Kuno Fischer\footnote{Geschichte der
neuern Philosophie. Zweiter Band. Mannheim 1855. p.~509.} has then
believed he could point out ``the contradiction'' in the Leibnitzian
``world-concept'' in the following manner: ``Of the countless worlds that
are possible God chooses the one he creates. The actual world is for that
reason metaphysically contingent and only morally necessary. But what are
those countless worlds that are possible in the divine understanding? One
thing can be apodictically maintained, that they all consist of monads,
for even the divine understanding cannot represent things to itself
otherwise than as monads. Must not now all monads be ensouled bodies,
representing forces, analogous beings? And the analogous beings must,
different
% --- p. 240 ---
\opage{240}as they are, form a gradual order, a perfect graduated series
in which there is not the slightest interruption. Every thinkable monad
is a thinkable difference, and this must, according to the law of
continuity, be actual. If we now assume that the possible worlds do not
consist of monads, then they are unthinkable, thus impossible in the
logical, and consequently also in the divine, understanding. If we assume
that they are monads which exist only in the divine understanding but not
in nature, then they are plainly \emph{lacking} in nature, then there is
here a \textit{défaut d'ordre}, then the actual universe is no perfect
graduated series, thus not the most perfect, not the best world. Our
syllogism runs: either those countless worlds are impossible, or, if they
are possible, then they must also be actual. In both cases there is
outside the actual world no other world possible''.

Against this criticism, however, objections must here be raised, if not
on the monads' behalf, yet on the possibilities'. From modality's
objective side it is indeed perfectly correct that the possibility which,
either because it has something contradictory in its concept, or because
it lacks and, according to the conditions of actuality, will forever come
to lack a condition, is impossible. But does the matter now also stand
thus when we consider possibility from modality's subjective side? By no
means. In the subjective sense it is so far from being the case that the
contradictory, and to that extent in the objective sense impossible,
should be an unthinkable, thus something impossible in the subjective
sense, that contradiction on the contrary lets itself be thought as a
sharply designated logical concept, and that so distinctly that the
objectively impossible becomes a subjectively possible.

``What does contradiction mean? An existent nothing is a contradiction; a
round square is a contradiction; a flying horse is also a contradiction.
Nonetheless these contradictions are of different kinds. A flying horse
one can represent to oneself, but not think. The unthinkable does not lie
% --- p. 241 ---
\opage{241}directly in the logical, that the predicate ``flying'' contradicts
the concept ``horse'', for this concept could indeed be extended; but in
this, that the whole concept thus extended stands in conflict with the
conditions and laws of experience. It is otherwise with the round square;
here the contradiction is seen immediately in the representation itself.
To represent the square as round is impossible: to think a round square
is, on the contrary, possible, insofar as it is possible to think the
predicate's \emph{conflict} with the subject. The contradiction is thus
not merely this, that the predicate's conflict with the subject makes the
judgment meaningless; but the contradiction is in a higher sense this,
that the contradiction itself, this unthinkable, nevertheless admits of
being thought''\footnote{Philos.\ Propæd. p.~43.}. If the absurd did not
admit of being assumed, if the unthinkable did not admit of being thought
at all, then there would be no insight, no grounding, no comprehending.
Upon what, then, does insight into what is acceptable rest? Manifestly
upon an insight into what is to be rejected. Upon what does the cognition
of the reasonable, of that which is in itself consistent, rest? Naturally
upon a cognition of the absurd, of that which is in itself conflicting
and inconsistent. How is it possible to think the thinkable except
precisely through the opposition to the unthinkable? But then, for the
sake of the opposition, for the sake of the thinkable, the unthinkable
must after all in some manner be thought. So supple is thought that it
thinks with ease its own negative, its own negation.

According to Aristotle one would indeed have to assume that the thought
of the false, the absurd, the imperfect was straightforwardly excluded
from perfect thinking; but this is an insipidity in Aristotelian
dogmatism, an insipidity whose consequence is that the Godhead is not
even able to think the possible. If now the Godhead cannot think the
possible, then neither can it properly think the actual, and what can
% --- p. 242 ---
\opage{242}it then think?\footnote{Grundideernes Logik I. p.~447.} It is
not the case that the thinking subjectivity should be ignorant of the
false, the unthinkable, in the same degree as it is knowing of the true
and the thinkable; on the contrary, in the same degree as insight into
the true rises, in the same degree insight into the false rises also; in
the same degree as the possibilities, the objective possibilities that
are determined to pass over into actuality, are seen through, in the same
degree are the merely subjective possibilities also seen through, those
possibilities that are determined to be excluded from actuality, since in
the objective sense they are sheer impossibilities.

As one of the most telling examples of the difference between the
impossible and the possible, the merely subjectively and the objectively
possible, the imaginary magnitude of mathematics may be adduced. Does the
imaginary magnitude belong to the possible or to the impossible? Since it
is a magnitude, it must naturally belong to the possible; it admits of
being assumed, of being designated, one can indeed calculate with it: it
cannot then straightforwardly belong to the impossible. But the imaginary
magnitude nevertheless stands in a determinate opposition to the real
one; this opposition would lose all significance if there were not in the
real something possible that is impossible in the imaginary. Let the
problem be, e.g., to determine the intersection between a straight line
and a circle. ``If we call the straight line's distance from the centre
$p$, the radius $r$, the chord that is cut off $c$, then we have, to the
question: how great is the chord? as answer: $c = 2\sqrt{r^2 - p^2}$. If
now $p > r$, $c$ becomes imaginary, or, as Newton expresses it,
impossible. The imaginary $c$ thus indicates the impossibility that the
line can cut off any chord from the circle. But why now is this
impossible? Because the line lies wholly outside the circle. But then
there is after all a line here, and we can from another side see
something real
% --- p. 243 ---
\opage{243}connected with the imaginary. For by transforming
$2\sqrt{r^2 - p^2}$, for $p > r$, into $2\sqrt{p^2 - r^2}\,\sqrt{-1}$, one
sees that the imaginary magnitude has a real component, and that this
component is precisely that piece of the tangent to the given circle
which is cut off by a circle tangent to the given line. --- In order now
to see quite clearly how the impossible and the possible are determined
by one another, the imaginary alternating with the real, we will consider
the adduced example a little more closely. Suppose the problem be put in
the following form: There are given two concentric circles, the one with
radius $p$, the other with radius $r$; let there be found the length of a
chord in the outer one which is tangent to the inner one. In the solution
of this problem we must then distinguish between:
\begin{center}
1)\enspace $p < r$, so that one obtains\\
$c = 2\sqrt{r^2 - p^2}$;\qquad (I)\\
2)\enspace $p > r$, so that one obtains\\
$c = 2\sqrt{p^2 - r^2}$.\qquad (II)
\end{center}
\noindent By multiplication with the imaginary magnitude $\sqrt{-1}$, one
can then easily reduce (I) to (II). From this it follows that (I) is
impossible when (II) is possible, and conversely that (II) is impossible
when (I) is possible. This mathematical matter is then at the same time
in perfect agreement with the logical, for $p$ cannot possibly, --- in
such cases dialectic knows very well what it has to concede to formal
logic, --- at one and the same time be both greater (to which (II) gives
the answer) and also (to which (I) gives the answer) less than
$r$.''\footnote{Om Begrebet Nøiagtighed med særligt Hensyn til Hr.\ Prof.\
Steens ``Bidrag''. Efterskrift til ``Philosophie og Mathematik'' af R.\
Nielsen. Kjøbenhavn 1860. p.~26--27.}

With the doctrine of what we shall here call the imaginary and the real
possibilities the doctrine of divine foreknowledge stands in the closest
connection. It has already been adduced under ``the logical
% --- p. 244 ---
\opage{244}ideal''\footnote{Grundideernes Logik. I. p.~235.} in what
sense the old theologians ascribed to God a conditioned knowledge
(\textit{scientia media [sc. de futuro conditionato s. futuribili,
hypothetica, simplicis intelligentiæ], qua Deus perspicit omnia, quæ,
positis quibusdam conditionibus, evenire potuissent}). To this there
attach themselves, naturally in a speculative manner, pronouncements in
speculative dogmatics. ``The contradiction that has been supposed to lie
between the concept of a free course of the world and the divine
omniscience'' --- so it runs in the doctrine of ``the free course of the
world and God's manifold wisdom'' --- ``rests upon a one-sided conception
of omniscience as mere \emph{foreknowledge}, upon a misapprehension of
the \emph{conditioned} in the divine decree \dots\ The final end of the
world's development, with the whole series of moments of development
necessary in and for themselves, must be thought to stand fast in the
eternal decree; but the actual execution of the eternal decree, the whole
fullness of the actual determinations of the course of the world, insofar
as these are conditioned by the freedom of the creature, can also be the
object only of a conditioned foreknowledge, that is, they can be
foreknown only as possibilities (as \textit{Futuribilia}), since other
possibilities may become actual''.

If speculative dogmatics only meant something by what it thus says, if
one could really bring it to lay scientific weight upon its free course
of thought and its own ``manifold wisdom'', it would now be an easy
matter to extort from it the following admission: ``If the
world-possibilities really depend upon the alternation of the conditions,
and the alternating conditions restrict foreknowledge, then the
consequence must be that every foreknowledge of the actual becomes wholly
impossible. On these terms it would not be possible for an omniscient
being to know exactly today what the weather will be tomorrow''. If, on
the other hand, the connection between the
% --- p. 245 ---
\opage{245}conditions is to be conceived thus, that he who knows the
state of the world at a given time can from it predict the alternation of
the conditions for all times to come, then ``the conditioned'' has set no
limit whatever to foreknowledge; but possibility is transformed into a
futuric necessity, the course of the world is unfree. What does
speculative dogmatics now communicate upon this head? Ambiguities!
According to speculative dogmatics ``the conditioned in the decree'' is
to restrict foreknowledge to a certain degree; according to orthodox
dogmatics the doctrine of the world-conditions is on the contrary to
extend the representation of foreknowledge. What orthodox dogmatics
teaches \textit{de Futuribilibus} is indeed not developed into concept,
but the assumption of a \textit{scientia, qua Deus perspicit omnia, quæ
positis quibusdam conditionibus evenire potuissent} nevertheless contains
significant hints toward a logical separation between imaginary and real
possibilities. The possibility which \textit{positis quibusdam
conditionibus} is imaginary can again \textit{positis quibusdam
conditionibus} become real, and conversely.

Since the opposition between the subjective and the objective, the
imaginary and the real possibilities is not absolute but relative, and
since this relativity can on account of the alternating conditions run
through innumerable gradations: it is seen that the reciprocal
determination of possible and impossible, possible and possible, possible
and actual must, in consequence of the two-sided modality of the
anticipating and preforming functions, be manifold. If there are in the
divine knowledge assertoric judgments in an infinite system, then there
is also an infinite system of problematic judgments: there is in the
circle of the possibilities infinitely much more than there is in
actuality. That Leibnitz has let the Godhead determine the actual world
by the exclusion of innumerable possible ones can from the standpoint of
the possibilities by no means be laid to his charge; what can be laid to
his charge on logic's behalf is the manner, the manner in which
% --- p. 246 ---
\opage{246}he lets the anticipation, the choice and the preformation of
``the best world'' take place.

That the Godhead's choice among innumerable possibilities acquires in
Leibnitz a strongly anthropomorphistic tinge derives, to be sure, in part
from this, that in his popular presentation he has had to accommodate
himself to the reader's power of comprehension and as far as possible to
avoid conflicts with the theologians; but, according to the fundamental
character of his philosophy, essentially from this, that representation
is constantly at odds with thinking. While creation is for thought a
fundamental relation, it is for representation a supernatural event. The
choice by which the Creator hits upon the best world forms in Leibnitz a
peculiar boundary-line between possibility and necessity, between the
abstractly subjective and the independently objective.

Prior to the choice subjectivity reigns one-sidedly; the Godhead's gaze
beholds as it were a sea of billowing possibilities, all is in suspense;
for every possibility is, however completely actuality may be
anticipated, a hovering determinateness; perhaps thus, perhaps thus! the
problematic functions are, under the anticipation, in infinite movement.

After the choice, on the other hand, all is changed, the best world is
now completely determined; all its possibilities are at one stroke
transformed into sheer peculiar dispositions, in their peculiarity
unalterable, dispositions that develop themselves with inner necessity.
From every side everything is now determined, calculated in advance and
fixed so punctually that the whole development is a becoming, a
continually unfolding necessity: here there is no room for possibility.

What in Leibnitz should have been developed into concept has remained
standing at a representation. Only representation can want to set from
eternity a finite limit between possible and necessary and then suppose
that the matter is thereby settled once and for all. The concept shows
that the relation between possible
% --- p. 247 ---
\opage{247}and necessary is an infinite reciprocal relation; but since
the opposition of the possible and the necessary is so absolute that the
possible can as little be made into a necessary as the necessary into a
possible, the development of the concept will demand a closer
determination.

\textit{γγ)} \emph{The Apodictic Judgment: Functions of Reflected
Determinateness.} The immediate determinateness, scepticism's
presupposition, turns over through reflection into indeterminateness; the
determinateness reflected in the indeterminateness transforms itself for
intuition into an immediately determinate essentiality and forms the
presupposition for dogmatism. That the reflected determinateness passes
for immediate makes the dogmatic reflection recognizable: different
points of departure condition different results; the conflicting results
introduce indeterminateness anew, every time it is supposed to have been
overcome. As possibility-philosophy dogmatism will develop a system of
objective essence-determinations; but since with all its eye for the
essence of possibility it is not able to see through the form of
possibility, it fixes its essence-determinations as finite
being-determinations, and must halt before the impossibility presented in
the antinomy. To the ontology of the possibilities dogmatism has,
however, furnished characteristic contributions, by seeking on the one
side their origin in the all-determining subjectivity, and on the other
side by demanding the objective connection of the possible with the
actual and the necessary. Although in this way a problem has arisen which
the dogmatic reflection is not able to solve, yet those conceptual
elements have been procured by whose working-up the solution of the
problem will be attained.

It is the form, the possibility-form, upon which everything essentially
depends in the working-up of the elements; but when the emphasis falls
upon form and development of form, possibility, actuality and necessity
will reflect one another as different forms of one form, particular
determinations in one totalized
% --- p. 248 ---
\opage{248}determinateness. Since now the totalized unity of the
form-determinations is itself form, there enters the form-question
characteristic of the modality of the functions, whether it is
possibility, actuality or necessity that furnishes the comprehensive
form-unity, or whether the unity itself is perhaps a still higher and
more concrete form: we shall then investigate the form-unity.

In the development of form possibility, the still indeterminate
determinateness, can be determined as unity of actuality and necessity.
The actual is, in essential opposition to the necessary, an alterable;
the necessary, in opposition to the actual, an unalterable. Each of the
terms thus opposed has at one and the same time its opposition, its
negative, within itself and outside itself; it has and is its own
opposite.

As the form of the alterable, immediate actuality has necessity, the
determinateness of the unalterable reflected into itself, outside itself:
the alterable is not the unalterable, the contingent not the necessary.
What then is actuality in its abstract immediacy, in its pure
contingency? Its positive being is to be determined negatively; the
alterable is the not-unalterable, and has in actuality this negative
determinateness for ground and foundation: the alterable is possible only
through the unalterable.

As the form of the unalterable, necessity is in itself negatively
determined; the unalterable can subsist only as the negation of the
alterable; but if the negation of the alterable is to continue, the
alterable itself must continue. A negation that has nothing to negate is
impossible. If the alterable can no longer be negated, because it has
vanished, the negation will instantly attack the unalterable. The
unalterable is possible only through the alterable.

But that the alterable is possible only through the unalterable is in
other words to say that the unalterable is the possibility of the
alterable; that the unalterable is possible only through the alterable
shows likewise that the
% --- p. 249 ---
\opage{249}alterable is the possibility of the unalterable. With the
doubling that the alterable itself contains the unalterability from whose
negation it proceeds, alterability is expressly determined as actuality;
with the doubling that the unalterable itself contains the alterability
upon whose negation its self-confirmation rests, unalterability is
expressly determined as necessity. That the actual has necessity for its
possibility means that necessity is possibility by means of the actual;
that the necessary has actuality for its possibility means that actuality
is possibility by means of the necessary. In this way, then, the
possibility-form is the form-unity of all forms. \emph{All is
possibility, all is possible}\footnote{Although it is only
form-determinations, pure determinations of reflection, that are dealt
with here, one must nevertheless not imagine that they are without
significance for actuality. What a different conception of the relation
between possibility, actuality and necessity may come to signify for
existence in its most decisive crises, S. Kierkegaard has masterfully
developed in his work ``Sygdommen til Døden''. --- ``The believer
possesses the eternally sure antidote against despair: possibility; for
with God all things are possible at every moment. This is the health of
faith, which resolves contradictions. The contradiction here is that,
humanly speaking, destruction is certain, and that there is nevertheless
possibility. Health in general is to be able to resolve contradictions.
So bodily or physically: a draught is a contradiction, for in a draught
cold and heat are disparate or undialectical; but a healthy body resolves
this contradiction and does not notice the draught. So too with
faith.~\dots\ --- The fatalist is in despair, has lost God and thus his
self. But the fatalist has no God, or what is the same, his God is
necessity; for as with God all things are possible, so God is this, that
all things are possible.~\dots\ In order to pray there must be a God, a
self --- and possibility, or a self and possibility in the pregnant
sense, for God is this, that all things are possible, or that all things
are possible is God; and only he whose being was so shaken that he became
spirit by understanding that all things are possible, only he has
involved himself with God. That God's will is the possible makes it
possible for me to pray; is it merely the necessary, then man is as mute
as the beast.'' Anf.\ Skr.\ p.~35--36.}.

% --- p. 250 ---
\opage{250}If now all is really possible, then it is also possible, by
the modality of the development of form, that the actual can constitute
the uniting middle of possibility and necessity, that in other words the
actuality-form can become the principal form. In its reflected opposition
to actuality the objective possibility is one with actuality; this
contradiction makes itself known in this, that the possible, although it
is not actual, must nevertheless, insofar as it is objective, have
existence in some actual or other, and thus be as an actual: only in the
actual is there an actual possibility. Only the actual possibility is
real, the unactual is imaginary; the imaginary possibility is in the
objective sense impossibility. If now possibility, in order to become
actual in one manner, must already beforehand be actual in another
manner, then it follows from this that possibility is a modal
determination, and thus only a side-determination in the actual:
actuality is immediately itself and possibility.

In immediate unity with its possibility the actual must necessarily
reflect necessity. Necessity is, as we have seen, the possibility of
actuality, and only through this possibility essentially different from
actuality; that actuality is in immediate unity with its possibility
means, then, that it integrates itself by means of possibility, that it
is in unity --- by means of possibility in reflected unity --- with
necessity. Through reflection actuality does indeed lose its peculiar
stamp, insofar as it is to be immediate; but the loss is made good by
this, that reflection sublates itself: reflection's self-sublation in
reflected determinateness is necessity. Through that immediacy with which
reflection begins and is broken off, through the immediate immediacy,
% --- p. 251 ---
\opage{251}actuality has taken possibility up into itself; through that
immediacy in which reflection is completed and therefore itself ends,
through the reflected immediacy posited by reflection, actuality has
taken necessity up into itself. In this way it is then the
actuality-form that is the form-unity of the forms: \emph{All is
actuality, all is actual.}\footnote{``For it is not so, as the
philosophers explain, that necessity is the unity of possibility and
actuality; no, actuality is the unity of possibility and necessity
\dots\ Already in relation to seeing oneself in a mirror it is necessary
to know oneself; for if one does not, then one does not see oneself, but
merely a human being. But the mirror of possibility is no ordinary
mirror, it must be used with the utmost caution. For of this mirror it
holds in the highest sense that it is untrue. That a self looks thus and
thus in the possibility of itself is only half the truth; for in the
possibility of itself the self is still far from being, or is only half,
itself. The question then is how this self's necessity determines it more
closely''. S. Kierkegaard. Sygdommen til Døden. p.~32--33. But however
willingly one must, from existence's and actuality's own standpoint,
concede that actuality is the unity of possibility and necessity, there
lies in this no obstacle to necessity's also being conceived, by a change
of standpoint, as the unity of possibility and actuality.}

It is immediacy, the immediate determinateness, with which reflection
begins; it is immediacy, the reflected immediacy, in which it ends. So
long as the actual can by reflection be converted into the possible, and
the possible again by the returning immediacy become the actual, so long,
in other words, as reflection and immediacy continue to designate an
alternation, the principal form will alternately assume now the stamp of
possibility, now that of actuality, and necessity will be a subordinate
form. But when the immediate itself has become reflected, and the
reflected immediate, the alternation ends, and necessity, the fundamental
form of the unalternating,
% --- p. 252 ---
\opage{252}shows itself. Why can the finite opposition of immediacy and
reflection not maintain itself? Because each of the opposition's terms is
equally alternable. Immediacy, which in the infinite just as in the
finite presupposes reflection, changes about and becomes reflected;
reflection, which in the finite as in the infinite presupposes immediacy,
changes about and is resolved into the immediate: the reflected becomes
immediate. Not even the separation between immediate and reflected
immediacy, a separation with which reflection has exerted itself to the
utmost, is able to maintain itself. At reflection's beginning the first,
and thus immediate, immediacy becomes a reflex in the beginning
reflection, and thus \emph{reflected}; with reflection's completion its
last reflex vanishes in the immediacy that has come to be, and is thus
existent, and is thus no longer reflected but \emph{immediate}.
Reflection may now turn and twist as it will: always it gets the same out
of the different and the different out of the same; always the same and
the same, that which in all alternations is unalternating, the unalterable
that masters all alterations, the necessary. As the absolute unity of the
possible and the actual, necessity is now the principal form: \emph{All
is necessity, all is necessary.}

That now each of the results of the development of form, with all its
many-sidedness, is nevertheless --- on account of the shifting modality
--- itself one-sided, strikes the eye. Were the result: all is actuality,
all is actual, the last and highest, then, as the empiricists assert, the
assertoric judgments would exclusively furnish the contingent foundation
of cognition; the logical functions would in innumerable modifications be
functions of immediate determinateness. Were the result: all is
possibility, all is possible, the highest of the development of form,
then, as the sceptics assume, the problematic judgments would constitute
the groundless ground of cognition, and all logical functions would be
functions of
% --- p. 253 ---
\opage{253}sublated determinateness. Were the result, finally: all is
necessity, all is necessary, the only rational one, then, as dogmatism
presupposes, the apodictic judgments --- axioms, definitions, etc. ---
would have to furnish the essential, fixed foundation of cognition; the
logical functions would then in all possible ways have to be modified
into functions of reflected determinateness.

The results obtained have in a certain manner the same validity. The
apodictic judgment will as little be able to displace or dominate either
the assertoric or the problematic as any of these will ever be able to
make itself the sole ruler of knowledge. The assertoric, the problematic
and the apodictic are in modality three forms under which the validity of
the judgment originally comes forward; but in the logic of the ideas ---
it lies in the essence of logic --- it is a matter of course that the
apodictic form must be the overarching one. For closer elucidation the
different phases of the apodictic are then brought out.

Insofar as the necessity-form arises through reflection's sublating and
arresting itself, the apodictic coincides with the logical: the apodictic
is \emph{logical}. Insofar as reflection can reach immediacy only by the
help of an immediate, and the immediate, in whose reflex reflection comes
to a stand, is as real negation opposed to reflection, the apodictic must
coincide with the antilogical: the apodictic is \emph{antilogical}. If
now the logically apodictic is not to come into conflict with the
antilogical, both forms must become modal determinations of one
absolutely overarching form, the rational: the apodictic is both as
logical and as antilogical \emph{rational}. Thus:

1) \emph{The Logical-Apodictic: essential necessity.} In the Hegelian
logic necessity is treated in three different places: in ``the objective
logic'' under ``actuality'', in ``the subjective logic'' under ``the
judgment'' (das Ur
% --- p. 254 ---
\opage{254}theil der Nothwendigkeit) and under ``the syllogism'' (der
Schluß der Nothwendigkeit); against such a distribution nothing can, in a
systematic respect, well be objected. Characteristic, on the other hand,
is it that while the development in the objective part of the logic is
detailed, acute and fruitful, the treatment in the subjective part is by
contrast short, aphoristic and barren. The ground of this is easy to see;
it is the misapprehension of subjectivity that avenges itself upon this
point also.

When necessity is elucidated by means of the judgment, one would surely
expect to see its significance from the subjective side, and this all the
more as its objective significance has already beforehand passed through
its logical development. But no! The judgment is objective, in the
concept of subjectivity, so objective that it objectively judges itself;
necessity must then present itself objectively a second time. Necessity
was now, as the necessity of the concept, to receive a step-by-step
development through the categorical, the hypothetical and the disjunctive
judgment\footnote{``\emph{Dieß ist die Nothwendigkeit des Begriffs}'',
says Hegel concerning the disjunctive judgment --- ``worin erstens die
Dieselbigkeit beider Extreme einerlei Umfang, Inhalt und Allgemeinheit
ist; zweitens sind sie nach der Form der Begriffsbestimmungen
unterschieden, so daß aber um jener Identität willen diese als bloße Form
ist. Drittens erscheint die identische objektive Allgemeinheit deswegen
als das in sich Reflektirte gegen die unwesentliche Form, als
\emph{Inhalt}, der aber an ihm selbst die Bestimmtheit der Form hat; das
eine Mal als die einfache Bestimmtheit der \emph{Gattung}; das andere Mal
eben diese Bestimmtheit als in ihren Unterschied entwickelt, --- auf
welche Weise sie die Besonderheit der \emph{Arten}, und deren Totalität,
die Allgemeinheit der Gattung, ist''. --- That these remarks are merely
adduced on the occasion of the judgment, and are by no means conceptual
determinations that are developed with inner necessity out of the
peculiar essence of the judgment, is obvious.}; but the remarks that are
adduced on the occasion of the judgment's form are so far from carrying
the development
% --- p. 255 ---
\opage{255}of the concept of necessity further that they on the contrary
serve to divert attention from the main matter. Only by modality can the
determination of necessity show itself to stand in essential dependence
upon the judgment. The dependence is known by insight: in insight the
subjective meets with the objective.

The objective that corresponds to insight is the evident; the evident
never presents itself in a merely immediate light: the light of the
understanding is always reflected. And yet in the reflected light
reflection is sublated, insight is in its reflecting immediate, in its
immediacy reflected. To insight's reflected immediacy there corresponds
the immediately evident; to the immediately evident the immediately
apodictic. The immediately apodictic has at first glance a striking
similarity to the assertoric. ``Water is fluid''; ``the straight line is
the shortest way between two points''; with regard to the immediate in
the assurance both judgments seem to fall under the same category, and
yet there is here a difference of sphere: the first is assertoric, the
last is apodictic. That the last judgment is essentially apodictic is
seen from this, that it is reflection, not sense-perception, that gives
us the certainty. We draw in thought innumerable lines, we compare them,
and now discern, without immediately seeing the innumerable lines, that
the straight one must be the shortest. Insofar as the insight rests upon
a comparing, the light in the evident is reflected. But a proper adducing
of proof nevertheless does not take place; the whole comparison rests
upon an estimation: the proposition is as axiom immediately evident.

That the self-evident rests essentially upon reflection, that no axiom
stands so fast that it must not fall before a dialectical attack, is at
the standpoint of dogmatism a logical mystery. And yet the experiment can
be made with any axiom whatever. ``Two magnitudes which are each severally
equal to one and the same third are
% --- p. 256 ---
\opage{256}mutually equal''; this proposition wavers on the instant, when it
is properly reflected through. Is there then really any contradiction here?
``In order to see what it actually is that is here spoken of, one must
expressly note that what is here spoken of is not things that have magnitude,
but magnitudes. If what were spoken of here were not the magnitudes
themselves, but things that have magnitude, e.g.\ an apple, a triangle, an
hour, etc., then there would certainly not be a trace of contradiction here.
That two triangles which are each equal to one and the same third are
mutually equal is self-evident. So long as we have to do with concrete
magnitudes, difference and plurality can very well subsist together with
equality of magnitude; but when pure magnitudes are in question, the matter
stands otherwise. A plurality of magnitudes presupposes with necessity that
the one must be capable of being distinguished from the other: without mutual
difference no plurality. But by what should pure magnitudes be mutually
distinguished from one another if not by their size? If the magnitudes are
mutually equal, then every glimmer of difference has vanished, and therewith
plurality has vanished as well: that several magnitudes are mutually equal
thus means, in other words, that they are not several, but one and the same
magnitude, which, however many times it be repeated, is nevertheless not
multiplied, but continues to be the selfsame''.\footnote{Forelæsninger over
``Phil.\ Propæd.'' Universitetsaaret 1860--61. p.~2--3.}

Should anyone hold that logical investigations of this kind are altogether
unfruitful, we need only recall how intimately the question of the difference
of pure magnitudes hangs together with the question of the
\textit{principium indiscernibilium}, a question that is decisive in the
assessment of the Leibnitzian philosophy. Should anyone hold that the dialectic
by which dogmatic axioms are shaken must absolutely be sophistical, it must
here be noted that it is not
% --- p. 257 ---
\opage{257}the negative result, but its interpretation and application, by
which sophistry is to be known. When an axiom has been shaken, the sophist
interprets this to mean that it is the apodictic judgment that has been
overturned, and uses the negative result for a play of thought. True dialectic
understands the matter otherwise. That the axiom can be shaken means that the
apodictic rests upon reflection; the wavering of the proposition is the
expression of the fact that reflection has been set in movement. The movement
is conditioned by abstraction. When the abstract proposition is set up as an
axiom, there must always be some accessory determination or other to hold it
in immediacy. Thus mathematics holds the various equal magnitudes apart from
one another by designations, by accessory representations of sums, products
(3 + 3 = 3.\,2 = 12\,:\,2), powers, etc. When now the dialectician lets
abstraction take away the accessory representations in question and sets
thought in unrest, he is fully conscious of what he is doing; for he knows
with himself that the whole operation is meaningless if he does not give back
to the proposition the immediacy that has been subtracted from it and thus
brings reflection to rest again. In this way it becomes evident that the
proposition is an apodictic judgment and, as apodictic judgment, a true axiom.

Sound human understanding relies, and rightly so, upon the immediately
evident; but when a separation is to be made between the assertoric and the
apodictic judgment, the task is not so easy. To test the apodictic by the help
of contradiction acquires significance only when one has first fully tested
the contradiction. How difficult it can sometimes be to point out and to test
a contradiction is seen from manifold examples, and reminds us, as concerns
the exertion of the ancients over their problems, in this connection
especially of the Liar (\textit{ψευδόμενος}). ``When someone says that he is
lying, does he then lie with this his statement, or does he speak the truth?
Here either Yes or No is required as an answer. ``Yes, he is lying!''
% --- p. 258 ---
\opage{258}--- And he says so himself, calls the thing by its right name,
says what is true, and therefore does not lie. ``No, he is not lying, he
speaks the truth!'' --- One must then believe what he says; now he says
precisely that he is lying, so he is not speaking the truth''. Whether Yes or
No be answered here, what is answered here is answered in an apodictic
judgment. The reflection-movement arises from this, that the concepts lie and
truth are incessantly exchanged for one another; they are two apodictic
judgments in mutual conflict, an antinomy in sophistical form.

That there is a sophism here is immediately evident. But when this immediately
evident is now at the same time to yield a corresponding insight, it must be
shown exactly wherein the sophistical consists; only thus can the sophism be
expressly annulled. There are two judgments which are here struck together
into one; a judgment about a judgment is fused with the judgment judged. He
who says: I am lying, judges about a statement which he either has given or
will give, and so far, therefore, about a judgment. If we call the judgment in
the statement x, the judgment about the statement a, it is then seen that the
speaker must apodictically be truthful in the one judgment when he is
apodictically a liar in the other. If the judgment x is a truth, then the lie
lodges in the judgment a; if the judgment a is a truth, then we have the liar
in the judgment x. He who himself says that he is lying must absolutely either
have said something, or in his statement refer to something, in which he is
not lying; he becomes a liar, then, by being truthful, and truthful by being a
liar. --- Instead of dismissing the sophism with contempt, we must rather
admire the ingenuity with which the knot is tied, and take it, not as a sign
of mental derangement, but as a speaking proof of the logical energy of Greek
thinking, when we hear that Chrysippus wrote six books upon this thought-riddle,
and that Philetas of Cos contracted a consumption through excessive exertions
to solve it. The immediately
% --- p. 259 ---
\opage{259}evident must be raised to insight; the apodictic must be seen to
be apodictic.

The assertoric judgment has the problematic outside itself and precisely for
that reason falls so easily forfeit to the problematic; the apodictic
judgment, on the other hand, has through reflection the problematic within
itself, has overcome it and masters it. That the apodictic is evident of
itself shows that the problematic has been overcome; that what is evident in
itself rests in reflection is seen from this, that the problematic, in order
to be overcome, must be presupposed: insight depends upon the transparency of
the problematic sublated in the apodictic.

Nor is it, then, dialectic alone that for the sake of grounding shakes the
axioms and gives the immediate up as prey to reflection; an exact science such
as mathematics does the same. For mathematics it can in several cases be
doubtful whether a proposition is to be set up as an axiom, regarded as a
definition, or proved as a theorem. The theorem and the axiom are both
apodictic judgments; the difference is that the former is mediately, the
latter immediately evident. But this difference thinking can sublate; for the
immediacy in the immediately evident is reflected, and the reflection in the
mediately evident is, through the illumination itself, immediate.\footnote{In
``Elementær Geometri'' by C.\ Ramus, Kjøbenhavn 1850, the proposition about
the straight line between two points is set up as an \emph{axiom}. ``An axiom
(\textit{axioma})'' --- it is said on p.~7 --- ``presents a truth as
immediately given, in that its necessity is evident in itself (e.g.\ only one
line can be the shortest between two points)''. By this we are referred back
to something preceding (p.~5), where it has already been said that the line
which indicates the shortest distance between two points is the straight line.
In ``Elementær Plangeometrie'' by L.\ Oppermann, Kjøbenhavn 1845, the
definition of the straight line is on the other hand given by itself (p.~4),
and the proposition: ``The shortest line between two points is straight'', is
then set up later on (p.~43) expressly as a \emph{theorem} with a detailed
proof.}

% --- p. 260 ---
\opage{260}In the apodictic judgment reflection has equalized the difference
between subjective certainty and objective essentiality. The assertoric
judgment is immediately subjective because it is immediately objective; the
apodictic is reflectedly objective because it is reflectedly subjective. In
what sense the logically apodictic rests upon the determinateness fixed by
reflection is seen from judgments in which the knower reflects upon his own
knowledge. In this way any judgment whatever can be determined by an apodictic
judgment. The blackboard is \emph{actually} black, therefore: the judgment
about the color of a blackboard is \emph{necessarily} assertoric. Chemical
attraction is \emph{perhaps} electrical, therefore: the judgment about the
identity of electrical and chemical attraction is --- at the present stage of
science --- \emph{necessarily} problematic. The objective essentiality of
knowledge is determined apodictically by knowledge itself: there is given, in
formal reflection, an apodictic knowledge of the assertoric and the
problematic just as well as of the apodictic. But the objectivity in the form
of knowledge is one thing, the objective in the content of knowledge another.
As concerns the content of knowledge, it is only the idea, the Absolute
itself, that can give true apodictic certainty.

When it seems as though there were in the relations of the finite much about
which apodictic certainty could be had without direct regard to the idea, it
must here be noted that the apodictic in such cases is conditioned, and the
conditioned in all reflections a reflex of its unconditioned. ``The critical
separation of the pure intuition without concept and the pure concept without
intuition is in its essence untrue; concept and intuition are moments for one
another reciprocally, their dialectical unity is their truth, in the idea they
are one. The critical separation of the categories of quantity, quality,
relation and modality
% --- p. 261 ---
\opage{261}has its significance in the finite, but is, when it is a question
of essence itself, idealess and untrue. To the critique's finite separation
there corresponds an equally finite connection; but the idea is as little a
sum of finitely connected categories as it is a category of reason for itself,
i.e.\ one existing over against the rest. The idea is, as idea, the inner
principial totality of the categories, the truth in the categorical
determinations. The critical separation of a thing in itself without
properties and a phenomenon of properties without a thing in itself is
idealess and untrue; the truth is that the thing is determined by its
properties and the properties by the thing: the infinite reciprocal
determination is grounded upon the idea. The critical separation of a subject
for itself without object and an object for itself without subject is, from
the side of essence, the highest idealessness. To deny the identity of the
subjective and the objective is to disown the idea and to make all cognition
of truth impossible.''\footnote{Philos.\ Propæd. p.~181--82.}

As concerns the visible world, there is an essential difference between seeing
things in the light and seeing the light as light; in the intellectual there
is a still greater difference between seeing truth piecemeal in the true
propositions and seeing the true propositions as moments in truth itself. The
truth in every true proposition, in every apodictic judgment, is a glimmer of
the light of the idea; just as no object can be seen in pure darkness, so
neither can any truth be cognized in pure idealessness: everywhere in the
evident, be it the immediately or the mediately evident, it is an illumination
of the idea that is seen in insight. But dispersed among the many
illuminations, distributed over the many apodictic judgments in their partly
unconditioned immediacy, partly finite conditionedness, the rays of the idea's
light cannot yet gather themselves in the focus in which the idea illuminates
itself; the relatively apodictic still lacks its Absolute.

% --- p. 262 ---
\opage{262}In this connection the relation between mathematical and logical
truths, between the apodictic judgments of mathematics and those of logic, can
be more closely illuminated. Singly and in their special connections the
mathematical propositions bear a far stronger stamp of apodictic validity than
the logical. The logical propositions find themselves, under the streaming
movement, in undulating transitions. This undulating, this dialectical
swinging and hovering, gives them, for the superficial observer, something
wavering, something slippery and untrustworthy, which does not seem able to
subsist together with the demand for apodictic certainty. It is otherwise in
mathematics. The mathematical propositions do indeed also bear the mark of
relativity, insofar as in their mutual connection they depend, the one upon
the other; but in spite of this dependence and reciprocal conditionedness each
proposition is nevertheless for itself so exactly determined, and gives singly
so incontrovertible a certainty, that the apodictic thereby acquires a
punctual omnipresence: it can literally be pointed out at every point. And yet
one notices in mathematics nothing of the idea as idea; with what right, then,
can it be said that every apodictic glimmer is a fulguration of the idea?

In the system of the apodictic judgments of mathematics the idea of knowledge
is present from first to last; but that the idea, despite its apodictic
presence, does not let itself be noticed as idea has apodictically its ground
in the peculiarity of punctuality. The punctuality with which each single
principal proposition is brought out and held fast for itself gives finitude,
mediacy and conditionedness so decisive a preponderance that the inner
infinity must keep itself hidden. Indeed, so foreign can the essence of
infinity be to a mathematician's consciousness that he even conceives the
concept of infinity as a sort of ``\emph{mixed concept}'' and asserts that in
mathematics there is no problem of infinity at all.\footnote{That thinking
mathematicians, however much weight they may otherwise lay upon the satisfying
results of the method of operation, have nevertheless never been able to get
the concept that is inconceivable to their thought (the concept of infinity)
separated in a satisfying manner from the method that rests upon this concept,
is shown by the history of mathematics from Leibnitz and Newton (Cf.\ Mathem.\ og
Dialekt.\ p.~90--91, p.~95 ff.) down to the present time. ``In the first
place, when the principle is unknown, we cannot indeed know whether all the
consequences that follow from it may likewise be admitted: we go as though
blindfold, and it might well happen that after laborious exertion we should at
last stumble upon something which displayed precisely the nature of the
principle, i.e.\ upon something we could not comprehend. In the second place:
is it not a shameful degradation of reason to let it follow a connection of
links, a series whose beginning and end it cannot see through?'' These
scruples of Savrien's sufficiently explain the unrest and tension in which the
analysis of the infinite has constantly held mathematical reflection
(Mathem.\ Dial.\ p.~55--56), whose problem --- as a prize question set by the
Berlin Academy in the year 1784 characteristically enough expresses it, ``wie
aus einer widersprechenden Voraussetzung so viele wahren Sätze haben können
hergeleitet werden'' --- will not cease to return. Forelæsninger over
``Phil.\ Propæd.'' fra Universitetsaaret 1860--61. p.~47--48.}
{} In a logic of the ideas, on the other hand, the
% --- p. 263 ---
\opage{263}single propositions cannot come forward with any special apodictic
independence, since every relative determination is a moment, a glimmer of the
Absolute. But precisely this reflected form of the apodictic moments, a form
transparent in reflection, makes the idea, the absolutely apodictic, visible.

The general principal forms of modality: possibility, actuality and necessity,
recommend themselves by turns to sound sense in such a way that now finitude,
now infinity, now the side of actuality, now the side of the idea is turned
toward the observer. At the stage of finitude it holds that there must
necessarily be something that is merely \emph{possible}, something that is
\emph{actual}, and something that is \emph{necessary}. Seen under the
standpoint of infinity, the difference between possible, actual and necessary
steps so far into the background that the one
% --- p. 264 ---
\opage{264}of the forms involuntarily gains power over the two others, and now
it is said: now that everything is possible, now that everything is actual,
now that everything is necessary. If the emphasis falls upon becoming
(Heraclitus), then everything is possible; if the emphasis falls upon being
and that which is, then everything is actual (the Megarians); if the emphasis
falls upon the essence of actuality reflected into itself, the infinite
substance, then everything is necessary (Spinoza). In the logical sense
necessity is overreaching, but the overreaching necessity is essentially
subjective.

2) \emph{The Antilogical-Apodictic: actual Necessity.} For the closer
determination of the essential or subjective necessity it is necessary to
determine the objective or actual; if the subjective necessity is logical,
then the objective must be antilogical. But in insight necessity must be known
by the apodictic; if there is an antilogical necessity, then there must be a
certainty which is at one and the same time both apodictic and antilogical;
there must, in other words, be an \emph{insight into the antilogical}. All
insight is essentially logical; an insight into the antilogical is then a hard
contradiction: an antilogical apodictic must at first glance appear to the
observer as a \textit{contradictio in adjecto}.

The objective, antilogical necessity is hard, the hard necessity is immediate;
but even in its hardest immediacy necessity is nevertheless essentially
reflected and stands, through reflection, open to insight. It is precisely the
objective necessity that with all its hardness steps mediatingly in between
the logical and the antilogical, and makes an insight into the antilogical
possible. The logical essentiality of necessity reflects itself in the
universal; its hard immediacy is seen in the singular. If the gaze is fixed
exclusively upon the singular, then necessity disappears; it is lost in a
contingency. The contingent is because it is, has its ground in itself, or
rather: it has no ground, because in all immediacy it is an existent ground
and consequently groundless;
% --- p. 265 ---
\opage{265}in this contingency resembles the hard necessity that is blind in
its hardness. The immediately necessary is a contingent, the contingent is an
immediately necessary: here the antilogical is seen. But precisely in that the
hard immediacy antilogically repels reflection from itself, this repelling
will essentially be an attraction. Subjectivity, which cannot reflect the
contingent through, can nevertheless reflect upon it, and thus make it the
object of a reflection.

As object of a reflection the contingent readily shows itself in its
difference from and unity with the objectively necessary. That the contingent
is in itself groundless then means that it has its ground outside itself; that
it is because it is means that it is because an other is; that it is because
an other is means that it, as an outward singular, stands over against another
outward singular. It is otherwise with the necessary. That the necessary is
because it is signifies that it is impossible for it to have its ground
outside itself, since it has its conditions and therewith its sufficient
ground in itself alone. Insofar as the necessary comes forward as the singular
case, it reflects itself in the universal necessity; as a necessary it has in
this reflection the universal, not outside itself, but in itself: the
universal inherent in the necessary cases is power. All manifestations of
immediate power are, in their immediate singularity, sheer cases, in their
outward reciprocity contingent; but power itself, the universal reflected in
the singular and yet in reflection immediate, is necessity, the \emph{actual}
necessity.

How the actual necessity, antilogical in the form of power, stands to insight
can be seen from the apodictic judgment. Upon the antilogical foundation the
apodictic judgment has so striking a likeness to the assertoric that in the
singular cases it is even difficult to keep them apart. ``This body is
actually heavy''; ``the body is necessarily heavy''; in the first case
attention is fixed upon the body
% --- p. 266 ---
\opage{266}in its immediate singularity, and the judgment must become
assertoric; in the latter case attention is fixed upon the corporeal essence
in the body, and the judgment must become apodictic, namely
antilogical-apodictic. The assertoric judgment is determined by the object,
the antilogical-apodictic by power. The object in its peculiarity is a word of
power; but the object's objective meaning, the powerful, universal thought,
makes the power-word a word of knowledge.

The great Kantian question: ``how are synthetic judgments \textit{a priori}
possible?'' can in this connection be paralleled with the question: how are
antilogical-apodictic judgments possible? That the synthetic judgments have a
predominantly empirical, and thus antilogical and so far assertoric character
is presupposed in the Kantian question; if the synthetic did not in the first
instance come forward \textit{a posteriori}, it could not be a problem to find
it \textit{a priori}. That the analytic and so far apodictic judgments are of
an aprioric character is likewise presupposed. Kant asks precisely after the
synthetic, not after the analytic, because the analytic judgments, as Hume has
already brought out, do not, on account of the formal essence of identity,
enrich cognition.\footnote{Analytische Urtheile (die bejahenden) sind also
diejenigen, in welchen die Verknüpfung des Prädicats mit dem Subject durch
Identität, diejenigen aber, in denen diese Verknüpfung ohne Identität gedacht
wird, sollen synthetische Urtheile heißen \dots\ \emph{Erfahrungsurtheile}, als
solche, sind insgesammt synthetisch. Denn es wäre ungereimt, ein analytisches
Urtheil auf Erfahrung zu gründen, weil ich aus meinem Begriffe gar nicht
hinausgehen darf, um das Urtheil abzufassen, und also kein Zeugniß der
Erfahrung dazu nöthig habe. Daß ein Körper ausgedehnt sei, ist ein Satz, der
\textit{a priori} feststeht, und kein Erfahrungsurtheil. Denn, ehe ich zur
Erfahrung gehe, habe ich alle Bedingungen zu meinem Urtheile schon in dem
Begriffe, aus welchem ich das Prädicat nach dem Satze des Widerspruchs nur
herausziehen und dadurch zugleich der Nothwendigkeit des Urtheils bewußt
werden kann, welche mich Erfahrung nicht einmal lehren würde. Dagegen ob ich
schon in dem Begriff eines Körpers überhaupt das Prädicat der Schwere gar
nicht einschließe, so bezeichnet jener doch einen Gegenstand der Erfahrung
durch einen Theil derselben, zu welchem ich also noch andere Theile eben
derselben Erfahrung, als zu dem ersteren gehörten, hinzufügen kann. Ich kann
den Begriff des Körpers vorher analytisch durch die Merkmale der Ausdehnung,
der Undurchdringlichkeit, der Gestalt \&c., die alle in diesem Begriffe gedacht
werden, erkennen. Nun erweitere ich aber meine Erkenntniß, und, indem ich auf
die Erfahrung zurücksehe, von welcher ich diesen Begriff des Körpers abgezogen
hatte, so finde ich mit obigen Merkmalen auch die Schwere jederzeit verknüpft,
und füge also diese als Prädicat zu jenem Begriffe synthetisch hinzu. Es ist
also die Erfahrung, worauf sich die Möglichkeit der Synthesis des Prädicats
der Schwere mit dem Begriffe des Körpers gründet, weil beide Begriffe, ob zwar
einer nicht in dem anderen enthalten ist, dennoch als Theile eines Ganzen,
nämlich der Erfahrung, die selbst eine synthetische Verbindung der
Anschauungen ist, zu einander, wiewohl nur zufälliger Weise, gehören. Kritik
der reinen Vernunft. Werke. Zweiter Band. Leipzig 1838. p.~42--44.}

% --- p. 267 ---
\opage{267}With the Kantian opposition between analytic and synthetic
judgments we further parallel the opposition between apodictic and assertoric.
The assertoric judgments are judgments of actuality, thus judgments of
experience, thus synthetic; the apodictic judgments are judgments of
necessity, thus judgments of reflection, thus analytic. But now it has shown
itself in the doctrine of the form of judgment that the analytic judgment
necessarily presupposes the synthetic, and conversely; thus in every judgment
the aprioric must necessarily presuppose the empirical, and conversely. The
judgment of actuality must, in order to be a judgment, reflect necessity, and
the judgment of necessity, insofar as necessity is objective, actuality; in
that the judgment of necessity reflects actuality, the analytic judgment
becomes empirical, antilogical; in that the judgment of actuality reflects
% --- p. 268 ---
\opage{268}necessity, the synthetic judgment becomes aprioric, logical. The
relation between the empirical and the aprioric in our judgments of experience
Kant determines in the following way: ``all judgments of experience are
empirical, but not all empirical judgments are on that account judgments of
experience. Judgments such as: the room is warm; sugar is sweet, wormwood is
bitter, are empirical; but as mere judgments of perception they have only
subjective validity, and are for that reason not judgments of experience. The
judgment of perception expresses only that two sense-sensations are united at
a given time in me, the sensing subject.'' This is, as we have seen above, a
conception of the assertoric judgment from the immediately subjective side. It
stands otherwise with the judgment of experience. ``What experience teaches me
under certain circumstances, it must teach me and everyone at every time''
\dots\ ``Therefore I utter all such judgments as objectively valid, as
e.g.\ when I say: the air is elastic, then this judgment is in the first instance
only a judgment of perception; I refer two sense-sensations to one another. If
I will that it shall be called a judgment of experience, then I demand that
the connection be placed under a condition which makes it universally valid.''
But in order that an immediate judgment of sense-perception may be able to
become a judgment of experience, it is required that the sense-perception be
subsumed under some ``concept of the understanding'' or other; ``the air,
e.g., falls under the concept of cause; this concept determines the judgment
about the air with respect to expansion as hypothetical. Thereby this
intuition is now represented, not merely as belonging to my sense-perception
of the air in my state or to a state for others' sense-perception, but as
\emph{necessary};'' the judgment becomes universally valid, and thus a
judgment of experience, through this, ``that certain judgments, which subsume
the intuition of the air under the concept of cause and effect, precede and
thereby determine the sense-sensations, not merely respectively by one another
in my subject, but on the whole with respect
% --- p. 269 ---
\opage{269}to the form of the judgment (here the hypothetical) and in this way
make the empirical judgment universally valid''.\footnote{Cf.\ Kant's kleinere
metaphysische Schriften. Werke. Dritter Band. Leipzig 1838. Prolegomena zu
jeder künftigen Metaphysik. p.~215 ff.}

From this it is seen that the Kantian judgment of experience falls between the
assertoric-objective and the antilogical-apodictic: in its unity with the
judgment of perception it expresses the assertoric, in its unity with the
category of the understanding it expresses the apodictic. But in three points
the deviation is discernible.

In the first place Kant must restrict the apodictic to the psychologically
subjective, since the aprioric with him has only psychological-subjective
validity. ``We must'', he says, ``separate the empirical laws of nature, which
always presuppose particular sense-perceptions, from the pure or universal
laws of nature, which, without particular sense-perceptions lying at their
ground, contain merely the conditions for their necessary union in experience,
and in regard to the latter nature and \emph{possible} experience are entirely
and altogether the same, and since the lawfulness in the latter rests upon the
necessary interconnection of the phenomena in an experience, without which we
could not cognize any object whatever in the sense-world, and thus upon the
original laws of the understanding, it does indeed sound somewhat strange at
first, but is nonetheless certain, when I say with respect to the latter:
\emph{the understanding does not draw its laws} (\textit{a priori})
\emph{from nature, but prescribes them to it}.''\footnote{Prolegomena
z.\ j.\ k.\ Metaphys.\ p.~240.}

In the second place the critique of reason mystifies the relation between
``the things of sight (\textit{Phænomena})'', which make up ``the
sense-world'', and ``the separate beings of the understanding
(\textit{Noumena})'', which are supposed to make up ``the world of the
understanding''. For the critique inculcates in the most emphatic manner
``that we do not at all know and cannot know anything determinate about these
pure beings of the understanding, since our pure concepts of the
understanding, as also our pure
% --- p. 270 ---
\opage{270}intuitions, go out only upon objects of a possible experience, and
thus upon mere beings of sense''. Since now ``the beings of the
understanding'' are assumed to be indeterminable, every mediation between the
universal categories and the individual objects, or, more exactly, between the
universal in the categories and the individual in the objects, is lacking.

In the third place the Kantian critique is not in a position to state how the
in-itself objective really stands to sense-perception and thinking; whether,
e.g., it first individualizes and determines the sensible through the pure
intuitions and thus fills out the forms of the understanding; or whether it
perhaps first determines the forms of the understanding and through these the
sense-impressions corresponding to the intuitions; or whether it at one
stroke, and thus equally originally, fills out and determines both the forms
of intuition and the categories. In the first case the unknown objective has
the greatest affinity to the sensible, in the second case to the intellectual,
in the third case it is equally akin to both.

The first deviation is removed in the following way. Reflected in the
ontological subjectivity, the aprioric has not merely subjective, but also
objective validity, not merely validity in knowledge, but also validity in
power. It is thus not the human understanding that prescribes laws to nature;
it cognizes, on the contrary, in nature the objectively aprioric validity of
the laws: to this cognition there corresponds the apodictic in the judgments
of experience. Every judgment of experience can now be both assertoric and
apodictic, according as the corresponding immediacy is \emph{merely immediate}
or \emph{reflected}. Determined by the object in all immediacy, the judgment
is assertoric, is from the subjective side what Kant calls a judgment of
perception (ein Wahrnehmungsurtheil), and becomes, when reflection is set in
movement, easily problematic. Determined, not by the object in its immediacy,
but by the necessity revealing itself in the object-relations, the judgment is
% --- p. 271 ---
\opage{271}apodictic. It is not the law as law that in the actual case
expresses the necessity; for the law as law presupposes a liberated,
hovering reflection in a determinate opposition to the singular case in
which it is fulfilled. Such a reflection does not take place here; the
necessity reflects itself in the case itself, and the case works upon the
insight with striking immediacy. When the table tilts, and glasses and
plates are on the point of falling, the mistress of the house raises the
alarm: ``Take care! it is going over, it absolutely must go over, if you
do not'' etc. Here we have an example of an antilogical-apodictic
judgment. The alarmed mistress of the house involuntarily receives an
impression of the conditioned necessity. She does not indeed yet see the
destruction, for it has not actually taken place; but she does see it
before her eyes, she sees it immediately and yet in reflection, for she
discerns \emph{the necessity}, ``if not'' \dots\ In such an insight there
does not fall first a glance at the general laws of nature, particularly
at the law of falling, then a glance at the conditions obtaining, and
finally an understanding that the law of falling will, in the case that
threatens danger, be sorrowfully fulfilled: such a mediacy, so finite a
parcelling-out, stands in ludicrous conflict with what is momentary in the
insight. And yet, despite the immediate and the momentary, all the moments
indicated by the parcelling-out are present in the reflection; for the
reflection is \emph{momentary}: it is the momentary reflection in which
the insight sees the necessity.

If the momentary reflection is confused with immediate sense-perception,
then the antilogical-apodictic judgment, the necessity-judgment about the
singular actual case, is confused with an assertoric judgment, an
immediate actuality-judgment, and the confusion is in full swing. If there
is no apodictic judgment about the singular actual case; if every
actuality-judgment is without more ado assertoric: then necessity in its
objectively actual presence is once and for all lost. We
% --- p. 272 ---
\opage{272}must then either, with Hume, renounce all actual necessity, or,
with Kant, let the human understanding prescribe its laws to nature.

The defect in ``the particular beings of the understanding'' now likewise
lets itself be removed. What indeed are ``the particular beings of the
understanding'' other than existence-categories, categorical
substantiality-forms, objectively a priori power-concepts, whereby the
thing in existential singularity shows itself as an individuation of its
universal. Such expressions as carbon, nitrogen, boron, silicon and the
like designate at one and the same time categories and existing materials.
As categories the ``particular beings of the understanding'' are
objectively a priori substantiality-forms (\textit{Noumena}); as
categorically determined materials they are existing substances,
sight-things in material immediacy (\textit{Phænomena}). The certainty
that the substance carbon exists in manifold repetitions of actually
existent carbon is antilogical-apodictic. The antilogical is determined by
the immediate carbon-existence, and in this immediacy points to the
assertoric; but the apodictic, which is determined by the
carbon-substance, rests upon momentary reflection and secures the insight
into the universally valid. Such judgments as: ``Carbon occurs in pure
state as diamond and as graphite, less pure as coal, anthracite'' etc.,
are from the standpoint of chemistry antilogical-apodictic. The same holds
of judgments that immediately concern the compounds of the materials.
``With chlorine and fluorine boron forms gaseous compounds which fume
strongly in the air; with nitrogen it forms a white, amorphous,
fire-resistant body which phosphoresces strongly with a greenish light
when it is heated in the air'' etc. If such judgments were not
\emph{apodictic}, they would be without scientific worth; if the apodictic
were not \emph{antilogical}, they would not be actuality-judgments, hence
not experience-judgments, not characteristic of an experience-science.

% --- p. 273 ---
\opage{273}It lies in the nature of the case that an actuality-judgment
which, in accordance with the antilogical immediacy, is assertoric, must
be able to become problematic; to that extent the experience-science does
not give out its judgments about the singular ``sight-things'' in the
singular cases as infallible either. But however many correctives the
judgments about the singular may undergo, the experience-science must
nevertheless of necessity --- hence apodictically --- hold fast the
universal validity of the ground on which the appraisal originally rests:
\emph{the assertoric judgments about the singular are individuations of
apodictic judgments about the universal}. Though the judgments about the
compounds of carbon, boron and chlorine in the singular cases, with the
conditions determined in the singular, may be manifoldly corrected and
altered, it is nevertheless apodictically certain that carbon is always
carbon, boron boron and chlorine chlorine; that is to say: every one of
the existing substances is, under all changes, like itself; the
unchangeable actualizes itself in its changes, the universal is realized
in the individual, the necessary reveals itself in the singular cases.

In this respect we can grant Kant his right when he insists that the pure
concepts of the understanding have no significance, ``wenn sie von
Gegenständen der Erfahrung abgehen und auf Dinge an sich selbst
(\textit{Noumena}) bezogen werden wollen'';\footnote{Prolegom. p.~232.}
but by thrusting out the \textit{Noumena} he makes precisely the unity of
the concepts of the understanding with the objects inexplicable. Concepts
such as: carbon, nitrogen, boron, chlorine etc. cannot possibly pass for
pure essentialities existing in themselves; so far is it from being the
case that the particular substance-categories, Kant's ``beings of the
understanding'', should have independent determinate being as ``Dinge an
sich'' (\textit{Noumena}) over against the existing materials, that the
universal category in its separation from the particular ones is much
rather itself a shadow. It is just the particular that mediates
% --- p. 274 ---
\opage{274}the union of the universal with the individual. The
unity-in-difference of the universal substance (the category) with the
existing substance (the object) by means of the particular substance
(``the being of the understanding''), the unity-in-difference of general
substance with carbon-substance in carbon, of general substance with
nitrogen-substance in nitrogen, of general substance with
sulphur-substance in sulphur, etc., is a necessary unity, one
antilogical-apodictic by its actual necessity. If, on the contrary,
reflection holds fast the subjectively a priori concept of the
understanding for itself, while sense-perception for itself is fettered by
the existing material, and ``the being of the understanding''
(\textit{Noumenon}), the particular substance, is excluded as something in
itself indeterminable, then actual necessity is misconstrued, and with it
the antilogical-apodictic.

The third defect in Kant is only a modification of the two foregoing; its
removal must therefore be regarded as a closer illumination of the
correctives already applied. The necessary unity of the a priori and the
empirical, the universal and the individual, the intellectual and the
sensuous is, categorically determined, a power-unity, and only as such an
actual necessity-unity. The categorical judgment here expresses the
equilibrium between the sensuous and the intellectual; this (a sensuous)
is carbon (a substance-category). Now while the assertoric judgment lays
immediate stress on the sensuous, without giving up the intellectual, the
apodictic judgment lays a decisive stress on the intellectual, without
however, as antilogical-apodictic, giving up the sensuous. Seen in the
light of both judgments, the existence in question is a word of power,
with its actual necessity an element in the categorical power-language. In
the actual necessity, the necessity for the momentary reflection in the
striking case, the intellectual is a sensuous, the sensuous an
intellectual.

It is a great error, an error which does not indeed
% --- p. 275 ---
\opage{275}directly, but nevertheless indirectly, stem from the
abstractions of formal logic, that what is necessary in actuality should
let itself be found out by induction and analogy, hence by a comparative
consideration of several singular cases, although each singular case
yielded only an assertoric certainty. In this way a certain number of
assertoric judgments should then at last be able to lead to an apodictic
one; how thoughtless! Actual necessity is in the judgment determined by an
original insight into the actual itself; the actual itself is, in
antilogical immediacy, the actual case: the apodictic certainty is in
actuality just as original as the assertoric. The reflection by which the
judging subjectivity immediately discerns the necessity in the actual case
is a \emph{momentary} one; from the momentary reflection, determined by
the essence-impression that is immediate in actuality, the
antilogical-apodictic has its origin.\footnote{How induction and analogy
unfold their significance on the ground herewith given will have to be
shown in a later connection.}

3) \emph{The Rational-Apodictic: absolute Necessity.}
As the necessity is, so is the insight; as the insight into the necessary
is, so is the certainty; as the certainty is, so is the apodictic. The
essential necessity, the reflection-necessity, is logically objective and
to that extent ontologically subjective; the actual necessity, the
impression-necessity, is objective, factual and to that extent
antilogical: the unity of the essential and the actual necessity must
therefore be a unity which has in itself the logical and masters the
antilogical; this unity presenting the absolute necessity is the
\emph{rational} one.

``The rational'' is indeed not seldom apprehended so one-sidedly that it
thereby acquires now ``the Irrational'', now ``the Suprarational'', now
``the Extrarational'' as its opposite; but for the idea of science these
oppositions are all vanishing,
% --- p. 276 ---
\opage{276}for the idea of science is absolutely rational. Hegel's words:
``the actual is the rational'', have often enough been misunderstood, and
the misunderstanding again and again refuted. That the actual is rational
by no means says that every fragment of actuality, every singular state,
event or action should directly satisfy reason. On the contrary! There is
in finitude, with its parcellings-out and distortions, infinitely much
that immediately stands in conflict with reason's eternal demand; but
mediately the finite, in all its conflict and confusion, belongs of course
to the actual development of reason's content. What immediately conflicts
with reason must mediately serve its purpose: reason is not merely in the
idea the infinite wisdom, but in actuality the overreaching power by which
the idea subjects all conditions to itself.

In what sense the absolute necessity is the form for reason's content is
seen from the way in which everything, all differences, all oppositions
and contradictions enter into the necessity. Thus the impossible in the
formal sense is a necessary just as well as the possible. For if there
were no boundary between possible and impossible, if everything were
equally possible and just for that reason equally impossible, then there
would be no necessity; but if the boundary is necessary, then the
impossible is --- for the boundary's sake --- just as necessary as the
possible. The unessential is a necessary just as well as the essential.
For if there were no difference between the essential and the unessential,
if everything were equally essential and precisely on that ground equally
unessential, then there would be no necessity; but if the difference is
necessary, then the unessential is, if for nothing else, then for the
difference's sake, necessary. What is necessary must also be essential:
the unessential is thus with necessity essential. In this way one could go
on to infinity.

% --- p. 277 ---
\opage{277}The interplay of forms under which necessity confirms and again
dissolves every relative determination has in actuality not merely formal
but also real validity. There must in the actual itself be something
unessential, in order that the essential may be. Insofar as the
unessential in this way becomes essential, it does indeed seem as if the
actual boundary between essential and unessential were effaced; but this
is only an illusion. The necessity that gives the unessential an essential
significance, the same necessity demands that the unessential shall
acquire significance precisely as an unessential. Still more sharply, and
in a more existential way, the necessary opposition between the necessary
itself and the contingent comes forward. It is e.g. impossible to drink
\emph{a} wine without drinking \emph{this} wine. But why just this one? It
might just as well be either this or this \dots\ or this! Which
``\emph{This}'' one chooses is, insofar as it merely comes to choosing a
wine, altogether indifferent: a ``This'', just a ``This'', is thus
contingent. But whichever one chooses, it is nevertheless altogether
unavoidable that one comes to choose a ``This''; a ``This'', just a
``This'', is thus necessary. The insight that the contingent is not
necessary is rational only in union with the insight that the contingent
is nevertheless necessary. That the contingent is necessary does not then
say that the opposition of necessary and contingent is to be given up, but
on the contrary that it is to be maintained. In this way every relative
determination has its determinate validity, i.e. its necessity: in the
totality of the relatively necessary moments the absolute necessity
consists; the absolute necessity is the form for the infinite reason.

Should anyone now ask whether one must not then, for reason's sake, demand
something irrational in the same way as one must, for the essential's
sake, demand something unessential, or for the necessary's sake something
unnecessary: then the
% --- p. 278 ---
\opage{278}answer does not lie far off. There must indeed be a difference,
an essential difference, between rational and irrational, just as between
necessary and unnecessary; but just as the difference by which the
unnecessary itself becomes necessary falls under the absolute necessity,
so too the necessity with which the irrational --- however much it
conflicts with the rational --- announces itself falls under the absolute
reason. Seen in the idea, the absolute reason can as little have anything
irrational against it as the absolute necessity can have anything
unnecessary outside it. Nor can it do, on the standpoint of the idea, to
give the rational an opposite in the antirational in a manner
corresponding to that in which the antilogical is set against the logical.

The opposition of the logical and the antilogical means that the logical,
however much it shows itself to reflection as an infinitely overreaching,
is nevertheless in its essence one-sided, just as pure knowledge is a
one-sidedness over against power. It is otherwise with the rational; here
knowledge is in unity with power, for power is in its essence just as
rational as knowledge.\footnote{Only on the standpoint of faith, hence in
a sphere that falls outside science, can there be talk of the
antirational; the antirational then becomes the distinguishing mark of the
paradoxical.} The infinite reason determines the logical by the
antilogical, and conversely. The certainty of the validity of the idea,
i.e. of the absolute reason, is rational-apodictic. Since now the
antilogical has its rational validity just as well as the logical,
actuality too in its difference from possibility and necessity,
possibility in its difference from necessity and actuality, necessity in
its difference from possibility and actuality, has rational-apodictic
validity in the infinite kingdom of reason. The assertoric judgments can
in the idea as little be transformed into apodictic ones as the
problematic can; but all three kinds of judgment hold with relative
% --- p. 279 ---
\opage{279}necessity; for the perfect knowledge the validity of them all
is rational-apodictic.

In order to discern how it is possible that systems of assertoric,
problematic and apodictic judgments can have validity not merely outwardly
in the objectification of the finite, but also inwardly in subjectivity's
self-objectification, we must above all remember that the ontological
self-objectification, which everywhere rests upon an objectification of
other, is not empty, but infinitely rich in content. The content is reason
itself, the idea's absolute content. Reflected in modality's three forms,
the idea, and with it subjectivity, will then assume the following
fundamentally different shapes. In relation to actuality, hence in the
system of the assertoric judgments, subjectivity is intuiting, and the
idea in actuality's form an idea of the phenomenon. In relation to
possibility, hence in the system of the problematic judgments,
subjectivity is self-determining, and the idea of the predetermined
possibilities a purpose-idea. In relation to necessity subjectivity is
thinking, comprehending, and the idea in necessity's form a universally
valid essence-idea. Since the idea is under all three shapes identical
with itself, it is a matter of course that modality's three forms must
presuppose one another in such a way that the one, by its determinate
appearance, always reflects the two others.

As regards the idea of the phenomenon, an immediate separation between the
phenomenon and the idea is presupposed in actuality. Here, in the actual
world, there must be able to be phenomena which fall contingently outside
the idea, as distinguished from phenomena in which the idea immediately
lets itself be seen. Precisely by the fact that the idea-phenomenon in its
ideal transparency steps into opposition to the actuality-phenomenon that
in its contingency is idea-less, there forms itself an assertoric judgment
about the beautiful, corresponding to the idea-phenomenon, hence to an
ideal form: ``this is actually beautiful'', and an assertoric judgment
about the unbeautiful,
% --- p. 280 ---
\opage{280}corresponding to the idea-less phenomenon that conflicts with
the ideal form: ``this is actually unbeautiful''. The conflict between the
idea-phenomenon and the idea-less actuality-phenomenon is equalized in
possibility: both phenomena can be altered in countless ways, and as the
idea, which by means of the possibilities masters, fulfils and penetrates
the idea-less actuality-phenomena to infinity, it satisfies the
subjectivity that intuits itself in all and all in itself. If the judgment
about the beautiful and the unbeautiful is, on account of the phenomenon's
immediacy, assertoric, then the judgment about beauty itself, i.e. the
judgment by which the judgment about the difference, the essential
difference, between beautiful and unbeautiful is originally fixed, is on
the contrary apodictic, rational-apodictic.

As regards the purpose-idea, the idea of the determinate possibilities, it
is self-evident that the all-capable subjectivity, by means of its power
over otherness, hence over immediate actuality, must live in an infinite
system of possibilities. But an infinite system of possibilities is only
thinkable in an infinite system of problematic judgments. ``God is this,
that all things are possible''; that means here: for the absolute self
nothing finite stands fast; all things are seen in infinite becoming, in
transitions, in possibilities. Through the infinite system of the
possibilities the divine subjectivity asserts its infinite
self-determining by a particular determining of other: it transforms the
actual into conditions and posits the possibility determined by the
anticipation as purpose. The rational purpose can thus be realized on
account of the possibilities; but it can in actuality also be partially
disturbed, distorted and hindered on account of the possibilities. The
judgment about the attainment of the purpose, of the good, is in finitude
conditioned and problematic; but in the idea itself rational-apodictic.
Just as the assertoric judgments on actuality's ground are individuations
of the apodictic that corresponds to the necessity of the phenomenal idea,
so
% --- p. 281 ---
\opage{281}too are the problematic judgments on possibility's ground
specifications of the apodictic that in divine insight corresponds to the
necessity of the purpose-idea.

As regards finally the essence-idea, the universal necessity-idea, the
thinking subjectivity, comprehending all in the unity of knowledge and
power, will objectify to itself all, the possible as well as the actual,
the eternal as well as the temporal, the inner and ideal as well as the
outer and real, in the word, by the word, as the word. The only thing
absolutely necessary for the absolute reason's self-revelation is the
word; in the word all is brought forth, by the word all is determined: the
universal necessity-idea is the idea of the word. What then is the object
of the judgment? A word of power! And the judgment itself? A word of
knowledge! And the judgment's validity? A word of actuality, of
possibility and of necessity! Or how should the identity of the judgment
and the judgment's object, and therewith the truth, be able to become
evident, if the object, despite its actual opposition to the judgment,
were not in the word itself essentially one with the judgment. In the word
as the word of power the intuiting subjectivity, with all power-actions
and judgments, enters upon the objective, antilogically ruling
actuality-necessity; in the word as the word of knowledge the thinking
subjectivity goes, through all actual oppositions, into the
logically-objective, ontologically-subjective essence-necessity; in the
reflected unity of the word of power and the word of knowledge the
all-comprehending subjectivity enters upon the absolute necessity, a
system of rationally-apodictic judgments in which all has its relative
truth because all has its relative necessity.

The judgment in and for itself is not the truth, for in the truth
knowledge must presuppose power; the object in and for itself is not the
truth, for in the truth power must presuppose knowledge. If now the word
in the judgment for itself is a word of knowledge, and the word in the
object for itself a word of power, how
% --- p. 282 ---
\opage{282}then is the unity of this word, doubled in itself and by the
doubling originally opposed to itself, to be more closely determined?

\begin{center}{\bfseries b)\quad Functions of the Syllogism: Modal Functions.}\end{center}
\addcontentsline{toc}{subsubsection}{b) Functions of the Syllogism: Modal Functions}

If one lets the judgment about an object sink down to simple
representation, and determines the representation not by a word but
immediately by the object's image, then one should at first glance suppose
that the agreement between representation and object would in this way be
secured, and the agreement thus secured the best guarantee for the
judgment's truth. But that a criterion of truth does not let itself be set
up on these terms, the ancient Skeptics have, particularly by their
objections to the Stoics' \textit{φαντασία
καταληπτική},\footnote{The definition of this otherwise somewhat ambiguous
category is given in Sextus (\textit{Adversus Logicos I. §~402.}) in the
following terms: \textit{ἤν γὰρ καταληπτικὴ φαντασία ἡ ἀπὸ ὑπάρχοντος, καὶ κατ᾿
αὐτὸ τὸ ὑπάρχον ἐναπομεμαγμένη καὶ ἐναπεσφραγισμένη,
ὁποία οὐκ ἄν γένοιτο ἀπὸ μὴ ὑπαρχοντος}, that is: the representation
which has its origin from the existing thing, and in perfect agreement
with this is so impressed and imprinted in the soul that it cannot
possibly derive from a non-existing thing. For brevity's sake we here use
the designation: the grasping representation.} sufficiently shown. The
discussions that Sextus Empirikus has preserved for us concerning this
point are in a high degree characteristic. ``There is'', it is said, ``no
representation grasping the object so peculiar that it may not derive from
another object than the one to which it is just in a given case referred;
this, besides the Skeptics, even the Academics, and then chiefly
Arkesilaus and Karneades, have striven from many sides to establish. The
assent (\textit{συγκατάθεσις}), which the Stoics call apprehension
(\textit{κατάληψις}) and grasping representation (\textit{φαντασία καταληπτική}),
is, according to Arkesilaus, only an empty word; for such an assent, if it
is given,
% --- p. 283 ---
\opage{283}must either be given by the wise man or by the fool. But if it
is given by the wise man, it is a knowledge (\textit{ἐπιστήμη}), and if it
is given by the fool, it is only an opinion (\textit{δόξα}). If the assent
we bestow on the grasping representation is to be the apprehension, then
this is nowhere at home; for in the first place it is not the
representation but the thought on which we bestow our assent; in the
second place there is no representation so true that it could not also be
false''.\footnote{\textit{Sext. Empirici Opera ed. Fabricius. Lipsiæ. Anno MDCCXVIII.
Adversus Logicos. I. §~143--54.}} To distinguish the true representation
from the false is not possible; for to wish to undertake such a
distinction by saying with the Stoics: the true or grasping representation
derives from an existing object, the false or non-grasping from a
non-existing one, is in vain, since there are delusions which in clarity
and vividness are not inferior to the true representations corresponding
to actuality. ``In his frenzy Herkules took his own children for
Eurystheus's, and acted in consequence of this delusion of his, i.e. in
the delusion he killed his enemy's children, but in actuality his
own''.\footnote{\textit{Adversus Logicos. I. §~406.}}

But even if one starts from the presupposition that the representation in
question has its origin from an actual object, it will nevertheless, the
Skeptic assumes --- and that from his standpoint, so long as reflection is
here directed merely upon a finite comparison between representation and
object, with perfect right --- be impossible to point out a representation
that should give so distinct an image of its object that by this image one
would be in a position to distinguish the object from all others. For the
forms can be quite alike, and the objects themselves nevertheless
different. ``I give'', says Karneades, choosing the example of two eggs
that are quite alike one another, ``the Stoic leave to try it alternately,
now with the one
% --- p. 284 ---
\opage{284}egg and now with the other, how far the wise man, when he sees
them, will be in a position to say which is
which''.\footnote{\textit{Adversus Logicos. I. §~409.}}

If one attempts to evade the difficulty by such evasions as that it is
after all not a matter of stating universal determinations which should
contain distinguishing marks for every singular case, that objects which
are so alike one another that they cannot be distinguished in the
representation itself by no means thereby lose the character of being
actual, etc.: then the Skeptic makes with great ease a leap from the
immediately singular to the abstractly universal; for his skill is, as we
have earlier seen, to leap from extreme to extreme in such a way that he
constantly leaps over the connecting middle term. If the difficulty in the
extreme of the singular was that the Dogmatist with his ``grasping
representation'' could not distinguish the one egg from the other, the
difficulty in the extreme of the universal now becomes that his ``grasping
representation'' can only be determined by the object which itself must be
determined by the representation.

With respect to the universal in the reciprocal relation of representation
and object, the Academic's objection runs that the Stoics make themselves
guilty of a diallelus. ``When, namely, we ask what a grasping
representation is, they say: such a one as derives from an existing
object, agrees with this, is stamped and imprinted in the soul, and of
such a constitution that it cannot derive from a non-existing
object.\footnote{\textit{Adversus Logicos. I. §~426--29.}} But since now
every determination must start from what is known, and we therefore ask
further: what is the existent? they return to the first determination and
say: existent is that which calls forth a grasping representation. Thus
the determination of the grasping representation points back to the
determination of the existent, and
% --- p. 285 ---
\opage{285}this again to the former; we thus obtain from neither the
object-representation nor the object any distinct concept. How far the
represented object has a merely apparent or both an apparent and an actual
validity cannot therefore be decided by any representation whatever; if
objects are to be appraised according to the non-grasping representation,
then this must --- what the Dogmatists themselves will not concede --- be
able to appraise the true. To take the grasping representation for the
something \emph{according to which} the true is to be appraised is
ridiculous; for since a determinate distinguishing mark for the grasping
representation is precisely what is sought, there must, for every
representation that is set up as grasping, constantly be required a new
representation for the confirmation of the first, and so on to
infinity''.\footnote{Cf.\ E.\ F.\ Nielsen. Den antike Skepticisme. p.~91--97.}

To remove these difficulties is impossible so long as the standpoint is
chosen exclusively in the finite; but in order that the infinite may come
along, the representation must be determined by the judgment, and the
opposition of judgment and object cleared up in the word. This is the
train of thought we have followed in the doctrine of the
judgment-functions. It begins with the nature-powerful thinking forming
and shaping language, so that in the developed language-forms it can know
and recognize itself. In the course of this development the subjectivity
working with language and thoughts has not merely objectified to itself an
actual outer world, but, since language presupposes and conditions an
inner ideal objectification, an inward world as well, a world of essential
content.

The essential thought-content articulated by language is in principle just
as originally objective as subjective; the identity of the subjective and
objective content is in the essence itself \emph{the idea}, but in the
form's immediacy \emph{the word}: the language- and thought-forming
functions have in this way become idea-forming. That the idea-forming
functions, from their
% --- p. 286 ---
\opage{286}peculiar point of departure begin with an immediacy as much a
force of nature as the language-forming ones do, is seen from the
formation of myth. The essential content of mythology exhibits moments
of logical originality, determinations of form in which an eternal
reason has found itself mirrored; we have seen the ideal results to
which the sublation of the mythical stage of consciousness can lead
with inner consequence. True, the Gnostic idea-functions could not, by
an ingenious union of imagining and thinking, give the absolute any
adequate form; true, the theosophical functions, which let thought
grasp the absolute absolutely only by slipping under it the form of
imagination and making the form formless, had to lose themselves in the
mystical; but at these transitional stages it nevertheless showed
itself ever more clearly what it means that thought grasps the idea by
the word. The mythical gods, the Gnostic aeons, the theosophical
essentialities are different idea-revelations, which are determined by
thought through the word, and which live and grow in the word.
Clarified into concept, the idea then also reveals itself at last as
the divine essence of the word, and the word as the eternal shape of
the idea.

The identity of word and idea is at the same time an identity of
judgment and content; the judgment absolutely is the word, the content
absolutely is the idea. But the identity was grasped only from the side
of form; its universal essentiality is as yet without actuality. From
the standpoint of actuality, then, one could not begin with a further
development of judgment and content; one had to begin with the sharp
opposition of judgment and object. When the object is set up
independently over against the judgment, one might almost fear that the
object existing in outward contingency would withdraw from every
determinate and reliable apprehension, so that the old quackeries with
\textit{φαντασία καταληπτική} would have to return. This awkwardness is
removed by the fact that the judging subjectivity, in an overreaching
manner, lets the judgment itself be determinative for the relation of
judgment and object. The most immediate independence that can be
% --- p. 287 ---
\opage{287}granted to the object is recognizable in this, that it is the
object which is assumed to determine the judgment; the consequence of
this assumption is that everything dissolves, both the
thought-determinations in the judgment and the reality-determinations
in the object. If, then, it cannot be the object that determines the
judgment, the question arises whether it is not perhaps, conversely,
the judgment that determines the object. On this assumption the
thought-determinations do indeed come into force, but so overreaching
is the power of ideality that the object succumbs, is dissolved into
concept and transformed into the content of the judgment. The object's
immediate independence over against the judgment has thus in a double
way shown itself to be a lack of independence: if without concept it is
to rest upon itself, it becomes a loose aggregate of properties without
substance, crumbles away into contingencies and collapses; if in virtue
of the concept it is to rest upon thought, and thus upon the judgment,
it is transformed into a phantom and vanishes like vapour.

If the object is to remain in force, there must be an ideal power in
the judgment; the ideal power in the judgment as logical judgment
presupposes a real power in the acting that corresponds to the
judgment, for originally knowledge presupposes power, and power
knowledge. The object in its immediacy is not concept, but its
existence has its ground in the concept; without ground no existence,
without concept no object: as the ground is, so must the existence be;
as the concept is, so must the object be. The object must thus have
ideal determinate being in its concept in order to acquire actual
existence in existence: the object must be anticipated; in the
anticipation the judgment reaches out beyond the object, just as
subjectivity reaches out beyond the objective. That the object comes to
be on objective conditions by no means signifies that it is torn loose
from the concept or arises alongside the concept through the
co-operation of thing-causes; no, on the contrary! The realization of
the object is precisely the realization of the concept: it is
% --- p. 288 ---
\opage{288}again the judgment that is overreaching; it is in virtue of the
antilogical judgment that the object is realized. In the antilogical
judgment thought is power, in the logical power is thought; and since
the expression of thought is originally the word, and the word is
articulated in the proposition, the object comes to be essentially
through a thought-proposition, actually through a power-proposition.

As a power-proposition reflected in a thought-proposition the object,
seen in the light of the concept, is more than an immediate thing; in
essence it is an objective judgment, a word, a word of power: judgment
and object stand to one another as the word of knowledge and the word
of power. In this unity-relation the opposition is preserved; for
object stands over against judgment as power over against knowledge.
But in the opposition the unity too is preserved; for the object is a
power-proposition, just as the judgment is a thought-proposition. The
real thought that is in the power-proposition corresponds in principle
to the ideal power that is in the thought-proposition; thus to infinity
the object is determined by the judgment, and the judgment again by the
object. Judgment and object are one in the word, since they are one in
the idea.

The unity of judgment and object is a unity-in-opposition; from this
unity-in-opposition the power of modality shines forth. The categories
actuality, possibility and necessity point not merely to three
different sides of existence, but they show at the same time, when we
look more closely, that these three sides are three totalities, so that
the standpoints alternate with the three judgments: everything (namely
everything existent) is \emph{actual}, \emph{everything is possible},
\emph{everything is necessary}. The three totalities thus form three
different spheres, on account of modality; but the three different
spheres are nonetheless side-determinations of one and the same
existence, on account of modality: judgment and object thus enter into
three different relations on account of modality. By this retrospect we
shall illuminate the modality-relations somewhat more closely.

% --- p. 289 ---
\opage{289}At the stage of actuality the object stands, although
essentially itself an objective judgment, in immediate independence
over against the subjective judgment; the judgment immediately
determined by the object is assertoric. If actuality so excluded
possibility that the latter arose only within subjectivity and obtained
a subjectively infinite play over against the former, then the whole of
logic would fall prey to skepticism, and the old game with
\textit{φαντασία καταληπτική} would repeat itself in new turns. But now
possibility, by virtue of modality, has objective being just as much as
actuality: ``everything is possible'' --- this is not an abstract
subjective judgment that has excluded objectivity, but subjectivity's
original judgment about everything objective. What is already
determined does not, insofar as it is determined, admit of further
determination: only what is still undetermined can be determined. If it
now stands fast, as was shown above, that no determination of existence
is thinkable without a corresponding determination of concept, no
determination of the object without a determining thought; if it stands
fast that the ground of everything objective must in principle be
determined subjectively and can be determined subjectively only insofar
as the objective is still problematic, then the overreaching
subjectivity must know the origin of all finite determinateness as
problematic: ``\emph{everything is possible}''; the problematic
judgment is not merely human, but divine.

In the problematic judgment the object is dissolved; it has gone under
in indeterminateness, been buried in uncertainty; the problematic
judgment does indeed have a content, but no object.\footnote{That at the
standpoint of finitude a problematic judgment can be passed upon an
object determinate in itself follows from this, that here, besides the
judgment, there is a representation present which places the object
itself outside the judgment.} If the judgment in the original sense is
to have both object and content, an object corresponding to the
determinate content, a content corresponding to the determinate object;
if, in
% --- p. 290 ---
\opage{290}other words, the unity-in-opposition of judgment and object is
to attain completed determinateness, then modality itself must be
completed in necessity, and the apodictic judgment become the
overreaching one. In virtue of necessity the identity of the judgment's
content and object is determined in two essentially different ways: the
identity is posited in the \emph{content}, and to this corresponds the
logically apodictic; the identity is posited in the \emph{object}, and
to this corresponds the apodictic in antilogical form. Both ways yield,
with all their opposition, two sides of the same. From the antilogical
side necessity is hidden in actuality, and still more determinately, in
contingency. If anyone will maintain that necessity, by coinciding with
contingency, must either make the contingent necessary, and thus
sublate contingency, or make the necessary contingent, and so, which is
self-contradictory, abolish itself: then he could surely profit by
considering whether the judgment ``the contingent is necessary'' does
not assert the significance of contingency just as fully as the
judgment ``the necessary is contingent'' asserts the significance of
necessity.

If necessity is present everywhere, in all cases and in contingent
phenomena, without however --- as the judgment ``the contingent is
necessary'' forbids --- transforming the contingent into semblance:
then necessity, by yielding to the case, respecting actuality and
itself assuming all the shapes of contingency, has precisely drawn
everything contingent under itself, and has thereby proved its infinite
superior power.

As the actual now stands to the necessary, so the assertoric judgment
stands to the apodictic. If anyone will deprive the assertoric judgment
of its peculiar validity and transform it into a reflex of the
apodictic, then necessity protests; if anyone will make the assertoric
judgment absolutely independent of the apodictic, then he precisely
cognizes the necessary in its peculiar character, baptizes it, so to
speak, with the baptism of necessity, and
% --- p. 291 ---
\opage{291}in this way lets the apodictic vouch for the assertoric.

As for the opposition of the problematic and the apodictic judgment, it
is at once evident that it coincides with the opposition of possibility
and necessity. Since possibility as possibility has its necessity,
possibility is itself a latent necessity, and from this it follows that
in the problematic judgment there is an apodictic moment, and that this
apodictic moment, far from annihilating the problematic, on the
contrary vindicates for it an ontological right, a validity divine in
the idea. The judgment offensive at first glance, \emph{everything is
necessary}, can thus be cognized as the logically overreaching one, and
the apodictic knowledge corresponding to necessity be held fast as the
ideal of logic, without the influence that necessarily belongs to the
assertoric and the problematic suffering in the least thereby.

If the articulating reciprocal determination of assertoric, problematic
and apodictic is fully acknowledged in the determination of judgment
and object, then it is seen from this as well what the identity of the
judgment's content with the judgment's object, seen from the standpoint
of necessity, originally signifies. In virtue of necessity the
judgment's object is identical with its content; for the apodictic
judgment is in its essence logical, and every determination of the
object is, in the logically apodictic judgment, transformed into a
thought-determination and thereby into a determination of content. But
in virtue of necessity the judgment's content is then also identical
with its object; for the apodictic judgment is in actuality
antilogical, indeed in the phenomenon assertoric, and every
determination of content a reflex of a corresponding determination of
the object.

It is in consequence of the overreaching necessity with its reciprocal
determinations of logical and antilogical that the all-capable
subjectivity transforms, by categorical power-language, a world of
actual objects into a world of objective
% --- p. 292 ---
\opage{292}judgments. How the concept thing stands to the concept
property is determined by a thought-proposition, thus by a
\emph{subjective} judgment; the relation between the existing thing and
the existential properties, on the other hand, is originally determined
by a power-proposition, a practical-antilogical proposition, thus by an
\emph{objective} judgment.

The investigation of the relation between judgment and object has thus,
with inner consequence, transformed itself into an investigation of the
relation between judgment and judgment. This investigation proceeds in
three directions; for what matters now is: partly to determine the
relation between judgment and judgment in the subjective sense, and
thus to gain insight into the connection of the subjective judgments;
partly to determine the relation between judgment and judgment in the
objective sense, and thus to gain insight into the connection of the
objective judgments; and finally to determine the relation between
judgment and judgment in the subjective-objective sense, and thus to
gain insight into the judgments' original connection in the idea.

To determine the relation between judgment and judgment is, so far as
the connection is concerned, to determine which judgment is the
original one and which the derived, that is, in other words, to derive
judgment from judgment, and thus to \emph{syllogize}: all syllogistic
functions are in the logical sense connection-functions. That the
category of connection forms the foundation is immediately discerned
from this, that the same judgment can, without anything being altered
either in form or in content, occur both as premise and as conclusion;
the difference between premise and conclusion is decided by the
connection. From this the relation between modality and connection is
seen as well; if the judgment is premise in one connection and
conclusion in another, then it is apprehended in the first connection
in a different way from the last; as connection-functions the functions
of the syllogism are categorically determined modal functions. Since
everything here depends on the connection, everything at the same time
depends on the mode, and thus on the mode of
% --- p. 293 ---
\opage{293}connection; for the mode determines the connection and the
connection again the mode.

In the categorical peculiarity of the modal functions lies the ground
of their division. To the subjective syllogism, the determination of
the mode of dependence of the subjective judgments, there correspond
functions of subjective connection, \emph{modal functions of
knowledge}; to the objective syllogism, the determination of the mode
of dependence of the objective judgments, there correspond functions of
objective connection, \emph{modal functions of power}. To the
idea-syllogism, the determination of the mutual dependence of the
subjective-objective judgments in reflection of the absolute itself,
there correspond functions of subjective-objective connection,
\emph{modal functions of truth}.

\medskip
\noindent\textit{α)} \emph{Subjective Syllogism: Modal Functions of Knowledge.}
\addcontentsline{toc}{paragraph}{α) Subjective Syllogism: Modal Functions of Knowledge}
\medskip

For a survey of the inner connections and separate modes of dependence
of the subjective judgments it is necessary from the outset to give
oneself an account of the point of departure, the goal and the method.
The point of departure, or more precisely the central point, the point
about which the whole development turns, is subjectivity itself, the
thinking, judging consciousness. The goal, the fundamental task to
whose solution the modal functions co-operate in their various
alterations, is an articulation of the content of the knowing
subjectivity. Since the content, apprehended as thought-content,
consists in a manifold of judgments, the knower can agree with himself
and possess the riches of his content in full transparency only insofar
as the manifold judgments are connected immediately or mediately into
unity, or at least insofar as they show themselves mutually consistent
through the peculiar mode of connection. With the task thus determined,
the procedure, the logical method, is at the same time indicated. Since
the mode of connection and of dependence is known from the syllogism,
the subjective syllogism will naturally assume a different form
according as the mode of
% --- p. 294 ---
\opage{294}dependence is determined partially by the connection between
two, between three, between several, or totally by the connection
between all judgments. If the mode of dependence is determined by the
connection between two judgments alone, the syllogism is
\emph{immediate}; if it is determined particularly by the connection
between three or more judgments, the syllogism is \emph{mediate}; if it
is determined finally by the original connection, which transforms the
judgments separated in representation and immediately set up for
intuition into moments in an all-sided system totalized by the
transparent unity of the mediate and the immediate, the syllogism is
\emph{immanent}. The transition from the immediate syllogism through
the mediate to the immanence-syllogism exhibits, as we shall see, a
progressive movement conditioned by the corresponding alterations of
the modal functions.

\textit{αα)} \emph{Immediate Syllogism: analytic Modal Functions.}
\addcontentsline{toc}{subparagraph}{αα) Immediate Syllogism: analytic Modal Functions}
Just as one can in a one-sided and arbitrary manner
restrict the logical judgment to the schema ``the singular is (is
subsumed under) the general,'' so too one can in a no less one-sided
and arbitrary manner dictate to oneself a schema according to which
every logical syllogism must always require a separation of two
extremes and a middle. When such a schema is to hold as the only
canonical one, it is undeniably a matter of course that in logic there
can no longer be any talk of ``immediate syllogisms.'' Hegel, who with
a view to this schema lets the judgment develop of itself into
syllogism, so that the two extremes correspond to subject and predicate
as the middle does to the copula,\footnote{``Der Schluß hat sich als die
Wiederherstellung des \emph{Begriffes} im \emph{Urtheile}, und somit
als die Einheit und Wahrheit beider ergeben. Der Begriff als solcher
hält seine Momente in der Einheit aufgehoben; im Urtheil ist diese
Einheit ein Innerliches, oder was dasselbe ist, ein Aeußerliches, und
die Momente sind zwar bezogen, aber sie sind als \emph{selbstständige
Extreme} gesetzt. Im \emph{Schlusse} sind die Begriffsbestimmungen wie
die Extreme des Urtheils, zugleich ist die bestimmte \emph{Einheit}
derselben gesetzt.'' Hegel. Die subjektive Logik. p.~118.} has to such a
degree
% --- p. 295 ---
\opage{295}overlooked the old doctrine of ``immediate syllogisms'' that
he has not even been willing to waste a word on getting it dismissed.

Nonetheless the doctrine of the immediate syllogisms is, when we look
more closely, just as indispensable in logic as the doctrine of the
mediate, and that for the good reason that the latter first acquire
their significance through the former. If we begin with the mediate
syllogism ``all men are mortal; Cajus is a man; thus Cajus is mortal,''
it is easily seen that the syllogistic function cannot restrict itself
to determining the conclusum ``Cajus is mortal'' after the two premises
have been set up beforehand. By degrading the function to so formal and
mechanical an act one misapprehends the logical significance of the
syllogism. The logical in the syllogism first unfolds itself when the
union of the subject and predicate of the consequent proposition is
expressly posited as a task, when it has in other words been brought to
consciousness that a proposition given in a determinate case has
validity only insofar as it admits of being presented as a consequent
proposition corresponding to determinate presuppositions. From this
standpoint the mediate syllogism will on all sides demand and
presuppose the immediate. The proposition ``Cajus is mortal'' is to be
determined as a consequent proposition, \emph{thus} it must have
premises; the proposition ``all men are mortal'' is the one premise;
\emph{thus} the proposition ``Cajus is a man,'' must be the other;
conversely! The proposition ``Cajus is a man'' is the one premise,
\emph{thus} the proposition ``all men are mortal,'' must be the other.
These are nothing but immediate syllogisms; but if a single one of
these immediate syllogisms fell away, what would then become of the
mediate one? Let one analyse whatever mediate syllo
% --- p. 296 ---
\opage{296}gism one will: provided only that one does not forget what and
how much originally belongs to the syllogistic function, it will surely
show itself that the mediate in every syllogism presupposes the
immediate.\footnote{So that it may not seem as though the commendation of
the immediate syllogism were an empty formality, and the syllogism
itself without influence on the actual progress of scientific
cognition, we choose at once an example from the exact sciences. If a
mathematician begins the solution of a problem of integration thus: it
is the expression $\frac{x^{n}dx}{\sqrt{a+bx^{2}}}$ that is to be
integrated, \emph{thus} it is the expression
$\frac{x^{n-1}}{b}\cdot\frac{bxdx}{\sqrt{a+bx^{2}}}$ that is to be
integrated, then everyone who reflects that a magnitude remains
unchanged by being multiplied and divided by an equal magnitude will
surely grant that here one begins with an immediate syllogism. But to
what end this syllogism? Well, that is seen from what for the expert
again shows itself as an immediate syllogism: the one factor, the
expression $\frac{bxdx}{\sqrt{a+bx^{2}}}$, is evidently the differential
of $\sqrt{a+bx^{2}}$; \emph{thus} one can integrate by parts, etc.}

That the discoverer of the immanent method has had no eye for the
pervasive significance of the immediate syllogism is remarkable enough,
since he otherwise has so sharp an eye for mediating transitional
forms; this circumstance, however, is by no means without a
characteristic influence on the method itself. If Hegel had not
regarded such logical movements as that from being and nothing through
becoming to determinate being, from determinate being and existence
through something and other to one (being-for-self), etc., as so
objective that they could be carried out by the concept itself, without
a comprehending subjectivity, then he would have seen from the outset
that these movements in actuality took place by judgment following upon
judgment in unbroken connection; but a series of development in which
the
% --- p. 297 ---
\opage{297}one judgment follows continuously upon the other will at the
same time show how the one judgment follows immediately from the other
for the free glance of thought; and that the one judgment follows
immediately from the other is precisely an expression of the immediate
syllogism. In a Hegelian logic one may choose whatever examples of
immanent thought-movement one will: on closer inspection it will surely
show itself that the transitions occur by immediate syllogisms. ``The
things that are equal,'' it says in Heiberg's ``\emph{logiske
System},'' ``thereby stand in relation to the common measure and
consequently also to one another, since otherwise they could not be
compared. But since all relation is difference, finitude, distance and
motion, all this is immediately present in the absolute unity and
difference of the original,'' etc. It is easy to reduce the indicated
thought-movement to a series of immediate syllogisms. The things are
equal, \emph{thus} they must have been compared; the things have been
compared, \emph{thus} they must stand in relation to the common
measure; the things stand in relation to the common measure,
\emph{thus} they stand in relation to one another. All relation is
difference, finitude, distance and motion, \emph{thus} finitude,
distance and motion are immediately present in the unity and difference
of the original, etc.

One will easily convince oneself that if in such a presentation a point
is discovered where the immediate syllogism does not hold good, then
the understanding of it fails as well. If anyone finds the syllogism
``the things are equal, thus they must have been compared'' untenable,
since, as it seems, there may well be many things both equal and
unequal that have not yet been compared: then he will also regard a
grounding such as this, that the equal things must stand in relation to
one another ``since otherwise they could not be compared,'' as
insufficient.

To be sure, the connection in which the several propositions occur can
shed a light over each one of them, so
% --- p. 298 ---
\opage{298}that transitions which are intelligible in the flowing
presentation can become unintelligible when the propositions are torn
in pairs out of the context in such a way that the one becomes premise,
the other conclusion. From the flowing presentation it is clear in our
example that here by ``the things that are equal'' are meant things
whose equality we cognize by comparison; whereas the immediate
syllogism ``the things are equal, thus they must have been compared,''
has what is offensive in it in this, that in the premise objective
equality is spoken of, \emph{thus} --- an immediate syllogism --- an
equality that is not seen to be able to rest upon the subjective
function of comparison. But the circumstance that the context as a
whole co-operates in the understanding of the singular does not alter
the matter; for the impression of the context now itself forms a
premise that has the understanding of the singular for its conclusion,
thus again --- an immediate syllogism. ``The context shows that a
thought equality is spoken of here, \emph{thus} --- an immediate
syllogism --- an equality that necessarily presupposes comparison.''

Since the context of a determinate passage forms the premise for an
immediate syllogism, the reader must not forget that the premise so
formed must itself become a consequent proposition to a more
comprehensive premise, and that this must come about by a new immediate
syllogism. ``In an objective logic there can be talk only of an
objective thinking, an objective comparison; \emph{thus} in this ---
belonging under the objective logic --- context there is talk only of
an objective comparing.'' One sees how the problem hidden in the
relation between equality and comparison shifts itself incessantly, by
the immediate syllogism, from one standpoint to another. The immediate
syllogism stands in the closest connection with what in the logical
sense is called glance, insight, survey; indeed it is itself the
divinatory advance of glance, insight and survey to mediate knowledge.

If it is asked how many essentially different forms
% --- p. 299 ---
\opage{299}the immediate syllogism can assume, then everything depends
chiefly on how one apprehends the relation between form and content on
the whole. If the form is treated by itself in abstract-mechanical
opposition to the content, then schematism arises; if the form is
developed in logical unity with the content, the difference of form
will rest upon a difference in the modal functions themselves; the
difference of form thus developed then leads us to gain insight into
the immediate syllogism's transition into the mediate.

1) \emph{Schematism: the Schema of the so-called immediate
Syllogisms.} To schematize is formal logic's strong side; the schema is
its ideal, the completion of the schema its flower and crown.
Everywhere that the logical formalism with its divisions and divisions
of divisions has really penetrated, the schema is unfolded and
ramified down to the microscopic. In exalted striving toward the
completion of schematism, formal logic has accordingly not been able to
restrict itself to a specification of the principal species of ``the
immediate syllogisms''; it has searched out the subspecies of each
species by itself, it has penetrated in earnest into the
singular.\footnote{A survey of the principal species of the so-called
immediate syllogisms might well afford a formal satisfaction, if one
avoided all prolixities and illustrated each species by an example. An
example such as: all birds lay eggs, \emph{thus} some birds lay eggs,
would then characterize the syllogism of subordination
(\textit{ratiocinatio subjectionis, vulgo subalternationis}); an example
such as: no animal is moral; \emph{thus}: every animal is a non-moral
being, would suit the identity-syllogism (\textit{ratiocinatio
æqualitatis, æquipollentiæ, conclusio ad æquipollentem}); an example
such as: all animals have voluntary motion, \emph{thus} there is no
animal that does not have voluntary motion, would be indicative of the
syllogism of opposition (\textit{concl. ad oppositam}), and an example
such as: all men are mortal, \emph{thus} some mortals are men, would
illuminate the syllogism of conversion (\textit{ratioc. per
conversionem}) and so forth.

But here formal logic cannot stop. According to its scholastic
principle the schema is a chief matter, and when the schema is a chief
matter, then what counts is perseverance and consistency: with a
principle one ought not to stop half-way. It is to be regarded as a
consistency respectable in its kind that the editors of formal logic,
after having for a time shown the weakness of yielding to the demands
of the age and restricting the schema, once more acknowledge the old
principle and unfold schematism. We shall now see how here, as regards
the immediate syllogisms, reflection is directed not merely upon the
species, but (Cf.\ Drbal. Lehrbuch der propædeutischen Logik. Wien
1865. p.~55--71.) also upon the subspecies.

Of conversion-syllogisms, e.g., there are now two species: proper
conversion-syllogisms and contraposition-syllogisms. Of proper
conversion-syllogisms there are again 4 subspecies: 1) conversion of a
general affirmative judgment (all birds lay eggs, thus: some egg-laying
beings are birds), 2) conversion of a general negative judgment (no
idler knows how to value time, thus: no one who knows how to value time
is an idler); 3) conversion of a particular affirmative judgment (some
dogs are faithful watchers, thus: some faithful watchers are dogs): 4)
conversion of the particularly negative judgment. For this conversion
there are no rules. To prove that there are no rules, one constructs
three geometrical figures, each consisting of two circles, which in the
first figure intersect one another, in the third figure fall outside
one another, while in the second figure the larger surrounds the
smaller. If one now considers the first figure, then some objects
falling under the sphere \textit{S} do not lie in the sphere \textit{P},
namely \textit{x}, \textit{y}, etc.; likewise there are objects belonging
under \textit{P} which do not lie in \textit{S}, namely \textit{n},
\textit{z}\dots. Besides, there are parts of \textit{S}, namely
\textit{s}, which lie in \textit{P}, and parts of \textit{P}, namely
\textit{p}, which lie in \textit{S}. Thus there are here four judgments:
1. some \textit{S} is \textit{P}; 2. some \textit{S} are not \textit{P};
3. some \textit{P} are \textit{S}; 4. some \textit{P} are not
\textit{S}. From this the conversion now follows: some \textit{P} are
not \textit{S}. E.g.\ some duties are not duties of compulsion (not an
object of compulsion), thus:
something that is an object of compulsion is not a duty. Some examples
are not good: thus something good is not an example\dots. To get any
such point out of it, one sets enormous schematic vehicles in motion,
seventeen pumping-engines to procure a single drop of water. Here too
it holds: \textit{parturiunt montes, nascetur ridiculus mus}.} It is the
% --- p. 300 ---
\opage{300}pure zeal for the pure form; no one shall say that in such
syllogisms as: what \emph{all} birds do, that must
% --- p. 301 ---
\opage{301}\emph{some} birds must do; what \emph{no} animals are, \emph{some}
animals cannot be; some beings of which what can be said of all men can be said
must be men, etc. --- too one-sided an emphasis is laid upon the logical
content; a suitable employment, on the other hand, of universally affirmative
and universally negative, particularly affirmative and particularly negative
judgments for suitable conversions may well be said to produce an excellent
effect, when the point is to get the schema as extensive and as widely
ramified as possible.

Were there really any sense in such a schematization, the schema would surely
at least have to be able to do service as a kind of register or table. But what
logical service should that be? As a register the schema is of no use. The need
for a register would presuppose that it mattered that the syllogisms, in order
to be rightly understood, should form a certain series and be set up in a
determinate order; but upon that no weight whatever falls here. The syllogism
that what \emph{all} birds do \emph{some} birds must do is evident by itself
without regard to series or register. It is wholly arbitrary, and for logical
insight completely indifferent, whether one will place the
equipollence-syllogism before the contradiction-syllogism, and this again
before the conversion-syllogism, or the reverse. As a table, then, the schema
is of no use either. Why does one draw up tables? Naturally in order to secure
the results of laborious observations, involved analyses, extensive
calculations, and the like. Thus one has star tables, logarithm tables,
chemical affinity tables, etc.

% --- p. 302 ---
\opage{302}But formal logic's immediate syllogisms are not results of extensive
calculations; one can really syllogize that what \emph{all} birds do
\emph{some} birds must do, without many observations or extensive calculations.

On every side formal logic's exertion with the schema stands in glaring
disproportion to the logical yield: the schema is extensive, the logical yield
$=$ 0.

2) \emph{The analytic Modal Functions.} The manner in which formal logic treats
the immediate syllogisms is in serious conflict with its own principle and
brings it at last into contradiction with itself. To an immediate syllogism
there belong, after all, two distinct judgments, of which the one forms the
premise, the other the consequent proposition. But how does this assumption now
stand to logic's supreme axiom, the so-called principle of identity? If there
are really two different judgments here, then the one cannot, without being
different from itself, contain the other. But if the one of two judgments
cannot be contained in the other, how then should the one be capable of being
derived from the other? If, on the contrary, it is not two different judgments,
but only two identical propositions, that we have before us, then the two
propositions give only one and the same judgment, expressed in two different
ways; but where there is only one judgment, there can impossibly be any
question of a syllogism. That ``All'' consists of ``Some'', that the
universally affirmative judgment contains the particularly affirmative but
excludes both the universally and the particularly negative, is surely
something that is self-evident in such a degree that there can be no room here
for any syllogism. It is thus in consequence of formal logic's own principle
that its doctrine of immediate syllogisms has at length been
rejected\footnote{``Diese angeblichen Schlüsse sind in Wahrheit keine Schlüsse,
sondern nur die Hervorhebung (Aussprache) dessen, was zwar nicht ausdrücklich,
doch aber unmittelbar im Obersatze mitgesetzt ist, so daß ich mir den Inhalt
desselben nur klar zu machen brauche, um es auch in ihm zu finden. Wo der
Mittelsatz fehlt, wo es also keiner Vermittelung zwischen der Prämisse und der
Conclusion bedarf, da bedarf es offenbar auch keiner Folgerung, d.\ h.\ da
findet die Function des Folgerns gar nicht statt, da ist auch überhaupt kein
Schluß vorhanden. Mit andern Worten: die Function des Schließens und somit der
Schluß als Ausdruck derselben besteht eben nur in der Ableitung der Conclusion
aus der Prämisse \emph{mittelst} des Mittelsatzes: wo also der Schlußsatz
\emph{nicht} abgeleitet, sondern nur ausgesprochen wird, daß sein Inhalt in der
Prämisse mit gesetzt ist, da findet offenbar auch kein Schluß statt.'' Ulrici.
Compendium der Logik. Leipzig 1860. p.~192. Let it be well noted that the
rejection is founded directly upon a conception of the principle of identity in
which there is in no respect any departure from formal logic.}.

% --- p. 303 ---
\opage{303}From the standpoint of objective logic, however, objection can also
be raised against syllogisms from one premise, although the difference here
develops out of the identity. It is urged from this standpoint that between
premise and consequent proposition there can here be only a
reflection-difference, and that such a difference is determined only by a
reflection, not by a syllogism. If we consider the matter abstractly
objectively, then the immanent transition from essence to phenomenon, from
ground to existence, from possibility to actuality, from cause to effect can
assuredly not deserve the name of syllogism; for each of the separated
categories is, when it is held fast for itself, so one-sided that, in order to
answer to its determination, it must reflect itself in its complement. The
advance indicated by the reflection is therefore just as much a retrogression;
it is a transition which is yet no transition. In that the thought-movement
goes forward from essence to phenomenon, it goes precisely from phenomenon back
to essence. Here, then, there is no distinct judgment about the essence from
which a judgment of its own about the phenomenon could be derived, for the
judgment about the essence falls in objective
% --- p. 304 ---
\opage{304}reflection necessarily together with the judgment about the
phenomenon; just as little can one think here of a distinct judgment about the
phenomenon from which a judgment of its own about the essence would let itself
be derived, since the judgment about the phenomenon must in objective
reflection necessarily fall together with the judgment about the essence.

These objections fall away when it is remembered, as was shown in an earlier
connection, that objective reflection must everywhere rest upon the subjective.
But the subjective in the reflection rests again --- by virtue of the word ---
upon the judgment. From this standpoint it becomes, then, not merely possible
but necessary to pass a distinct judgment about the essence in order to be able
to pass a distinct judgment about the phenomenon. Here there are two questions
for subjective reflection to answer: what is essence? what is phenomenon? That
these questions must be answered by different judgments is incontestable. Since
now the judgment that is passed about the phenomenon is just as much
determinative for the one that is passed about the essence as the reverse, the
judging subjectivity would find itself in an infinite hovering if the one of
these judgments were not expressly posited as the original, the other as the
derived. If, e.g., the judgment about the phenomenon is posited as the original
and to that extent immediate, then the judgment about the essence must direct
itself accordingly; it must as consequent proposition be derived by analysis
from its presupposition. The presupposition is premise, the consequent
proposition conclusion; the phenomenon is apprehended \emph{thus},
\emph{therefore} the essence must be apprehended \emph{thus}: here is an
immediate syllogism.

What this means must surely strike the eye when the question is how far one may
syllogize from possibility to actuality and from actuality to possibility.
Formal logic is, in accordance with what has once been agreed, just as certain
that every syllogism from actuality to possibility must be true:
\textit{ab esse ad posse esse valet consequentia}, as it is convinced that
every syllogism from possibility to actuality must be false: \textit{a posse
esse ad esse non valet conse%
% --- p. 305 ---
\opage{305}quentia}. That a syllogism from actuality to possibility can
nevertheless just as well be false (the actual has been, but is no longer, a
merely possible, since possibility has been superseded by actuality), as a
syllogism from possibility to actuality can be true (the possible in the
objective sense must become actual, for otherwise the possible becomes
impossible), is easy to see. Everything depends, in the given cases, upon a
distinct judgment about the possibility and a distinct judgment about the
actuality.

The significance of the immediate syllogism becomes still more striking when we
reflect that the objective categories determine one another only objectively,
in that they show themselves determinative for the objects concerned. Only when
the object is cognized as determined by its category does the reflecting
subjectivity become judging; but in the objective relations the one categorical
determination must necessarily draw the other after it, the one judgment lead
to the other: here is an effect (the treetops are moving), \emph{therefore}
there must here be a cause (the wind); here is a phenomenon (the pallor and the
flashing eyes), \emph{therefore} there must here be a ground (strong emotion);
here are determinate means (implements of agriculture), \emph{therefore} there
must here be a determinate purpose (agriculture). By apprehending the immediate
syllogisms from the objective, practical standpoint one will convince oneself
that they, far from being superfluous, play an important part in the
subjectivity's orientation in the actual surrounding world. What would merely
sensuous seeing and hearing be without a corresponding manifold of immediate
syllogisms? Or how should one explain to oneself the fact that there can be
sense-deceptions, optical and acoustic illusions, if there were no immediate
syllogisms?

% Errata: p. 305 --- printed "De", corrected to "Det"; translated as corrected.
Since now an immediate syllogism can in a given case just as well be false as
true, logic must adduce a criterion by which the true can be distinguished from
the false. It is the connection between the two judgments, premise and
consequent
% --- p. 306 ---
\opage{306}proposition, upon whose determination the matter essentially turns;
let one then investigate how far this connection lets itself be reduced to a
hypothetical judgment or not! If this reduction can be carried out, then the
syllogism is reliable; if not, it is unreliable. ``This is a plane triangle,
thus the sum of the angles $=$ two right angles''; this syllogism is true, for
``\emph{if} it is a plane triangle, \emph{then} the sum of its angles must
necessarily be two right angles.'' If, on the other hand, someone syllogizes
thus: ``The family was well this morning (8~o'clock), \emph{therefore} it is
well still (12~noon)'': then it is easily seen that such a syllogism is not
founded upon necessity, but only upon probability, and thus cannot be entirely
reliable. In theory this is recognizable by the fact that premise and conclusum
do not stand to one another as antecedent and consequent clause in a
hypothetical judgment; in practice it is recognized by the fact that the
syllogism has a reassuring effect only under ordinary circumstances, but not
when there is danger afoot.

The analytic modal functions must accordingly aim not merely at stating the
relation between the immediate syllogism's premise and conclusum in general,
but also at determining, for every peculiar case, how far the syllogism's
foundation is necessity or probability, and, if it is a probability, at stating
its degrees. Should it in particular cases turn out that the foundation is not
a being one, but hovers between the poles of probability and necessity, then it
is the task of the judging subjectivity to draw a boundary-line between
\emph{syllogism} and \emph{conjecture}, an often very fine boundary-line which
both in theory and in practice is of decisive significance for the thinking
that advances from judgment to judgment.

In order to apprehend the significance of the analytic modal functions in its
whole extent, let one represent to oneself the subjectivity in relation not
merely to a manifold of immediate phenomena, but, as judging subjectivity, to a
manifold
% --- p. 307 ---
\opage{307}of immediate, logical, antilogical, categorical, assertoric
judgments. How should it now be possible to bring about unity in such a
consciousness, if it were not possible, by an analyzing reflection, to derive
one judgment from another, if, in other words, it were not possible to make
immediate syllogisms. Only by virtue of the immediate syllogism does an
essential connection arise between judgment and judgment, only by virtue of the
immediate syllogism is there continuity in advancing thinking. A continuously
advancing thinking is an intellectual movement; but such a movement can surely
go in the most different directions.

That there can here be question of a plurality of analytic modal functions
rests then chiefly upon this, that not only can another judgment be derived
from every judgment, but that from one and the same premise essentially
different syllogisms can even be drawn. ``Cajus is a man, \emph{therefore}
Cajus is mortal; Cajus is a man, \emph{therefore} Cajus is immortal; Cajus is a
man, \emph{therefore} in need of sleep,'' etc. Each of these syllogisms
acquires its peculiar significance from the peculiar direction in which the
thought moves: the physiologist syllogizes otherwise than the ethicist, the
ethicist again otherwise than the aesthetician, etc. ``But'', someone might
ask, ``are there now also different functions here in the logical sense? It
really seems as if the thought functioned in the same way whether it is here
syllogized: \emph{therefore} C. is mortal, or: \emph{therefore} C. is immortal;
or: \emph{therefore} C. is in need of sleep: insofar as the one Therefore is
different from the other, this difference surely concerns only the content, not
the form.'' We draw attention to the fact that every essential
content-difference must always reflect itself in a corresponding
form-difference. Wherever the content-difference withdraws itself from the
form-difference, the content sinks down to being mere bare material. If the
difference between a physiological, an ethical and an aesthetic content is more
than a material-difference, then the difference between a physiological,
% --- p. 308 ---
\opage{308}an ethical and an aesthetic syllogism is also an essential
form-difference. The difference itself does not concern logic, but that there
is here a difference, an essential form-difference, does concern logic. That
formal logic does not know the essential form-difference is explicable; formal
logic looks to the result and overlooks the thought-movement, it holds to the
schema and forgets the form. If the whole attention is concentrated upon the
conclusum that follows upon the immediate syllogism's \emph{Therefore}, then
the one Therefore is quite certainly like the other, there is in the function
itself no noteworthy difference. But whence now does that come? It comes from
this, that one overlooks the thought-movement, the advance in which the
syllogism-act itself marks a turn, and looks only to the result, whether, e.g.,
the concluding proposition is affirmative or negative, universally negative or
etc., difference-determinations which one-sidedly suit the schema, the
mechanical form torn loose from the content.

Concentrated in the \emph{Therefore} which, according to the schema, supports
itself exclusively upon the one premise, the thought-movement is reduced to a
vanishing minimum, an element in which its \emph{Whence} and \emph{Whither} has
become wholly unrecognizable. All \textit{S} are \textit{P},
\emph{therefore}: some \textit{S} are \textit{P}, here is --- under the aspect
of a syllogism --- an infinitely small movement from All being to Some being, a
movement so small that one may even dispute whether it is progressive or
retrogressive. It stands with the movement-differentials in the different
\emph{Therefore}s belonging to the immediate syllogism as with the different
differentials in ``the analysis of the infinite.'' The elements \textit{dx},
\textit{xdx}, \textit{adx}, \textit{cosxdx} and the like can, when they are
regarded as existent, be taken to be equal, since they are all infinitely near
0; but when they are integrated, then one sees \emph{Whence} and
\emph{Whither}, then one sees the \emph{movement}, and then the difference is
thoroughgoing. So too with the logical movement-differential, the thought
enclosed in \emph{Therefore}; it acquires its logical significance from the
whole
% --- p. 309 ---
\opage{309}thought-movement. In virtue of the thought-movement it makes an
essential difference whether it is a physiological, an ethical or an aesthetic
\emph{Therefore} that we have before us. To the fundamental difference in the
immediate syllogism's \emph{Therefore} there corresponds the difference in the
analytic modal functions.

3) \emph{The Transition of the immediate Syllogism into the mediate.} It lies
in the concept of the immediate syllogism that its consequent proposition must
be derived from one premise; but it by no means lies in its concept that it
must in every case be absolutely immediately evident. That this demand cannot
be set up as universally valid is a direct consequence of the fact that what is
instantly clear to the one may be obscure to the other. What one man sees at
first glance another is perhaps able to grasp only after many preparations and
explanations\footnote{This difference is often of great comic effect. In
``Hovmesteren i Knibe'' Ane can see at once from the forgotten hat (thus by an
immediate syllogism) that there has been a lady in the tutor's room; but it
costs her considerable effort before she can make the phenomenon evident to
Hans.}. Talent for a particular science rests upon a particular eye, but the
particular eye is precisely recognizable by the ease and sureness with which
the man gifted in a given direction makes immediate syllogisms. Objectively, or
in and for itself, the demand cannot hold either; for in the first place there
is no syllogism so absolutely objective that it could syllogize itself; in the
second place there is no syllogism so absolutely immediate that it must not,
when reflection is awakened, at once pass over into the mediate\footnote{In a
peculiarly naive way the mediacy in the immediate syllogisms also comes to
light in this, that for every such syllogism a proof is expressly demanded. One
proves --- only think! --- the validity of a syllogism such as this: it is true
that all animals have voluntary motion, thus it is false that some animals do
not have voluntary motion. The proof runs (Drbal. Lehrbuch der propädeutischen
Logik. p.~59) thus: ``If the judgment: all animals have voluntary motion is
true, then the judgment some animals have voluntary motion must also be true.
If it now be assumed that the judgment: some animals do not have voluntary
motion were likewise true, then, since the judgments have the \emph{same
particular} subject, the propositions: some animals have voluntary motions and
precisely \emph{these same} (animals) have no voluntary motion, would be true
at one and the same time, which conflicts with the axiom of contradiction.''
How mediate the immediate syllogism must become through such a conduct of proof
is further seen from this, that the proof of its validity expressly grounds
itself upon the proof of the validity of the subordination-syllogism. This
proof runs thus: ``If the universally affirmative judgment: All \textit{S} are
\textit{P}, is valid, then the particularly affirmative judgment: some
\textit{S} are \textit{P}, must also be valid, for the opposite would, it is
said, be `unreasonable'.'' That all these conducts of proof are a fiasco must
be granted, for he who cannot at an immediate glance see the contradiction in
some animals not having voluntary motion, once it is given that all animals
have voluntary motion, is a complete idiot; but a complete idiot does not
become wiser by being referred to the axiom of contradiction.\par For an
assessment of formal logic's formal obscurity the conduct of proof is
characteristic insofar as it hereby becomes manifest that it has not duly made
clear to itself the concept of the immediate syllogism. If the immediate
syllogism is to be known by the immediately evident: where then does logic mean
to go with its proofs? If the syllogism's immediacy is to be known by its
having only one premise, why then is it restricted to vacuous forms and
examples?}. ``The line from \textit{a} to \textit{b} is straight, thus it is
% --- p. 310 ---
\opage{310}the shortest way between \textit{a} and \textit{b}''. Is this
syllogism now evident by itself or is it not? Subjectively regarded, both
conceptions hold, as has been shown in the foregoing. The syllogism is
immediately evident, for it denotes an axiom; yet no, the syllogism is not
immediately evident, for the alleged axiom is really a theo%
% --- p. 311 ---
\opage{311}rem, and a theorem requires proof. Objectively regarded, the
syllogism is, on the contrary, unconditionally immediate; for upon what does
its immediacy rest? Solely and alone upon this, that its consequent proposition
is derived from a single premise.

From this it follows again that the immediate syllogism can not only demand a
rich and many-sided analysis, but can also have a rich scientific content.
Syllogisms such as these: ``The daily motion of the sun and the stars is only
apparent, \emph{therefore} it must be the earth's motion that is actual''; ``a
stone thrown from a high tower deviates towards the east, \emph{therefore} the
earth moves from west to east'', undeniably require for their right
understanding a multitude of reflections and intermediate determinations; but
for that reason they are, according to the concept laid down, none the less
immediate. That the conclusum rests upon a single premise makes the syllogism
immediate.

The immediate syllogism is a finite form to which an infinite content lets
itself be reduced. The result that is science's last and richest must, so
surely as it comes about by a conclusion, be just as capable of being referred
back to a single premise as that which is drawn from the very first and
simplest axiom. The premise of the last conclusion must indeed, when it is
analyzed, contain a manifold of premises; but the demand is that the many
premises with their conclusions be reduced to one premise. It is an ever richer
reduction which, in that it gives the premise an ever richer content, furnishes
the foundation for the analytic activity of the modal functions. If great and
content-rich knowledge-results could not be concentrated in immediate
syllogisms, a true and comprehensive science would be impossible.

Were each single one of the premises lying in the content-rich premise to be
brought forward, gone through and explained separately every time, then science
would have to begin constantly over again at the treatment of every new point,
and could in this way never come to an end with its
% --- p. 312 ---
\opage{312}investigations. One would, e.g., hardly get very far in geometry if
one could not support oneself upon such immediate syllogisms as: ``The triangle
is right-angled, thus the square on the hypotenuse is equal to the sum of the
squares on both the legs;'' ``the figure is a circle with radius $=$~\textit{r},
\emph{therefore} its area is $= \pi r^2$''; ``\textit{A} is the area of the
whole spherical surface, thus \textit{A} is equal to the product of the
diameter and the circumference of the great circle'', etc.

Since now the premise in an immediate syllogism can contain as great a manifold
of mediately won results as one will, the transition from the immediate to the
mediate syllogism is therewith at the same time motivated: the immediate
syllogism is at bottom mediate. Were there a syllogism so immediate that it did
not let itself be resolved into a mediate one, then its premise would have to
be an absolute axiom evident by itself; but that such axioms are not to be had
has been shown in the foregoing.

The immediate syllogism is not merely the presupposition of the mediate, but
also its result, and can as result be resolved again into its presupposition. A
syllogism such as: ``Cajus is a man, thus mortal'', evidently rests upon the
tacit presupposition that every man is mortal; but when this presupposition is
expressly brought forward, the syllogism acquires two premises: every man is
mortal, Cajus is a man, thus --- here is a mediate syllogism.

\textit{ββ)} \emph{The mediate Syllogism: synthetic Modal Functions.} While the
analytic modal function, as regards the terms of the syllogism, holds to
identity, the synthetic holds on the contrary to difference. Only by a peculiar
co-operation of thinking and representing can there be brought about here an
equilibrium of analysis and synthesis which makes the mediate syllogism
possible. If in the often-cited example: ``All men are mortal, Cajus is a man,
thus Cajus is mortal,'' one gives analyzing thinking a one-sided
% --- p. 313 ---
\opage{313}preponderance over the synthesis conditioned by representation, then
the syllogism loses its significance. This Hegel has done, and has for that
reason been unable to acknowledge the syllogism. ``The major premise'', so runs
his objection\footnote{Die subjektive Logik. p.~151.}, ``is correct only
insofar as \emph{the consequent proposition is correct}; if Cajus happened not
to be mortal, the major premise would not be correct. The proposition that is
supposed to be the consequent proposition must already be immediately correct
for itself, since otherwise the major premise could not comprise all singulars;
before the major premise can pass for correct, one must ask beforehand whether
that consequent proposition is not itself an \emph{instance} against it''.

From the objection it is seen that Hegel lets the analytic modal function
absorb the synthetic and thereby sublate itself. If Cajus is a man, then, in
order to be able to say what every man is, one must be able to say beforehand
what Cajus is. In this way the identity of what Cajus is and what every man is
is reduced to a tautological paraphrase: ``Cajus is what every man is, when
every man is what Cajus is.''

This formalism ceases, however, when finite thinking is brought into the right
relation to intuition and representation. If the individual Cajus is set over
against the concept man, so that the individual as subject is for the
representation what the concept as predicate is for the thought: then it is
seen that the subject, the single man, must have the universally human
properties in an individual manner. Finite thinking, conditioned by the
representation, now apprehends this in such a way that the individual's
properties fall into two systems: systems of the universally human and systems
of the individual ones; the individual properties being then assumed to hold
either exclusively for \emph{this} one or at all events only for \emph{some}
singulars.

% --- p. 314 ---
\opage{314}Only when the analysis has been carried out in this way is there
room here for a corresponding synthesis.

In accordance with the separation of general and particular properties in every
single man, the total cognition then divides into two partial cognitions. On
the one side a property common to all men is apprehended and determined. This
universally valid property is seen, under the aspect of thinking, to lie in
man's essence; it is to that extent the category, the concept as concept, upon
which the universal validity is founded. ``Man is --- it lies in human nature
--- mortal.'' But finite thinking is particularly conditioned by the
representation; in that the universally valid property is brought forward, the
gaze fastens not upon the pure concept, but upon the individuals in which the
concept is realized, thus not upon man, but upon men. Thus we have in the major
premise a cognition for itself: ``\emph{All men are mortal}''. How do we now
come to the minor premise? Evidently not by a continued analysis, but by a
synthesis. In the cognition that all men are mortal, the cognition that Cajus
is a man is by no means comprised. Cajus might, e.g., be either a horse or a
dog or an elephant; indeed, there is in logic nothing even to hinder Cajus, if
one will have it so, from being a stockfish.

If one will turn the relation about and say that it is really the cognition
uttered in the minor premise that comprises the cognition of the major premise,
this too is a misunderstanding. It is assuredly impossible to cognize Cajus as
a man without seeing him as sharing in certain properties universally valid for
human nature. But from this it by no means follows that representing thinking
would need to reflect through all the properties of human nature in order to
apprehend the single individual as a man. On the contrary! it is here expressly
presupposed that the apprehension takes place in the well-known immediate
manner, namely by intuiting the hu%
% --- p. 315 ---
\opage{315}man figure. The minor premise is therefore, in the particular
cognition, just as independent of the major premise as the latter is of the
former. In that the two distinct cognitions are brought together, the syllogism
acquires its significance, and the modal functions show themselves to be
synthetic.

The cognitions uttered in the premises of the immediate syllogism are, with all
their respective independence, by no means matters of indifference to one
another. On the contrary! in that they are brought together, that is, in that
the synthesis is carried out, it shows itself that by their union they form a
new cognition. The consequent proposition, the new cognition, can be grounded
in the premises in different ways; how the grounding is more precisely
constituted in the particular cases must be decided by the analysis. For it
depends upon how the elements of the syllogism, (\textit{syllogismi elementa}),
the terms which in premises and consequent proposition make up subjects and
predicates, stand to one another. If particular weight is laid upon the
formulae conditioned by the shifting position of the elements, we have
\emph{the Syllogistic}; if the question here is of the original fundamental
determinations, then everything essentially depends upon a \emph{Reduction} of
the \emph{syllogistic Forms}; if finally the distinct premises are apprehended
as form-giving moments of a content continuously developing itself in
syllogism-systems, then \emph{the Dissolution of the Syllogistic} is logically
motivated.

% Errata: p. 315 --- printed "Syllogistisk" and without the item number, corrected to "1) Syllogistik"; translated as corrected.
1) \emph{Syllogistic.} The mediate syllogism is according to its concept the
act of thinking which grounds a judgment's validity upon a union of two or more
other judgments, and is designated in particular by the
Aristotelian-scholastic name \emph{Syllogism}. The Syllogistic, or the doctrine
of the syllogisms and their application, then begins schematically with a
distinction between the \emph{simple} syllogism (\textit{syllogismus simplex}),
which has only two premises, and the \emph{composite} (\textit{s.\ compositus}),
which forms itself by the connection of a series of simple syllogisms
(polysyllogism). The simple, non-composite syllogism
% --- p. 316 ---
\opage{316}is again divided into the \emph{categorical}, whose premises are both categorical judgments, the \emph{hypothetical}, at least one of whose premises must be a hypothetical judgment, and the \emph{disjunctive}, at least one of whose premises must be disjunctive. --- Schematization is once more in full swing\footnote{It is a rich harvest that scholastic syllogistic offers to the sharp-witted, whose thinking-mechanism is especially assisted by rubrics and schemata. If we designate, in the categorical syllogism, the conclusion's subject by \textit{S}, its predicate by \textit{P} and the middle term by \textit{M}, then the following schema gives us the four figures:
\begin{center}\begin{tabular}{c|c|c|c}
First Figure & Second Figure & Third Figure & Fourth Figure\\
\textit{M}---\textit{P} & \textit{P}---\textit{M} & \textit{M}---\textit{P} & \textit{P}---\textit{M}\\
\textit{S}---\textit{M} & \textit{S}---\textit{M} & \textit{M}---\textit{S} & \textit{M}---\textit{S}\\ \hline
\textit{S}---\textit{P} & \textit{S}---\textit{P} & \textit{S}---\textit{P} & \textit{S}---\textit{P}\\
\end{tabular}\end{center}
If one now moreover lays particular weight on the division of judgments into universally affirmative (\textit{a}) and particularly affirmative (\textit{i}), after the two first vowels in the verb \textit{affirmo}, into universally negative (\textit{e}) and particularly negative (\textit{o}), after the vowels in the verb \textit{nego}, then one obtains, by letting each of these 4 kinds of judgment in each of the 4 figures serve alternately as major and minor premise, 64 possible combinations; and when one thereupon applies the three golden rules: \textit{ex mere negativis nihil sequitur} (no \textit{M} is \textit{P}, no \textit{S} is \textit{M} yields no syllogism), \textit{ex mere particularibus nihil sequitur} (some \textit{M} are \textit{P}, some \textit{S} are \textit{M} yields no syllogism) and \textit{conclusio sequi debet partem debiliorem præmissarum}
(when one of the premises is negative, the conclusion becomes negative; if the one premise is particular, the conclusion becomes particular), one will find that only 19 of the combinations yield true syllogisms. These 19 combinations are then applied in the 4 figures in such a way that each figure thereby obtains its valid moods. --- Thus, e.g., the first mood in the first figure is \textit{Barbara: M (a) P, S (a) M, ergo: S (a) P} (\emph{everything} immaterial is imperishable, \emph{all} absolutely singular beings are immaterial, thus: \emph{all} absolutely singular beings are imperishable). The third mood in the first figure is \textit{Darii: M (a) P, S (i) M, ergo: S (i) P} (\emph{all} acids color blue vegetable matters \emph{red}, \emph{some} oxides are acids, thus: \emph{some} oxides color
blue vegetable matters \emph{red}). In the third figure, on the other hand, the fourth mood is called not \textit{Darii} but \textit{Datisi}, a remark which, as one sees, is very important: \textit{M (a) P, M (i) S, ergo: S (i) P} (\emph{all} fungi are cryptogams, \emph{some} fungi are parasitic upon other organisms, thus: \emph{some} things parasitic upon other organisms are cryptogams).

If one has once said \textit{A} --- this is an acknowledged saying --- then one must also say \textit{B}. If one really wishes to lay scientific weight upon the three Aristotelian figures, then it is quite in order that one should also demand respect for the fourth, the Galenic. Why should the fourth figure not, considered as figure, be just as respectable as the three others? True enough, the fourth figure has not quite so many moods as the third, for the third figure has six moods: \textit{Darapti, Felapton, Disamis, Datisi, Bocardo, Ferison}; but it has nevertheless more than either the first or the second. These namely have only four moods each, the first: \textit{Barbara, Celarent, Darii, Ferio}, the second: \textit{Cesare, Camestres, Festino, Baroco}, whereas the fourth has five valid moods, very important for ``science,'' namely: \textit{Bamalip, Calemes, Dimatis, Fesapo, Fresison}. Thus: if the fourth figure is to go, then away with them all! but if one is in a position to hold on to the fourth figure, one is also in a position to hold on to them all.

Only by holding on to all four figures will one rightly be able to grasp the deep logic in the ingenious verse that derives from Petrus Hispanus (he became, though we do not precisely venture to maintain that it was for the verse's sake, afterwards Pope under the name of John XXI, and died in 1277); the verse runs as follows:
\par\medskip
\hspace*{2em}\textit{Barbara, Celarent primæ, Darii, Ferioque.}\\
\hspace*{2em}\textit{Cesare, Camestres, Festino, Baroco secundæ.}\\
\hspace*{2em}\textit{Tertia grande sonans recitat Darapti, Felapton,}\\
\hspace*{2em}\textit{Disamis, Datisi, Bocardo, Ferison. Quartæ}\\
\hspace*{2em}\textit{Sunt Bamalip, Calemes, Dimatis, Fesapo, Fresison.}
\par\medskip\noindent
If one now takes note --- for in so acute a verse there is much to take note of and exceedingly much to remark upon --- that not only the vowels \textit{a, e, i, o}, but also the initial letters \textit{B, C, D, F}, are significant, in that they call to mind what is analogous
in the kinds of syllogism belonging under the different figures (\textit{Bamalip}, e.g., analogous with \textit{Barbara}), then one is inspired by the logical ``conscience''; and if one moreover becomes aware --- and here there is much to become aware of --- that even the consonants \textit{s, p, c, m} acquire a deeper meaning through this, that they according to the verse:
\par\medskip
\hspace*{2em}``\textit{S vult simpliciter verti, P verte per accid.}\\
\hspace*{2em}\textit{M vult transponi, C per impossibile duci;}''
\par\medskip\noindent
indicate what alterations must be undertaken with each of the kinds of syllogism in the three last figures, when they are to be reduced to the first figure, then one feels ``\textit{testimonium spiritus}''.}
.

% --- p. 317 ---
\opage{317}Although the detailed schematization of the multiplicity and mutual difference of the forms set up also,
% --- p. 318 ---
\opage{318}as regards the mediate syllogisms, shows itself to stand in disproportion to the logical yield, there is here nevertheless, on account of the synthesis, a plurality of intermediate determinations and middle terms which makes it explicable that thinking, reflecting upon its own laws, can find itself impelled to seek an exhaustive determination of cases and rules.

Already the categorical syllogism's three elements: \emph{the middle term} (\textit{terminus medius}), the common third to which the conclusion's subject and predicate must each for itself be referred, \emph{the minor concept} (\textit{terminus minor}), the conclusion's subject, and \emph{the major concept} (\textit{terminus major}), the conclusion's predicate, must present points of combination and awaken interest in a schematization. If one now takes note of how \emph{the major premise} (\textit{propositio major}), the premise in which the major concept is connected with the middle term, and \emph{the minor premise} (\textit{p.\ minor}), the premise in which the minor concept is united with the middle term, are formed in different cases: then one easily discovers that the middle term can change place in the premises and thereby give these a different character.
In the syllogism: ``all \emph{men} are mortal, Caius is a \emph{man}, thus Caius is mortal,'' the middle term is subject in the major premise, but predicate in the minor premise; in the syllogism: ``no rectangle is \emph{equilateral}, all squares are \emph{equilateral}, thus no square is a rectangle,'' the mid%
% --- p. 319 ---
\opage{319}dle term is predicate in both premises. The question then arises, how many cases are there really here, and what may these different cases signify? Once one has found out that through the middle term's different position four syllogistic figures can be formed, then the cases belonging under each of the figures, which show themselves now fit and now again unfit for an adequate syllogism, will provoke the question of the different ways in which one can syllogize in each single figure. The schematism shows to that extent a striving to secure the cases and to bring the special under the universal; but the endeavor is fruitless, the movement illusory; instead of immersing itself in its own essence, thought carries on a combinatorial play with the forms that vary upon the surface.

2) \emph{Deduction of syllogistic Forms.} There holds of the mediate syllogism what has already been said of the immediate: the development of form must go as one with the logical development of content. But the development of form, a deduction of the syllogism's particular forms, is conditioned by a corresponding \emph{reduction}, a leading back of the manifold to the beginning itself, the simple first.

Among the logicians who in more recent times have aimed at a rational reduction of the mediate syllogism's particular forms, we must especially name Ulrici\footnote{Compendium der Logik. Leipzig 1860. p.~183--194.}
. ``In the universal's relation to the singular,'' he says, ``in the genus's relation to the species, as in the species' relation to its specimens, thus in the concept of the judgment, there lies immediately this, that what holds (is predicated) of the universal must also hold of the singular comprised under the universal, and that what holds of the genus must also hold of its species, subspecies and specimens.'' In consequence of this presupposition Ulrici reduces the syllogism's necessity to the first logical axioms: ``\textit{A} = \textit{A}, of like holds like, what
% --- p. 320 ---
\opage{320}holds of the genus must hold of its species and of their specimens; this is the universal axiom for all \emph{positive} syllogisms; further: \textit{A} is not = \textit{non A}, of like the unlike does not hold, and of the unlike not the like, what does not hold of the genus cannot hold of its species and their specimens; this is the universal axiom for all \emph{negative} syllogisms''.

``From the universal axiom for all positive syllogisms it follows: 1) that when a whole species is subsumed under a general predicate-concept or higher genus-concept, then every one of the species' specimens is also subsumed under the same concept. With the syllogisms formed according to this norm the first so-called syllogistic figure is implicitly given. From the same universal axiom it follows: 2) that a subject which is to be subsumed under a higher genus-concept must necessarily be subsumable under some one or other of the species-concepts comprised in it, an axiom which belongs to the second syllogistic figure, in which it finds especially a negative application. From the axiom it follows further: 3) that when a species is subsumed (allows itself to be subsumed) under two different genera, each of these must also in part be subsumable under the other. The syllogisms formed according to this norm belong to the third figure. Finally it is: 4) a consequence of the general axiom that when a species-concept is to be subsumed under a genus-concept, and this again under a higher genus-concept, then some specimens belonging to the last concept must also be subsumable under the species-concept; and from the axiom for all negative syllogisms, that when the genus-concept under which the species-concept is subsumed does not belong under a higher genus-concept, no specimen of this last can be subsumed under the species-concept; as also that, when the species is not subsumed under the lower genus-concept, some specimens of the higher genus-concept must also be excluded from
% --- p. 321 ---
\opage{321}the species-concept. It is the syllogisms formed according to this norm that are designated by the fourth figure\footnote{Under his Hegelian standpoint Heiberg has in a very simple and natural manner deduced the three figures with rejection of the fourth, the non-Aristotelian one. After having indicated the system of the three syllogisms by a) \textit{E}---\textit{S}---\textit{A}, b) \textit{S}---\textit{E}---\textit{A} and c) \textit{E}---\textit{A}---\textit{S}, according to the position that the singular, the particular and the universal each time occupy, he continues: ``Since the predicate is always of greater extension than the subject, in the first syllogism every preceding term is subject, and every succeeding one predicate. But in the second syllogism the middle term is subject, and the extreme terms predicates, and in the third the extreme terms are subjects, and the middle term predicate. Or: in the first syllogism the term which is first predicate becomes afterwards subject; in the second one subject is common to two predicates; and in the third one predicate is common to two subjects. But in the conclusion of all three syllogisms the first term is subject, and the third predicate. [Here it is to be remarked that with the completed system of syllogisms subject and predicate become of equal extension; the relation between species and genus disappears; the general is not of greater extension than the singular. Herein already lies the transition to objectivity]''. Specul.\ Logik. §~155.}
''.

But however clear and judicious this deduction must be said to be in its main features, it nevertheless suffers from formal logic's old, well-known fundamental defect, that difference is held outside identity, and contradiction is left standing as a piece of logical lumber. If the equation \textit{A} = \textit{A} is to express identity, and identity so expressed is to be principle, then logic will be unable to bring it as far as a judgment, let alone as far as a syllogism. Even if one were, with Ulrici, to regard the proposition in which a singular as singular is subsumed under the general as being the only logical judgment, it is nevertheless clear that by such a proposition one goes far beyond pure identity. Could a glance at the judgment not at once teach one that the identity \textit{A} = \textit{A} is a tautology, and a tautology insufficient to ground a judgment, then
% --- p. 322 ---
\opage{322}an attempt to deduce and determine the syllogism, in particular the mediate one, ought nevertheless to be able to do so. It is reciprocal determinations of genus, species and individual about which the whole matter turns. How is it now possible to think the species as middle concept between genus and individual, without seeing this middle term on the one side identical with the individual, on the other side with the genus? But the genus is of course different from the individual; the middle term is thus identical with two different things and, despite this doubling of identity, identical with itself; here are three identities in one identity: such a triune identity has in truth its difference not outside itself, but in itself.

If one overlooks the unity-in-difference without which the middle term could not possibly unite the extreme terms, then one will also remain standing at a very one-sided conception of the syllogism-function itself. This is precisely the case with Ulrici. Instead of cognizing the fundamental function as it comes forward in its particular forms, i.e. as it actually functions, he holds fast to its abstract likeness with itself; function = function, just as \textit{A} = \textit{A}, and in this way obtains a function that cannot function\footnote{``Es giebt \dots\ vier \dots\ Schlußfiguren, die sich insofern logisch unterscheiden lassen, als sie aus dem allgemeinen Grundsatz alles Schließens in seiner positiven wie negativen Form, der zugleich die Norm für die Gestaltung aller Schlüsse ist, unmittelbar sich ergeben. Sie bilden indessen, wie von selbst einleuchtet, keineswegs vier verschiedene ``\emph{Arten}'' von Schlüssen. Denn sie sind keineswegs Modificationen der logischen \emph{Function} des Schließens, sondern nur --- im Grunde unerhebliche --- Specialformen der allgemeinen \emph{Gestalt} der Schlüsse''. Anfr.\ Skr.\ p.~188.}
.

That difference belongs just as originally in the ground, and thus in every fundamental function, as identity; that identity itself is in the ground difference, has ``the grounding reflection'' in its place\footnote{Grundideernes Logik. I. p.~93--99.}
{} already demonstrated; that a fundamental
% --- p. 323 ---
\opage{323}function, when it is held fast in abstract generality, cannot be applied at all, and thus must remain in an empty being without being able to reach intellectual existence, can even be proved outright mathematically.

$F(x) = y$ is every possible function and precisely on that ground indeterminate: with such a function no reckoning can be done. Insofar as the function is \textit{in abstracto} as fundamental function, it is, by being in the ground, only in possibility. ``In relation to the reduction possibility is ground; in relation to the deduction the ground is possibility. The identity of ground and possibility is seen from the disjunctive judgment's different forms: \textit{A} is neither \textit{B} nor \textit{C} nor \textit{D} \dots\ $f(x)$ is neither $ax$, nor $x^a$ nor $a^x$ \dots: the particular determinations are dissolved in the universal, and have thereby gone to ground; \textit{A} is, in the ground, both \textit{B} and \textit{C} and \textit{D} \dots\ $f(x)$ is, in possibility, both $ax$, $x^a$ and $a^x$ \dots\ Insofar as what is particular in the ground's content lets itself appear, it is reflected in possibility; where one case is possible, at least two cases must be possible, for if the possible case's opposite were not also possible, the case itself would be necessary: possibility has alternating terms. If now one of the terms is to be posited as existent, then the disjunctive judgment in its usual form shows: \textit{A} is either \textit{B} or \textit{C} or \textit{D} \dots, $f(x)$ is either $ax$ or $x^a$ or $a^x$ \dots\ the tension between the alternating terms\footnote{Mathematik og Dialektik. p.~50--51.}
''.

What in this respect holds of the mathematical functions holds in corresponding fashion of the logical ones. It is not possible to apply the syllogism's fundamental function without specifying it according to one of the four figures. If the judgment: ``Caius is mortal,'' is demanded as conclusion, and the judgment: ``all men are mortal,'' is given as the one premise, then the judgment: ``Caius is a man,'' must surely be the other premise; it is then by no means arbitrary,
% --- p. 324 ---
\opage{324}but altogether necessary, that the syllogism should correspond to the first figure. If the judgment: ``some animals living in the water are mammals,'' is demanded as conclusion, and the judgment: ``all whales live in the water,'' is given as the one premise, then one must necessarily have either: ``all whales are mammals''; or: ``some whales are mammals,'' for the other premise, and it is then a matter of course that the syllogism falls under the third figure. If now the syllogism-function is recognized, as was shown above, not only by the act of thinking through which a conclusion is derived from given premises, but also by the synthetic activity of thought through which the premises are found out and set up, then it lies in the nature of the matter that the syllogism-function is modified with each of the syllogism's typical ``special forms.'' For the synthesis of the premises takes place under other conditions and in another manner when the middle term is to be subject in the major premise and predicate in the minor premise, than when it is to be predicate both in the major and in the minor premise, and in the latter case again under other conditions and in another manner than when it is either to be subject in both premises or predicate in the major premise and subject in the minor premise. If one overlooks this, the act of syllogizing must sink down into a thoughtless mechanism; but if one sees this, one will also see that the synthetic modal functions must change with the changed synthesis.

3) \emph{The Dissolution of the Syllogistic.} But what significance do the mediate syllogisms really have for knowledge and science? It has been objected against the syllogistic, especially in more recent times, that it finitizes the forms of knowledge, parcels out cognition, and, through the syllogisms' ``well-known convenience for proving the most diverse, indeed even contradictory things,'' promotes scepticism. The syllogism's arbitrariness lies in this, ``that the one who conducts the proof can choose the middle term (\textit{terminus medius}) at pleasure. For since every thing presents several points of view for consideration, and the more of them the more con%
% --- p. 325 ---
\opage{325}crete it is, it can by any single one of these be joined in syllogism with the most diverse, indeed even opposite things. \dots\ E.g., if one wishes by syllogisms to prove, first that man is an animal, next that man is not an animal, then one says in the first case: ``Man has senses, organs, circulatory, digestive, nervous system, etc., in short, all the conditions for animal life. Whatever has these conditions is an animal. Thus man is an animal.'' --- In the second case one says: ``The animal is not a rational being. Man is a rational being. Thus man is not an animal\footnote{J.\ L.\ Heiberg. Speculativ Logik. §~148. Anm.~1.}
''.

What is true in this objection is however obscured by the fact that syllogism and proof are without more ado run together into one. Here, as far as the syllogism is concerned, there is only a question of cognition's formal side, not of its real content. The essential development of form must indeed also be a development of content, but it is then a logical, not a real content that is being thought. It is the articulating determination of the connection of judgments, and thus of cognitions, that gives the syllogism, the mediate as well as the immediate, its essential significance. That now%
{} a \textit{methodus mathematica}, which binds advancing cognition to a chain of continuous syllogisms, cannot be commended in philosophy, is a settled matter; but the syllogism's significance as a scientific reduction-formula ought not on that account to be overlooked. The Anselmian proof of God's existence, which, when it is, if we may so say, seen in foreshortened perspective, allows itself to be reduced to the following syllogism: ``the most perfect being (\textit{id quo majus cogitari non potest}) must have, not a merely thought, but also an actual being; God is the most perfect being; thus God has, not a merely thought, but an actual being,'' is indeed, considered as proof, unsatisfactory; but as a reduction-%
% --- p. 326 ---
\opage{326}formula for an infinitely rich content the syllogism is nonetheless of pervasive significance. For the major premise contains the whole of ontology, the minor premise the whole of the philosophy of religion, and the syllogism itself the insight that both sciences in conjunction bring forth. Certainly the single syllogism, and still less the single judgment, cannot give the richer content an exhaustive, secure and all-sided designating form; but for all that it is nonetheless of importance that great totalities can, for punctual illumination, be concentrated in a single syllogism, indeed in a single judgment.

It is further alleged as a defect in the use of syllogisms that an infinite series is established which science cannot sublate. ``Since the two judgments contained in the syllogism,
\textit{E}---\textit{S} and \textit{S}---\textit{A}, are presented as isolated propositions, as premises, i.e. as something that does not lie immediately in the syllogism but precedes it, it is required that each of these propositions shall previously have arisen as a concluding proposition in a preceding syllogism, or, as one says, that it shall have been proved. This presupposes two preceding syllogisms, but since the premises of these must likewise have been proved in what precedes, these two thus presuppose four preceding ones, these again eight preceding ones, and so on in an infinite series, which in its 21st member would comprise about a million and a half syllogisms, and thence rapidly rise to billions and trillions \dots\ To believe oneself in possession of the truth because one finds oneself in a system of syllogisms is still more unreasonable than to believe oneself free because one is permitted to walk about in a cage; for even the smallest birdcage is larger in relation to the whole earth than the cognition brought forth by syllogisms is in relation to that which they presuppose\footnote{Heiberg. Speculativ Logik. §~148. Anm.~1.}
''.

But as regards the objections in this direction also,
% --- p. 327 ---
\opage{327}it must be remembered that the demonstration of a misuse proves nothing against the true use. It goes with the syllogistic on the whole as with the finitizing of thought. If the finite were not present in the development, the infinite could not be comprehended at all. If the immanent method can treat \emph{reflection's} endless series as conditions for the finite's transition into the infinite, why then should it not be able to treat \emph{the mediate syllogism's} endless series in the same way; if the series of reflection allow themselves to be sublated, why should the syllogistic series not also allow themselves to be sublated?

There is in the whole of Hegelian logic not a section whose content cannot with ease be reduced to syllogistic form; but if the content of the immanent development can be reduced to syllogism, then the syllogism must surely also be reducible to immanent development. Let one begin, e.g., with the beginning: ``In a system there can be only one beginning that is the true one; logic is a system, thus in logic there can be only one beginning that is the true one. It lies in the nature of the beginning that it must be contentless; only ``pure being'' is contentless; thus logic must begin with pure being. ``Pure being'' is absolutely contentless; the absolutely contentless is nothing; pure being is thus nothing.'' --- In this way one can then go on through being, essence and concept, and transform even the very last syllogism into a scholastic syllogism%
\footnote{That one may along the way have use for all four syllogistic figures, and, when it finally comes to performing feats of skill, even get occasion to apply both \textit{Darapti} and \textit{Bamalip}, both \textit{Ferison} and \textit{Fresison}, is a matter that is understood of itself.}
.

But what do we now learn from this? In the first place, that the syllogistic form is an abbreviation characteristic of the reduction, an abbreviation quite suited to gather together and
% --- p. 328 ---
\opage{328}preserve the results of a development already traversed.
In the second place, that the manner in which the results won are apprehended and presented in the relatively independent judgments which make up the syllogism's premises corresponds exactly to the manner in which the immediate premises are again dissolved and surrender their content to a development of the concept liberated from syllogistic constraint, a development of the concept that is then identical with the logical development of the matter. One need only consider how, e.g., the two premises: ``pure being is absolutely contentless, and the absolutely contentless is nothing,'' are laid out, grounded and clarified through in the immanent development of the concept, and one will soon discover how the syllogistic series allow themselves to be sublated. In the third place, that every logical development of content, whether it otherwise proceeds in the immanent or in the syllogistic form, must necessarily pass through a system of judgments. But a system of judgments is impossible without a corresponding system of syllogisms. What indeed is the syllogism other than the logical expression for a system of judgments, the immediate syllogism for a system of two judgments, the mediate for a system of three judgments, the chain-syllogism for a system of several judgments, etc.

``Many,'' says Heiberg\footnote{Specul.\ Logik. §~148. Anm.~1.},
``stand in the mistaken belief that when one adduces grounds in one's discourse, these are cryptic syllogisms. But now one gives grounds, and then stands in reflection; now, when one stands in ideality, one brings forward syllogisms, but never scholastic syllogisms. When one says: ``Caius is imperfect, because he is a man,'' that is no syllogism. One states the ground of his imperfection. His human nature is the essence whose phenomenon is imperfection.'' If these and similar misunderstandings did not hang together in the innermost way with the Hegelian mystification of the judgments' and
% --- p. 329 ---
\opage{329}the syllogisms' metaphysical objectivity, we should not dwell upon them further; but there is in the manner of viewing here set up a crookedness which for the principle's sake absolutely must be corrected. The crookedness is that the syllogisms, like the judgments, are to figure objectively at a stage where they have not even yet become subjective. If one says: ``Caius is imperfect, because he is a man,'' one certainly sees the imperfection grounded in human nature, insofar as one sees objectively; but the proposition: ``Caius is imperfect,'' is surely nevertheless a judgment, and the proposition: ``Caius is a man,'' likewise. In order to get sight of the syllogism, one must see to getting sight of the judgment; in order to get sight of the judgment, one must not only see objectively, but also subjectively. The grounding: ``Caius is imperfect, because he is a man,'' evidently rests upon the syllogism: ``Caius is a man, thus imperfect,'' an immediate syllogism which reflection with infinite ease transforms into a mediate one: ``Every man is imperfect, Caius is a man, thus'' etc.

Once it is cognized that judgment and syllogism must be determined subjectively, and thus originally presuppose a judging and syllogizing subjectivity; once it is cognized that every thought, whether it otherwise stands in ``reflection'' or in ``ideality,'' can be determined as thought only insofar as it is determined in the word by a judgment: then it will also be seen that the connection of thoughts is one with the connection of judgments, and the connection of judgments, even where the syllogistic form is not expressly brought out, essentially one with a system of immediate and mediate syllogisms.

\textit{γγ)} \emph{The immanent Syllogism: dialectical Modal Functions.} How the judgment as judgment becomes form for the absolute content has already been elucidated under ``Functions of the Judgment.'' Since now the single judgment first obtains its significance in the connection with other judgments, and reflects the connection the more deeply the richer its content is,
% --- p. 330 ---
\opage{330}thus in an infinite depth when it has the idea, the absolute itself, for content: it follows from this that only an infinite system of judgments will be capable of containing the absolute. But an infinite system of judgments is, when the middle terms are emphasized, \textit{eo ipso} an infinite system of syllogisms. As the form for the absolute develops itself into a system of syllogisms, it is a matter of course that all parcelling-out and finitude must be sublated. The single premises then become moments for the corresponding syllogism, as each single syllogism becomes a moment for a relative totality of syllogisms, and each relative syllogism-totality again a moment for the totalities' pervading total-syllogism. In order to be able to become a moment, each of the particular terms, be it premise, single syllogism or syllogism-totality, must have a particular determinate being; for the moment arises only through a disappearance of the particular determinate being. Through the reciprocal determination of the terms' moment-being and immediate determinate being that thus sets in, the modal functions become dialectical.

Insofar as the premises of the single syllogism give up their particular determinate being and become moments in the act of syllogizing, the syllogistic is abolished: the syllogism is \emph{immanent}. The immanent syllogism can thus be a moment in a totality of syllogisms and a syllogism-totality for itself. As the absolute's subjective form the immanence-syllogism presents three principal sides, in that it corresponds originally to \emph{subjectivity itself}, to the development of \emph{subjectivity's content} and to \emph{the free necessity} with which subjectivity possesses its content. We consider each of these principal sides by itself.

1) \emph{The immanent Syllogism as Form of Subjectivity itself.} That subjectivity's absolutely original syllogism must correspond to an absolutely original judgment can be expressly demonstrated. If the schema \textit{E}---\textit{S}---\textit{A} is to hold for an original syllogism of subjectivity, then surely every one of
% --- p. 331 ---
\opage{331}its members have absolute originality. E, the singular, cannot
then signify a certain singular, this or that singular, a singular among
several others, but singularity itself, singularity as principle. But
singularity as principle is selfhood: \emph{the absolutely singular is
the absolute self}. A, the general, can then for its part just as little
be an abstraction of something common to many, an abstract general, for
an abstract general is only a relativity; the absolutely universal
is\footnote{Grundideernes Logik. I. p.~181--192.} infinitely concrete:
\emph{the absolutely universal in all concretions is the absolute self}.

What holds abstractly in the uttermost abstraction, namely that the
infinite opposition is an infinite unity, and conversely, the same holds
concretely in the completed concretion: the singular is absolutely one
with the universal only by being absolutely opposed to the universal. If
now the unity of the singular and the universal cannot be immediate,
since it rests upon opposition, and the opposition not immediate, since
it rests upon the unity: then it follows of itself that here there must
be a mediation. Only by the determination of this mediation can the
relation between reflection, judgment and syllogism be more closely
fixed.

In reflection the absolutely singular is indeed separated from the
absolutely universal; here there is posited in one and the same selfhood
a thoroughgoing opposition between self and self. But this opposition
vanishes again for reflection as quickly as it is posited, for the self
is after all itself, from whichever side it be seen: the self-opposition
is a reflex in the unity, just as the self-unity for its part is a reflex
in the opposition. It is otherwise in the original judgment: that the
self in the absolute sense is both subject and predicate means that the
opposition of subject and predicate is just as ineradicable as the unity.
While reflection incessantly levels out the
% --- p. 332 ---
\opage{332}opposition of opposition and unity, the judgment sets itself
with decisive emphasis against this all too fleeting levelling, and
expressly demands the subsistence of the opposition in the totalized
unity. Since now it is in an equally absolute sense a serious matter with
the identity and the opposition of the members, the opposed members, the
subject (the absolutely singular) and the predicate (the absolutely
universal), form extremes whose middle term, S, \emph{the absolutely
particular}, shoots up out of the ground of identity. The original
judgment develops with unrefusable consequence into immanent syllogism.

According to the immanent syllogism, subjectivity is at once the absolute
presupposition of the original judgment and its result. The judging self
cannot possibly be derived from the judgment; the judgment presupposes a
judgment-act, the judgment-act a judging selfhood: upon the judging
selfhood rests the absolute priority of the self. But as the first
logical expression of the original act of selfhood, the judgment is not
directed upon this or that finite object, does not have this or that
relative content; no, the original content of the judgment, of the
original judgment, is subjectivity itself, the judgment-act a first
absolutely necessary act of self-positing, a primitive deed whereby that
which judges in the concept originally divides and apportions between
selfhood and otherness\footnote{``Das Urtheil (ἀπόφανσις, als
Bestandtheil des Schlusses πρότασις = \textit{propositio}) ist gleichsam
der richterliche Bescheid, welchen das Denken über das Verhältniß zweier
Begriffe ertheilt; daher ist das Wort ``\emph{urtheilen}'' nicht, wie
Einige behaupten, von ``\emph{Ur- und Theilen}'' abzuleiten, sondern es
kommt von dem mittelhochdeutschen Worte ``\emph{urteilen}'';
niederdeutsch: \emph{ordelen} (wovon ``Ordalie'') \emph{aussagen}, zu
Recht erkennen, rechtlich entscheiden.'' Drbal. Lehrbuch der
propädeutischen Logik. p.~29. If it really came down, in a logic, to
bringing linguistic and dialectical etymologies into agreement, it would
surely not be difficult to show that ``\emph{ordelen}'' is at bottom
``urtheilen'' and ``urtheilen'' ``ordelen''.}.

% --- p. 333 ---
\opage{333}The judgment in which eternal subjectivity utters its essence
is not, as Fichte assumes, a tautological proposition: I = I; such a
tautological proposition is an uttermost abstraction of psychological
thinking, but no \textit{in concreto} original judgment. In the
concretely original judgment the self, the absolutely singular, is as
subject set over against the self, the absolutely universal, as
predicate. The thoroughgoing opposition, subjectivity's self-opposition,
is posited and determined by this, that the subject in the judgment is
itself the judging subject and the predicate of the judgment the original
concept in which the subject grasps itself as object.

That the fundamental act whereby subjectivity makes itself subject and
predicate, so that the predicate becomes the objectified subject, is
neither an immediate self-mirroring nor an abstract reflecting, but a
logical judgment, acquires decisive significance for the determination of
the immanent syllogism. In order to utter its essence in the judgment,
and by this utterance to separate itself into subject and object,
subjectivity must necessarily hover over its own self-separation and to
that extent posit itself, not merely as the \emph{singular} in the
subject of the judgment and the \emph{universal} in the predicate, but
likewise as that which in the judging selfhood is identical with subject
and predicate and yet different from both, thus as the
\emph{particular}. It is by this threefold self-separation (E--A, E--S,
S--A) that subjectivity transforms the original judgment of its essence
into an original immanent syllogism.

But subjectivity's threefold self-separation, which by a continued
hovering of the self-hovering selfhood can be doubled and repeated to
infinity, would nevertheless sink down into an infinitely empty
formalism, if there were not given here in subjectivity itself a means
whereby extremes and middle could keep themselves separate in their
mutual transition. But such a means is \emph{the Word}. If the Word were
not, the subjective thinking, and therewith subjectivity, would eternally
consume
% --- p. 334 ---
\opage{334}itself; if the Word were not, the immediate other of intuition
would stand as an alien, inexplicable object upon which mute subjectivity
would have to stare itself blind; if the Word were not, the self, the
singular, would be so shone through by the particular and the universal
that the pure light of reflection, the indeterminate clarity
indeterminable in itself, would at last vanish into differenceless unity
and lose itself in pure darkness.

By mediating between thinking and intuiting, and thus making principled
self-determination possible, the Word puts the immanent syllogism into
force. In the Word thinking has its elementary determinations of
otherness, its infinite store of images and representations. In the
actual world, in the infinite realm of the realities, there is not an
object, no whole, no part, not a thing, not a property, without the
Word's being able, by image and representation, to take up and render
what is peculiar therein. Enriched with reflexes of all possible real
determinations, the Word now has that ideality whereby in reflection it
gives up its immediacy and sacrifices itself for thought so inwardly and
absolutely that the thinker through the Word thinks himself in all things
and all things in himself, and comprehends his subjectivity as \emph{the
absolutely universal}.

The absolutely universal is indeed a thought; but it is not a poor
thought, for thought lives upon the Word, and the Word has an infinite
richness. The absolutely universal is indeed, considered for itself, a
word; but it is not an empty word, for the Word that nourishes thought is
enriched by thought, and thought is infinitely rich.

In the thinking enriched by the Word, subjectivity comprehends itself
under all three shapes, that of singularity, that of particularity and
that of generality, but is nevertheless, inasmuch as the universal
syllogizes itself together with the singular through the particular, as
yet posited only in the one-sided form of thinking, of the universal:
subjectivity is thinking. The opposition of subject and object is indeed
thought as a thoroughgoing, absolute
% --- p. 335 ---
\opage{335}opposition, but it is nevertheless as yet only a thought
opposition; the unity of subject and object is indeed thought as an
absolute unity determined by the particular selfhood --- at once hovering
over the extremes and yet dwelling within them --- but for all that it is
nevertheless only a thought unity. The whole immanent syllogism has, by
the one-sidedness of the thinking that thinks through all things, itself
become an extreme.

Over against selfhood in its extreme of the universal there now stands,
not as before the extreme of the singular, but a selfhood-extreme of the
particular. This change of extremes is decisively determined by the Word.
In thinking the Word is negative, its immediacy burns away in reflection;
only the manifoldness of illuminations, shadings, gleams and reflexes
makes up the richness of the subjectivity that thinks itself in all
things and all things in itself. That reflection cannot now consume this
richness and plunge subjectivity into such poverty that it looks like a
shadow of abstraction, a psychological spectre of I = I, is prevented by
the Word; for the Word, which is negative in its unity with thought, is
positive in its unity with intuition. The Word must indeed sublate all
sensuous immediacy and to that extent effect a transformation of all
sense-images; but the pictorial elements that are consumed by thinking
are transfigured in intuition and determined by the Word. Insofar as the
Word is negative, the image passes through the representation over into
thought; but insofar as the Word is positive by means of intuition,
thought is resolved into the representation, and the representation into
a transfigured image, a thought-vision, an intellectual intuition. In
infinite manifoldness of visions, images and representations the
intuiting subjectivity is all that it is and can be in the universal as
well as in the singular, in the Word, by the Word, as the Word. By its
positive side in intuition and its negative side in thinking the Word is
the particular. It is not a single extreme of subjectivity, but the whole
subjectivity in extremes and middle, that is transfigured by the Word.
The thinking subjectivity, objectifying itself in the Word,
% --- p. 336 ---
\opage{336}is thus determined as \emph{the particular}: by means of the
particular as object, the self as subject is the universal by being the
singular, and the singular by being the universal.

And yet there is here still an indeterminateness, a wavering, a touch of
formalism, which threatens to hollow out the syllogism and to transform
the doctrine of subjectivity into logical phraseology. ``How can
subjectivity, as that which is particular in the Word, at one and the
same time both hold itself in opposition to the singular on the one side
and to the universal on the other, despite this opposition likewise be
identical with the singular as well as with the universal, indeed by this
identity bring it about that the singular is itself the universal, the
universal the singular?'' Such a question is, from the standpoint of
formal logic, not merely natural but, in the present connection,
perfectly justified. By insight into subjectivity and the immanent
syllogism the matter can however be cleared up.

All determination of subjectivity is essentially self-determination, but
a self-determination that presupposes another and is conditioned by
another. That the complete determination of subjectivity assumes a
threefold form has its ground precisely in this, that the determination
of the other is to be resolved into the determination of the self. On the
level of relativity and finitude this is now an impossibility. The
determination of the other stands, in the psychological sense, as a
barrier, a hindrance to self-determination. Face to face with the other,
the mighty barrier, self-determination is impotent, abstract and formal.
The impotence is then to be seen upon thought and word alike, for pure
self-thinking is an empty reflecting, and the word a name that is capable
of nothing over the object. Referred to the subjectivity that is formal
and abstract in its limitation, the immanent syllogism with its extremes
and middle must naturally also become formal. The difference between the
self-thinking and the thought self, between the self-objectifying and the
objectified subjectivity, is here a
% --- p. 337 ---
\opage{337}difference of word, thus only a difference of name: the
determinations are nominal. But when the talk is not of the
psychological, but of the ontological, not of the relative, but of the
absolute subjectivity, nominalism must vanish.

It is the absolute ideality of the almighty Word upon which the emphasis
falls. Already in the doctrine of the idea-forming functions the thinking
subjectivity was grasped as the origin of the Word, the intuiting
subjectivity as the Word itself, and the subjectivity determining itself
by thought and Word as the developed unity of both. But so long as the
judgment had not passed over into syllogism, the knowledge-word with its
ideality and the power-word with reality were still so finitely separated
that power one-sidedly ran upon otherness, while knowledge ran upon
self-determining and selfhood. In the progressive development of the
functions, what was one-sided in the opposition was gradually sublated;
in the passage of the judgment into syllogism the ideality of knowledge
and of power was expressly posited in the Word. Now that the Word thus
acquires subjective, ideal significance from power and knowledge alike,
its mighty ideality, capable of all things and overreaching in all
things, comes to light.

What the Word determines concerning subjectivity's threefold
self-separation in the threefold form of the immanent syllogism, it
determines not merely on behalf of knowledge but also on behalf of power,
not merely on behalf of selfhood but also on behalf of otherness:
subjectivity's determination by the Word is its own absolute
self-determination, for the origin of the Word is, in its opposition to
the Word, identical with the Word.

If we now let \emph{the middle term} indicate the difference between the
three syllogisms of the syllogistic system, then the schema E--A--S
becomes designative of subjectivity as \emph{the absolutely universal},
E--S--A of subjectivity as \emph{the absolutely particular}, and S--E--A
of subjectivity as \emph{the absolutely singular}. That every one of the
members of the total syllogism is itself a syllogism, and every one of
the particular syllogisms in turn a
% --- p. 338 ---
\opage{338}moment in the total syllogism, is conditioned by the Word; we
shall now see how the Word works as the ideality of power.

The thinking subjectivity uses the Word, as has been said, negatively. In
the immediate word of knowledge and of power the ideal richness of
all-existence is enclosed; but the immediacy is sacrificed to thought.
What it means that the Word in the infinite thinking must be used
negatively can be elucidated from yet another side. It does indeed make
an essential difference whether, with such expressions as: the singular,
the particular, the universal, we reflect upon the word or upon the
thought. If we reflect one-sidedly upon the word, then we have here, in
the singular, the particular and the universal, three different words.
These three different words are, in all superficiality, three different
appellations; each appellation then requires, if it is to signify
anything, a peculiar representation, an intuition, an image of what is
named: by the immediacy of the word we are led away from the thought over
into the circuit of images, intuitions and finite representations. It is
otherwise when reflection is here expressly made upon the thought, the
thought made almighty by the negativity of the Word. The three words are
now three ideal power-determinations, three substantial concepts; in this
sense it is with the Word as with light: we see the objects in the light
and forget the rays of light; we see the thoughts in the Word and forget
the appellations.

Subjectivity, the thinking One (E), receives through the Word's negative
working no images, no visions, no finite representations, but an
incessant vanishing of countless images, visions and representations,
countless particular thoughts. In particular thoughts subjectivity posits
an essential difference between the thinker, the thinking and the
thought; by this difference it thinks itself particularly. The
subjectivity thinking by the Word (E) and the subjectivity thought in the
Word (S) now form substantial extremes; but in the negative functions of
the Word, in the infinite thinking, these extremes have their essential
middle. Insofar as the infinite thinking is the middle of the extremes,
thus the working form in which the thinker is himself the thought,
subjectivity is determined by determining itself as the infinite unity
that thinks through all things,
% --- p. 339 ---
\opage{339}that is capable of all things in thought, thus determined as
\emph{the absolutely universal} (A). The schema E--A--S is satisfied.

By the infinite self-thinking subjectivity may indeed be said to brood,
sunk deep in its absolute unity; the unity is an infinite deep, but the
deep is clarity: magical visions, gnostic aeons, dream-images of βυθὸς
and σιγὴ, προπάτωρ and ἔννοια, vanish before the energy of thought in the
immanent syllogism. The energy of thought is the Word's negative working.
In its negative working the Word is hidden; but the hidden Word must
become manifest, its negative working pass over into the positive.

The thinking subjectivity, sunk deep in itself, which keeps the Word
hidden, is \emph{the origin of the Word}. The Word does not come to the
aid of thought from without; the Word that is determinative for thought
is originally determined by thought itself: the Word is with necessity in
thought, it is born of thought. Born of thought, the Word first acquires
positive significance as word: the thinking subjectivity becomes
intuiting in reflection, and the schema of the syllogism is changed into
E--S--A.

With the reflected immediacy of the Word \emph{the particular} comes into
force. Thinking must indeed, as was shown above, presuppose both the
particular and the singular; but in thinking as thinking the said
extremes can work only negatively, through reflexes; only when the
immediate emerges from the stream of reflection, and thought is coloured
by the light of intuition, does the difference break forth, and the
particular come into its right. With the stamp of immediate intuitability
the particular is likewise determined as the objective, the outwardly
manifold; the intuiting subjectivity (E) now fixes its gaze upon the
particular (S) and sees a manifoldness of visions. But they are not
sense-images; they are visions in words, a multiplicity of particular
judgments, that stands before the gaze of subjectivity. In each of these
judgments the word is particular; the immediacy is in the subject, the
reflection in the predicate: subject
% --- p. 340 ---
\opage{340}and predicate are determined particularly by particular words;
the unity of immediacy and reflection must then also be determined by a
particular word. The particular word is relative; the relativity lies in
the manifoldness. By its relativity the word of selfhood has taken up
otherness; to otherness objectivity attaches; the word particularized in
objective determinations thus stands as subjectivity's other over against
subjectivity itself: by the particular word the intuiting subjectivity
itself becomes a particular (S).

If subjectivity is set sharply over against the Word, then it shows
itself clearly how the particular comes forth; it is not merely
subjectivity that becomes particular by means of the Word, it is, in
accordance with the original unity of Word and thought, just as much the
Word that becomes particular by means of subjectivity. If subjectivity is
in the Word \emph{subjectively} two-sided by being itself and its other,
then the Word is in subjectivity \emph{objectively} two-sided by being
itself and its other; if subjectivity, with its relative acts of
thinking, particularized by relativity, is an absolute unity, then the
Word, with its relative elements, particularized by relativity, is
likewise an absolute unity. With its gaze fixed, not upon the Word in
itself, but upon \emph{the Word's other}, upon the word
\emph{particularized} in manifoldness and difference, the intuiting
subjectivity sees only \emph{its other}, not \emph{itself}: the
particularity is on both sides relative. By losing itself in the
relative, the intuiting subjectivity would then not merely become
particular over against its other in the Word, but likewise particular
over against itself in the thought; the intuiting subjectivity would, in
other words, not be thinking, nor the thinking intuiting; by a finite
separation of thought and intuiting, the infinite subjectivity would sink
down into finitude and impotence; but that is a contradiction, indeed an
infinite self-contradiction.

In order to sublate the self-contradiction, the intuiting subjectivity
must therefore raise itself above the relative, transfigure
% --- p. 341 ---
\opage{341}thinking by intuition, intuition by thinking, and grasp itself
absolutely by comprehending the Word. As subjectivity is, so is the Word;
as the Word is, so is subjectivity. In the absolute Word the thinking
subjectivity intuits its own image. There is in this image an objective
immediacy, for the thinking subjectivity is \emph{intuiting}; but
subjectivity's image is subjectivity itself, thus a non-pictorial image,
an almighty thought-image, the subjectivity objectified in the Word and
as the Word. In the absolute self-objectification, therefore, thinking
goes together with intuition, the general (A) together with the
particular (S); herewith the syllogism has reached its completion in the
form of particularity.

``But'' --- one might ask --- ``is there now also any true logical
necessity that the syllogism should repeat itself in the form of
singularity as well? It really does seem as though everything had already
been said here that can in any way be said, from the general standpoint
of subjectivity, about the immanent syllogism. If it is given from the
beginning that subjectivity itself is the absolutely singular, the
absolutely particular and the absolutely universal alike, then it must
surely be an art as easy as it is unfruitful to transpose the three
members in three different ways, to justify these transpositions
calculated upon the middle term by suitable paraphrases, and to let
Hegelian combinations of E, S and A step into the place of
\textit{Barbara, Celarent} and \textit{Baroco}''.

It is however no empty schematism, but a clear exposition of the absolute
energy of the Word and of thought, that is aimed at here in the present
connection. If it came down merely to the designation, the linguistic
expression, the categorical formal side, then the word energy has, in
Aristotle and after him in the Neoplatonists, the Scholastics and
Leibnitz, been grasped, determined and explained from many sides. But
when it is not the word energy, but the energy in the Word, at the
cognition of which the aim here essentially lies, definitions,
% --- p. 342 ---
\opage{342}paraphrases and formal explanations will not suffice. The
primitive energy of the Word, the energy in the divine Word, is
determined by a unity-in-opposition of knowledge and power: in the Word
and by the Word knowledge is different from power, in the Word and by the
Word knowledge is identical with power. From this it follows that, when
the unity is posited, the opposition breaks forth with irresistible
force: in this irresistibility is the energy; conversely, when the
opposition is posited, the unity raises itself with overreaching,
irresistible superiority: in this superiority the energy is seen. Since
now the Word, in the absolute sense, is identical with subjectivity, the
energy of the Word must likewise be the energy of subjectivity; but
subjectivity's absolute energy is its absolute self-determination.

By absolute self-determination the subjectivity identical with the Word
then grasps itself in all things, all things in itself, and thereby
becomes absolutely determined as \emph{the absolutely singular}. In the
determination of subjectivity as the absolutely singular the Word works
at once negatively and positively. It works negatively; the absolute
self-determination is negative energy, for the self can be the absolutely
singular only by being the absolutely universal. But since the universal
in the form of thinking has its mighty opposition in the intuition of the
particular, the self-determination becomes recognizable by the energy
with which the negative word of thinking masters and sublates immediacy.
Yet the Word can work negatively for thinking only by working, with equal
energy, positively for intuition. The absolute self-determination is
positive energy, for the self can be the absolutely singular only by
positing itself against everything other and everything other against
itself, thus only by becoming the absolutely particular. The negative
energy with which subjectivity excludes the other in order to posit
itself is a force-measure for the positive energy with which subjectivity
posits the other in positing itself, and conversely. In energetic unity
of negative and positive working, the absolute self-determining is an
absolute
% --- p. 343 ---
\opage{343}willing: only the willing subjectivity can in the proper sense
be said to be \emph{the absolutely singular}. With the absolutely
singular as middle term, the immanent syllogism is developed into being
the adequate form of ontological subjectivity.\footnote{``But'', the
attentive reader will perhaps ask, ``is this not a masked attempt to
explain the Trinity, and thus to encroach upon that Church privilege of
speculative dogmatics which excludes all philosophical concurrence?'' We
answer that the rational forms with which dogmatics, in order not to
restrict itself outright to a compilation of scriptural passages, is wont
to adorn its doctrine of the Trinity are all borrowed, sometimes even
scraped together, from various philosophical systems. An all-sided
development of the immanent syllogism as subjectivity's original form
falls, on the contrary, wholly and entirely under logic; here nothing is
borrowed from without, least of all from dogmatics. What in the objective
doctrine of the Trinity is genuinely theological, and to that extent
belongs exclusively to dogmatics, philosophy, and especially logic, will
take good care not to covet. To this belong such ``dogmatic
elucidations'' as, e.g.: ``If it can be said that the one God as it were
looks into the world with three faces (τρία πρόσωπα), then it must also
be said that these faces are not merely turned outward toward the world,
but that they are also turned inward toward themselves, that they see
themselves in mutual mirroring. Otherwise they would not be the
revelation of God's true inwardness, but only deceiving masks''. That is
an important, a very important ``dogmatic elucidation'', for from it one
sees that each face must, in order not to be ``a deceiving mask'', at one
and the same time be turned both inward and outward, and consequently be
a double face. In order not to have ``three deceiving masks'' by merely
having three outward-turned faces, the Triune of dogmatics must therefore
have, at the very least, six faces. Further: ``the whole Trinity stands
in one present Now, three eternal flames in the one light''. Further:
``the God revealed to himself, to whom all is revealed, \emph{has} not an
eye, but he \emph{is} eye''. Thus: the Triune of dogmatics \emph{has} six
faces, but not an eye. For he \emph{is} eye, a triune eye: three flames
in one light, six faces in one eye! To get such an eye, with such flames,
in such a light, with such faces, into science --- for that supernatural
powers are required, indeed for that are demanded faith and
\textit{testimonium spiritus}.} In
% --- p. 344 ---
\opage{344}this form subjectivity has rest, for extremes and middle are
simultaneous; but this rest is an eternal movement, for extremes and
middle are posited and transposed to infinity.

2) \emph{The immanent syllogism as form for the content of
subjectivity.} In the absolute subjectivity the content is itself
absolute; the difference between subjectivity and content is thus a
difference in the absolute itself, a difference between absolute and
absolute. But the difference in the absolute itself is one with
relativity. Since now the absolute must absolutely have the relative as
an immanent determination, the difference between subjectivity and the
content can rest solely upon the twofold relation in which the relative
stands to the absolute: \emph{either the relative is resolved into the
absolute, or conversely the absolute into the relative}. If the relative
is resolved into the absolute, then the immanent syllogism gives us, as
was shown above, subjectivity itself; if on the contrary the absolute is
resolved into the relative, then the syllogism gives, as shall be shown
in what follows, the immanent form for the content of subjectivity.

The difference in the relative's relation to the absolute is recognizable
upon every one of the members of the syllogism. That the singular,
insofar as the relative is resolved into the absolute, must become
subjectivity itself, has been shown; but when it is now the absolute that
is resolved into the relative, what then becomes of the singular? A swarm
of singulars! That the universal, insofar as the relative is resolved
into the absolute, must become subjectivity itself, has been shown; but
when the absolute is now resolved into the relative, what then becomes of
the universal? A plurality of universals! And what, on the same terms,
becomes of the particular? A multiplicity of particulars!

In subjectivity itself, the absolute subjectivity, there is no finite
difference whatever between genus, species and individual. Only with
relativity does the separation come. It is relativity that in the genus
effects the separation of the many
% --- p. 345 ---
\opage{345}species, and in each species again the separation of the many
individuals. Sprung from the absolute, relativity is principled, indeed
itself a principle: the principle of relativity is the principle of
particularity; otherness in thought and word is the energy of the
principle.

If the principle of particularity did not, through otherness, work
through the infinite content of subjectivity, then the subjectively
objective, the differenceless universal that levels out all thoughts, all
concepts and forms, would rest in infinite uniformity and condense itself
into eternal darkness. But when the principle of particularity sets the
universal trembling, when thought shines in the Word, and the Word, with
the light of thought in the colour-gleam of relativity, disperses the
darkness and opens an entrance for cognition: what visions then present
themselves to intuition, what shapes will the gaze of thought then meet
in the realm of the ideas? Surely it can hardly be the ``playing
possibilities'', ``plastic prototypes'', ``magical actualities'' of ``the
speculative dogmatics'', a ``heavenly host'' of those ``visions'' which
cannot be content with a ``revelation \textit{ad intra}'', but are in
importunate tension to slip out and be ``released to a revelation
\textit{ad extra}''? By no means! In order to see such visions,
speculation in company with faith must first take in a quantum of
ecclesiastical \textit{spiritus}, for in a sober state thought sees
something quite different. ``Magical visions'', speculative
hallucinations are not fitting for a divine subjectivity, and poorly
enough for a human one. With a host of such visions individuality may
perhaps become a fermenting mystic, but not a poet, still less a thinker.

The visions that arise for the poet, the prototypes that hover before the
thinker, have, as influxes from the source of the ideas --- however
strongly the stream may run --- always the distinguishing mark that they
come with the Word, order themselves into thoughts and are transfigured
in judgments; even in the mythical consciousness the visions must, for
all their magic, subordinate themselves to the power of the Word. But
visions that subordinate themselves to the power of the Word do not have
the
% --- p. 346 ---
\opage{346}``creaturely'' obtrusiveness, that apocryphal importunity, of immediately ``desiring to be released into a revelation \textit{ad extra}''; no, on the contrary! they clarify themselves precisely into revelation \textit{ad intra}, for the word of knowledge in the deep of the ideas is revelation \textit{ad intra}.

In the realm of the ideas, in the deep of the Godhead, there is nothing fermenting and nothing immediately being: all is determined by thought and the word, all consists in infinite reflecting, judging and syllogizing. In the infinite judging and syllogizing the three fundamental determinations that hold for the absolute totality of the content --- the singular, the particular and the universal, or: the individual, the species and the genus --- have the same original validity. It is necessary for the thinking subjectivity to comprehend its content in the form of the universal, since it is itself the absolutely universal; it is necessary for the intuiting subjectivity to apprehend its content in the form of the particular, since in its reflecting intuiting it is itself the particular; it is necessary, finally, for the willing subjectivity, determined by self-determining towards another, that it must represent its content to itself as articulated into nothing but determinate singulars, since it is itself the absolutely singular. The content of subjectivity is thus not parcelled out, so that one part by itself should fall under the universal, another under the particular, a third under the singular; by such a parcelling-out the Absolute would, instead of being absorbed into the relative, be lost in nothing but finite relativities. But when the whole undivided content, in itself indivisible, is to be distributed among the three forms, then the contradiction that the undivided nevertheless underlies a fundamental division can be resolved only in this way, that two of the forms alternately constitute extremes which have the third for their middle, thus only by a system of immanent syllogisms.

In the course of the system's development the movement now goes: partly from the singular through the particular to the universal, whereby the syllogism-series become categorical, and the emphasis falls upon \emph{subsumption}; partly from the universal through the
% --- p. 347 ---
\opage{347}particular to the singular, whereby the syllogism-series become disjunctive, and the emphasis falls upon \emph{disjunction}; finally it is discerned what it means that the movement goes in both directions at once, so that the syllogism-series are totalized, and the decisive emphasis falls upon \emph{individuation}.

\emph{Categorical Immanence-Syllogisms: The Subsumption-Problem.} In formal logic there smoulders a remarkable conflict between reflection and the judgment, a conflict that concerns precisely \emph{being subsumed} under something and \emph{being} something. According to the judgment the individual is its species, the species its genus. But according to reflection the individual cannot possibly be its species, nor the species its genus; for the species, which has a greater extension than the individual, is what the singular individual is not; the genus, which has a greater extension than the species, is by its greater extension what the species is not, and the species, which has a greater content than the genus, is by its greater content what the genus is not. For finite reflection the singular cannot possibly be the particular, any more than the particular can be the universal.

Although reflection in formal logic is otherwise eager enough to secure to the judgment its general form, that the copula \emph{is} shall not be confused either with the verb \emph{is} or with the reality \emph{is}, it nevertheless reserves to itself the right to decide what meaning may lie in the copula of the judgment. The judgment reduced to its typical form says straight out: the singular is the universal; this is a horse, the individual is its species; the horse is a pachyderm, the species is its genus; but reflection assures us that it cannot possibly be meant so. The singular is not the universal, but the meaning is that it is subsumed under the universal; the individual is not species, but the meaning is that it is subsumed under the species etc.

It is characteristic that the judgment, when it literally answers to the form determined by logic itself, cannot be understood literally. Every time the judgment assumes its logical stamp,
% --- p. 348 ---
\opage{348}the paraphrase ``\emph{is subsumed under}'' is displaced by the copula ``\emph{is}''; every time reflection steps in with its explanation, ``\emph{is}'' is paraphrased into ``\emph{is subsumed under}''. Out of this tension there arises, then, \emph{the subsumption-problem}.

In formal logic the subsumption-problem cannot come forward for sheer subjective reflecting; in ``speculative'' logic it cannot come forward for sheer objective being. From the standpoint of finitude and of one-sided concepts it is certainly a contradiction that the subject should be its predicate. Here there is indeed not merely a difference, but a \emph{finite} difference between the singular, the particular and the universal, between thing and concept, between individual, species and genus. Now in order properly to bring out the finite difference without giving up the unity, emphasis is laid upon subsumption. It is otherwise from the standpoint of the idea. In the idea there is indeed also difference, a pervasive difference between the individual and the species, between the species and the genus; but in the pervasive difference finitude is sublated. The individual is in an \emph{individual} manner what the species is in a \emph{particular}, and the genus in a \emph{general} manner. There are no other and no more conceptual determinations in the individual than there are in the species: the identity of individual and species is absolute, the individual \emph{is} its species. But the conceptual determinations, which in the individual are all individual, singular-moments in the singular, are in the species all particular, reflection-moments in the peculiar thought-image: the difference between individual and species is absolute, the individual \emph{is subsumed} under its species.

As regards next the relation between the species and the genus, it might well seem at first glance as though the species contained more and consequently other determinations than the genus; but this way of looking at it has validity only for abstraction. Seen in the idea, the species has the same totality of determinations as the genus; the difference rests solely upon the mode of
% --- p. 349 ---
\opage{349}determination.\footnote{Cf.\ Grundideernes Logik. I. p.~218--20.} In opposition to the genus as the universal, the species is the particular, and therefore the species-determinations have the form of \emph{particularity}, the genus-determinations the form of \emph{generality}. If we let, e.g., the categories swine, tapir, horse and elephant represent the species, and the category pachyderm the genus, then it is easy to show that the species cannot possibly contain more determinations than the genus, inasmuch as the genus, just as well as the species, is not an empty schema but a relative absolute. The horse is not a pachyderm and then something more, but the pachyderm itself in a peculiar shape. Nor, then, do the horse and the elephant stand to one another in such a way that they should have a complete likeness in being pachyderm and, alongside this, an unlikeness in the form-determinations that fell outside being pachyderm. No, the likeness is pervasive, since the unlikeness is pervasive. There is indeed a difference between the foot-form of the elephant and that of the horse, but this difference falls precisely in the genus, just because it falls in the species: the elephant's foot is the pachyderm's foot just as well as the horse's foot is. The genus is indeed different from its species in that the species are mutually different from one another; but by this difference the genus is at the same time identical with its species. If, then, reflection is here directed upon the identity, the species \emph{is} its genus; if reflection is here directed upon the difference, the species must \emph{be subsumed under} the genus.

The double-sided relation between individual, species and genus finds its logical development and confirmation in the categorical immanence-syllogism, which we may therefore also call \emph{the subsumption-syllogism}. When the judgment has posited abstract unity, and thereby abstract difference, between the singular and the universal, then the syllogism introduces the particular as middle term, and by this middle term reconciles subsumption with being. In this respect it makes an essential difference whether a judgment is apprehended immediately as a judgment or as a conclusum. If we say
% --- p. 350 ---
\opage{350}``the horse (the individual) is a pachyderm'', ``the elephant (the individual) is a pachyderm'', then we are tempted to set up the two immediate judgments as premisses and to draw the inference that the horse is what the elephant is, that the horse is elephant and the elephant horse. It is otherwise, however, when each of the judgments adduced is apprehended as a concluding proposition, as a conclusum. The horse (the individual) is then not without more ado a pachyderm, but is subsumed under the genus pachyderm; the individual horse is not without more ado the species horse, but is subsumed under the species. If the individual were immediately subsumed under the genus, subsumption would one-sidedly displace being; but inasmuch as the species steps in between in such a way that the individual is subsumed under the species, and the species under the genus, being emerges out of subsumption, for the individual \emph{is}, as we have seen, essentially its species, the species essentially its genus.

The immediate judgment ``the horse (the individual) is a pachyderm'' is then, when we look more closely, a consequent proposition torn loose from its premisses. In reflected unity with the premisses the judgment as consequent proposition will run thus: this individual is, as horse, a pachyderm; thus, as regards the remaining judgments, this individual is, as swine, a pachyderm; this individual is, as elephant, a pachyderm, etc. But in this way the species-difference is laid into the genus. The pachyderm is not without more ado pachyderm; for the pachyderm horse is different from the pachyderm elephant, from the pachyderm swine, etc. The genus is thus identical with itself only in that it posits and deposes the species, just as the species is identical with itself only in that it posits and deposes the individual; the universal is identical with itself only in that it posits and deposes the particular, just as the particular is identical with itself only by positing and deposing the singular. But that the genus deposes the species, and the species in turn the individual; that the universal deposes the particular, and the particular in turn the singular, is logically grounded in disjunction.

% --- p. 351 ---
\opage{351}\emph{Disjunctive Immanence-Syllogisms: The Disjunction-Problem.} Formal logic's schema: \textit{A} is either \textit{B} or \textit{C}, now \textit{A B}, thus not \textit{C}; or: now \textit{A} is not \textit{C}, thus \textit{A B} --- yields a meagre return for cognition so long as it is held fast in empty abstraction. As form for the idea, the absolute content, the disjunctive syllogism acquires, on the contrary, an infinite significance. If the subject (\textit{A}) is the universal-absolute, then the terms of the predicate (\textit{B} and \textit{C}) must be particular-absolutes. The particularity is expressly indicated by the disjunction: either \textit{B} or \textit{C}, and to that extent falls outside \textit{A}; for \textit{A}, immediately as \textit{A}, is not different from itself, not a particular in itself: \textit{A} is only what \textit{A} is; in \textit{A} the particular is not. And yet, according to the disjunctive judgment, the disjunction falls precisely within \textit{A}. It is precisely \textit{A} as \textit{A} that is either \textit{B} or \textit{C}. Here the judgment is again in contradiction with reflection. In expressly saying that \textit{A} is either \textit{B} or \textit{C}, the judgment presupposes that \textit{A} is at bottom both \textit{B} and \textit{C}, namely in the one case \textit{B}, in the other \textit{C}. Made anxious by the contradiction, reflection would now gladly transform the disjunctive judgment into a mere possibility-judgment, a judgment that determines possibility by limiting and specifying its alternatives. Outside the idea, in the idealess multiplicity of empty forms, it may now indeed be much the same whether one takes the disjunctive judgment for a paraphrase of the problematic or sets it up by itself, for it is certain that disjunction itself one does not understand. Seen in the idea, disjunction is absolute, and the disjunctive judgment's difference from the problematic is recognizable precisely by the syllogism. The distinguishing mark of the disjunctive judgment is just this, that according to its original significance it forms the premiss for as many positive-negative syllogisms as the predicate contains alternating terms, an incontestable proof that disjunction is categorical. If the judgment is apprehended merely as the expression of possibility, then here, however many terms the predicate may contain, there can be had
% --- p. 352 ---
\opage{352}the premiss for only a single positive syllogism. \textit{A} is either \textit{B} or \textit{C} or \textit{D}, that is to say: \textit{A} is perhaps \textit{B}, perhaps \textit{C}, perhaps \textit{D}. Whence come the three alternatives? From uncertainty alone! The alternatives are only subjective, not objective. In an objective sense there is here only a single term; when this is posited, the rest are excluded, their nothingness has become manifest; if, conversely, the rest are excluded, then the remaining term must be posited; with this the matter is settled. The positive in the syllogism is then at once separated from the negative in such a way that, to the multiplication of the forms, two separate syllogisms arise, a positive one by the so-called \textit{modus tollendo ponens}, and a negative one by \textit{modus ponendo tollens}. It is otherwise in the Absolute; here disjunction is what the judgment says it is, objective, i.e. logically objective. The Absolute, undivided in itself and indivisible in the finite sense, is by the disjunction in principle separated into peculiar, mutually exclusive absolutes: either \textit{B} or \textit{C} or \textit{D}; each of the particular terms has the same objective validity. The positing of every term is reflected in the negation of the rest, and conversely; therefore each term must be grasped by a positing of its own and demand a syllogism of its own. Since the positing of each single term reflects itself in the negation of the rest, each syllogism by itself is at once both negative and positive: \textit{modus ponendo tollens} reflects itself in \textit{modus tollendo ponens}, and conversely. In the positing \textit{A} is \textit{B} the negation reflects itself, thus \textit{A} is neither \textit{C} nor \textit{D}; conversely, in the negation: \textit{A} is neither \textit{C} nor \textit{D}, the positing reflects itself: thus \textit{A B}.

And what, then, is the essential result of these syllogisms? nothing but categorical judgments: \textit{A} is \textit{B}, \textit{A} is \textit{C}, \textit{A} is \textit{D}; these judgments, resting upon disjunctive syllogisms, are now the expression of the fact that the universal is a particular.

With regard to subsumption we sought a grounding of the judgment: the particular is a universal; with regard to disjunction we have found just the reverse, namely that the
% --- p. 353 ---
\opage{353}universal is a particular; subsumption and disjunction still fall apart from one another, a sign that the problems are not yet completely solved. The solution of the problems is made difficult by the fact that syllogisms and general reflection-movements are only too easily confused. If we say: ``All is general, the general is all'', then it makes an essential difference whether what is to be thought here is a reflection-proposition or a system of syllogisms. As a reflection-proposition it is a paraphrase of the ground's first proposition, ``All is at bottom one'', that we have before us, and the movement corresponding to the proposition is reduction. When the outer multiplicity of existence is led back to the unity of the ground, the ground-unity is the concretely universal; but as a logical function reduction will herewith at the same time signify that the manifold is dissolved in the ground-unity, that the existent goes to the ground. It is otherwise, however, when the reflection-movement is developed into syllogism. Instead of reduction, subsumption then holds. Subsumption dissolves nothing; it preserves all: the singular does not go to the ground in the particular in order that the particular may in turn go to the ground in the universal; no, the singular is preserved in the particular in that the particular is preserved in the universal. The syllogism is essentially a syllogistic union, and the syllogistic union a categorical expression of the fact that the extremes together with their middle are completely preserved. Insofar as the syllogism ``All is general'' is expressly apprehended as a syllogism, the syllogism is itself a system of categorical immanence-syllogisms, indeed a system of systems to infinity. The movement from the multiplicity of the singular through the particular to the universal is a movement from periphery to centre, a syllogism-movement that in absolute knowledge can run through innumerable intermediate stages. A mere proposition such as ``this (the individual horse) is a living being'' already contains, when it is to be reflected as consequent proposition in its corresponding syllogism, a multitude of essential middle terms between subject and predicate: ``this is
% --- p. 354 ---
\opage{354}a horse, the horse is a pachyderm, the pachyderm is a hoofed mammal, the hoofed mammal is a mammal, the mammal is a vertebrate, the vertebrate is an animal, the animal is an ensouled being; the ensouled being is a living being; thus: this individual is a living being''. If I say: ``this (the writing-desk) is an existent,'' what a distance does there not then seem to be here between the individual in the subject and the universal in the predicate, when we consider that \emph{this} (the house) and \emph{this} (the star) and \emph{this} (the Emperor of China), in short, all individuals, things and properties, must be subsumed under the existent. If determinate being is now posited as the universal, the existent on the other hand as the designation for every determinate individual, any singular whatever, then it is easily discerned that the way from the existent (the subject) to determinate being (the predicate) is at once infinitely long and yet infinitely short; a straight and yet infinitely ramified line.

The thinking subjectivity, which has its whole content in the general, sees all the singular at one glance, because at one glance it sees the infinite system of syllogisms in which each singular, through its middle terms, goes together with the universal.

But now it would indeed be impossible that an infinite system of syllogisms should be thinkable without a corresponding system of middle terms. What, now, is the middle term in categorical determinateness? Always a particular! Whether we choose the schema \textit{E}---\textit{S}---\textit{A} or \textit{S}---\textit{E}---\textit{A} or \textit{E}---\textit{A}---\textit{S}, the middle is, in the character of middle term, always the particular, for in its difference from the extremes it reflects the extremes' relative unity. But thereby it also comes about that the extremes, in their relative unity with the middle, come to reflect their mutual difference, that the extremes, in other words, become particular themselves by virtue of the particular in the middle. But when the extremes become particular by being extremes, since the middle is particular by being middle, then in this way all the content-forms become particular: all is a particular, the particular is
% --- p. 355 ---
\opage{355}all. If we now confuse the syllogism with a reflection-movement, then the reflection-movement is an immanent but ambiguous conversion of reduction into deduction, and conversely. By reduction all, the singular together with the particular, is dissolved into the general; but by taking up into itself all determinations of the singular and the particular, the general must itself coalesce with the totality of particulars. But just as the universal has no other subsistence than its reflection in the particulars arising from the dissolution of the singulars, so these particulars have no other subsistence either than the reflection in the universal determined by that dissolution. The ground-reflection is a mutual dissolution. The unity of the general is dissolved and goes to the ground in the multiplicity of the particular, inasmuch as this multiplicity is dissolved and goes to the ground in the unity. To this there corresponds, then, the ground's second main proposition: ``All is at bottom different.'' With difference the emphasis falls upon deduction, and by deduction the particular arises; but with unity the emphasis falls again upon reduction, and the particular is drowned in the ground.

If pure reflection can carry it no further than to a development of the proposition that all is at bottom a particular because all is at bottom a general, and conversely, then it is at once evident that the complete determination of the particular must rest upon the syllogism. In the syllogism, i.e. in an infinite system of syllogisms, the intuiting-reflecting subjectivity sees the particular everywhere: the general grounds itself in grounding the particular, and by virtue of the particular the singular as well. The movement that thus goes from the general through the particular to the singular, a syllogism-movement from centre to periphery, is determined by disjunction, thus by a system of systems of disjunctive syllogisms. An example: ``The existent is in its generality both a living and a non-living (\textit{A} is both \textit{B} and \textit{C}), but in its
% --- p. 356 ---
\opage{356}particularity either a living or a non-living (\textit{A} is either \textit{B} or \textit{C}). The non-living is \emph{posited}, thus the living is \emph{excluded}; the living is \emph{posited}, thus the non-living is \emph{excluded}. The living and the non-living are both, by mutual exclusion and corresponding reciprocal determination, posited as particulars. By virtue of the carried-through disjunction the universal is concluded and closed off in itself in such a way that the particular is at the same time singly concluded and closed off.''
% The 1866 text never closes the quotation opened at ``The existent (p. 355); closed here, where the sense requires it.

Since disjunction in each of the particular systems must now be continued right down to the extreme of the singular (the living being is either plant or animal, the animal either mollusc or articulate or vertebrate, the vertebrate either etc. down to the individuals), we must recognize that in propositions such as ``the existent is (in the present case) either a horse or an ox'', the way from the subject (the existent) to the predicate (horse) must, for the judgment's sake, be infinitely short, and for the syllogism's sake infinitely long.

\emph{The Totalization of the Syllogism-Systems: The Individuation-Problem.} Whence, then, has the individual its first and proper origin? Realism (\textit{univers. ante rem}) answers: from the universal, and exerts itself to whip the concepts up to such a concentration of the particular in the singular that it may look to reflection as though the individual in its immediacy really were an intellectual condensation of the universal. In opposition to realism, nominalism (\textit{univers. post rem}) starts from the presupposition that the in itself individual, far from having any proper origin, is much rather itself the original, and then, by the aid of abstraction, gives it the appearance that it is the universal that is a derivative. The realists misuse \emph{disjunction}, the nominalists \emph{subsumption}: the former lose themselves in an individualization that never reaches the individual itself; the latter hold fast to an individual that must exclude every individualization. From both sides, then, the
% --- p. 357 ---
\opage{357}particular too is misapprehended; for the originally particular presupposes that the universal is just as much determined in principle by the individual as the latter is by the former, that the universal and the individual have, therefore, the same originality.

Since all three fundamental forms are equally original, and the content in each of them identical with its form, the universal, the particular and the individual must have the same absolutely ideal reality. Here there can then just as little be talk of universalizing the individual as of individualizing the universal. It follows, on the contrary, from the principle of relation that the three ideal realities must absolutely subsist by virtue of one another: the relatively universal subsists by virtue of the particular in original unity with the singular; the particular subsists by virtue of the singular in original unity with the universal, and the singular by virtue of the universal in original unity with the particular. The fully developed relation is thus syllogism. But a subsisting in syllogism is a subsisting in ideal energy, a subsisting in infinite subsumption and disjunction: the articulating reciprocal determination of subsumption and disjunction is individuation.

The individual does not come to be by finite individualization; it \emph{does not come to be}, for it \emph{is}. And yet the individual cannot subsist without individuation, for individuation is the completed syllogizing, and the individual subsists in eternal syllogism. Here, in the determination of the relatively absolute, there is not an individualizing activity alongside a universalizing one; for individuation, which gives the individual subsistence, at the same time gives the universal a being different from the individual and precisely thereby individual; only a ratiocination-philosophy that splashes about in words can imagine that it is really speaking of something when it speaks quite generally of the ``all-universalizing'', the ``all-individualizing'' etc.

If now each of the fundamental forms is, by virtue of
% --- p. 358 ---
\opage{358}individuation, a totality developed in itself which embraces the whole content of subjectivity, then the extremes and middle of the totalized syllogism must necessarily constitute \emph{three immanent regions}, and each region by itself must in its form embrace the whole infinite system of categorical and disjunctive syllogisms. To the region of the universal there corresponds, then, subjectivity as thinking; to the region of the particular, subjectivity as reflecting-intuiting; to the region of the singular, subjectivity as willing. Just as little as the absolute subjectivity can be said to give up its unity because its opposition of thinking, reflection-intuiting and willing separates itself into three phases, just as little can the three immanence-regions bring about a splitting of the original reason-content, as though the content were stratified and the ideal world parcelled out into heterogeneous provinces. As pervasive as the opposition is between thinking, reflecting intuiting and definitively determining willing, just as pervasive is also the opposition of the total syllogism holding for the content, according as it takes place in the region of the universal, of the particular, or of the singular.

From the standpoint of the general survey, the region of the universal may most nearly be referred to what Neoplatonism has apprehended on the one side as the pure light (τὸ φῶς), on the other side as the divine reason (ὁ νοῦς). But although the most acute among the Neoplatonists, Plotinus e.g., determined the Absolute itself as the infinite energy\footnote{Εἰ οὖν τελειότερον ἡ ἐνέργεια τῆς οὐσίας, τελειότατον δὲ τὸ πρῶτον, πρῶτον ἂν ἐνέργεια εἴη· ἐνεργήσας οὖν ἤδη ἐστὶ τοῦτο, καὶ οὐκ ἔστιν ὡς πρὶν γενέσθαι ἦν· ὅτε γὰρ οὐκ ἦν πρὶν γενέσθαι, ἀλλ᾽ ἤδη πᾶς ἦν. Ἐνέργεια δὴ οὐ δουλεύουσα οὐσίᾳ καθαρῶς ἐστιν ἐλευθέρα· καὶ οὕτως αὐτὸς παρ᾽ αὐτοῦ αὐτός. Πλωτῖνος. Cf.\ Kirchner. Die Philosophie des Plotin. Halle 1854. p.~35--46.}, they were nevertheless not able to see through the relation between ``the
% --- p. 359 ---
\opage{359}One'' and ``the Reason'', since they were not yet able to discern the identity of the activity of reason with the immanent syllogism\footnote{``Aus der reinen Einheit geht zunächst der Geist des Universums hervor, dem alle einzelnen Geister entspringen \dots\ Das erste Sein (οὐσία) und die erste Thätigkeit (ἐνέργεια) lebt er in ruhender That (ἐνέργεια ἑστῶσα), in sich ruhend, aber um jenes sich bewegend (ἐνεργῶν). Man darf die göttliche Vernunft nicht mit der menschlichen verwechseln, die einem äußeren Sein gegenübersteht. Es ist von keiner Vernunft die Rede, die durch Berechnungen und Schlüsse zur Wahrheit gelangt, sondern von der, die Alles hat und Alles ist und ewig bei sich selbst bleibt''. Kirchner. Anf.\ Skr. p.~46.}.

If we refer ``the first being'' (οὐσία), the ``in itself reposing'' ground-energy (ἐνέργεια ἑστῶσα), to the syllogism-region of the universal, then we easily find the expression for the absolute unity, exalted above all particular determinations, in the thinking selfhood, the subjectivity identical with itself in all thinking. The self that thinks its content thinks, by the negative functions of the word, the singular and the individual just as sharply and distinctly as either the particular or the universal. The thinking self thinks, by its negatively determining word, the totality of all syllogism-systems through all individuations in all directions, without the multiplicity in the least degree either infecting or weakening the original unity. It is precisely the transparency of terms and syllogisms corresponding to the syllogism-energy that in this region makes the realm of reason an eternal realm of light.

If we next refer ``the unfolding of the highest unity'', in which the Neoplatonists place the origin of ``the ideas'', to the region of the particular, then we become able to exchange the veiled, figurative and half-mystical intimations for pure logical determinations. To the intuiting-reflecting gaze every determination presents itself with necessity
% --- p. 360 ---
\opage{360}as a particular. To speak of the particular's stepping out of the universal, of the separation of the coloured rays out of the pure light, of the passage of the uniform into the multiform (τὸ ποικίλον), is to speak of thinking's passage into reflecting intuiting, a passage that takes place by the word, inasmuch as the word exchanges the negative functions for positive ones. The immediacy of intuition that thereby arises introduces into reflection precisely that tinge of finitude in which the Neoplatonists saw a characteristic difference between an intellectual matter (ὕλη νοητή) and a multiplicity of changing forms\footnote{Among the several passages occurring in Plotinus which Kirchner adduces, the following (Anf.\ Skr. p.~48) is especially characteristic: Εἰ οὖν πολλὰ τὰ εἴδη, κοινὸν μέν τι ἐν αὐτοῖς ἀνάγκη εἶναι, καὶ δὴ καὶ ἴδιον, ᾧ διαφέρει ἄλλο ἄλλου. Τοῦτο δὲ τὸ ἴδιον καὶ ἡ διαφορὰ ἡ χωρίζουσα ἡ οἰκεία ἐστὶ μορφὴ· εἰ δὲ μορφὴ, ἔστι καὶ τὸ μορφούμενον, περὶ ὅ ἡ διαφορὰ. Ἔστιν ἄρα καὶ ὕλη ἡ τὴν μορφὴν δεχομένη καὶ ἀεὶ τὸ ὑποκείμενον.}.

Under the region of particularity there belong now genera and species (γένη and εἴδη). This is, however, not to be understood as though it were an indeterminate multiplicity of ``visions'' and ``plastic prototypes'' that involuntarily emerged out of the eternally billowing sea of the reflections; no, individuation's determinate difference between genus and species is in all genera and in all species an express syllogism-determination. By the syllogism-determination the individual is everywhere present in the species, since the species is present in the genus; the peculiar individual (the individually particular) syllogizes itself together, by virtue of the species (the particular in the special sense), with the genus (the particularly universal).

That the Neoplatonists with Plato take the ideas (εἴδη) for ``archetypes'', and that Plotinus with Aristotle regards every εἶδος as an individual essentiality (οὐσία), can in a certain sense be approved; but the fundamental error is that the individual essentialities are apprehended in a merely immediate and therefore superficial
% --- p. 361 ---
\opage{361}manner. The conception is equally superficial whether with Plato we
let these essences rest in eternal unconditionedness, or with Aristotle lose
ourselves in a forming of matter. In order to master its matter the forming
essence must at the same time hover above the matter. The singular essence
hovering above its matter cannot be determined individually without being
determined particularly; cannot be determined particularly as species without
being determined particularly as genus; cannot at one and the same time be
determined as individual, species and genus except by an immanent syllogism.

If we finally bring the whole content of subjectivity under the region of
singularity, the region of the individual, then it must become properly evident
to us what individuation originally has to signify. Neoplatonic reflection does
not, so far as the principle is concerned, get beyond εἶδος; but εἶδος belongs
essentially to particularity: Neoplatonism does not reach punctual
determinateness, but remains in speculation. It is a resolute step on Herbart's
part that, in order to get out of speculation without entangling himself in the
metaphysical contradictions of empiricism, he lays ``the reals'' at the
foundation, and by this punctuality means to secure for himself the absolutely
individual. But as intellectual atoms the reals are themselves burdened with
contradictions far too glaring for them to be able to render the
immanence-syllogism superfluous. Leibniz with his monads is, so far as
individuation is concerned, plainly the one who comes nearest to the true. The
monad is really an individual which, instead of excluding the species and the
genus, has on the contrary drawn the genus and the species into itself: the
region of singularity may to that extent well be conceived as a region of
monads. The willing subjectivity is, through the definitive opposition of the
determinate self and the determinate other, \emph{exactly representing}; the
representations, the ideal objects of the absolute positings, undeniably look
like monads.

But as pure products of representation without species and genera the monads can
neither be thought nor intuited. If the
% --- p. 362 ---
\opage{362}monad is to be thought, it must be absolutely determined as the
universal. The monad is then according to its concept one in all; the monadic is
in all monads like itself. The disjunction: A is either B or C or D, must be
transformed into: A is either A\textsubscript{1} or A\textsubscript{2} or
A\textsubscript{3}, for A\textsubscript{1} is as monad essentially identical with
A\textsubscript{2} and with A\textsubscript{3}. The disjunction \emph{is
thought}; but since the difference remains at the negative, particularity cannot
come forward in particular. That the disjunctive judgment becomes premise for a
syllogism, that A by exclusion of A\textsubscript{2} and A\textsubscript{3} is
expressly posited as A\textsubscript{1}; by exclusion of A\textsubscript{1} and
A\textsubscript{3} as A\textsubscript{2}; and by exclusion of
A\textsubscript{1} and A\textsubscript{2} as A\textsubscript{3}; that thus
A\textsubscript{1}, A\textsubscript{2} and A\textsubscript{3} can all be subsumed
under A, shows however that the subsumption is likewise \emph{thought}. The
thought subsumption is just as negative a determining as the thought
disjunction; the individual is a negation of the universal, because the
universal is a negation of the individual. But when all monadic determinations
are negative, since they are thought and as thought essentially universal, then
the thought monad is itself negative in all individuations: \emph{the monad is
the universal; the universal is monad}.

If the monad is made an object for reflecting intuiting, then the particular
comes especially into force. The disjunction: A is either B or C or D,
transforms itself into an expression for an articulation of particularity: S
(the monad in its particularity) is either S\textsubscript{1} or
S\textsubscript{2} or S\textsubscript{3}. The particular within particularity, or
the express difference between S\textsubscript{1}, S\textsubscript{2} and
S\textsubscript{3}, now at the same time gives the subsumption its peculiar
significance: S, the generally particular, under which S\textsubscript{1},
S\textsubscript{2} and S\textsubscript{3}, the specially particular, are
subsumed, is determined as genus: \emph{the monad is genus, the genus is monad};
the specially particular are at the same time determined as species: \emph{the
monad is species, the species is monad}.

If finally the monad through original representing is expressly determined
monadically, then it is a matter of course that the
% --- p. 363 ---
\opage{363}monad as individual has subsistence only through the monad as species
and the monad as genus. For absolute representing presupposes reflecting
intuiting with the same necessity with which it presupposes universal thinking.
If it is necessary that the monad be represented, then it is also necessary that
it be reflected; thus: if it is necessary that it be conceived as individual,
then it is also necessary that it be conceived as species and genus. Individual,
species and genus are monadically different and monadically identical in one and
the same syllogism. But just as monadic difference maintains itself in the
extreme of the individual by means of the acts of representation, the absolute
positings of the willing subjectivity, so monadic identity maintains itself in
the extreme of the universal by means of thinking, the system of subjectivity's
negative functions.

The monad as the individual and the monad as the universal thus constitute two
regions; but the two regions are extremes which have the monad in the form of
particularity for their middle. Truth, individuation carried through, lies just
as little in the extremes by themselves as in the isolated middle, just as
little in the individual by itself as in the particular or in the universal; the
truth in the content of subjectivity is identical with the truth in subjectivity
itself, for truth is a syllogism, a syllogism totalized in systems of
syllogisms.

3) \emph{The immanent Syllogism as Form of free Necessity.} In individuation the
reciprocal determination of selfhood and otherness has found a characteristic
expression. All individual determining is originally self-determining; but all
self-determining is in turn a determining of other by other. If reflection here
falls upon selfhood, then the individual determination is an expression of
\emph{freedom}; if reflection here falls upon otherness, the individual
determination is an expression of \emph{necessity}. Just as freedom corresponds
essentially to subjectivity,
% --- p. 364 ---
\opage{364}because it mirrors selfhood, so necessity corresponds to objectivity,
because it corresponds to otherness.

Insofar as we here look to what is subjective in thought, judging and
syllogizing, weight is here laid upon self-determination, and thereby upon
freedom. Insofar as we here look to the content of the thought, the judgment and
the syllogism, thus to that by which thought, judgment and syllogism are
essentially determined, we here look to the objective, thus to necessity. So
long as the necessity by which the thinking subjectivity is originally
determined, and the freedom with which subjectivity in all thinking determines
itself, are not cognized in their origin in principle, the relation between the
Godhead and the ideas will not be able to be cleared up.

When Plato in aesthetic naïveté now speaks of God as having created the ideas,
now again lets the creating God look to the ideas as to eternal prototypes, then
by the first expression freedom, by the last necessity, is indicated;\footnote{Plato, or rather the Platonic Socrates, sometimes expresses himself in such a way that God must be assumed to have created the ideas. Ζωγράφος δή, κλινοποιός, θεός, τρεῖς οὗτοι ἐπιστάται τρισὶν εἴδεσι κλινῶν. Ὁ μὲν δὴ θεός, εἴτε οὐκ ἐβούλετο, εἴτε τις ἀνάγκη ἐπῆν μὴ πλέον ἤ μίαν ἐν τῇ φύσει ἀπεργάσασθαι αὐτὸν κλίνην, οὕτως ἐποίησε μίαν μόνον αὐτὴν ἐκείνην ὃ ἔστι κλίνη· \textit{Platonis Dialogi Lipsiæ MDCCCLX.} Πολιτείας ί, p.~290. On the other hand in Τίμαιος (cited edition p.~333): εἰ μὲν δὴ καλός ἐστιν ὅδε ὁ κόσμος ὅ τε δημιουργὸς ἀγαθός, δῆλον ὡς πρὸς τὸ ἀϊδιον ἔβλεπεν·} but when the higher unity of necessity and freedom is explained to the effect
that freedom persuades necessity, then the logical fails. By his naïve
utterances Plato has indicated a problem upon whose solution philosophy can only
begin to work when it has grasped it in its whole depth. In the Middle Ages the
problem acquired a special significance in the Schoolmen's dispute over the
relation between the idea of the good and the
% --- p. 365 ---
\opage{365}divine will. According to the Thomistic proposition: \textit{Deus
vult, quia bonum est}, it is the idea that determines the will, the objective
that moves subjectivity, necessity that is the ground of freedom; according to
the Scotistic proposition: \textit{Bonum est, quia Deus vult}, it is conversely
the will that determines the idea, subjectivity that governs the objective,
freedom that dissolves necessity. How the problem even in the most recent times
moves through ethics and the philosophy of religion is seen among other things
from the way in which Feuerbach has set idea and God against one another.
Instead of making, e.g., the idea of justice a predicate of God as subject, one
ought, according to Feuerbach's direction, to make the divine the predicate and
the idea of justice the subject; instead of the religious judgment: ``God is
just,'' one ought to posit the ethical judgment: ``Justice is divine.'' Not only
in ethics and the philosophy of religion, but in aesthetics too the problem must
come forward, when the question of the reality of beauty is to be answered in
the last instance. But that the problem stirs so strongly in the more concrete
spheres of the doctrine of ideas contains precisely a summons to investigate its
origin in logic; for if it is to be solved from the ground up \textit{in
concreto}, it must first have been solved \textit{in abstracto}.

That the opposition in principle between freedom and necessity has its logical
origin in the opposition between the thinker and the thought, between the
subjective and the objective in thinking's own essence, cannot be doubted. All
thinking presupposes with necessity a thinking subjectivity; but the thinking
subjectivity must then also with necessity think something. If now the thinker
becomes in this way just as dependent on the thought as the thought on the
thinker, then the reciprocal dependence and with it necessity does indeed seem
to be the original; but the necessity so formed is the preformation of freedom.
The logical essence of subjectivity is thought; the thinker determines
% --- p. 366 ---
\opage{366}in all thoughts and in all thinking essentially himself: every act of
thinking is originally an act of freedom. The thought has indeed objectivity and
in objectivity necessity; but what is objective in the thought is so far from
weighing upon the thought that it is on the contrary a vehicle for the energy of
thinking and thereby for freedom.

Although freedom is just as essentially different from necessity as the
subjective is different from the objective, yet neither has freedom any finite
limit in necessity, nor necessity in freedom. Everything happens originally with
freedom, because everything happens with necessity, and conversely. Every
thought in which subjectivity posits its freedom has as conclusion its necessary
premise in an already presupposed thought, and thereby comes to lean upon
necessity. The presupposed thought, which exercises the power of necessity over
the one that follows, is itself a product of freedom conditioned by a preceding
necessity. Seen immediately, the premise is free, the conclusion necessary; but
since every premise in the infinite circuit is itself conclusion, and every
conclusion premise, freedom is in every syllogism the subjective moment of
necessity, necessity the objective moment of freedom: necessity is free, freedom
necessary.

That the immanent syllogism is the adequate form of free necessity (\textit{libera
necessitas}) is now to be known by this, that it is able to express the absolute
opposition of freedom and necessity in absolute identity. The reconciliation of
the absolute opposition with its corresponding identity can, according to what
has preceded with us, just as little be sought in an immediate being as in a
finite becoming; everything depends on the energy with which movement and rest
infinitely presuppose and sublate one another.

In this idea-movement the various systems of syllogisms belonging under
``relation,'' the hypothetical, categorical and disjunctive, can enter. From the
standpoint of the absolute,
% --- p. 367 ---
\opage{367}in which subjectivity itself penetrates its content, the hypothetical
syllogism will serve to exhibit the relation between necessity and freedom. ``If
A is B, then C must be D,'' by this judgment the judging subjectivity has
acknowledged necessity and yet reserved freedom to itself. The reservation lies
in the antecedent; how far the antecedent is to be assumed depends exclusively
on how far subjectivity itself will posit it. The syllogism's \textit{p. minor}:
``Now A is B, thus'' \dots\ designates an act of judgment which begins with
freedom and ends with necessity. The determination of the individual, the
non-necessary, is in the realm of ideas an unconditional act of freedom; that
this act of freedom in turn yields an element for necessity is seen from this,
that the premise which expresses freedom acquires significance only through that
in the conclusion which expresses necessity.

It is the relation between subjectivity and its content that is to be more
closely determined by the idea-movement. Insofar as the movement goes from self
to self within subjectivity, from the self, the universal, from the self, the
particular, from the self, the singular, and back through the doubling of the
middle, the movement is central: in the central movement necessity is a moment
for freedom. Insofar, on the other hand, as the movement takes place in the
content, freedom is a moment for necessity. The movement in the content is, like
the central movement, threefold: a movement from the center (the region of the
universal) through middle and periphery back to the center; a movement from the
middle (the region of the particular) through center and periphery back to the
middle; finally a movement from the periphery (the region of the singular)
through center and middle back to the periphery. Under these movements
disjunction and subsumption alternate with one another infinitely. The movement
from the center through the middle to the periphery takes place by disjunction
--- in systems of disjunctive syllogisms; the movement from the periphery
through the middle to the center takes place by subsumption --- in systems of
categorical syllogisms. The movement from middle
% --- p. 368 ---
\opage{368}through center to periphery and from periphery through center to
middle brings with it a crossing of individuations in the most various
directions.

By making the immanence-syllogism the form of free necessity or, what in this
connection is the same, of necessity's unity with freedom, we have, superficially
regarded, done nothing but develop a form for forms, a form of forms; for
freedom is indeed form just as much as necessity. It can then be emphasized once
more that the logical development of form must always yield a logical
development of content; every pervasive determination of form is in the absolute
itself a determination of content.

That necessity and freedom are two mutually opposed forms of the absolute is, in
other words, to say that the absolute as the necessary and the free is
originally opposed to itself; that a unity-form for necessity and freedom is
here sought then acquires the significance that a unity-form is here sought for
one and the same, self-opposed content. What this means becomes properly evident
to us when we now, instead of treating the content from the side of form,
reverse the relation and treat the form from the side of content, instead of
letting the content reflect itself in the form, let the form reflect itself in
the content.

When the form, the absolute form, is reflected in its absolute content, we have
the idea. If the question is raised: is the idea essentially form or content?
still more distinctly: is the idea the content which has the absolute for its
form, or is perhaps the absolute the content which has the idea for its form?
then it is at once noticed that a finite separation of form and content is here
impossible. If the idea as idea is to be absolutely content, then it transforms
itself into form; if it is to be absolutely form, then it transforms itself into
content. If we grasp the idea as content in the form, the three fundamental
movements
% --- p. 369 ---
\opage{369}in the preceding development of form will give us the idea in the
shape of three fundamental ideas: \emph{truth, beauty, goodness}.

It is then in the first place evident that the form-movement in its syllogism
from center to center, from the universal through the particular and singular to
the universal, A--(S, E)--A, must give us the content of the idea in the form of
generality and necessity, thus \emph{the idea of truth}.

It is in the second place evident that the form-movement in its syllogism from
middle to middle, from the particular through the singular and universal to the
particular, S--(E, A)--S, must give us the singular transfigured in the universal
and the universal individualized in the singular; an equilibrium of the free and
the necessary so reflected in particularity that the free is necessary and the
necessary free then corresponds categorically to \emph{the idea of beauty}.

It is in the third place evident that the form-movement in its syllogism from
periphery to periphery, from the singular through the universal and particular
to the singular, E--(A, S)--E, must give us the absolute content in the
overarching form of singularity, of self-determination and freedom, thus the
will-idea, freedom itself, the \emph{idea of the good}.

How adequately the syllogism corresponding to the first fundamental movement
corresponds precisely to the idea of truth must strike the eye. That truth is
one, that it has its essence in the universal, the necessary, the eternal, is
generally acknowledged. But it is then also generally acknowledged that
everything has its truth, that there are many and many kinds of truths, that
eternal truth is not empty, but infinitely rich in content. It is likewise
acknowledged that truth, although it rests in eternal being and governs
everything with the irresistible power of necessity, nevertheless enters into
the temporal, involves itself with what is becoming; for when everything in all
eternity has its truth, then time, becoming and what is becoming must also have
their truth. It is finally acknowledged that truth,
% --- p. 370 ---
\opage{370}though in essence unlimited, yet has an actual limit in untruth,
illusion, the lie, etc.

These heterogeneous determinations all allow themselves to be reconciled and
brought into agreement in the syllogism. In showing how the universal itself, by
developing the particular and singular as middle, divides into extremes, the
syllogism shows at the same time what it means that from the one absolute truth
manifold relative truths can proceed, and that these relative truths, held fast
each by itself immediately and finitely, can become untruths, illusions,
deceptions and the like.

It is in the particular and singular that truth has its infinite wealth of
determinations of content. Thus every special science develops truth from its
own side, in its peculiar system of cognitions. But from the standpoint of
finitude it is impossible to get beyond the separation of extremes and middle.
The movement does indeed take place from the universal through the singular and
particular to the universal; but the middle remains standing in outward
independence; the extremes do not allow themselves to be joined together in a
syllogism. The universal, which constitutes the first extreme and from which one
sets out, is conceived in hovering intuitions, empty abstraction-concepts, with
which one proceeds to the investigation of facts and phenomena, thus of the
singular. The intuition of the singular is then to be resolved into its
particular determinations, so that these particular determinations may in turn
be taken up into the universal; the movement takes place from periphery to
center: the method is \emph{analytic}. But the singular and particular do not
allow themselves to be resolved; the universal, which is the goal of the
endeavor, remains standing as a last extreme; the goal, the concretely
universal, is reached only in an approximation. The universal thus attained can
next be held fast in its generality as a first foundation, the last extreme can
be made the first. The particular and singular are then anticipated in this
first foundation in such a way that the universal by inner self-division can
yield, enclose and hold together its
% --- p. 371 ---
\opage{371}content's heterogeneous parts; the movement goes from center to
periphery: the method is \emph{synthetic}\footnote{In essential agreement with the Hegelian model Heiberg (Speculativ Logik. Prosaiske Skrifter. Første Bind. p.~354--62) has sharply and aptly set forth analytic and synthetic cognition of the idea, namely under the presupposition that the one who cognizes is not a divine but a human subjectivity.\par
``Cognition, as life's renewed immediacy, is itself immediate in its beginning, namely as \emph{intuition}, and thus stands under the schema of the first (immediate) syllogism, E--S--A, whose significance is that the subject, by means of cognition, is joined together with objective truth.~\dots\ But the immediacy of intuition must, like every other, be sublated in its opposite, which is its dialectical development. Hereby arises \emph{analytic cognition}, which stands under the schema of the second syllogism, S--E--A, whose significance is that cognition itself, by means of the subject, is joined together with objective truth or the general; i.e., that the cognizing subject makes its cognition an object, or places itself between its cognition and the latter's object, and in this position sees the object joined together with cognition. This happens by \emph{analyzing} cognition, i.e., by means of separation and abstraction bringing it into the form of the abstract-general. The form thus attained is \emph{the definition}, at which analytic cognition remains standing, as at its goal.''\par
``The immediate certainty, contained in intuition, of the subject's identity with the object in cognition is sublated in the analysis, where the subject, as the connecting middle between its cognition and the latter's object, thereby at the same time separates these, so that their union takes place merely in the subject. The sublation of the immediate certainty has only the determination of being itself sublated, and of returning to its origin. This happens in synthetic cognition, which stands under the schema of the third syllogism, E--A--S, whose significance is that the object itself is the uniting middle which joins the subject together with its cognition. Just as, therefore, analytic cognition is the dissolution of itself into the form of the abstract-general, so synthetic cognition goes the opposite way, and brings the abstract into the form of the concrete; i.e., it constructs the object which the analysis regarded as given. And just as analytic cognition ends with the definition, so synthetic cognition begins with it.''\par
The defect of this presentation, however, is that it is not shown here how it becomes possible for the cognizing subjectivity to complete either the analysis or the synthesis and by this completion to appropriate the content of the idea.}.

% --- p. 372 ---
\opage{372}The opposition between analytic and synthetic method remains in
reflection and cannot be raised to immanent syllogism; it belongs to finitude
and does not reach the absolute. That the analytic method has the idea only in
approximation is recognizable from the disproportion in which the individual,
according to this method, comes to stand to the particular and the universal.
The individual with which one here begins is an immediately concrete existence,
a thing with many properties. What now does the analysis do? It exchanges the
thing for its concept, and determines the concept by a definition. The thing,
that which immediately exists, falls however outside the concept, as the
antilogical falls outside the logical; even if the universal determinations won
by the analysis correspond exactly to the singular determinations in that which
exists, there is still here a yawning chasm between existence and concept:
existence is sensuous and concrete; the concept is intellectual and abstract.

That the thing, the individual, is by abstraction separated from its concept is
expressly seen from the definition. For it is not the individual but the
particular, not the thing but the concept, that is defined. In the definition's
\textit{genus proximum} and \textit{differentia specifica} the relation between
the universal and the particular is indeed reflected, insofar as the concept is
determined as species, and the species as individualized genus; but since the
genus is an abstractum, and the species, instead of appearing particularly as
middle term between the genus and the indi
% --- p. 373 ---
\opage{373}vidual, sinks down to a logical mark appended to the concept of the
genus: the definition cannot appear as syllogism, but takes on the character of
an \emph{identical judgment} in the proper sense\footnote{Heiberg in his speculative logic (Pros. Skr. 1ste Bind p.~356) has indeed endeavored to determine the definition as a syllogism; but the endeavor ends in much ambiguity, since he even three times changes his standpoint. First it is said that the definition, just as much as the judgment, must not be ``regarded merely subjectively, as a human operation of thought, but objectively.'' Nonetheless it is still to be ``the cognizing subject'' which, even when the definition is objective, constitutes the middle term of the syllogism. ``Just as every thing is a judgment, a syllogism and a system of syllogisms, so too every thing is a definition in the cognizing subject, namely a subsumption of the particular under the general, by means of the singular; i.e., cognition's subsumption under the abstract-general, by means of the cognizing subject.'' This way of looking at it is then finally altered to the effect that ``the cognizing subject'' drops away as middle term and the singular is represented by ``the sublation of the subsumption''; the whole explanation ends, despite the careful exemplification, in something altogether inexplicable. ``Furthermore,'' so runs the explanation, ``it is by means of its singularity that the particular is subsumed under the general, so that the definition is in itself a syllogism standing under the schema of the second figure: S--E--A. Every definition therefore has three members: 1) that which is to be defined, S; 2) the subsumption under the general, A; and 3) the sublation of this subsumption, in that it is determined as singularity, E \dots\ E.g. the definition: `Man is a rational being,' contains the three members: 1) man; 2) a being (subsumption under the general); 3) but a being which is rational (sublation of the subsumption, in that the general is determined as singular).'' The obscurity could hardly be greater. For how now does the singular in the sublation of the subsumption stand related to the singular as the cognizing subject? Are there here, in the sublation of the subsumption, two singulars or only one? How can ``the general be determined as the singular'' in the predicate, and the predicate become identical with the subject, when the subject itself is not the singular but the particular?}.

% --- p. 374 ---
\opage{374}That the synthetic method for its part is likewise unable to vanquish
finitude and appropriate the idea is equally evident. If the analytic method,
far from mastering the subsumption and reaching from the individual through the
particular to the concretely universal in the system of syllogisms, entangles
itself in abstraction and ends with the definition: then the synthetic method
begins with the definition, not with the concretely universal in the system of
syllogisms, moves not through disjunctions in principle but through formal
divisions, and thus reaches not the individuals, the concrete existences, but
only theorems and formal proofs. The definition, which was to determine the
\emph{universal}, falls outside the division, which was to give the
\emph{particular}; the division in turn falls outside the theorem, and an
immanent total syllogism shows itself to be impossible.

In that finite knowledge, whichever way it otherwise prefers, the analytic or
the synthetic, misses the necessary, pervasive reciprocal determination of
individual, species and genus, it misses the truth and entangles itself in an
untruth. If truth is determined as the necessary unity of thinking and being,
concept and reality, then this unity must itself be expressly determined as
syllogism-unity. The Aristotelian denial of the ideal independence of genera and
species deviates just as much from the truth as the Platonic hypostatizing.
There are indeed species and genera which do not have eternal validity (the
species: whirlwind, of the genus: wind; the species: sandheap, of the genus:
heap); but what does that prove? That the individuals corresponding to such
species and genera (this whirlwind, this sandheap) likewise do not have eternal
original validity! Seen in the light of the idea, individual, species and genus
have the same necessity, though in different form, the same independence, though
in different regions. Every idea is with necessity individual, species and genus
in absolute opposition, and precisely therefore individual, species and genus in
absolute
% --- p. 375 ---
\opage{375}unity: the self-unity of the idea is an immanent syllogism-unity.

That the individual has its necessity in species and genus, the species its
necessity in individual and genus, the genus its necessity in individual and
species, is a proof of freedom's unity with rational necessity, and of this
unity's original unity with truth. For freedom shows itself in this, that each
of the syllogism's members preserves its independence; but since the
independence of the extremes is one with the independent middle, and conversely:
the reciprocal independence is identical with the reciprocal dependence, and
freedom is posited as a moment in necessity. Upon freedom's complete merging
into necessity truth depends.

But upon what then does untruth depend? Untruth can assume very different forms.
It can show itself nominalistically in this, that an unconditional independence
is ascribed to the individual and to what is individual at the expense of the
species and the genus; it can show itself realistically in this, that genera and
species obtain independence at the expense of the individuals; it can show
itself in an overvaluation of thinking at the expense of intuition, of intuiting
reflecting at the expense of representation, etc. But what especially dazzles
and confuses finite subjectivity is attendant representations of sensuous
immediacy, of finite existential determinateness.

By mixing in attendant representations of the sensuously concrete and the
intellectually abstract, finite knowledge is constantly entangled in the
illusion that the individual is temporal and only the universal eternal, that
individuals must perish and only the species, the genera or still more universal
essence-forms endure. Hereby arises an intolerable conflict between the actual
and the true. The individual is the actual; but the individual is in actuality
perishable, and the actual to that extent something untrue. The universal (the
species, the genus, etc.) is in essence the imperishable and to that extent the
true.
% --- p. 376 ---
\opage{376}But the universal does not exist; it has its being in thought, not in actuality: the true is thus not actual, and the actual not true; so hard is the contradiction.

In the immanent syllogism, by contrast, the individual and the particular are enclosed from both sides by the universal. With the singular and the particular as middle term, the idea of truth forms out of its universal two extremes whose syllogistic union is precisely the true concrete unity; in this way truth gains power over actuality: all the moments of actuality receive their truth.

Just as the universal would be one-sided and untrue if in the immanent syllogism it did not expressly have the particular as its opposite, so the idea of truth, the idea of essential objectivity, would itself become one-sided and untrue if it did not have the idea of beauty, the idea of phenomenal objectivity, as its absolute opposite. That the absolute opposition must be grounded upon absolute unity will in the present connection further become evident from this, that the idea of beauty is true only by being opposed to the idea of truth.

From two sides, namely, the idea of beauty becomes obscured, in that it either straightway melts together with truth and loses itself in necessity, or by immediate freedom is separated from truth and abandoned to contingency. The Neoplatonists, Plotinus in particular, identify the principle of beauty with beauty itself, the beauty exalted above all beautiful forms, a beauty which nevertheless cannot itself be beautiful, since it is different from the beautiful. ``That which gives beauty is itself without shape; for what we call shape is something in itself formless which has received form, and thus has a share in beauty, but is not itself beauty''.\footnote{Ἀρχὴ δὲ τὸ ἀνείδεον, τὸ μὴ μορφῆς δεόμενον, ἀλλ᾽ ἀφ᾽ οὗ πᾶσα μορφὴ νοερά. Καὶ τὸ κάλλος αὐτοῦ ἄλλον τρόπον καὶ κάλλος ὑπὲρ κάλλος· οὐδὲν γὰρ ὂν τί κάλλος; ἐράσμιον δὲ ὂν τὸ γεννῶν ἂν εἴη τὸ κάλλος. Δύναμις οὖν παντὸς καλοῦ ἄνθος ἐστὶ κάλλους καλλοποιόν. Plotin. Cf. Kirchner. Anfr.\ Skr.\ p.~36--39.}
 It is, in
% --- p. 377 ---
\opage{377}other words, the thought beauty, beauty in the form of truth, that is here posited as the highest. But when it is now the idea as truth, not specifically as beauty, that gives beauty, so that ``having a share in beauty'' must become the same as having a share in truth: how then does the principle of beauty stand to the beautiful shape? As the shape is formed there must, besides the form, be something that receives the form, namely the matter. If the matter were not formed, the shape could not come forth; but matter as matter stands in a contingent relation to the \emph{determinate} form, and of course also to the \emph{beautiful} form. There must then, according to the Neoplatonic view, be something contingent in the beautiful shape; but since something contingent steps in between the beautiful shape and the principle, beauty itself has sunk down into the contingent and therewith into the perishable.

``Macht' ich doch, sagte der Gott, nur das Vergängliche schön''. These words, with which the poet lets Zeus console ``perishable beauty'', contain an essential violation of the idea of beauty. Only earthly beauty, the beauty which in the changing phenomena must engage itself with the material and the contingent, can be perishable; the idea of beauty itself cannot possibly be perishable. The element of otherness which beauty originally demands cannot fall outside the idea; it must with necessity be a moment inherent in the idea itself. And yet reflecting intuiting posits an essential difference between self-determination and determination by another. The original unity of idea and shape which expresses beauty is precisely a reflected
% --- p. 378 ---
\opage{378}unity of self-determination and determination by another, a reflected unity of freedom and necessity; a reflected unity so thoroughgoing that reflection itself dissolves into transfigured immediacy. The shape must be determined, for the idea of beauty must have its shape; if the determination is abstract self-determination, then the shape is arbitrary; if it is \textit{in abstracto} a determination by another, then the shape is contingent. Against arbitrariness and contingency stands necessity. In the separation of the abstractly arbitrary, the abstractly contingent and the abstractly necessary, reflection is made finite; in the finitized reflection unbeauty is seen, but in unbeauty the particular is again seen.

The unbeautiful and the beautiful have thus this in common, that they both rest upon the particular. There is a particularity which is unbeautiful, since it falls outside the idea, and a particularity which is beautiful, since it expresses the original unity of idea and shape and is thus identical with the idea. That the particular expresses the idea is conditioned by this, that the actual and the contingent dissolve into the necessary; by this dissolution freedom comes forth: the idea-shape is determined with freedom, since it is determined with necessity.

When there is talk of beautiful objects (a beautiful horse, a beautiful flower, a beautiful woman), it is easily seen that there is here no talk of different phenomena which each for itself have their share in one and the same beauty, but of different species of beauty. And yet beauty in all that is beautiful must surely be like itself: how then does beauty stand to the different species of beauty? Every philosophical system which remains in reflection without reaching the syllogism must certainly declare this problem insoluble. For an object to be able to be called beautiful there is required, in the first place, that it be an individual reflection of its species, for the species is its prototype; in the second place, that the species itself must have
% --- p. 379 ---
\opage{379}unambiguous determinateness in opposition to other species, and thus be a typical individuation of the genus. Unbeauty can in this respect come forth in two ways: partly in that the individual stands in an aesthetic disproportion to the species (a deformity), partly in that the species itself is an aesthetic ambiguity (a toad).

That the many species of beauty are now assumed to be particular revelations of beauty itself, the one and same beauty, reflection in its one-sidedness apprehends in such a way that beauty is absorbed by truth, freedom by necessity. That the different species of beauty step forth independently as particular types, divine prototypes, the reflection that is unsure of the syllogism interprets in such a way that intuition has set itself up against thought, that the beauty which is identical with truth wants to separate itself from truth. The separation can moreover take place in two ways: either in that the divine prototypes are assumed to be eternal beings, independent in themselves, and the liberation of the particular must then rest upon an eternal contingency; or in that the Godhead, absolute subjectivity, arbitrarily determines the particular types: ``God creates the ideas'', so that the liberation is arbitrary.

These abstractions and contradictions disappear when the idea is expressly apprehended as syllogism. In the syllogism, not in the immanent syllogism of thinking but in that of intuition, the gaze is fixed in particular upon the particular, and the particular expressly determined as idea-shape, as type (παράδειγμα).

The idea-shape is with necessity; for the universal and the individual must --- this is the idea's inescapable demand --- reflect themselves in the shape in such a way that the universal, with the individual as its moment, becomes genus (γένος), the individual, with the universal as moment, becomes species (εἶδος), and the species the middle in which the universal is syllogistically united with the individual. In this respect it holds that subjectivity does not arbitrarily produce, but originally intuits, the
% --- p. 380 ---
\opage{380}eternal types determinate in themselves (δῆλον ὡς πρὸς τὸ ἀΐδιον ἔβλεπεν).

But the idea-shape is then also with freedom. The idea can take on the form of intuition only in that subjectivity intuits; but the intuiting subjectivity can be determined to the intuition of another by another only in determining itself. In countless intuitions, thus in countless shapes, the universal can unite syllogistically with the individual. Why now precisely in \emph{this}, and not in another shape? Here there is not first a shape objective in itself, independent of intuiting, by which the intuiting subjectivity ``has to direct itself''; no, the validity of the form is objective only insofar as it is subjective; the necessity with which the idea-shape is intuited is originally determined by intuition itself, thus by subjectivity, thus by freedom (οὕτως ἐποίησε μίαν μόνον αὐτὴν ἐκείνην).

The immanent syllogism is the free and yet necessary form in which the subjectivity that intuits the particular produces and preserves the eternal types of beauty.

As regards, finally, the third fundamental movement, the movement from the singular through the particular and the universal back to the singular, thus from periphery through middle and center to periphery, it shall now be shown that the syllogism corresponding to it is a will-syllogism which originally determines the idea of the good.

The willing self must absolutely will something; since selfhood is determined as will, and the will as the subjectively singular, otherness is with that determined as the will's content, the content again as sheer ideal objects, a multiplicity of objectively singular things. Since the beginning is here made from the periphery, the beginning is made here with the hard opposition of singular against singular, of the willing self against the object of the will. The hardness in the opposition is to be apprehended from two sides. Insofar as the represented object of the will stands against the will, it falls outside the will and remains standing as an
% --- p. 381 ---
\opage{381}insurmountable barrier. Over against the barrier the will is empty, and the empty will impotent. Insofar as the object of the will, by contrast, is posited and determined by the will itself, it falls under the will and cannot possibly hold itself up as a barrier. The barrierless will, unrestricted in itself, now determines its content singly in each single case, as it itself wills; the will's content, the good, is then from this standpoint absolutely determined by an absolute arbitrariness: \textit{bonum est, quia Deus vult}.

But absolute arbitrariness is, when we look more closely at it, an intolerable contradiction. The power of arbitrariness was to show itself precisely in this, that the will which masters all barriers expressly posits and determines the objectively singular, and does with this whatever it wills. It is the will that in its absolute arbitrariness gives its content objective determinate being in the objectively singular; it is the will that in its absolute arbitrariness can recall this singular and make it into nothing. What follows now from this? From this follows with inescapable consequence that the singular has no objective validity; that the will, by consuming its objective opposite in all single objects, consumes its content and feeds upon itself. Absolute arbitrariness, which was to be an expression of the will's omnipotence, becomes, essentially regarded, an expression of its impotence.

In order that the contradiction may be resolved, the willing self must grasp itself as thinking. Since willing goes together with thinking, the thinking will becomes an objective universal will, a will which in each and all lets the content's own essential development hold good. But the content's essential development is a development in the form of thinking; in the completed form of thinking the objective content is by the immanent syllogism absolutely determined as truth. It is thus truth that takes on the form of the will, since the universal will takes on the form of thinking. In the form of thinking the universal will is of course absolutely determined by its content, for the content is
% --- p. 382 ---
\opage{382}expressly posited as the content of truth. From this standpoint it is thus by no means the will that arbitrarily determines the good, but conversely the good that essentially determines the good will: \textit{Deus vult, quia bonum est}.

That this resolution of the contradiction is nevertheless also unsatisfactory is equally striking whether we consider first the content of the good or its form. As regards the content, the good is straightway exchanged for the true, since the will, as regards the form, has straightway become exchanged for thinking: the universal will is a pure thought-will. As thought-will the will is contemplative and theoretical; the contemplative essence of the theoretical will is recognizable in this, that instead of determining its content it lets itself, point for point, be determined by the content. But a will whose universal determining consists in letting itself be determined is a passive, will-less will, thus a will which is no will. By veiling the contradiction one can nevertheless still make an apparent advance, by letting the will, after it has by virtue of its unity with thinking become universal will, at the same time by virtue of its unity with intuition become a particular will. It is then in this way not only the idea of truth but also the idea of beauty that is immediately apprehended as a will-idea and made one with the good.

% Errata: p. 382 — the 1866 errata list keys this correction to "282", a misprint for 382; καλοκαγάθία corrected to καλοκἀγαθία.
The illogical amalgamating of the true, the beautiful and the good is in general characteristically marked by an aesthetically superficial apprehension of divine wisdom. Not enough that the Greeks in all naivety identified the beautiful with the good (καλοκἀγαθία);
``speculative dogmatics'' continues still, in its logical impotence, to do the same. ``What speculation calls the idea, the world-forming thought, is designated in Holy Scripture as the wisdom that played eternally before the countenance of the Most High, and in which he had his delight. And it is described not merely as an inner mirror-image in God, but also as the effecting, the all-forming
% --- p. 383 ---
\opage{383}thought. For wisdom (the idea, the divine Sophia, the heavenly Virgin, as the Theosophists have called it) is the \emph{master-workwoman} of all things'' etc. ``The forming thought'', which points toward the idea of truth, ``the heavenly Virgin'', which points toward the idea of beauty, ``wisdom'', which points toward ``an infinite purpose of will'', thus toward the idea of the good --- these eternal hands on the clock of logic ``Christian speculation'' rakes together, in the most Christian of meanings, into one box.

By emphasizing the will's unity with thinking and intuiting, so much is nevertheless already won, that the objectively singular, the will's ideal object, can now form the middle term in which the willing subjectivity\footnote{The logical-ontological opposition of the three fundamental ideas: the true, the beautiful and the good, is of decisive significance for life as well as for science. If the idea of truth did not have its independence over against the good, science would never be in a position to assert for itself a principle independent of the practical view of life; if the idea of the good did not have unconditional self-validity in opposition to the idea of truth, ethics and religion could not possibly assert for themselves a particular principle independent of science. The Middle Ages held to the principle of faith and cowed science; it made the idea of truth straightway dependent upon the idea of the good: that was barbaric. The present age goes to the opposite extreme; it holds to science and wants by the means of knowledge to dissolve faith; it makes without further ado the idea of the good dependent upon the idea of truth: the latter misunderstanding is just as striking as the former.}
 lets the universal unite syllogistically with the particular. Since the middle term, the objectively singular, is by the will itself determined as its express object, the two extremes, thinking's universal and intuition's particular, are with that determined as the objective content of the thinking and intuiting will. By the extremes necessity is then determined, by the middle freedom is determined; by the syllogism itself the unity is determined, the unity decisive for the idea of the
% --- p. 384 ---
\opage{384}good, the unity of necessity with freedom in the form of freedom.

\medskip
\noindent\textit{β)} \emph{Objective Syllogism: Modal Functions of Power.}
\addcontentsline{toc}{paragraph}{β) Objective Syllogism: Modal Functions of Power}
\medskip

It was shown in ``the Logic of Subjectivity'' from various sides that the Hegelian ``syllogism'' was too objective to be subjective, and again too subjective to be objective. The fundamental error was: 1) that the difference between the singular, the particular and the universal was posited and sublated in an altogether too fleeting and just therefore superficial manner;\footnote{Grundideernes Logik. I. p.~530--32.}
2) that the extremes' unity mediated by the middle term was apprehended as a unity of being, instead of being apprehended as an essence-unity determined by thinking and effecting, i.e. as a movement-unity;\footnote{Anfr.\ Skr.\ p.~216.}
 3) that, since every regard for the energy of movement was lacking, no true difference could be shown between the syllogism in the ideal system, the ontologically subjective syllogism, and the syllogism in the real system, the ontologically objective syllogism,\footnote{Anfr.\ Skr.\ p.~224.}
 so that the completion of the syllogism in teleology received only a schematic, not a significance essential for the content.\footnote{Anfr.\ Skr.\ p.~224--28.}
 The result was that the character of the terms and their mutual relation had to be stated, in each sphere and for each system, in a manner exactly corresponding to the energy of the functions.\footnote{Anfr.\ Skr.\ p.~433--34.}

In the present connection, where the task is to illuminate the objective syllogism from the standpoint of the logical functions, the objection to the objectivity of the Hegelian syllogism will essentially run thus: that the character of the objectification is overlooked and the function missed. The Hegelian proposition: \emph{All is a syllogism}, lets itself in an immediate objective sense
% --- p. 385 ---
\opage{385}be defended just as little as the proposition that \emph{All is a judgment}. To be sure, the existing object is a judgment insofar as it is in essence a word of power; but thus it shows itself only when its immediate objectivity is again objectified: the logical functions of the word of power are all functions of a doubled objectification. For the immediate objectification the objects in their actual relations and connections are facts, not concepts and judgments, still less syllogisms. By the doubled objectification the objects become objective judgments, antilogical propositions in the categorical power-language: to the categorical power-language, the system of objective judgments, correspond the interpreting functions.

That the objective syllogisms can as little as the objective judgments be apprehended and determined without a doubling of objectification and the functions corresponding to such a doubling, is a matter of course, and lets itself be shown without difficulty. When the objective syllogism, as with Hegel, is to be apprehended exclusively from the standpoint of being and determined by indicative functions, it sinks down into an empty schema, a schema whose objective significance comes moreover to rest upon subjective reflection; whereas the system of the three syllogisms, even with the form to which it is restricted in the Hegelian logic, receives essential significance when it is borne, not by indicative, but by explicative and interpreting functions. We shall illustrate this by examples.

Of what he calls the product of the formal mechanical process Hegel expresses himself in the following way: ``\emph{Unmittelbar} ist das Objekt \emph{vorausgesetzt} als Einzelnes, ferner als Besonderes gegen andere, drittens aber als Gleichgültiges gegen seine Besonderheit, als Allgemeines. Das \emph{Produkt} ist jene \emph{vorausgesetzte} Totalität des Begriffes nun als eine \emph{gesetzte}. Er ist der \emph{Schlußsatz}, worin das mitgetheilte Allgemeine durch die Besonderheit des Objekts mit der Ein%
% --- p. 386 ---
\opage{386}zelnheit zusammengeschlossen ist; aber zugleich ist in der Ruhe die \emph{Vermittelung} als eine solche gesetzt, die sich \emph{aufgehoben} hat, oder daß das Produkt gegen dieß sein Bestimmtwerden gleichgültig und die erhaltene Bestimmtheit eine äußerliche an ihm ist''.\footnote{Wissenschaft der Logik. Zweiter Theil. p.~189.}

It is now easy to show how this utterance changes its character according as it is indicative, explicative or interpreting functions upon which it is assumed to rest. If one starts out from Hegel's own presupposition, that ``the product of the mechanical process'' is to be held fast in immediate objectivity, then such an object is in antilogical actuality an external object, not a word of power, not a judgment, still less a syllogism. As existing object the object is wholly indifferent to the process whose product it is posited to be. First the object is presented as though its existence rested upon a mediation expressed in a syllogism; next it is presented as a product indifferent to this mediation, thus as an object whose existence is indifferent to the alleged syllogism. The mechanical object is thus a syllogism in the sense that it is no syllogism; what is crying in this contradiction is only muffled by the ambiguous turn: ``aber zugleich ist in der Ruhe die \emph{Vermittelung} als eine solche gesetzt, die sich \emph{aufgehoben} hat.'' If the ambiguity is to be removed, then the old confusion of objective movement of the matter with subjective movement of thought shows itself again. That a mechanical object comes forth as the product of a mechanical process, an aggregate e.g. of a process of aggregation, is logically in order; just as it is then also in order that such a process can be at a standstill when the product is finished, and ``the mediation'' in this rest be to that extent posited as one that ``has sublated itself''. But this does not concern at all the process of which there is talk here. The mechanical
% --- p. 387 ---
\opage{387}process, which is here to count as an objective syllogism, is so entirely a subjective thought-process that it concerns the intellectual element in the objectification, without in any way concerning the actual object.

If, by contrast, the object's concept is set against the object, so that reflection is expressly directed upon the concept, then the functions become explicative, and then the Hegelian interpretation will also show itself in a better light. What is said straightway of the mechanical object, that it is ``\emph{presupposed} as a singular'', that it is ``as a particular against others'', that it ``as indifferent to this particularity is a general'', are really three chief sides of the mechanical object's concept, and the comprehending subjectivity, which applies the concept's measure to the object itself, is from this standpoint then also quite rightly called upon to let the concept's three chief sides separate with an immediacy corresponding to ``mechanism'', and thereupon to let the chief sides thus separated unite syllogistically as two extremes with a middle, and to that extent to give a syllogism corresponding to the object's concept, thus an objective syllogism. But the concept of the syllogism in this syllogism cannot notwithstanding be completed unless reflection is expressly directed upon the satisfaction which the concept cognized in the syllogistic union of the extremes affords the comprehending subjectivity, i.e. unless subjectivity knows its comprehending of the given object to be identical with the comprehending that produces the object, and thus through the object corresponds with the producing subjectivity; but on these conditions the functions become interpreting.

The original connection between the objective syllogism and the interpreting functions becomes still more evident when we, instead of a mechanical formal process, choose a chemical one for treatment. If an objective syllogism requires only a syllogistic union of two extremes with a middle, every chemical process will be able to serve as an example. If oxygen is the negative and hydrogen the positive extreme,
% --- p. 388 ---
\opage{388}then their affinity is the essential middle, and water the product which marks the completed syllogism. In this way the chemical process is interpreted by Hegel: ``Er (der Proceß) beginnt mit der Voraussetzung, daß die gespannten Objekte, so sehr sie es gegen sich selbst, es zunächst eben damit gegen einander sind; --- ein Verhältniß, welches ihre \emph{Verwandtschaft} heißt \ldots\ Die Mitte, wodurch nun diese Extreme zusammengeschlossen werden, ist \emph{erstlich} die \emph{ansichseyende} Natur beider, der ganze beide in sich haltende Begriff. Aber \emph{zweitens}, da sie in der Existenz gegen einander stehen, so ist ihre absolute Einheit auch ein \emph{unterschieden} von ihnen \emph{existirendes}, noch formales Element; das Element der Mittheilung (Hegel is thinking here in particular of water) \ldots\ Das Verhältniß der Objekte ist als bloße Mittheilung in diesem Elemente einer Seits ein ruhiges Zusammengehen, aber anderer Seits ebenso sehr ein \emph{negatives Verhalten}, indem der konkrete Begriff, welcher ihre Natur ist, in der Mittheilung in Realität gesetzt, hiermit die \emph{realen Unterschiede} der Objekte zu seiner Einheit reducirt werden''.\footnote{Wissenschaft der Logik. Zweiter Theil. p.~202--203.}

But however much value one must set upon the dialectical development of the concept which runs through this whole part of the Hegelian logic, it is in the first place evident that it rests upon explicative functions; its deviation from the development of the matter can be seen by looking to the presentation that is given in chemistry. Next it is obvious that the result of the chemical process cannot possibly be presented as a syllogism (Conclusum) so long as the development of the concept does not go beyond the merely explicative. For the chemical product to be transformed into a syllogism, it must expressly be presupposed that the result is anticipated in its presuppositions, that it comes forth as something originally intended, that the process with its extremes and middle first, upon the attainment of the result,
% --- p. 389 ---
\opage{389}shows itself in its proper significance; but with this apprehension the syllogism becomes teleological, and the functions interpreting.

The movement of the concept in Hegel's logic from ``mechanism'' through ``chemism'' to ``teleology'' exhibits an interesting mixture of the true and the false, insofar as it rests upon a confounding of movement of the matter, objective movement of the concept and movement in subjective comprehending, thus, as regards the activity of thought, upon a confounding of indicative, explicative and interpreting functions. The presentation of ``chemism'' closes with the following words: ``Der Begriff, welcher hiermit alle Momente seines objektiven Daseyns als äußerliche aufgehoben und in seine einfache Einheit gesetzt hat, ist dadurch von der objektiven Aeußerlichkeit vollständig befreit, auf welche er sich nur als eine unwesentliche Realität bezieht; dieser objektive freie Begriff ist der Zweck''.\footnote{Anfr.\ Skr.\ p.~208.}
 Can there be any clearer utterance of the fact that the movement of the matter in chemism has been, not a real movement, but a movement of the concept, and the whole understanding built upon explicative functions? The presentation of teleology begins thus: ``Wo Zweckmäßigkeit wahrgenommen wird, wird ein Verstand als Urheber derselben angenommen, für den Zweck also die eigene, freie Existenz des Begriffes gefordert''.\footnote{Anfr.\ Skr.\ p.~209.}
 Can there be any clearer utterance of the fact that understanding of the purpose in things rests upon an understanding of the understanding that produces things, that thus the functions decisive for teleology are interpreting? Here there are not first a couple of systems of objective, mechanical and chemical syllogisms which in an immanently objective manner form the transition to and complete themselves in teleology. No, the matter stands the other way round! Here the beginning is made at the lower stages with a tacit exchange of the object's objectification for a corresponding objectification of the object's
% --- p. 390 ---
\opage{390}concept, thus with an exchange of indicative functions for explicative. By explicative functions, since attention is constantly fixed upon the concept-object, the thinker then discovers in the mechanical and chemical processes, not indeed proper syllogisms, for syllogisms in the proper sense presuppose a judging and syllogizing subjectivity, but yet something that resembles syllogisms, namely syllogistic unions of extremes with a middle; by pursuing these syllogism-analogies with increasing emphasis upon the concept-objects and the corresponding explications, the thinker at length reaches a turning point, when a light dawns upon him and it becomes clear what the concept of the objective syllogism properly is: that it is the end, the aim, the intention that gives the objective syllogistic union of the extremes the stamp of a syllogism. With his gaze fixed upon the intention and ``the purposive'', the thinker finally makes the discovery that what his understanding has been occupied with in all the lower instances is not the actual object, not even the object's objective concept for itself, but the concept's original meaning, the divine understanding that utters itself in concept and object: every objective syllogism is without difference teleological.

With this presupposition we shall then in what follows develop the objective syllogism's \emph{formal concept}, \emph{its relation to the word of power and to the categorical power-language}.

\noindent\textit{αα)} \emph{The Formal Concept of the Objective Syllogism.}\\
The difference between the subjective and the objective syllogism is originally determined by the difference between functions of knowledge and functions of power. The objective syllogism is determined by objective judgments, the objective judgments by the categorical power-language. The functions of the subjective syllogism are all indicative; the difference between indicative, explicative and interpreting functions, a difference determined by judgment and object, can first receive its full significance in the objective syllogism. Here the indicative functions then correspond
% --- p. 391 ---
\opage{391}to the first immediate objectification, the objectification
whereby the object is apprehended as object; to the doubled
objectification, the objectification whereby the object's concept becomes
object, namely intellectual object, there correspond the explicative
functions: by the explicative functions there arises a system of objective
judgments. Finally, inasmuch as the objectification goes determiningly
back upon the objectifying subjectivity, which seeks not merely the
concept in the fact, but in the fact's concept the original meaning, it
appears that the functions of the objective syllogism are essentially
interpreting.

If one overlooks the peculiar character of the functions as interpreting,
then it shall not be possible to exhibit a single example of a true
objective syllogism. Let one choose the solar system, the relation between
soul and body, the state with its citizens, government and laws, in short
whatever example one will: so long as the chosen example is treated
immediately objectively, i.e. so long as one looks immediately at the
existing terms and the relations obtaining among them, it will be
impossible to get beyond categories such as whole and parts, cause, effect
and reciprocal action. Here, in the reciprocity of the factors, there are
no objective judgments, and where there are no objective judgments, there
are no objective syllogisms either. The highest one brings it to is to
exhibit a totalized unity-in-opposition, a reciprocal action of three
totalities, in which the one totality furnishes the unity when the two
others form the extremes.\footnote{Trendelenburg makes several apt remarks
with respect to Hegel's examples; as regards in particular the state as a
system of three syllogisms, he says: ``Nach dieser Ansicht wächst das
Ganze dadurch kräftig zusammen, daß das Besondere, Einzelne und Allgemeine
wechselsweise und gegenseitig Grund und Folge wird. Daß sich die
Thätigkeiten des Ganzen und der Theile innig durchdringen, das bildet
allerdings die Selbsterhaltung des organischen Ganzen. Will man die
zusammenwirkenden Glieder das Allgemeine, Besondere und Einzelne nennen:
so hat auch das im Sprachgebrauche einen Grund. Aber man verwirrt die
Sache, wenn man das Analogon eines Schlusses bildet'' (Logische
Untersuchungen. Zweiter Band. p.~278--79). But when he then sets about
giving an account of wherein the unsatisfactory element really consists,
he plainly goes astray. He makes the syllogism, namely, into so typical a
form that no union of extremes and middle deserves the name of syllogism
unless it be in every way arranged according to the rules for a syllogism.
That the state now cannot become a syllogism is evident, ``denn die
Bedürfniße sind nicht als Art der Allgemeinheit des Staates, noch die
einzelnen Bürger als Individuen oder Art eines Geschlechts den
Bedürfnissen subsumirt''. For the same reason the end with its means
cannot be conceived as a syllogism either. ``Im Zwecke \dots\ soll das
Mittel den \textit{terminus medius} bilden, durch den sich die subjective
Vorstellung mit der Objectivität zusammenschließt. Die drei Termini des
Schlusses wären in einem einfachen Beispiele: deutlich sehen wollen das
eine Extrem, das optische Glas der Mittelbegriff, das wirkliche deutliche
Bild das andere Extrem. Mag man hier vergleichungsweise sagen, daß sich
das Subject mit der objectiven Welt, der es seinen Zweck abgewinnt oder
einbildet, zusammenschließt: dies Bündniß ist noch kein logischer Schluß.
Wie will man in den genannten drei Terminis das deutlich sehen dem
optischen Glas als den Umfang dem Inhalt unterordnen? oder gar das
wirkliche deutliche Bild den deutlich sehen wollen so subsumiren, wie
sonst der Unterbegriff in den logischen Umfang des Oberbegriffs fällt? Man
kann doch die wirkliche Ausführung nicht als eine Art der vorgestellten
betrachten. Wenn der reale Schluß, wie er von Hegel in die Objectivität
eingeführt ist, wirklich dem logischen entspräche: so müßte er sich in die
vollständige Form eines Syllogismus fassen lassen'' (Anf.\ Skr.\ p.~276).}

% --- p. 392 ---
\opage{392}From the standpoint of the interpreting functions it will on
the contrary appear that any example whatsoever, be it at first glance
ever so unfavorable, provided only it can furnish
% --- p. 393 ---
\opage{393}two extremes and a middle, lets itself be conceived as an
objective syllogism. We choose one of the examples Trendelenburg has
employed against Hegel. Let sulphuric acid and lime be the
\textit{terminus medius}, thus the particular, whereby the universal, the
concept gypsum, closes itself together with the singular, the existence
gypsum. Here, be it noted, no proof is to be given that gypsum is now also
in a finite sense an objective syllogism; here it is only to be shown what
it means to consider the formation of gypsum under the point of view of
syllogism. From the teleological standpoint of the interpreting functions
the existence gypsum is then expressly conceived as the product upon whose
formation the preceding formation of sulphuric acid and lime is laid out;
\emph{the meaning} of the sulphuric acid and the lime is to be sought in
the gypsum. In the concept the gypsum is anticipated; with this
anticipation the beginning has taken place. The beginning is the
syllogism's first premise; for the concept of the gypsum could not
possibly be determined unless the determination of sulphuric acid and lime
were included in the anticipation. Sulphuric acid and lime in particular
existence now furnish a middle, not in an abstract objective sense, but in
virtue of the relation in which the particularly existent are set on the
one side to the determinate beginning, on the other side to the end, the
attainment of the result, the actualization of the anticipated concept.
The concept occurs here three times and in three different ways: 1) as
anticipated concept, the universal, the concept in subjective possibility,
2) as concept in the objective elements, the particular, the concept in
objective possibility, and 3) as realized concept, the singular. It is,
according to what has preceded, incontrovertible that every concept which
is to be realized must pass through these three phases.

No concept can be realized immediately. Were a concept to be realized
immediately, it would have to happen in one of two ways: either its
subjective being in the anticipation would have to be transformed, by an
absolute power-language, thus without conditions and middle terms, into
objectively actual determinate being, which
% --- p. 394 ---
\opage{394}would then presuppose a miracle; or the concept would have to
begin, without any anticipation, by being so sunk into the objective
elements that it was realized in blind objectivity, and came forth solely
as result, as immediate product. But the product is as product only
object, not concept. As result of the union of the sulphuric acid with the
lime the gypsum is a chemical product, a salt; but this salt is in its
immediacy no concept. In order that the concept may come to light in the
gypsum, the gypsum as the product of the union must reflect itself in the
union itself and in the united constituents; in this reflection
subjectivity, and with subjectivity the concept determining the
constituents and their union, shines through: only on the presupposition
that there is here a concept determining in subjective form, and that the
elements subjected to the conceptual determination have objective
existence, can there be any question of a concept's realization.

But if decisive emphasis is now laid on this, that it is the anticipated
concept, determined by the anticipation and determining its objective
elements by a system of executive functions, that is realized, then the
realization is thereby at once determined as syllogism. The objective in
the syllogism is that the process ceases, that it comes to an end, that a
result appears; the subjective, which determines the process of
realization as a concept-movement and gives it the stamp of a
syllogism-act, is that the subjective concept meets itself in the
objective product; through the middle the end goes back to its beginning;
here there is no interruption, but a completed circuit, a realized
purpose, a syllogism. How closely akin the objective syllogism is to the
actual closure, and this again to the attained aim, can even be seen from
turns of speech such as: ``when the end is good, then everything is
good'', \textit{respice finem}, and the like. In these turns of speech
sound sense has with vigorous one-sidedness emphasized now the objective,
now the subjective side, and precisely thereby shown that
% --- p. 395 ---
\opage{395}the matter must be seen from both sides. As regards the first
saying, the objective may indeed be emphasized so immediately that the
good end which makes everything good becomes an adventurous occurrence, a
happy outcome which suddenly brings those sorely tried to forget all
adversities. But such an end, such an outcome is only a play of chance, a
whim of fate, but no syllogism. That there is here no syllogism, but only
an interruption, a stoppage, an arbitrary end, is evident from the
objective as well as from the subjective side. Objectively regarded, the
lack of syllogism is recognizable in this, that the interruption of the
misfortunes can itself be interrupted again, and the happy end show itself
as the beginning of new misfortunes. Subjectively regarded, the lack of
syllogism is recognizable in this, that there is here no concept, no plan,
that is realized. If the end is in truth to be good and to make everything
good, then it must in development and explanation correspond to its middle
and beginning; it must then, in other words, be a syllogism. For what
indeed is the beginning? The anticipated concept! And the middle? The
concept's objective elements! And the end? The concept realized in the
elements! In the actualization of the concept the closure is a
syllogism-act, an ideal power-act, a logical-antilogical act of the
syllogizing subjectivity.

Since the syllogism, the true objective syllogism, is to gather and unite
what is scattered in its middle and laid out in its beginning, it is seen
from this that in poetry as well as in actuality it must belong among the
great tasks to reach a syllogism. In what relation of dependence the true
syllogism must stand to the true beginning is discerned from the
well-known turns of speech: ``Well begun is half done''; ``a good
beginning never gets a bad end'', and the like. It is precisely the
remarkable thing about these turns of speech that when they are taken from
the side of finitude and regarded superficially, they are only half-true,
indeed even untrue. In finite existence beginning, middle and end fall,
precisely in consequence of the parcelling-out
% --- p. 396 ---
\opage{396}characteristic of finitude, so far apart from one another that
the beginning may, as it is said, promise well, and yet in the course of
the matter most sorrowful accidents may occur, just as in the morning it
may promise to become fine weather, and yet, when noon comes, the sky may
be overcast, and the day end with storm and tempest. When the relation
between the end and the beginning is on the contrary seen from the side of
essence, as in an objective syllogism, the turn of speech is perfectly
true; the beginning is in the objective syllogism nothing other than the
anticipated end. In the result determined by its presupposition, in the
end preformed by its beginning, nothing can come forth but what the
presupposition and the beginning themselves contain. The good beginning is
the disposition toward the good end; if the disposition is developed, if
the beginning reaches through the middle to the syllogism, the end must
absolutely become good: it is in this sense literally true that a good
beginning can never get a bad end.

In the turn of speech \textit{quidquid agis, prudenter agas, et respice
finem} the word \textit{finis}, like the Greek \textit{τέλος} as well,
means at one and the same time both end and purpose. But that the end is
purpose, and the realized purpose end, means precisely that the end is
syllogism, and the syllogism teleological. There are thus not, as Hegel
assumes, abstractly objective, mechanical and chemical syllogisms which,
themselves falling outside the teleological, are to introduce teleology;
for every syllogism, however objective it may be in its facts, is yet
syllogism only by virtue of the subjective and teleological in the
syllogism-act itself.

If we now, after these elucidations, set the objective syllogism beside
the subjective, the difference and opposition between them is so prominent
that the way in which Hegel has sought to maintain their identity must be
regarded as untenable. Nonetheless there is here in the opposition a
unity; it is a unity-in-opposition of knowledge and power, to which we
% --- p. 397 ---
\opage{397}are again referred. Just as the functions of the subjective
syllogism have their ground in a knowledge identical with ideal power, so
the functions of the objective syllogism have their ground in a power
identical with knowledge. The subjective syllogism-act is a theoretical
thought-act with premises of logical judgments; the objective
syllogism-act is, from the standpoint of the producing subjectivity, a
practical thought-act with antilogical judgments for premises. From the
standpoint of the contemplating subjectivity the syllogism expressed in
the production is determined by interpreting functions. In the subjective
syllogism what matters is a mediate knowledge of the concept expressed in
the consequent proposition of the premises: the subjective syllogism is
exclusively theoretical. The objective syllogism, on the contrary, is of
both practical and theoretical significance. For the producing
subjectivity the syllogism is practical, for the contemplating, on the
contrary, theoretical: absolutely understood, the producing subjectivity
must be contemplating, the contemplating producing. It is the mediate
actualization of the anticipated concept on which everything practically,
and an insight into this actualization on which everything theoretically,
depends.

In the features adduced the opposition has been emphasized; it will now be
a question how far the identity lets itself be exhibited on the ground of
this opposition. The task is to investigate how far the subjective
syllogism lets itself be brought under determinations that originally hold
for the objective, and conversely.

The point of view of the objective syllogism is, as we have seen,
decisively teleological; this point of view can in a formal sense hold for
the subjective syllogism as well. If the connection of subject and
predicate in the consequent proposition is set as a task, then the
solution of this task is indeed the syllogism's purpose, and the whole way
of viewing it to that extent teleological. If it is asked here, to keep to
the hackneyed example, whether the sub%
% --- p. 398 ---
\opage{398}ject Cajus lets itself be united with the predicate mortal,
then the premises are means, the consequent proposition end. But this is,
as said, only formal.

% Errata: p. 398 — printed "realiset", corrected to "realiseret"; translated as corrected.
The point of view of the subjective syllogism is decisively logical, and
by reduction to its finite schema syllogistic. The logical and syllogistic
is recognized, particularly in the categorical syllogism, by this, that
the subject in each of the premises is subsumed under its predicate, so
that the consequent proposition becomes a necessary expression of the
unity of the double subsumption. That the subject Cajus is subsumed under
the predicate mortal is an expression of this, that Cajus is subsumed
under that (man) which is subsumed under mortal. The question now is
whether such a system of subsumptions also fits the objective syllogism.
By the application of explicative and interpreting functions it will do
well enough. In the formation of the gypsum the determinate salt-species
(the gypsum) closes itself together with the genus (oxygen salt),
\textit{S}---\textit{E}---\textit{A}; but since the premises are
antilogical, the subsumption must take place in reverse order. According
to the first premise of the realization, the genus can be actualized only
in its species; the actualized species is thus to be subsumed under an
actualization of the genus. According to the second premise of the
realization, the species can be actualized only in its individuals; the
actual individual (the specimen of gypsum) is thus to be subsumed under an
actualization of the species. Syllogistically: a realization of the
gypsum-concept is a particular realization of the universal salt-concept;
in the actual gypsum the gypsum-concept is realized; thus: in the actual
gypsum the universal salt-concept is realized in a peculiar manner.
Practically: first the genus-concept, then the species-concept, and then
the actualization of the concept!

Forced and far-fetched as this conception is, it yet shows this much, that
one cannot, as formal logic is so inclined to do, delimit subsumption in a
mechanical manner. If I conceive, e.g., the proposition: gypsum is a
chemical combination of sulphuric acid and lime, as a definition of
gypsum, then the predicate has indeed no greater extension than
% --- p. 399 ---
\opage{399}the subject, and to that extent there can here be no question
of subsumption; but the proposition can also be seen from another point of
view; I can, namely, conceive the gypsum in a singular representation by
itself, and regard the combination of the sulphuric acid with the lime as
one of many possible combinations, thus as the particular. Finally, there
is then nothing either to hinder the combination particularized in the
concept from being set over against its actualization in such a way that
the actualization is seen as the general, under which not merely the
combination of lime and sulphuric acid, but all chemical combinations, if
it is true that they are to count for more than possibilities, must be
subsumed. But when this has been discerned, it shall for the rest be
conceded that, just as the teleological point of view is inadequate for
the subjective, syllogistic syllogism, which lacks the objective
intermediate determinations, so the syllogistic subsumption is forced and
inadequate for the objective syllogism, specially grounded in objective
elements.

If the objective syllogism is conceived not from the contemplating but
from the producing subjectivity's point of view, then its premises become
antilogical judgments. The gypsum, as anticipated concept \emph{the
universal}, is then determined by the peculiar combination of the sulphur,
the oxygen and the calcium metal, \emph{the particular}; this combination
is in turn determined by its actualization into \emph{the singular}, the
existing gypsum. Yet in this respect too there is no lack of analogy with
the subjective syllogism. The subjective syllogism is, namely, as we have
seen, conditioned not merely by the subsumption, but also by the
disjunction. \textit{A}, the universal, is either \textit{B} or
\textit{C}, the particular --- for the particularity lies in the still
undecided reciprocal determination; by negative determination, i.e. by
excluding the one term: \textit{A} is not \textit{D}, the particularity
becomes restricted and thus individualized into the singular: thus
\textit{A}, the universal, is \textit{C}, the singular in its existing.

If the antilogical movement could not be reproduced logically, there could
be no objective syllogism; the objective syllogism-act must necessarily
reflect itself subjectively in
% --- p. 400 ---
% Errata: p. 400 — printed "blot Række", corrected to "Række blot"; translated as corrected.
\opage{400}the producing subjectivity, but the subjective reflection
hovering over the production is logical. If the logical movement could not
be reproduced antilogically, the transition from the theoretical to the
practical functions, from functions of knowledge to functions of power,
would be impossible. In the anticipation everything is logical, and yet
all logical determinations are so antilogically turned that what is really
antilogical in the execution is, as regards movements and transitions,
preformed in the logical. Here, in the master-workman's knowledge, there
is not a series of merely logical judgments beside or running parallel
with a series of antilogical ones; here there is besides a system of
judgments, namely logical-antilogical judgments, which furnishes a
connecting link for the extremes. To the system of the logical-antilogical
judgments there now corresponds a system of \emph{logical-antilogical
syllogisms}, an intermediate system in which the subjective and objective
syllogisms are unceasingly converted into one another.

What this conversion signifies can be elucidated by a closer consideration
of the relation between extremes and middle term. Is it necessary that
extremes and middle term must in every syllogism, objective as well as
subjective, constitute three totalities? It seems at first glance as
though the question must be answered with no. Manifold premises of
subjective syllogisms are indeed judgments with one-sided predicates, pure
property-predicates, as e.g. all men are mortal. Likewise the terms in
many of the premises holding for objective syllogisms are so elementary
and one-sided that they cannot at all present themselves as totalities. If
the concept of the gypsum is set over against the existing gypsum, one may
indeed say that the extremes are totalities; but how can the middle term,
falling into two elements, be said to constitute a totality? The
difficulty becomes no less if we, after Hegel's example, set the sulphuric
acid as the one extreme, the lime as the other, the affinity as middle,
and the gypsum as the actual syllogism.

% --- p. 401 ---
\opage{401}As regards the simultaneous and the successive too, misgivings
arise. Inasmuch as the subjective syllogism is conceived as a syllogism,
the premises are formed successively, the conclusion with its consequent
proposition follows after: the succession is unmistakable. If the
subjective syllogism is on the contrary transformed into an
immanence-syllogism, then the succession disappears; the movement becomes
rest, and the terms of the syllogism simultaneous. Something similar is
seen in the objective syllogism. If we choose the formation of gypsum for
an example, the terms are successive; if we choose equilibrium, they are
simultaneous.

By closer investigation of these doubtful operation-phenomena the
conversion of the logical into the antilogical and of the latter into the
former will gradually be exhibited. As regards now the relation between
the simultaneity and succession of the terms, to begin with this, it will
soon appear that the opposition of subjective and objective syllogism is
grounded in identity. The succession in the subjective, the ideal,
presupposes its simultaneous; for that which gradually becomes evident
becomes evident as that which already was; the becoming of truth, i.e. its
successive coming to consciousness, reflects itself in its eternal being.
What subjectivity gradually comes to know, it comes to know \textit{a
priori} as something that beforehand \emph{has been} and therefore is in
eternal knowledge. The succession in the objective, the real, is indeed so
finitized that it even becomes temporal; but precisely because the
successive is parcelled out in the schema of temporality, simultaneity is
outwardly realized in the schema of space. Whether, then, we consider the
syllogism from the subjective or the objective, from the logical or the
antilogical point of view, we see in both cases a reciprocal determination
of the simultaneous and the successive.

That the reciprocal determination in the logical is the opposite of the
reciprocal determination in the antilogical goes without saying; but it
lies in the essence of the syllogism that the identity is mediated by the
opposition. The total movement in the logical goes from the ideal through
% --- p. 402 ---
\opage{402}the real back to the ideal; here there are, in other words, two
ideal extremes with a real middle term: \emph{the anticipation is in this
light more than a judgment, it is a syllogism}.

The first extreme of the anticipation-syllogism consists in the totality
of logical determinations which originally determine the anticipated
concept. But since the anticipation is itself determined by objective
considerations, the totality of the elements objective for these
considerations will furnish the syllogism's real middle. The real middle
is seen at first glance not as middle, but as extreme. From each of the
objective elements there goes out an antilogical reflex which from the
objective side is just as determining for the anticipating judgment as the
logical moments are from the subjective side. Notwithstanding, it is not,
for the objective view, reality in its immediacy as reality in which one
here acquiesces. The regard for reality is only a vehicle for the
development of the concept: the antilogical determinations are only
reflexes, gleams and color-tones in the logical. The movement thus goes
through the real back to the ideal; the last extreme, the concept
articulated by the antilogical moments, is now, by virtue of the real
middle, different from and yet one with its ideal beginning: the
anticipation-syllogism is completed.

That the fully developed anticipation must absolutely end as syllogism
could already be foreseen in the doctrine of judgment. In the reciprocal
determination of the disjunctive and the hypothetical judgment, decisive
for the anticipation, particular weight was laid not merely on the choice
of objectively determined elements, but, what is here from this point of
view decisive, on the influence which the choice of the objective elements
had to exert, not on the actualization of the concept, for that goes
without saying, but on the anticipation-concept's own anticipated
determinateness. In the real middle of the anticipation-syllogism there
thus goes a stream of movements from
% --- p. 403 ---
\opage{403}the ideal toward the real, from the real toward the ideal with
overreaching ideality, thus a movement from the logical toward the
antilogical, and conversely, with an overreaching of the logical. The
anticipation-syllogism is thus, when it is seen in its middle, a complete
expression of the subjectivity-movement in which the logical unceasingly
turns in order to reproduce the antilogical, and thereby to close itself
together with itself.

If we fix our gaze on the actualization, is it not then striking that
\emph{the realization of the concept is more than a judgment, that from
its objective side it is a complete syllogism}? That the anticipated
concept, fully developed in the anticipation, however forms not, as it
might seem at first glance, the syllogism's first extreme, but its middle,
must in the present connection be expressly emphasized. So long as we
still considered the actualization of the concept, in unity with its
anticipation, only as one single objective syllogism, the
anticipation-concept, i.e. the anticipated concept, naturally had to
become the syllogism's first extreme. Now, on the contrary, after the
anticipation has been determined as a syllogism by itself, the
actualization as a syllogism of its own must necessarily go out from
actuality and return to actuality; the realization-syllogism must, in
other words, have two real extremes with an ideal middle.

The ideal middle is precisely the anticipated concept. For what indeed is,
in an objective sense, the first and most important thing for him who has
a concept, a plan, that he will realize? Surely never to sit and stare at
the plan and rehearse which and how many demands it really was that lay in
the concept; by incessant repetition of such considerations the producing
subjectivity gets nowhere with the objective. If the producer is not to
remain hanging in anticipating reflections like the fly in the web, then
he must, as he proceeds to the execution, instantly and immediately lay
hold of the objective elements, choose his material; here, then, is the
% --- p. 404 ---
\opage{404}first extreme. In its outward objective existence the material
is not merely wholly foreign to, but, as it seems, even wholly
heterogeneous with, the concept that is to be realized in it; now the
movement begins from the material toward the concept, from the first
extreme toward the middle.

Let us consider this movement somewhat more closely. The extreme is real,
the middle is ideal; if it is to be possible to reach from extreme to
middle, it must then be possible to convert the extreme's reality into
ideal reflexes, i.e. to dissolve its antilogical determinateness into
logical determinations. Inasmuch as a beginning is here made with
immediate reality, a beginning is here made precisely with a system of
logical judgments. The material consists, e.g., of iron, stone and wood;
how is it now possible for the producer to make anything out of such a
material if he does not from his point of view know its constitution; for
upon the constitution the usability precisely depends. But an exhaustive
cognition of the material's constitution presupposes a system of logical
judgments. Every constituent, every one of the particular substances
furnishes here the subject for a multitude of predicates to which the
producer must necessarily have regard in the execution of his plan.

The system of predicates in which the material's concept is developed now
meets with the system of predicates by which the anticipated concept is
determined; upon this meeting depends the validity of the ideal middle.
Through the logical essentiality of the predicates the purpose-concept
goes into its objective elements, becomes determining for their treatment;
the movement thus goes from the ideal toward the real, from the logical
toward the antilogical. That reality of the elements which in the form of
material furnished the syllogism's first extreme now corresponds to the
reality, the actualized concept, which forms the last extreme. By virtue
of the ideal middle the last extreme is identical both with the concept
itself and with the first extreme, is, in other words, the first extreme's
identity with
% --- p. 405 ---
\opage{405}the actualized concept. The doubling of movement in the ideal
middle is again an example of this, that the antilogical in every
objective syllogism reproduces the logical, in order that the logical may
reproduce the antilogical.

It began with the realization of the concept showing itself as one
syllogism; the one syllogism has now by immanent doubling become two
syllogisms; what is thereby gained? The objective syllogism has developed
its subjective moment! The contradiction, that the objective syllogism as
function of the syllogism-act must necessarily be subjective, is dissolved
in the doubling: the anticipation-syllogism has shown how the objective as
moment is taken up into the subjective; the actualization-syllogism has
shown how the subjective becomes a moment for the objective. But the
subjective can --- in the syllogism --- become a moment for the objective
only on the express condition that the objective has become a moment for
the subjective. The two syllogisms, the anticipation and the
actualization, cannot then be set over against one another as mutually
independent; the anticipation has significance only as presupposition for
the actualization, and the latter significance only as consequent
proposition of the former: the \emph{two syllogisms thus form extremes for
a third}, which maintains the validity of both by taking both up as
moments.

The third syllogism, enclosing the two particular syllogisms, must again
have its middle and be recognized by the middle. If now the
anticipation-syllogism with its real middle is the subjective extreme, the
actualization-syllogism with its ideal middle the objective, then it lies
in the nature of the matter that the middle corresponding to the syllogism
of these syllogisms must be the concrete unity of the two particular
middles. But each of the two particular middles is as middle itself the
whole particular syllogism, the real middle the whole
anticipation-syllogism, the ideal the whole actualization-syllogism; the
two middles form by their \emph{difference} the extremes of the third
syllogism. But from this
% --- p. 406 ---
\opage{406}follows again that the two middles must by their
\emph{identity} at the same time form the middle of the third syllogism.

And what, then, is all this to say? It is to say that the objective
syllogism is a system of three syllogisms, that the motion in this system
is infinite rest. The motion goes from extreme to extreme; seen finitely,
the rest passes over into motion. But every extreme is itself the whole
syllogism: seen infinitely, the motion passes over into rest. The
subjective goes from extreme to extreme through the real middle, the
objective from extreme to extreme through the ideal middle: the subjective
and the objective go in opposite directions. But the ideal middle is real,
the real is ideal: the subjective and the objective, with opposite
direction, nevertheless go in the same direction. \emph{Immediately} this
is a contradiction; but in the objective syllogism everything is
\emph{infinitely mediate}, and the contradiction removed.

Inasmuch as the logical in an infinite manner reproduces the antilogical,
the thought-act has from all sides gone together with the power-act;
knowledge has according to its reality become identical with power.
Inasmuch as the antilogical in an infinite manner reproduces the logical,
the power-act has at all points gone together with the thought-act; power
has according to its ideality become identical with knowledge. The
teleological identity of thought-act and power-act is clarified through.

If now the form-development thus carried through is applied to the
determination of the finite extremes, however these may be placed in
relation to the respective middle, then the totality in all three terms is
easy to discover. As regards the example of the gypsum, there can here of
course be no doubt that the real concept gypsum is a totality, and the
actual gypsum likewise. The doubt falls upon the middle alone. But now
consider the following: ``The concept gypsum is a --- by combination of
sulphuric acid and lime --- peculiarly determined real concept''; here we
have the anticipation-syllogism with the two ideal extremes and the
% --- p. 407 ---
\opage{407}real middle. ``The elements: sulphur, oxygen and calcium metal
unite --- in accordance with the demand of the gypsum-concept --- into
actual gypsum''; here we have the actualization-syllogism with the two
real extremes and the ideal middle. If we now, in immanent reflection of
the syllogism's essence, let the ideal and the real middle with their
mutual opposition form extremes, and by their identity again the middle
for a new syllogism, then the totalization of the syllogism's terms is
carried through on all sides. The same result comes out when we apprehend
the gypsum-syllogism in such a way that the sulphuric acid becomes the one
extreme, the lime the other, and their affinity the essential middle.

From a chemical standpoint it would, to put it plainly, be galimatias if
anyone asserted that either the sulphuric acid by itself or the lime by
itself was at bottom gypsum; but from a logical point of view there is
here, as regards the gypsum-syllogism, something else to take into
account. Not the sulphuric acid by itself as existing acid, but the
sulphuric acid as an extreme of the gypsum-syllogism is what here in a
one-sided manner represents the concept's totality. Only in relation to
the whole syllogism can the single term become an extreme. Inasmuch as the
sulphuric acid is posited as extreme for the gypsum, its corresponding
extreme, the lime, is expressly presupposed. The system of properties
holding for each extreme acquires significance solely thereby, that both
systems of properties are reflected in one system, thus thereby, that two
one-sided totalities are reflected in one all-sided totality; thus
thereby, that this one all-sided totality is particularly reflected in two
one-sided totalities determined not merely \emph{by} one another, but also
\emph{for} one another. This way of viewing then serves at the same time
as confirmation that the objective syllogism is essentially teleological.
As extreme in the gypsum-syllogism the sulphuric acid has not merely
certain properties, \emph{so that} it in combination with the lime
\emph{can}, but also expressly, \emph{in order that} it under certain
conditions \emph{shall} be able to form gypsum. Where such a teleology
neither holds nor can hold, there it has (as e.g.
% --- p. 408 ---
\opage{408}in chemistry) no significance either to speak of objective
syllogisms; where teleology on the other hand holds, there the objective
syllogism holds too.

\textit{ββ)} \emph{The Word of Power in the objective Syllogism.} What the
word of power signifies, both according to its ideality and according to
its reality; in what relation the anticipating, executive and explicative
functions must be thought to stand to the power-word's utterance in
judgment and object; in what manner the ideal power-functions as
imperative, indicative and interpreting enter into the determination of
the idea of judgment and object; how the anticipating judgment develops
itself into anticipating syllogism: these and similar principal points
have in our foregoing already been treated and illuminated from so many
sides that a further working-out in the same direction would only look
like unfruitful and wearisome repetitions.

The development of the objective syllogism has already been brought to the
turning point where the logical functions now meet with the logic of the
functions. In the logic of the functions the objective syllogism acquired
particular significance thereby, that ``the system of the
totalities''\footnote{Cf. Grundideernes Logik. I. p.~434 ff.} had to be
particularly determined by systematic alterations in the energy of the
corresponding functions. By the essential alteration in the energy of the
functions the objective difference between the mechanical, the
mechanically-dynamical and the organic syllogism then became recognizable
in such a way that the objective in the transition from the one of these
syllogisms to the other thereby at the same time let itself be seen. Here,
under the point of view of logical functions, the objective syllogism in
its relation to the power-word must be apprehended and determined in such
a way that it can be discerned what difference there is between referring
the objectified concepts, realized in the objectification, to the
psychological and to the ontological subjectivity.

% --- p. 409 ---
\opage{409}So long as the objectified concepts, actualized in the objects,
are, as regards the comprehending, referred exclusively to the
psychological subjectivity, the finite separation between concepts and
objects, between a priori and empirical concepts, between thinking and
being will never be able to be removed. At this standpoint the one who
cognizes can indeed, by subjective syllogisms supported upon facts, work
his way forward to objective results; but in the objective results finite
knowledge can never reach the objective syllogisms. In order to clarify
the reality of the objective syllogisms and thereby sublate the finite
parcelling-out of a priori and empirical cognitions, logic must refer the
objectified concepts, realized in the objectification, not to the
psychological, but to the ontological subjectivity, overreaching in
knowledge and power. An illumination of the reality of the objective
syllogisms then serves at the same time to clear up the logical element in
the difference between the real cognition of philosophy and that of the
special sciences. What this signifies we shall in what follows strive to
show; we choose as examples \emph{gravity}, \emph{affinity} and \emph{the
organic process}.

1) \emph{In the Word of Power Gravity is an objective Syllogism.} Both in
logic and in the experience-science an express difference must be made
between the concept of gravity and actual gravity; but in the
experience-science the difference is determined from a quite other point
of view than in logic. What does physics teach about gravity? ``Every body
which is supported by another'', says physics,\footnote{Naturlærens
mechaniske Deel af H.~C. Ørsted. Tredie Udgave ved Holten. Kjøbenhavn
1859. p.~154.} ``exercises a pressure upon the same, but when it is not
supported it falls perpendicularly down toward the earth's surface. This
property of bodies is called gravity''. Gravity is thus according to this
explanation a \emph{property}. Immediately after
% --- p. 410 ---
\opage{410}it reads\footnote{Anfr.\ Skr.\ p.~155.} ``Gravity is not a
force peculiar to the earth, but the consequence of a general attraction
which takes place between all bodies''. From this point of view gravity
thus becomes determined by two new categories, the category \emph{force}
and the category \emph{consequence}, a consequence whose effective ground
is a \emph{general attraction}. When here in conclusion it is added ``that
the force which on the earth we call gravity is the general force-bond by
which the great universe is held together'',\footnote{Anfr.\ Skr.\
p.~174.} then by this ``general force-bond'' the category gravity glides
imperceptibly over into the category power; gravity is thus a
\emph{power}.

That gravity really lets itself be apprehended from all three sides, as
\emph{property}, as \emph{force} and as \emph{power}, without the one of
these apprehensions at the standpoint of the special science coming into
conflict with the other, shall willingly be conceded. But the formation of
the real concept in the special science is one thing, its development in
logic another. When logic sets about investigating how the empirical in
the natural-scientific concept of gravity stands related to the a priori,
difficulties present themselves which at the stage of finite knowledge it
is impossible to remove.

In order to explain the universal, the necessary and to that extent a
priori in gravity, one subsumes the gravity-phenomena occurring on the
earth under ``the general attraction''; in order to explain the general
attraction, one appeals to the force of attraction inherent in all
material parts; in order to explain the general force of attraction, one
refers to the property common to all bodies; but in order to explain this
property, what does one then do? The judgment about the common property of
bodies, the judgment that
% --- p. 411 ---
\opage{411}all bodies are heavy belongs, according to what we have at an
earlier stage\footnote{See above p.~264--75.} set forth, among the
antilogical-apodictic judgments.

How the empirical and to that extent contingent in the
antilogical-apodictic judgments conflicts with the necessary and to that
extent a priori is then recognizable enough in mechanical physics'
explanation of gravity. If a property common to all bodies were really to
let itself be apprehended as the ground-cause of the gravity-phenomena,
then this property would have to let itself be derived \textit{a priori}
from the concept of matter. But to the concept of matter as matter there
belongs, according to the doctrine of experience, only this, to fill
space. ``The space-filling in bodies we call \emph{matter}, the bounded
filled space a \emph{body}''. But from the mere concept of space-filling a
general force of attraction will never let itself be derived with logical
necessity. If it were to be possible to derive the general attraction from
the concept of matter, then it would surely also have to be possible to
derive the law of attraction from the concept of matter. But how
precarious it is at the standpoint of finite knowledge to want to
determine the law of gravity \textit{a priori} is then seen clearly
enough, among several examples, from the well-known dispute between Buffon
and Clairaut. Buffon is right in setting out from the assumption that the
law of gravity must correspond to a concept of reason holding with
objective necessity, thus to a concept whose validity is \textit{a
priori}. But the untenability of his logical and metaphysical as well as
mathematical arguments shows that he is wrong in wanting to construe
nature's a priori concepts of reason \textit{a priori}.\footnote{On the
dispute's logical moments cf. Forelæsn.\ over ``Phil.\ Propæd.'' fra
Universitetsaaret 1861--62. p.~116--118.}

A real concept like gravity cannot be pieced together out of heterogeneous
constituents, so that an objective presupposition and a subjective
thought, an empirical element and an a priori
% --- p. 412 ---
\opage{412}category, a contingent content and a necessary form should by
ingenious combinations be able to satisfy the highest demand of knowledge
and of science. At all points thought demands the a priori, at all points
intuition stops at the empirical. The heterogeneous constituents out of
which a concept of experience at first glance seems to be composed must in
the comprehending itself be transformed into homogeneous moments: in the
concept of experience as concept all determinations are a priori; in the
concept as object, the existence of the concept of experience, all
determinations are empirical\footnote{If one overlooks this, one will
hardly be able to escape the consequence of the following discussion:
``A. You will then surely grant me that the a posteriori on the whole
stands related to the a priori as sense-perception to thinking? B. Of
course. A. The more of the sensible there is in a concept, the more it
contains of an a posteriori. B. That is understood. A. But the sensible
is, as we have seen, an unthinkable, something that is not there for
thought, that does not let itself be thought? B. That must be granted.
A. The more of the a posteriori a concept contains, the more of the
unthinkable must it thus contain; that is to say: a concept is unthinkable
in the same degree as it is a posteriori; that is to say: a posteriori
concepts are unthinkable''. Forelæsn.\ over phil.\ Propæd. fra
Universitetsaaret 1861--62. p.~107--108.}.

The unity of actual gravity and the concept of gravity is a
syllogism-unity, a unity of extremes which are themselves syllogisms. But
it is, be it noted, the ontological, not the psychological, the divine,
not the human subjectivity that masters such a syllogism. The subjective
extreme is the power-word's ideality; in this extreme all
gravity-determinations are anticipated knowledge-determinations: the real
concept with terms and relation is borne \textit{a priori} by imperative
functions. The objective extreme is the power-word's reality; in this
extreme all gravity-determinations are realized
existence-determinations: the concept is empirical, and the empirical an
objectification by indicative functions.

% --- p. 413 ---
\opage{413}From the standpoint of finite knowledge the syllogistic union
of the extremes cannot, as has been said, be cognized; but precisely in
the manner in which they are set over against one another, human reason
shows its natural urge to maintain them both.

If thought directs its gaze toward the subjective extreme, then the
explanation runs thus: ``Whether we consider gravity in its simplest
utterance as attraction between two bodies separated in space, or from a
higher point of view as the force-bond that holds the world-systems
together, it is to us equally unthinkable that gravity should rest upon an
objective property contingently inherent in matter. But if gravity in its
objective determinate being is not to rest upon a merely factual property,
contingent in its facticity, then it is surely incontrovertible that an
almighty reason, a divine wisdom, must have implanted this property in
matter. Since it is now further altogether unthinkable that a fundamental
property should be able to be implanted in a matter already existent
beforehand, it follows from this at the same time with necessity that the
same divine reason which has determined matter's fundamental property must
have brought forth matter itself. Matter is thus created together with
gravity. The general attraction is the phenomenal expression for the
thinking-together of the parts in the divine knowledge. The ``general
force-bond'' is an almighty thought-bond. We cannot comprehend it, but our
reason is compelled to assume it''.

If thought on the contrary fixes its gaze upon the objective extreme, then
the explanation goes in the opposite direction: ``However we may apprehend
the concept of gravity in and for itself, so much is certain, that gravity
is a fact. As fact the actuality of gravity, its objectivity, is
indisputable. But if gravity really is objective, the mutual attraction of
the material parts cannot possibly be explained by a reason hovering above
thinking the parts together, bringing them together by subjective acts of
thinking. It lies in the nature of matter, it lies in the parts
themselves, that they mutually attract one another.
% --- p. 414 ---
\opage{414}As fact gravity has its ground solely and alone in matter's
actual existence. In the actual existence the ground is one with the
objective necessity. Against an objective necessity there is no
ratiocinating; in the objective necessity lies precisely the ground why
gravity as a property common to all bodies does not let itself be
grounded: the concept of gravity is and remains a posteriori''.

This antinomy, insoluble as it might seem to a critical philosophy, finds
in a simple and natural manner its solution in the objective syllogism, a
solution conditioned by the psychological point of view being exchanged
for the ontological. For the ontological subjectivity each of the extremes
has the same original validity. The a priori does not in time precede the
empirical; the imperative functions do not come to operate for themselves
alongside the indicative; the middle of the extremes is not brought about
by finite combinations of subjective and objective elements; no, the
middle is itself the articulated identity of the extremes: all the terms
of the syllogism are simultaneous. The syllogistic union is an infinite
reciprocal determination of the extremes' peculiarity, thus of the a
priori with the empirical, of the subjective with the objective, of
freedom with necessity, of the power-word's ideality with its reality, a
reciprocal determination which is grounded precisely upon the doubling
that both extremes of the syllogism are themselves syllogisms.

2) \emph{In the Word of Power Affinity is an objective Syllogism.} In the
general gravity-syllogism all finite succession is excluded: extremes and
middle come forth as at one stroke. Precisely because gravity as property
is a general property inherent in matter as matter, there can here, as
regards the realization of the concept of gravity, be no thought of finite
succession. But how does it now stand with the concept of affinity? If it
is established that the ontologically powerful subjectivity cannot by any
whatever of its
% --- p. 415 ---
\opage{415}imperative functions force upon the objectively existent matter
a single property, but must indicatively apprehend and determine each
single property as objective expression for matter's own objective
essence: then it follows of itself that the properties corresponding to
affinity can as little have been arbitrarily impressed upon matter as
gravity.

None the less the realization of the concept of affinity presents
peculiar principal aspects characteristic of the objective syllogism.
Gravity is in its essence uniform; affinity, on the contrary, multiform.
So far is it from being the case that the actual difference between the
weight and specific gravity of bodies of different kinds should mark an
essential difference in their gravity, that it, according to physics, on
the contrary serves as proof that all bodies are equally heavy. It is
otherwise with affinity. Though hydrogen and iridium are equally heavy,
notwithstanding that their specific gravity is so different (hydrogen's
0.06, iridium's about 23.0), one nevertheless cannot say that their
affinity is the same, seeing that they enter into very different
combinations. From one and the same general gravity the particular
weight-relations can, under given conditions, be derived; but from the
general affinity the particular affinities can, however much weight one
might otherwise attach to the conditions, in no way be derived.

In the doctrine of ``the logic of the functions'', in particular of
``functions of the mixed ground'',\footnote{Grundideernes Logik. I.
p.~316--329.} it has already been shown how impossible it is, under the
functions' concrete-physical reciprocal relation, to derive the particular
affinities from the fundamental essence of a general affinity. One can
indeed set up a theory which teaches that the affinity in every substance
must always at the same temperature have the same intensity; but what
intensity the affinity in the given substance, at the given temperature,
actually has is an empirical determination which as
% --- p. 416 ---
\opage{416}little lets itself be derived from the general concept of
affinity as from the general concept of substance.

The empirical in the concept of affinity points, however, in a quite
peculiar manner to the a priori. Just as, namely, the property of the one
thing in general first comes into view thereby, that it, under the
reciprocal action of things, reflects itself in the property of another
thing, so it holds in a peculiar sense that a substance's peculiar
affinity can be seen only insofar as it, in the combination, renders
recognizable the corresponding affinity in another substance. The
different relations in which carbon unites with hydrogen show us at the
same time the different relations in which the hydrogen unites with the
carbon: the affinity of the carbon becomes manifest by that of the
hydrogen, inasmuch as that of the hydrogen becomes manifest by that of the
carbon. The facts belonging hereunder exhibit a remarkable reciprocal
determination of contingent and necessary, of possibility and disposition.

Apprehended according to their phenomenon, the chemical substances of
different kinds, metals and metalloids, are surely quite independent of
one another. Or how should it be able to be shown that the existence of
the one substance, if it is true that it really is an elementary
substance, should with necessity presuppose the existence of the other?
Oxygen could after all very well be thought to be there without hydrogen,
hydrogen without nitrogen, etc. Insofar now as the oxygen existing for
itself is independent of the hydrogen existing for itself, the properties
of the oxygen and of the hydrogen fall quite outside one another: the
properties of the oxygen do not concern the hydrogen at all, those of the
hydrogen not the oxygen; the two systems of properties are thus contingent
for one another. If we next consider affinity, in that we seek to
penetrate into the substances' ``nature and essence'', what do we then
discover? Why, then we discover that the chemical properties of the
substances, far from being indifferent and contingent to one another, are
on the contrary determined by one another and for one another, and thus
stand in a mutually necessary reciprocal relation. What the oxygen in its
combination with the hydrogen
% --- p. 417 ---
\opage{417}signifies rests upon the hydrogen; what the hydrogen signifies
is again determined by the oxygen: if oxygen and hydrogen were not
originally determined for one another, water would never have come to be.

What the reciprocal determination of contingent and necessary (actual
necessity is here thought of) is with respect to the phenomenon, the
reciprocal determination of \emph{possibility} and \emph{disposition} is
again with respect to the essence. Referred to its possibility the
peculiar product is a contingent or, more exactly, contingently-necessary
result; referred to its disposition the product is a realized purpose. For
the objectively rational consideration the chemical products will on the
whole appear now as contingently-necessary results, now as realized
purposes. That the observer can give up neither the one nor the other of
these ways of viewing, and yet cannot point out their mutual boundary
either, shows again how the psychological subjectivity oscillates between
the extremes of the objective syllogism, without mastering the syllogism
itself.

If we let the so-called sound human understanding (or ``sound reason'')
choose the starting point of its reflections first in the subjective, next
in the objective extreme, then the conflicting ways of viewing indicated
in the discussion of the a priori and the empirical in the concept of
gravity will develop themselves into further-reaching consequences
answering in particular to the concept of affinity. If the gaze is thus
fixed one-sidedly upon the subjective extreme, the consideration will go
in the following direction: ``No doubt it can for our limited knowledge in
manifold cases appear as if the chemical products were only
contingently-necessary results of affinity-forces taking place and
determined by the constitution of the substances and the conditions of the
surrounding world; but if we look more closely, then we cannot possibly
withdraw ourselves from the impression that what in the fact looks like
the product of a possibility conditioned from without shows itself on
closer reflection to be an actual utterance of divine plan and
disposition. The possibility of the oxygen's combination with the hydrogen
is only
% --- p. 418 ---
\opage{418}thinkable under the presupposition that the same divine wisdom
which has \textit{a priori} determined the chemical constitution of the
oxygen has also determined that of the hydrogen. But if the chemical
properties of the named substances are determined by a reason exalted
above both, if, in other words, the peculiarity of the oxygen is fixed,
not merely so that it really \emph{can}, which only hints at
\emph{possibility}, but also, which plainly testifies to
\emph{disposition}, so that it under given conditions expressly
\emph{shall} unite with the hydrogen: who will then be able to deny that
the same divine understanding which has determined the oxygen for
combination with the hydrogen has by this combination at the same time
predetermined, and thus intended, the formation of water. How
characteristic every disposition to chemical combination must be utters
itself plainly in the manifoldly varying opposition of acid and base:
these combinations are acids, because those are bases. By the antecedently
fixed separation of acids and bases not merely the possibility of a great
number of salt-formations is given; but by the separation of acids and
bases predetermined in the disposition it is expressly laid out for the
formation of determinate salts.'' It is Leibnitz's doctrine of
possibilities, dispositions and predetermined harmonies which by
reflections upon the subjective extreme of the objective syllogism will
again and again find confirmation.

If the gaze is fixed just as one-sidedly upon the objective extreme, the
consideration goes of course in the opposite direction: ``As unthinkable
as it is that gravity by a divine choice, an almighty arbitrariness, can
be forced upon matter, just as unthinkable is it that such a thing should
be able to happen with any other property. Not in accordance with divine
deliberation, calculation and predetermination can now these, now those
more or less purpose-serving properties be grafted into matter; for matter
with its actual and possible properties is at one stroke what it is and as
it is: here neither purpose nor choice nor arbitrariness holds. The
% --- p. 419 ---
\opage{419}difference between possibility and disposition is a difference
of reflection in the observing subjectivity, and has nothing to do with
the objective. No doubt the opposition of acids and bases plays an
important role in nature's ``economy'', and to that extent belongs
essentially in the order of things. But what has the objective order of
things to do with such categories of subjectivity as anticipation,
preformation, disposition and mechanically arbitrary predetermination?
There is in nature an order of things, not because there is a particular
ordering principle, but because a disorder of things is impossible. Nature
is an infinite whole which subsists in its parts. As the parts are, so
must the whole also be, and conversely: as the whole is, so must the parts
be. Had we another whole of nature, then we should also have other parts:
nature's order rests precisely upon this, that the parts cannot do
otherwise than correspond to the whole, and the whole to the
parts.''\footnote{Cf. Phil.\ Propæd. p.~189.} It is empiricism with its
facts and its factual necessity that will again and again seek its point
of support in the objective extreme of the objective syllogism. But
neither abstract idealism with its a priori nor abstract realism with its
empirical will be in a position to master the syllogism itself.

For insight into the relation between the power-word and the objective
syllogism, on the other hand, everything in the foregoing has been
prepared. With the word of power in its ideality all the specific
determinations of affinity are \textit{a priori} posited and presupposed;
with the word of power in its reality the same determinations are
empirically realized. The misunderstanding of the predetermined harmony in
the subjective extreme is removed in the doctrine of the anticipating and
preforming functions. The misunderstanding of the empirical
determinateness in the objective extreme is removed in the doctrine of the
objective possibilities and their dependence upon subjectively determining
functions. The middle of the syllogism is
% --- p. 420 ---
\opage{420}thus just as original as its extremes, and the extremes just as
original as the middle: the a priori and the empirical interpenetrate one
another at all points.

3) \emph{In the Word of Power the organic Process is an objective
Syllogism.} For the ratiocinating observer the two extremes of the
objective syllogism fall apart from one another like two separated
hemispheres; if the one is above, then the other is below the horizon of
cognition.

Absorbed in the subjective extreme, finite knowledge does indeed demand
the objective; but since objectivity is not permitted to strike out into
its extreme, its essential validity, its relative independence, will never
be able to be maintained. Instead of having its thing-causes, its grounds
and possibilities in itself, everything objective on the contrary gets its
last cause and therewith its original ground and possibility outside
itself, namely in subjectivity.

Impressed by the objective extreme, even to such a degree that it cannot
so much as notice it as extreme, finite knowledge is of course so taken up
with facts, sense-perceptions and empirical presuppositions that it has
not a glance to spare for the originally subjective. If the objective
world with its facts is not itself an extreme, but something absolute in
itself: how then should it be able to tolerate, let alone demand,
subjectivity as a corresponding extreme? A subjectivity is indeed required
in order to consider the objective and form for itself a concept of the
actual; but in actuality the objective is itself just as indifferent
toward the observer as toward the consideration and the concept won by the
consideration.

From neither of the extreme points of view is it possible to solve the
problem of objectification. In the subjective extreme the objectifying
subjectivity is intuited; but the objectification lacks its objective. In
the objective extreme objectivity is intuited, that which is objectified
in actuality, but the objectifying subjectivity is only a shadow. Only in
the objective
% --- p. 421 ---
\opage{421}syllogism's all-sided development can the problem of objectification find its solution.

\emph{Gravity} and \emph{affinity} have, as shown above, this characteristic about them, that, insofar as they fall under the horizon of finite knowledge, they each in its own way categorically hold the syllogism's extremes so separated that the empirical investigator must direct his gaze exclusively toward the objective, and let the speculative thinker grope after the subjective. With the organic process the phenomenon has become another. The fact now itself gives an equal mirroring of both extremes: the organic process is in its phenomenon an objective syllogism. As objective syllogism the organic process is nevertheless manifestly misinterpreted both by empirical physiology and by immanent speculation.

As a special science physiology holds to the objective extreme, naturally not as extreme, but as plain and simple objectivity; it compares the organism and its single parts with a machine: ``The heart resembles a pumping-apparatus furnished with suitable valves, the arteries elastic pipelines'' etc. The organic process is explained by mechanical and chemical functions: the consideration is altogether objective. If now the activity of the organs is determined exclusively by their ``anatomical, physical, chemical and thence derived or co-operating properties,'' then it is indeed a matter of course that ``the special anatomy,'' the doctrine of the structure of the tissues, ``physics'' and ``chemistry'' must furnish the foundation for physiology. But now, as has already been shown in ``the Logic of the Functions,'' an organic function can in no way be derived or directly explained from merely mechanical and chemical activities. Physiology's explanations then also run out at last upon this, that science can at bottom explain nothing. ``We do not know,'' says the physiologist, ``what life is.'' But if it really stands thus, that physiology cannot lead to insight into what life is, then neither can it lead to insight into what an orga%
% --- p. 422 ---
\opage{422}nism is, and then neither to insight into what an organic process is.\footnote{``Life,'' says the physiologist, ``must be constituted as the organization is. A single cell's life is \emph{the sum of this cell's activity} and is simple in the same proportion as the organization is simple. The life of an animal with a highly developed organization is the sum of all the forces that stir within it, and is complicated in the same proportion as the organism is complicated.'' The one-sidedness shows itself in this, that the physiologist, who in this manner makes life into ``\emph{a consequence of the organism},'' forgets that the ideality appealed to is grounded in an essential difference, that life is a consequence of the organism only insofar as the organism itself is a \emph{consequence of life}.
Forelæsn.\ ov.\ ``Phil.\ Propæd.'' fra Universitetsaar 1860--61. p.~223.}

The organic process is not clarified in immanent speculation either. Speculation is indeed in the right when, against the doctrine of experience, it maintains that the concept as concept is precisely realized in the organism; but it is in the wrong in behaving as though it could really prove that the concept realizing itself in the organism was the same as the concept which, at the stages of mechanism and chemism, struggles with the conditions and works itself up through lower presuppositions. By this misinterpretation, pointed out in the foregoing from several sides, the speculative philosophy of ``the concept'' comes into just as skewed a position toward the organic process as physiology.

Physiology's misdirection is, as far as the concept is concerned, this: that the organic process should be merely a result of the co-operation of mechanical and chemical functions. Speculation's misdirection is, as far as the fact is concerned, this: that the organic process should with irrefutable consequence let itself be developed and determined in logic as a naturally necessary expression of the concept identical with its content.\footnote{``Daß das lebendige Individuum sich in sich selbst gestaltet, damit spannt es sich gegen sein ursprüngliches Voraussetzen, und stellt sich
als an und für sich seyendes Subjekt der vorausgesetzten objektiven Welt gegenüber. Das Subjekt ist der Selbstzweck, der Begriff, welcher an der ihm unterworfenen Objektivität sein Mittel und subjektive Realität hat; hierdurch ist es als die an und für sich seyende Idee und als das wesentliche Selbstständige konstituirt, gegen welches die vorausgesetzte äußerliche Welt nur den Werth eines Negativen und Unselbständigen hat. Sein Trieb ist das Bedürfniß, dieß Andersseyn aufzuheben, und sich die Wahrheit jener Gewißheit zu geben''.
Hegel. Die subjective Logik. p.~255--56.}

% --- p. 423 ---
\opage{423}The organic process is indeed in Hegel's logic apprehended and determined as objective syllogism. But the apprehension is so one-sided that what is here posited as a system of objective syllogisms, namely the system of the organic processes in the organism, on closer inspection lacks completeness and falls one-sidedly under the objective syllogism's objective extreme. How matters stand with this we can easily elucidate somewhat more closely.

Already in the Logic of the Functions\footnote{Grundideernes Logik I. p.~434--446.}
{} the organism is, as ``the third fundamental totality of objectification,'' determined by an objective syllogism: S\,--\,A\textsubscript{o}\,--\,O\textsubscript{o} is the syllogism's schema. If we consider this schema somewhat more closely, then it is clear that the organic process takes place in O\textsubscript{o} and is referred to A\textsubscript{o} only insofar as O\textsubscript{o} itself proceeds from A\textsubscript{o}. Since now O\textsubscript{o} in the system of all its syllogisms expressly designates the objective syllogism's objective extreme, it follows from this that the organic process is indeed an objective syllogism, but a partial syllogism which is taken up as a moment in the objective extreme's total syllogism.

This illumination now has a force that reacts back upon the gravity- and affinity-syllogism. For just as the gravity-syllogism in its separate phenomena goes in under the system of the abstract-mechanical functions, and thus in under ``the first fundamental totality of objectification,'' so does the affinity-%
% --- p. 424 ---
\opage{424}syllogism, specified in its chemical appearances, go in under the system of the mechanical-dynamical functions, and thus in under ``the second fundamental totality of objectification.''

And what then do we learn from this? We learn from this that the objective syllogism is just as much to be developed in systems of total and partial syllogisms as the subjective-immanent one, that the functions of power are to that extent just as dialectical as the functions of knowledge. The dialectic of the power-functions can then be developed by a closer determination of

\textit{γγ)} \emph{The objective Syllogism and the categorical Power-Language.} --- The objective syllogism is a teleological act, an actual utterance of the comprehending subjectivity that presents its concept in the object, an utterance which, when it has been determined by interpreting functions, can rightly be said to give the object its proper significance, and to that extent to establish that there really is meaning in the thing, since originally something was meant by the thing.

The error, however, lies near at hand, that the subjective regard should now all at once acquire an illogical, arbitrary preponderance over the objective. One could indeed reason thus: ``If every thing is a realized concept, every realized concept determined by an objective syllogism, and every objective syllogism an utterance of the divine reason's aim, its meaning about the thing and with the thing: what a swarm of aims, subjective intimations and divine meanings must there not then be found in the world of things! If it is now to be the highest goal of human knowledge to trace out, discover and determine these meanings, then there is surely no doubt that we must soon come back to a teleology, a teleology in truth theological, which is just as profound and in its kind just as meaningless as that which put an Athanasius Kircher in a position --- in the age of ``lithotheology,'' ``petinotheology,'' ``grasshopper-theology'' --- to discover 6561 proofs of God's existence.''

% --- p. 425 ---
\opage{425}That this reasoning rests upon a misunderstanding is easy to demonstrate. In the first place, a separation has indeed been made at an earlier stage, an essential, categorically determined separation between the divine meaning with things and men's opinions about things. In the second place, what holds at a lower stage of the essential and the necessary holds at a higher stage of the teleological. That the essence of things pervades things, that everything which happens in the world of things happens with a certain necessity, can neither exclude the determination of anything unessential nor sublate the significance of anything contingent. In a corresponding manner it stands with the teleological. Because everything in the world of things is subordinated to teleology, since everything in the world of things is subordinated to the objective syllogism, it does not follow from this that every single thing in all immediacy should, as a specially attained aim, have a special teleological significance.

What therefore the difference between necessity and contingency is for reflection, that the opposition between end and means is for the objective syllogism. In systems of objective syllogisms the all-capable subjectivity has uttered its knowledge's content in the word of power: a rich content, a manifold word! But for the understanding of this word of power it is above all necessary to make a difference between \emph{main-thoughts}, \emph{middle-thoughts} and \emph{side-thoughts}.

In each of ``the fundamental totalities of objectification'' a divine \emph{main-thought} utters itself; such a main-thought is, with respect to form, objectively expressed in the total syllogism's objective extreme. Each main-thought is, as far as the singular and the particular are concerned, developed in a sum of characteristic \emph{middle-thoughts}: the determinate total syllogism consists in a system of peculiar particular syllogisms necessary for the totality. Finally, in that the particular syllogisms enter into finitude, take up the contingent and are altered in the external, these
% --- p. 426 ---
\opage{426}alterations, at once essential and yet unessential, may well be said to express the system's \emph{side-thoughts}.

``But,'' one will perhaps ask, ``is this explanation really more than an ingenious application of interpreting functions; will the doctrine of the categorical power-language really be able to exercise any fruitful influence upon science?'' We answer: Yes! The doctrine of the categorical power-language is of decisive significance in the determination both of philosophy's goal and of its relation to the special sciences. We shall at once test its significance upon the three fundamental totalities of objectification, the mechanical, the mechanical-dynamical and the organic.

1) \emph{The mechanical Totality in categorical Meaning}. In a totality such as our planetary system a concept complete in itself has been realized. Herein lies this, that this concept must present itself as a system of facts, and that the system of facts must conversely let itself be explained in a concept. From the standpoint of physics the facts must be seen as the original, the concept as the derived. The functions, which here are referred to the psychological subjectivity, are indicative and explicative, but not interpreting: physics' apprehension of the power-concept is and must be unteleological. From the standpoint of the philosophy of nature, on the contrary, the concept as concept must be the original in the divine sense, the system of facts the derived; the functions, which here are referred to the ontologically powerful subjectivity, are not merely indicative, but also imperative, not merely explicative, but also interpreting: the apprehension of the power-concept in the philosophy of nature is and must be teleological.

The strife between mathematical physics' mechanical and immanent speculation's teleological apprehension of the planetary system has come to utterance in a quite unreserved and striking manner in Hegel's philosophy of nature.

Hegel's speculative objections to Newton's \textit{principia
% --- p. 427 ---
\opage{427}philosophiæ naturalis mathematica}\footnote{Vorlesungen über die Naturphilosophie. Werke 7ter Band. 1. Abtheil. Berlin 1842. p.~98. Cf.\ Forelæsn.\ ov.\ ``Phil.\ Propæd. Universitetsaar 1861--62.'' p.~57--66.}
{} run thus: \textit{a}) that Newton's mathematical presentation of the Keplerian laws is a distortion of Kepler's theory built upon experience, \textit{b}) that Newton's mathematical proofs at bottom prove nothing, and finally \textit{c}) that the mathematicians' doctrine of the parallelogram of forces, introduced by Newton, and what hangs together with it, is without application to the motion of the heavenly bodies and their mutual relation. And what then will Hegel put in place of the mathematical analysis of the mechanical forces? Syllogisms, a system of objective, teleological syllogisms!

The colossal blunders of which Hegel on this occasion has made himself guilty allow themselves to be explained solely by this, that the relation between the objective syllogism and the categorical power-language has been altogether unclear to him.

Hegel has a sharp eye for the three-sided form of the objective total syllogism, but when he has got the total syllogism's threefold schema set up and, on the occasion of each figure, has brought forward some general remarks, then he is finished. Upon physics' observations and calculations he can in no way enter; the mathematical determinations of concepts are to such a degree incompatible with his absolute concept that he plainly combats them, and that even, which has in no way been to philosophy's advantage, without having understood them.

Mathematical physics is for its part consistent; it considers the whole planetary system, not as a system of objective syllogisms, but as a result proceeding from given presuppositions and conditions; it is then in good agreement with itself when it apprehends the system's single
% --- p. 428 ---
\opage{428}parts from the same standpoint. Hegel on the contrary is in the highest degree inconsistent. He sees the system in its totality under the standpoint of objective syllogisms; but, remarkably enough, he does not see that what holds of the total must also hold of the particular; he does not see that the total syllogism can consist only in a system of particular syllogisms.

If the motion of the planets, of the secondary planets and of the comets about the sun as central body utters a main-thought realized in the objective syllogism's objective extreme, then there is nothing either in the way of the earth's motion about the sun being interpreted as a middle-thought corresponding to the main-thought, and the moon's motion about the earth as a side-thought corresponding to the middle-thought. That the categories main-thought, middle-thought and side-thought must however be infinitely relative is unmistakably seen from this, that each of the middle-thoughts can in turn be posited as a main-thought with corresponding middle-thoughts and side-thoughts.

This way of seeing yields a consequence very important for the philosophical development of the real concept. The consequence is that the analyses of the particular syllogisms in the philosophy of nature can go into detail just as well as mathematical physics' analyses of the hypothetically necessary results. Indeed, what is more, by this and only by this way of seeing does the philosophy of nature become able to appropriate the yield of the mathematical-physical analyses, a yield essential and indispensable for idea-cognition.

The dialectical reciprocal determination of main-thoughts, middle-thoughts and side-thoughts is namely not to be settled according to a speculative estimate; such an estimate is too indeterminate and subjective to make decisions in science. The dialectical must here necessarily have the exact as its foundation. The task of the philosophy of nature will thus consist in showing how mathematical physics' hypothetically necessary results let themselves in detail be transformed into significant particular syllo%
% --- p. 429 ---
\opage{429}gisms, in that these particular syllogisms in turn show themselves to be moments in a corresponding total syllogism.\footnote{What is meant by this can provisionally be elucidated by a couple of examples. ``If a material point is acted upon by an accelerating force directed toward one and the same fixed centre, and is imparted an initial velocity, then the areas described by the radius vector in equal times will be equal.'' If this judgment is made the major premise in a formal subjective syllogism with a categorical minor premise, we obtain a \emph{hypothetically necessary result}. When on the other hand this result, in its connection with the main-thought uttered in the totality, shows itself as a realized middle-thought, it is transformed into an objective particular syllogism, takes on a teleological stamp, and thus yields a proposition belonging to the categorical power-language. The same transformation can naturally be undertaken with all similar results. ``If the attraction takes place according to Newton's law, the orbit must become a conic section;'' ``if the initial velocity is small in proportion to the force of attraction, the conic section must become an ellipse''. ``If different material points move in ellipses about the same centre, then, since $\frac{\mathrm{T}^2}{a^3}$ is constant, the squares of the periods of revolution will be related as the cubes of the semi-major axes.'' In that these hypothetically necessary results are transformed into objective particular syllogisms, each of them is seen to be a special disposition belonging under the total disposition A\textsubscript{z},
an intimated aim, a proposition of its own with its meaning determined by the power-language.}

Physics' logical functions are essentially explicative, those of the philosophy of nature interpreting: the categorical power-language, which has transformed the concepts into facts and precisely thereby assigned physics its task, has at the same time laid a divine meaning in the facts, and thereby set a task for philosophy.

2) \emph{The second Fundamental Totality of Objectification is determined by the categorical Power-Language}. ``In that geology presupposes geognosy, it must necessarily take its beginning in the finite, in actuality's piecework. To
% --- p. 430 ---
\opage{430}collect and examine minerals, to characterize rock-types, to point out the relation between loose blocks, intrusive masses, Neptunian, Plutonic and volcanic formations etc. belongs to such a degree under the category of piecework that on a preliminary consideration one would hardly suppose that such heterogeneous constituents could be worked together into a scientific whole, still less that their constituents should lead to a cognition of the idea. Nevertheless geognosy has become a science; the more this science masters the facts, the more thoroughly it works up the elements of finitude, the more it points beyond itself, for its content is of the idea \dots\ Seen from geognosy's standpoint the earth is only a pieced-together whole; from geology's standpoint the earth-body is on the contrary to be apprehended as an elementary organism, an essential whole whose parts have gradually developed out of an inner disposition. In the originality of the disposition we then have again the principial, in the parts' separation and the whole's planned articulation the typically teleological, i.e. we have, if only the cognition of the infinite mediation were complete, the idea realizing itself in the formation of the earth''.\footnote{Phil.\ Propæd. p.~186.}

The reciprocal determination of geognosy and geology shows, when the categories are taken in the sense given above, how natural science's all-sided working-up of an elementary organism must with inner necessity lead from the abstractly objective, non-teleological way of seeing over into teleology. For it is not possible for the geologist to judge the relation between the conditions in the different periods of the earth and to characterize these as a series of lower and higher stages of development, unless he presupposes a kind of final end, a state of maturity toward which the whole development strives as toward a conclusion. But the fundamental relation between the physically objective and the teleological the geologist will never be able to
% --- p. 431 ---
\opage{431}see through: that which dazzles him is again the categorical power-language.

In virtue of the categorical power-language geology's concept is realized in actual facts; in the system of facts thing-causes are seen everywhere; but the ground-cause identical with the purpose-cause is lost. How hollow and fantastic the philosophy of nature must become when, in order to remedy empiricism's one-sidedness, it springs off from the thing-causes (\textit{causæ efficientes}) and stares exclusively at the ground-cause (\textit{causa finalis}), can be seen from Hegel's explanation. ``Wie die Luft sich als allgemeines Element ihre Spannung aus der Erde nimmt, so das Meer seine Neutralität. Die Erde dünstet gegen die Luft aus, als Meer; gegen das Meer aber ist die Erde der Krystall, der das überflüßige Wasser aus sich abscheidet, in Quellen, die sich zu Flüssen sammeln. Aber dieß ist, als süßes Wasser, nur die abstracte Neutralität, das Meer dagegen die physische, in die der Krystall der Erde übergeht. Der Ursprung der unversiegbaren
Quellen darf also nicht auf mechanische und ganz oberflächliche Weise als ein Durchsickern
dargestellt werden, so wenig als nach der andern Seite das Entstehen der Vulkane und heißen Quellen; sondern wie die Quellen die Lungen und Absonderungsgefäße für die Ausdünstung der Erde sind, so sind die Vulkane ihre \emph{Leber}, indem sie dieß Sich an ihnen selbst Erhitzen darstellen.'' That which dazzles the speculative philosopher of nature is again the categorical power-language. He sees in the totality of facts not merely a concept, but in the concept a meaning. The meaning he quite rightly seeks in the aim; thus in his eyes the earth's springs and volcanoes get their existence from their significance and their significance from the service they render the earth-body: they stand to the earth-body as organs to organism. But in the determination of the single thing the speculative philosopher forgets that the word in which the concept's meaning is objectively uttered
% --- p. 432 ---
\opage{432}is a word of power, and that the word of power is proclaimed in actual facts, reciprocal actions of efficient causes.

It is ``the household of nature'' upon whose interpretation both the natural scientist and the speculative philosopher of nature, each from his side, are bent. But from both sides the interpretation will miscarry so long as the relation between hypothetically necessary results and objective particular syllogisms is not clarified.\footnote{How geognosy's hypothetically necessary results let themselves in cognition be transformed into particular syllogisms has been shown above by examples from chemistry.}
{} Only in the reciprocal action of particular and total syllogisms does the power-language become intelligible, and the expression ``the household of nature'' come to be apprehended, not as a figurative turn of phrase, but as a natural-scientific concept.

3) \emph{The third Fundamental Totality of Objectification is determined by the categorical Power-Language}. That the living organism must be apprehended as much under the standpoint of purpose-causes as under that of thing-causes is generally acknowledged. Leibnitz's words: \textit{finalem causam non tantum prodesse ad virtutem et pietatem, in Ethica et Theologia naturali, sed etiam in ipsa Physica ad inveniendum et detegendum abditas veritates}, will surely, if not the physicists proper, then at least botanists and zoologists willingly endorse. In the interpretation of a vertebrate's bone-structure, of its teeth, claws etc. the zoologist must necessarily support himself upon teleological presuppositions. By this interpretation the special science and philosophy thus seem to approach each other, each from its own point of departure. That the approach cannot lead to a fusion is however settled by this, that the special science, which apprehends the a priori through the empirical, can never reach the ground-cause through the thing-causes, while philosophy, which according to its method must begin with the ground-cause determined by facts, never
% --- p. 433 ---
\opage{433}by any conceptual development whatever will be able to make the relation between the ground-cause and the actual thing-causes \textit{in concreto} transparent.

What then in this connection is philosophy's significance? That with a clear cognition of the limits of human knowledge it at the same time knows and cognizes a divine knowledge, a knowledge which, in its unity with power, utters its content in the word of power, in a categorical power-language. That philosophy really does know and cognize a divine knowledge can naturally not be seen from its immediately presupposing, and by supposing functions appealing to, the unity of God's omniscience and omnipotence; but it is seen from this, that philosophy, after many groping attempts, has at last become conscious of its method, in that it has become conscious of how it, without mistaking the limit of human knowledge, will be able to develop and make clear the concept of a divine knowledge\footnote{Cf.\ Grundideernes Logik I. p.~230 ff.}
.

The organism, the living organism, is of all forms of existence that in which the reality of the categorical power-language is most clearly uttered. Human knowledge is able to understand this language only insofar as philosophy, under advancing conceptual development, is able ever more deeply, more thoroughly, more all-sidedly to give an account of knowledge and power in their pervading opposition and yet original unity. With the general unity as foundation the opposition of knowledge and power is developed in the doctrine of \emph{judgment and object}; with the developed opposition as foundation the unity-in-opposition in the form of power is developed and, by the categorical power-language, articulated in \emph{the objective syllogism}. But with this the objective syllogism is at the same time developed into a syllogistic union with the subjective.

% --- p. 434 ---
\medskip
\noindent\opage{434}\hspace*{1em}\textit{γ)} \emph{Idea-Syllogism: Modal Functions of Truth.}
\addcontentsline{toc}{paragraph}{γ) Idea-Syllogism: Modal Functions of Truth}
\medskip

That the expression: functions of truth, cannot in the present connection be specially determined by the contradictory opposition to functions of untruth, but by the ontological opposition to functions of knowledge and functions of power, is surely understood of itself. As form for the absolute subjectivity with its content and free necessity the subjective immanence-syllogism is an absolute form for the absolute itself, an adequate idea-form, and to that extent an idea-syllogism. Knowledge's content is from this standpoint identical with the power-word's ideality, and the idea-syllogism's subjective functions are modal functions of knowledge. As form for the power-word's reality in actual facts the objective syllogism is an adequate form for the development of existence's absolute, and for its part likewise an idea-syllogism. The syllogism's content is knowledge's object; knowledge's object is objectified in a categorical power-language; the syllogism's functions are with respect to this determined as modal functions of power.

During the development both of the subjective and of the objective syllogism there showed itself, however, the remarkable doubling, that each of the extremes itself became a syllogism. This doubling repeats itself, and the integrating of one-sidednesses in every direction, conditioned by the doubling, is now a distinguishing mark for modal functions of truth. Thus in this connection subjectivity's all-sidedly developed immanence-syllogism is expressly posited as the idea-syllogism's subjective extreme, while the objectivity-syllogism with real concepts, imperative, executive and interpreting functions etc. for its part yields the idea-syllogism's objective extreme. Since both syllogism-extremes are idea-syllogisms, the concrete idea must necessarily itself yield the middle in which they are syllogistically united. Only by this syllogistic union is the concept \emph{idea-syllogism} all-sidedly apprehended and determined.

% --- p. 435 ---
\opage{435}But the concept is as yet determined only in general; the concept fixed in its generality demands a determinate development going through the particular and the singular. Under this development the doubling will then consist in this: that the whole subjective-objective idea-syllogism is posited in its subjective extreme and characterized as \emph{ideality-syllogism}; that next the whole, objective-subjective idea-syllogism is posited in its objective extreme and characterized as \emph{reality-syllogism}; that finally the ideality-syllogism's subjective and the reality-syllogism's objective extreme are syllogistically united in an all-explaining middle, which thus forms a \emph{conclusion}.

\textit{αα)} \emph{Ideality-Syllogism: the Pre-existence of the Ideas.} If in logic the question is raised, which of the two categories: \emph{being} or \emph{existing}, is in the stricter sense to be posited as the original one for the idea, the answer may well demand a closer consideration. Insofar as the idea as idea is the universal, being seems undeniably to have to be preferred to existing; for being expresses the general. One would then have to say: the idea is, but it does not \emph{exist}; for the general is, but it does not exist. But now the idea as idea is not merely the absolutely universal; it is, according to what we have in the foregoing elucidated from so many sides, just as much the absolutely particular and the absolutely singular. Insofar then as the idea is determined as the absolutely singular, there is every reason to lay emphasis upon the idea's existence: the idea is, because it \emph{exists}.

And yet, immediately regarded, the category existence is just as inadequate for the absolute idea as the category being. Being is ambiguous, because it runs just as much upon the abstract as upon the concrete, just as much upon the non-self-subsistent as upon the self-subsistent, just as much upon the unessential as upon the essential: being is indifferent to the difference between abstract and concrete, essential and unessential, self-subsistent and non-self-subsistent. This indifference stands in
% --- p. 436 ---
\opage{436}decisive conflict with the idea's demand; for the idea is and wills to be, not the abstract, but the absolutely concrete, not the unessential, but the absolutely essential, not the non-self-subsistent, but the absolutely self-subsistent: being is inadequate for the idea. What being lacks, existence possesses. By the category existence there is expressly aimed at a subsistence for itself, a particular, self-subsistent determinate being: that which exists is not merely something existent, but a concrete singular essence subsisting for itself. And yet the category of existence has, immediately understood, side-determinations that do not let themselves be united with the idea as idea. A thing exists just because it is a thing: the thing and that which exists are identical expressions. To make the idea into something that exists is, in other words, to make it into a thing; but the idea is just as little a thing as the thing is an idea.

So long as being and existing fall apart from one another in a finite manner, neither of these categories will satisfy the idea's claim to all-sided determinateness; in order to become adequate, they must become identical; in order to become identical, they must surrender side-determinations for a higher category. The higher category, enclosing being as existing and existing as being, can be designated by an \emph{absolutely reflected subsistence}.

In the absolutely reflected subsistence being is present; for that which subsists is something that is. In the absolutely reflected subsistence existence is likewise present; for that which subsists is something self-subsistent, and the self-subsistent is something existing in and for itself. The formerly inadequate being has now, by being taken up into an absolute subsistence, become adequate, for it has come to run upon decisive self-subsistence; the formerly inadequate existence has likewise, after having given up its immediacy and been transformed into absolutely reflected subsistence, become adequate, for the idea has its subsistence, its absolute self-subsistence, in absolute reflecting.

But how are we now to determine the absolute reflecting? According to our foregoing, as an infinite energy! The idea \emph{is},
% --- p. 437 ---
\opage{437}that is to say: it \emph{subsists} in infinite energy; the idea \emph{exists}, that is to say: its infinite energy is one with \emph{its absolute self-subsistence}. The infinite energy is, subjectively seen, an infinite \emph{comprehending}, \emph{judging}, \emph{syllogizing}. In the infinite comprehending, judging, syllogizing the idea is at one and the same time both the comprehending, judging, syllogizing subjectivity and also the content of the concept, of the judgment, of the syllogism: the energy is that the idea's self-separation rests upon its absolute self-unity, and this again upon the former.

Subjectivity has in its threefold being, as thinking, intuiting, willing, a threefold existence: the absolute subjectivity exists in three existences. That the three existences are in the matter itself more than three modifications of one and the same existing can be seen from subjectivity's threefold content. The threefold content is expressly determined by the three absolute ideas: the true, the beautiful, the good. The thinking subjectivity exists with the energy of the universal in eternal truth, the intuiting subjectivity with the energy of the particular in eternal beauty, the willing subjectivity with the energy of individuality in the eternal good. Here there are in the idea three centers, three regions: the ideal world (\textit{κόσμος νοητὸς}) is originally divided into three realms. But the threefold existence is, in the syllogism-unity of subjectivity and the ideas, nonetheless one subsistence reflected into itself, one energetic being, one single infinitely concrete existing.

The ideal world with its infinite system of ideas and idea-syllogisms is in its reflected subsistence a relative-absolute existence; it is by the energy of objectification determined as the idea of objectivity, the idea of the ideas in objective form, that is to say as \emph{world-idea}. As the word of the Godhead the world-idea is one with and yet different from the Godhead itself. It is one with the Godhead, for in all objectification of another subjectivity originally objectifies itself. It is different from the Godhead, indeed opposed to the Godhead, as
% --- p. 438 ---
\opage{438}otherness in its absolutely pervasive relativity can be opposed to the original selfhood that determines itself in all and all in itself.

As the idea of objectivity the world-idea has its selfless selfhood in the principle of relativity: the objectivity-idea is in its objective essence the idea of otherness. Herein lies it, that the one original world-idea must have its reflected subsistence in systems of relative ideas. According to the ideality of the power-word the systematic unity of the relative ideas constitutes \emph{the ideal system}; according to the reality of the power-word it constitutes \emph{the real system}. The ideal system now furnishes the idea-syllogism's subjective extreme, the real system its objective extreme. Since, however, the idea-syllogism's subjective extreme is itself a syllogism, an \emph{ideality-syllogism}, it follows from this that the whole real system with its power-determinations and antilogical moments must already, in an ideal-concrete manner, be anticipated in the ideal idea-syllogism.\footnote{On the original opposition and teleological unity of the ideal and the real systems cf.\ Grundideernes Logik I. p.~216--229.}

For the all-sided determination of the ideality-syllogism, the ideal idea-syllogism, it is thus necessary that the totalities in the objective extreme be transparent in the subjective: the objective extreme's real totalities are anticipated and, by the anticipation, co-determining for the idea-syllogism's subjectively absolute content. If the real system's co-determining reflexes are thought away, the ideal system is only an empty abstraction. The ideality-syllogism is at one and the same time extreme and syllogism. As syllogism it has, by integrations from reality, the whole world-idea for its a priori content; as extreme it has the infinite richness of its content in absolutely reflected subsistence, in a being identical with eternal working. The anticipated real system is by the ideality-syllogism determined at one and the same time both as \emph{possibility} and as \emph{outer actual%
% --- p. 439 ---
\opage{439}ity}, as possibility by virtue of the anticipation, as outer actuality in the character of objective extreme.

By the objective extreme's syllogistic conjunction with the subjective the world-idea, the world of the ideas, is infinitely concrete, an actual world \textit{ad intra}. But since the anticipation of the objective extreme has validity only by virtue of the extreme's own objective reality, the world of the ideas can have subsistence in energetically reflected existence only under the inviolable presupposition that the real world of things has an outward immediate and material existence. In the inward-going, reflected existence all becoming in all activities, formations and movements bears the stamp of \emph{unalterable being}: the ideal existence expresses the eternal. In the outward-going, immediate and material existence, on the contrary, all being and determinate being takes on the stamp of finitude, parcelling-out, outward alteration, \emph{unceasing becoming}: the real existence expresses the temporal.

As necessary as it now is that a syllogism's extremes --- far from being indifferent to one another --- must as extremes have their subsistence in mutual opposition and reciprocal tension, just as necessary is it then also that the ideal and the real existence must assert their validity by opposition and tension.

But how is the ideal existence now to be determined by its decisive opposition to the real existence? As eternity-existence, thus as pre-existence! In opposition to the ideal existence the real existence in temporality and becoming is then to be apprehended as a derived existence, an existence in which every existing term has its finite beginning. In the tension between the ideality-syllogism's two extremes lies ``the birth of time''.\footnote{On the theologians' quackeries in determining the beginning of time and of the world cf.\ D.\ Strauß. Den christelige Troeslære. Oversat af H.\ Brøchner. Første Bind 1842. p.~557--71.}

% --- p. 440 ---
\opage{440}Every real existence begins in time; but time itself does not begin. If the real existence presupposes a disposition, then the disposition again presupposes a pre-existing idea. \emph{The pre-existence is partly absolute, partly relative.}

In opposition to the real existence's manifold worlds and edifices of the world the eternal world-idea has \emph{an absolute pre-existence}. But the edifices of the world that are in actuality different could nevertheless impossibly, if it is true that each singular one is to develop out of its principle according to its type, proceed from the absolutely pre-existing world-idea, if particular ideas, answering to particular worlds, did not by a principial disjunction develop out of the fundamental idea itself: the ideas so developed have in relation to the particular worlds of the real development not an absolute but \emph{a relative pre-existence}.

How it must originally stand with the relativity of the pre-existence and with the principial disjunction is seen clearly enough from the relation between \emph{idea and disposition}. That an edifice of the world develops out of a disposition, that a terrestrial globe develops out of a disposition, that an organism develops out of a disposition, means, according to our foregoing, that the edifice of the world, the terrestrial globe and the organism each in its own way actualize an idea. What the disposition is in objective being, that the idea is in subjective reflecting; the idea is thus one with and yet different from the disposition. The unity-in-difference is determined by this, that the disposition is the idea dwelling in the real beginning, the idea the disposition completed by preformation and pre-existing in the anticipation. The identity of disposition and pre-existing idea is, when it is more closely determined, a categorical expression for the ideality-syllogism's original middle. In the present connection it is especially a matter of making clear what essentially lies in this middle.

In the first place, then, the middle expresses that the extremes must hold good by their mutual opposition, for if the opposition were not, then the middle would surely be superfluous. In the
% --- p. 441 ---
\opage{441}second place the middle expresses that the mutually opposed extremes must be identical; for if the identity were not, then the middle would be without significance. We will now begin by reflecting upon the opposition by itself and seeing what follows from it. The subjective extreme will then, by its opposition to the objective, present freedom in a one-sided manner, while the objective from its side will in a just as one-sided manner present necessity.

According to the freedom-extreme the real system's principial totalities with their dispositions would then have to be immediately referred to relatively pre-existing ideas, and each singular one of these ideas be straightway determined by imperative functions of the divine subjectivity. Thus, as regards e.g. the organism and its types, not only the emergence of organic life upon the terrestrial globe, not only the typical fundamental difference between plant- and animal-organism, but also every one of the typical forms within the plant- and the animal-kingdom would have to be referred singly and immediately to a divine act of will. This is the extreme in the subjective. Whether one otherwise prefers to represent the matter to oneself in such a way that the divine subjectivity finds the pre-existing types already present in its realm of ideas and then, with these in view, forms the different dispositions of real existence, or prefers to hold that the divine will first brings forth the paradigms and thereafter the corresponding dispositions, is wholly indifferent to science; for thought fails: the concept has gone out, the objective extreme is violated.

According to the necessity-extreme not only would the pre-existing ideas fall away as idle, deuteroscopic fantasies, but even the dispositions characteristic of the real system would necessarily have to give up their principial significance and transform themselves into bare possibilities; thus a world of possibilities without disposition! It is incredible what exertions the science of experience, indeed even philosophy, has made in order to come to terms with such a world of possibilities without disposition.
% --- p. 442 ---
\opage{442}Where the disposition is lacking, there must naturally be peculiar circumstances and favorable conditions to cover the lack. Thus organic life upon the earth is by no means to have its origin from a divine disposition, but solely from a possibility that supervenes under peculiar circumstances and favorable, especially favorable, conditions. Since now no possibility can be actualized without sufficient conditions, while every sufficiently conditioned possibility is always actualized with necessity, necessity reigns. This is the extreme in the objective. How one otherwise will picture to oneself the peculiar circumstances and favorable conditions, whether one will grant the divine knowledge a contemplative role restricted to indicative functions alone, or straightway reject every representation of a divine knowledge, is indifferent, wholly indifferent; for the concept has gone out, the subjective extreme is violated.

By their mutual opposition the extremes are thus posited as one-sided freedom against one-sided necessity. Hereby the middle, the totalized identity, is expressly determined. That freedom in the abstractly subjective is one-sided means, in other words, that it must have its real possibility in the objective necessity; that necessity in the abstractly objective is one-sided means that it must have its essential ground in the subjective freedom. Freedom in the ideal system is thus to be integrated by necessity in the real system, and conversely: here is the middle. Reality weighs with necessity's pressure upon freedom; but by self-determining reaction freedom gets the better of necessity. It is necessity's pressure and subjectivity's reaction that give the syllogism's middle its peculiar tension. Upon this tension rests the teleological unity of the real and the ideal system; upon this unity rests the reciprocal relation between real and ideal disjunction; upon this reciprocal relation rests the idea's divine productivity.

% --- p. 443 ---
\opage{443}In a real world with possibilities without disposition nature is without productivity. As little as the ``harvestless'' sea can be said to have productivity because, on condition of wind and storm, it can bring forth high waves, just as little will a real world with possibilities without disposition be able to be said to have productivity because, on condition of substances with determinate melting-points and a changing temperature, it can bring forth solid, fluid and gaseous masses. Only with the principial disposition does productivity come. Water's transition from the drop-fluid state to ice cannot in the categorical sense be called productive bringing-forth; in the emergence of organic life, in the coming-to-be of every new organic type, nature on the contrary displays an essential productivity. Every productive form is determined by its peculiar disposition; every peculiar disposition is again determined by a real disjunction: either plant or animal, either mollusc or articulate etc. The productive transitions rest upon sheer totality-separations, sheer principial disjunctions. In such disjunctions nature utters itself productively.

And yet productive nature does not, properly speaking, speak out of itself, as little as any language either speaks out of itself or itself speaks. No, it is subjectivity that speaks in and by language, in that it speaks out of itself: the divine subjectivity speaks in and by nature, in that it speaks from itself. The real system's productivity is at all stages of the formations an utterance of the word of power, a categorical power-language in disjunctive judgments and disjunctive syllogisms: the principial disjunction is real.

But the principial disjunction could impossibly be real if it were not originally ideal; it could impossibly take place in the real system if it did not originally take place in the ideal system; it could impossibly enter into necessity if it did not originally spring from freedom. Were the divine selfhood not infinitely productive, it would be weighed down by otherness and stifled by objectivity.

% --- p. 444 ---
\opage{444}It lies in selfhood's essence that it can be determined by another only in that it determines itself; it lies in the word of power that it could impossibly be real if it were not originally ideal; it lies in the disjunction that it cannot have reality in a disposition if it does not have ideality in a pre-existing idea. Upon the relation between the determined disposition and the determining idea rests the relation between the ideal or inward-going and the real or outward-going productivity. The pre-existing idea, the particular word of the power-ideality, is not idle: it bears the objective totality's disposition, in that it is itself borne by subjectivity; it effects the disposition's real development in the objective and determines the word in the categorical power-language of bringing-forth, in that it receives from subjectivity its content-moments, answering to the antilogical power-determinations of the real development, developed in reflected subsistence.

The outward-going productivity answers to necessity, the inward-going to freedom; but in the middle of the extremes --- and the middle is here the decisive thing --- the ideal productivity is necessary by virtue of the real, and the latter free by virtue of the former. But all productivity expresses, according to its concept, that subjectivity at one and the same time both acknowledges and masters its objective. In productivity freedom is necessary, and necessity free; but since the real productivity depends upon the ideal as the determined disposition upon the determining idea, freedom has in the free necessity necessity itself under it.

That the earth's development, geologically seen, points to a disposition means then, logically understood, that it takes place in consequence of a pre-existing idea; here one may thus with full right speak of \emph{the earth-body} itself as elementary organism, of \emph{the earth-body's disposition} and of \emph{the earth-body's} idea. The development's beginning is preformed by this, that the disposition in the totality of the objective elements is all-sidedly determined by the idea reflected in subjectivity's productive energy.

% --- p. 445 ---
\opage{445}If the disposition were not borne by the idea, it would vanish as a pure contradiction; for what a wonderful demand is there not laid in a disposition! In the disposition the whole is, essentially seen, to be finished from the beginning, and yet the disposition requires precisely that an actual development must take place: in the earth's original disposition there lies at one and the same time everything earthly, and yet this earthly all is first to come forth successively, and to that extent to come to be \emph{by degrees} in the course of the development. This contradiction can be resolved only by this, that the same content is at one and the same time both \emph{outwardly}, thus in the reality of the power-word as disposition, and \emph{inwardly}, thus in the ideality of the power-word as pre-existing idea.

``But,'' someone might ask, ``does not the same contradiction now arise in the idea as in the disposition? What is gained, after all, by moving the contradictory demand out of the disposition and into the idea?'' We must on this occasion impress it that the pre-existence is categorically determined as subsistence in energetic reflecting, as reflected subsistence: in reflected subsistence the idea has all the moments of the content at one and the same time; the energetic reflection is in this subsistence anticipation, and in the anticipation that which is actually to come must surely be essentially present. What the anticipation is in the idea, that \emph{the drive} now is in the disposition. The outward-going productivity, the actual development of the earth, shows us only the drive, not the anticipation; and yet all formative drive in the objective is only a word of power that confirms the energy of anticipating in the subjective.

From the energy of anticipating it is then further seen that the development of totality in the outward-going productivity must point for point answer to the idea-development itself in the inward-going. The pre-existing earth-idea has indeed from the beginning all the moments of development in reflected unity, and the real development can to that extent not enrich the idea with anything whatever; but since the idea-moments are themselves anticipated, the real development in the outer totality, in the earth-body itself, must neces%
% --- p. 446 ---
\opage{446}sarily be the condition for a corresponding ideal development in the earth-body's idea: the executive functions reflect themselves in the anticipating ones. That the earth at a certain stage of development attains, as it is said, a certain maturity means then, logically speaking, that the pre-existing earth-idea has attained a certain determinate energy.

Whence then, at the terrestrial globe's relative stage of maturity, do the new dispositions come, the dispositions namely to plants and animals? Naturally from a preceding idea-development. In relation to the paradigms for organic life the earth-idea is pre-existing; in relation to the organic individuals the life-ideas are pre-existing. The free necessity with which the typically different life-ideas step forth out of the earth-idea has in the foregoing already been determined by the doctrine of the immanent subsumption- and disjunction-syllogisms and the individuations answering to them. We have still to consider the reciprocal relation between the individuations and the fundamental idea's infinite pregnancy.

It lies in the nature of the matter that the world-idea, the absolute idea of otherness, of objectivity and of relativity, must be infinitely pregnant. The otherness-form is finite; the content infinite: in every one of the finitely determined forms the content thus yields an infinite surplus. This disproportion persists at all the intermediate stages of formation, until the otherness-form can again be taken up into the selfhood-form. In the first syllogism, the syllogism fundamental for the edifice of the world, which according to the ideality of the power-word has the fundamental idea itself for its ideal, and according to the reality of the power-word this determinate actual world for its material extreme, there lies the first intermediate idea in the form of particularity, the idea of \emph{this kind of} edifice of the world. This intermediate idea, derived from the fundamental idea but determining for the actual world's disposition, an anticipation-idea which over against its own determinate real world has relative pre-existence, now assumes the determinate form of particularity, thus a form with determinate limitation, determinate boundary; it is then with this
% --- p. 447 ---
\opage{447}form to be an individuation of the fundamental idea, is to hold its whole content, thus an infinite content in a finite form: here there must come to be an infinite surplus, the particular idea must be infinitely pregnant.

With the fundamental syllogism the form of productivity is peculiarly bounded. However often the same kind of syllogism be repeated, however many finite edifices of the world be added to those already given: the productivity in the mechanical form will constantly suffer from the one-sidedness that the idea, as well after its ideal as after its real development, remains in otherness and, by virtue of otherness, in finite opposition to subjectivity.

That the subjectivity which is master of its content masters otherness precisely thereby, that otherness is driven to an extreme, rests in selfhood's deep upon productivity's inexhaustibility. But the inexhaustible productivity rests again upon this, that a higher form constantly relieves a lower, until the otherness-form at last is resolved into selfhood's adequate form. In the real system every new form of productivity is determined by a new disposition, thus in the ideal system by a new idea. The new idea, decisive for the new disposition, is not created out of nothing, nor does it rest upon a divine whim, but proceeds by ideal disjunction out of an already given foundation-idea.

The foundation-idea can expressly be determined as the immediately preceding syllogism's \textit{terminus medius}. If we consider the syllogism: \textit{A} (the world-idea in its generality) --- \textit{S} (the particular world-idea) --- \textit{E} (this actual, determinate solar system), then the particular world-idea, which furnishes the syllogism's middle, is precisely that foundation-idea out of whose deep the idea for the new productivity, the planet-idea answering to the system, proceeds. The new productivity beginning with the planet-idea now forms the extreme for the new syllogism: \textit{A}\textsubscript{2} (the planet-idea in generality) --- \textit{S}\textsubscript{2} (the particular planet-idea,
% --- p. 448 ---
\opage{448}the earth's idea e.g.) --- \textit{E}\textsubscript{2} (this actual determinate earth). By this syllogism productivity already makes an essential advance, for otherness now begins to enter into selfhood: as an idea realizing itself the actual earth is an elementary selfhood.

That the form of productivity, however, cannot be brought to a close with the syllogism \textit{A}\textsubscript{2}---\textit{S}\textsubscript{2}---\textit{E}\textsubscript{2}, that \textit{S}\textsubscript{2} must, so to speak, teem with surplus, is a straightforward consequence of the selfhood-drive lying in otherness. For the selfhood-drive in the new form \textit{S}\textsubscript{2} must then again become foundation-idea; in other words: the new disposition in the earth-body, the disposition to organic life, must presuppose a pre-existing life-idea developing itself out of the earth-idea. With the syllogism: \textit{A}\textsubscript{3} (the life-idea in earthly generality) --- \textit{S}\textsubscript{3} (the life-idea in earthly particularity, plant- and animal-life) --- \textit{E}\textsubscript{3} (the life-idea realizing itself in the individual), and the systems contained in this system, selfhood at last wins the victory over otherness, and productivity attains its completed form.

What otherness is in the outer infinity, that selfhood is in the inner; from this it is seen that the form of productivity can first be completed with the selfhood developed into a self-conscious self; only self-consciousness with its inner infinity can be the adequate form of the infinite productivity: the starry heaven is too small to hold what can dwell in a human soul.

How necessity in the infinite productivity unites itself with freedom, in what relation knowledge's indicative functions, which let everything rest upon objective power-consequences, originally stand to subjectivity's imperative functions, which let everything depend upon self-determination, now shines forth from our foregoing.

In stepping forth out of its ideal foundation the new idea steps forth according to a principial disjunction. By a principial disjunction the life-idea has proceeded from the pre-existing
% --- p. 449 ---
\opage{449}earth-idea, this again from the pre-existing planet-idea, etc. The principial disjunction is in its energy essentially conditioned by a corresponding subsumption. Seen outwardly, the subsumption presupposes the already existing individual (this existing thing is a horse, this existing thing is an elephant); the subsumption goes from existence to category, from individual to idea.

It is otherwise when the ideality itself, the inward, is here looked at. That energy of the anticipation upon which the disjunction, the original separating-out of the new idea, each time rests cannot indeed posit the two terms of the syllogism without presupposing the third, cannot possibly posit the foundation-idea as the universal (the genus), and the new idea as the particular (the species), without presupposing the individual as individual; but the presupposition is here an anticipation. By the anticipation-energy it is thus just as much the species and the genus that become determined by the individual as conversely the individual that becomes determined by species and genus. If now individuation requires, as we have shown above, that the immanent syllogism's three terms must in the principle itself have the same originality: then it follows with necessity from this that the individual in which the species is to be realized must by anticipating subsumption be an individual-idea co-determining for the species-idea: only in reflection of its individual-ideas is the species-idea a paradigm.

According to the subsumption the pre-existing idea becomes determined by objective considerations, that is to say: it becomes determined as disposition; here there is, e.g. at this stage of the terrestrial globe's development, a possibility of organic life, and thus this objective possibility is now in and by the life-idea to be transformed into a disposition. But life upon the earth is impossible without a living shape; the living shape must be an earthly-individual shape, according to the objective conditions, either this or this shape. Once the ideal disjunction has the possibility of several dispositions, the energy of the anticipation posits this, precisely this
% --- p. 450 ---
\opage{450}shape in individual determinateness; the presupposed individual can now be subsumed under its species, and the species-idea pre-exist as paradigm.

Insofar, then, as it is the subsumption conditioned by the real system that motivates the ideal disjunction, the subsumption indeed points to necessity, the ideal disjunction to freedom; but since the subsumption is just as originally motivated by the disjunction as the latter by the former, freedom goes together with necessity: freedom posited in necessity's form is disposition; necessity posited in freedom's form is pre-existing idea. Unity of disposition and pre-existing idea is a syllogism-unity, that unity of the idea-syllogism which completes ideality.

\textit{ββ)} \emph{Reality-Syllogism.} \emph{The Real Existence of the Ideas.} In the ideality-syllogism, the idea-syllogism's subjective extreme, otherness is a moment transparent in selfhood; reality is here a reflex in ideality; the word of power is transfigured in the form of knowledge; productivity \textit{ad extra} is there only as presupposition for productivity \textit{ad intra}. If we now attempted to hold fast the subjective extreme alone by itself, thus not as extreme, but as a syllogism self-subsistent for all eternity, then it would undeniably be an abstraction, and this abstraction would give us precisely that relation in God which Hegel\footnote{Philosophie der Religion. Zweiter Theil. Berlin 1840. p.~248.} has designated as ``a play of love with itself, in which it does not come to serious otherness, to separation and discord''. If the subjective extreme is to have validity, it must originally have its objective extreme, precisely that extreme in which it does come to ``serious otherness, to separation and discord''.

If speculative dogmatics knew the idea-syllogism, it would be able to spare itself all objective trouble over deriving nature's corruptibility from the Devil's principle and the Fall%
% Ellipses: the Danish of these pages prints a literal "..."; set as \dots\ here,
% as in the rest of the English.
% Afslutning is rendered "closure" throughout (as at pp. 394--95); Nielsen's own
% term Conclusion is what "conclusion" renders elsewhere in this volume.
% --- p. 451 ---
\opage{451}. Of course, the derivation is indeed of a piece with it. ``It
is'', it is said in the objective treatment of ``Death and the
Corruptibility of the Whole Creation'', ``difficult to deny that there are
in nature phenomena which themselves raise the problem of the incursion of
evil into nature. We shall not merely appeal to the fact that what stirs
in the darkness of the human soul hisses through the foliage of the forest
in the coiling of the serpent, and howls through the desert in the
bloodthirst of wild beasts. But we shall point in particular to the fact
that there is revealed in nature an enigmatic contradiction against
nature's own inner teleology, a contradiction against nature's own inner
purposiveness \dots\ This disturbance of nature's own ends we are unable
to cognize as grounded in a true necessity; natural phenomena which in
themselves bear a disturbing principle we are unable to acknowledge as
normal''.

It does indeed look solemn at a distance when the theological inquirer
takes the non-theological logician to task and submits it to him to
consider whether he ``is confident that he can show that the redemption of
fallen human nature does not also comprise the redemption of thought, that
this is not Christianity's intention, to restore true humanity completely,
but to let essential elements --- and thought surely belongs among the
essential elements --- remain outside its sanctifying and perfecting
activity''. Yes, in such language the talk of the redemption of thought
even sounds edifying; but is it not strange all the same? Every time the
``redeemed thought'' really has to set about solving a problem, how very
poorly the scientific redemption fares! When a Jacob Bøhme, although he
acknowledged ``that a No was required if the Yes was to be able to reveal
itself, and that opposition was necessary if life was to be felt'',
nevertheless did not believe he could explain to himself ``the wondrous
regime of present nature, the palpability and
% --- p. 452 ---
\opage{452}obscurity of matter, the struggle and strife of all elements
and creatures against one another'', without the presupposition ``that
Lucifer's arrogance had set fire to the pure etherial essence of the
original creation, whose slag forms the present world'':\footnote{See
D. F. Strauß. Den christelige Troeslære. Oversat af H. Bröchner.
Kbhvn.\ 1843. p.~13.}
{} then it is the task of science in such a mysticism to distinguish what
belongs to imagination and what belongs to thought. But upon such a
distinguishing the thought redeemed according to the direction of
dogmatics will hardly enter. But where then is the sign of thought's
redemption? For that the dogmatician is unable to cognize the natural
disturbance of nature's own ends as ``grounded in a true necessity'' can
surely not be set up as a proof of thought's redemption.

If we regard the actual nature present to our senses as the objective
extreme of the idea-syllogism, an extreme which is itself a syllogism,
indeed, as has been shown in our foregoing, a system of syllogisms; if we,
as regards this reality-syllogism with its systems, lay due weight both on
\emph{the natural otherness} and on \emph{the natural selfhood}, then it
may well succeed in cognizing the natural disturbance of nature's own ends
as grounded in true necessity, and in interpreting the aesthetic-ethical
parallel between that ``which stirs in the darkness of the human soul''
and that which alarms us in the hissing of the serpent and the howling of
the beast of prey.

As regards, then, in the first place, otherness, its elementary character
can come into strife with selfhood and mark corruptibility only at the
point where life begins, hence first in the plant organism. If the
terrestrial globe is regarded by itself, even the greatest upheavals will
mark no disturbance: for elementary selfhood, hurricanes and earthquakes
are just as adequate as sunshine and dead calm. In relation to the
organism, by contrast, there begins the tension between
% --- p. 453 ---
\opage{453}principle and conditions, between idea and actuality. ``In the
plant's natural shape the unity-in-difference of the whole and its parts,
the inner and the outer, the total individual and its partial individuals
\dots\ is all-sidedly determined by the articulating opposition of
selfhood and otherness \dots\ If we ask which of the two powers is the
stronger, selfhood or otherness, a dialectic of power and impotence will
show itself on both sides. Essentially understood, selfhood, which in
principle bears otherness within itself, is the original, the higher, and
to that extent the more powerful; otherness, which has its integrating
element outside itself, is by contrast the derived, the lower, and to that
extent the impotent, the weaker. But if now the emphasis falls, as
actuality likewise demands, on the phenomenon, the factual, the immediacy
that strikes one in the fact, then the relation reverses itself, then
otherness, which has its strength precisely in the outward immediate, must
necessarily show itself as the original, the stronger, and that not merely
\emph{although} it is, but \emph{just because} it is that which according
to its essence is the derived, the weaker; likewise selfhood, whose
validity rests on reflection-energy, must necessarily present itself as
the derived, the abstracted, and that not merely \emph{although} it is,
but just \emph{because} it is that which at bottom is itself the original,
that which in objective reflection is eternal and
incorruptible''.\footnote{Forelæsninger over ``Phil.\ Propæd.'' fra
Universitetsaaret 1860--61. p.~317--18.}

Corruptibility, which shows itself in many ways in the plant world, is
raised in the animal world to the higher power of suffering, pain and
death. This happens not by any intervention of Lucifer, but by ``a true
necessity''. The ensouled individual is in its immediate self-sensation
its own end; for this end the actual surrounding world must furnish means
and conditions, and thus be determined as the subordinate. But the
surrounding world, which itself lies subject to elementary
% --- p. 454 ---
\opage{454}otherness, resists the self, and so the conflict begins.
``Every animal species is, like every plant species, determined to live in
its own way. But the ensouled lives on higher terms than the unensouled;
the unensouled lives only insofar as it is, in all immediacy, a living
thing; but the ensouled is a living thing which expressly lives its own
life. To live one's own life for oneself is to \emph{try} life; to try
life is, in the immediate sense, to \emph{enjoy} and to \emph{suffer}
\dots\ The goal is not enjoyment of life for itself, as little as it is
suffering; for the same thing that makes it possible for the individual to
enjoy also makes it possible for it to suffer; actual enjoyment
presupposes the possibility of an actual suffering, just as conversely
actual suffering presupposes the possibility of a corresponding
enjoyment''.\footnote{Anfr.\ Skr.\ p.~392.}

Already the struggle thus characterized between selfhood and elementary
otherness will be sufficient to explain how it comes about that there can
be revealed in nature ``an enigmatic contradiction against nature's own
inner purposiveness.'' ``But'', the dogmatician with ``the redeemed
thought'' continues: ``natural phenomena which in themselves bear a
disturbing principle we are unable to cognize as normal \dots\ A theodicy
which would here reassure us with the general representation of finitude
and of life's necessary development through oppositions would only evade
the problem. For both true and false development moves through
oppositions, and the question is precisely whether all the oppositions
that natural life shows us are normal''. We answer ``the redeemed
thought'' that the objective extreme of the idea-syllogism naturally is
and remains normal as extreme; but that it is a contradiction, indeed a
gross contradiction, to demand of an extreme what can first be attained in
a middle. Let ``the redeemed thought'' for once make a serious scientific
attempt at throwing the real extreme away, and it will
% --- p. 455 ---
\opage{455}soon see where, despite all declamations about the normal and
the abnormal, the true and the false opposition, it finally ends up.

The strife between selfhood and otherness, growing in the animal world,
reaches its maximum in human life. Inasmuch as consciousness, which in the
animal is only immediate sense-consciousness, doubles itself in the human
being into self-consciousness, by this doubling not only joy and enjoyment
but also pain and suffering are infinitely doubled: death, which merely
strikes the animal, comes to be for the human being. With the human
being's awakening self-consciousness there awakens spirit, the idea of the
will with freedom's demand in freedom's form. Natural selfhood is egoism;
the strife between selfhood and otherness that runs through the whole of
nature is then transformed in the human being into a struggle of freedom
between egoism and spirit. For spirit, egoism is a guilt; the intimation
of guilt is anxiety. Seen from this standpoint it must surely be evident,
without Lucifer's help, to everyone whose thought is not so ``redeemed''
that it has become quite loose, how it can come about that something which
hisses in the foliage and something which howls in the desert can be an
alarming reminder of something ``which stirs in the darkness of the human
soul''.

It is the attractive thing about original mysticism that it everywhere
sees nature in the light of spirit and spirit in the light of nature; but
when the reflection of the understanding begins, the ambiguity must go.
What shines as in holy twilight with the metallic lustre of primitivity
can neither be replaced nor reproduced by the aesthetic polish of a
superficial ``thought-redemption''. That a Jacob Böhme, by one of spirit's
strong flashes of lightning, can believe he sees the corruptibility of the
whole creation explained in the red glow of Lucifer's flames, we can
understand; we can be involuntarily gripped by what is profound in it. But
how feeble and ambiguous the explanation becomes when the ``redeemed
thought'', by objective grounds in an objective manner, seriously
undertakes to furnish
% --- p. 456 ---
\opage{456}proof that world-egoism is the cosmic principle, and the cosmic
principle the devil in becoming.

Revelation and science have, with all their opposition, this in common,
that both reject the ambiguous: an ambiguous revelation is just as great a
\textit{contradictio in adjecto} as an ambiguous science. In order to do
away with the ambiguity, one must from both sides make a separation
between nature and spirit. But when the separation has taken place, how
deep is the gulf that divides faith from knowledge! For it stands fast,
provided only one will not let science's ``redemption'' play conjuring
tricks with faith, that pain and death have been in existence for millions
of years before the human being appeared on earth. And it stands fast,
provided only one will not let faith's ``redemption'' play conjuring
tricks with knowledge, that death is the wages of sin, and nature's
corruptibility a guilt of spirit.

That the reality-syllogism is an extreme has beguiled believing
speculation, which in consequences of the immediate power of otherness has
thought to trace the works of the devil; that the reality-extreme is
itself a syllogism consisting in systems of syllogisms has beguiled
immanence-philosophy, which by letting the ideal be absorbed in the real
has taken a syllogism-extreme for the completed syllogism.

\textit{γγ)} \emph{Closure: The Completion of Idea-Existence.} With the
fully developed middle of the ideality- and the reality-syllogism, and
with the two extremes' syllogistic union as determined by the middle,
the idea-syllogism itself at last reaches its closure. The succession
hereby indicated is, however, grounded in a corresponding simultaneity;
the idea-movement is rest, rest itself a movement. For the case is not
that the perfect middle should now abolish the two imperfect extremes,
that totalized one-sidedness of ideality and of reality on which dualism
originally rests; no, on the contrary! as true and perfect middle it must
precisely let the extremes stand in force in all their one-sidedness.
Another
% --- p. 457 ---
\opage{457}matter is that the middle, which is itself a syllogism, a
system of syllogisms, indeed a system of syllogism-systems, can again
unfold extremes which in their perfection are themselves middle terms,
inasmuch as the perfect middle terms are themselves extremes.

It is \emph{otherness} (the limit, the real negation, the difference)
which hereby presents itself in a new light. In the ideal system otherness
is ideal, and the negation immanent; the differences are formed by the
reciprocal determinations which the word originally carries with it in the
infinite thinking, intuiting and willing of ontological subjectivity. In
the real system otherness is real, by the immediate manifestation of power
even sensuous and material. The differences lose themselves in superficial
diversities, temporal parcellings-out, outward and contingent relations.

If now it really were impossible, as has so often been asserted in ``the
philosophy of the concept'', that real otherness could be raised to a
higher level than that which characterizes this actuality present to us,
then the ideals could indeed, as it is said, never be attained, then idea
and existence would have to find themselves in eternal discord. But a
transformation of otherness, and therewith a sublation of corruptibility,
is, logically seen, so far from being an impossibility that it is on the
contrary an inescapable necessity. The discord between ideal and actuality
which is so characteristic of the objective real extreme is by no means a
sign of the idea's relative impotence; it is much rather a proof of the
absolutely overreaching power of idea-filled subjectivity.

A feeble idea naturally has feeble extremes with a weak middle and a
feeble syllogistic union; hence immanent speculation is always so
anxious about whatever might resemble ``a dualism''. The powerful idea,
divine in essence, by contrast has powerful extremes; its manifestation in
a vigorous dualism is only the presupposition of the infinite superiority
with which it carries the syllogism through; the completion of
idea-existence in concrete closure must necessarily have an all-explaining
% --- p. 458 ---
\opage{458}middle, a last \textit{terminus medius} that masters all
extremes. Since now the all-explaining middle must be a realm of ideas
just as much as each of the extremes, the three \textit{termini} of the
idea-syllogism have thus with inner consistency transformed themselves
into three distinct worlds: an ideal world, a real world and an ideal real
world.

``But'' --- sober criticism asks --- ``where are we getting to now? If
logic is really to occupy itself with illuminating the ideal real world,
that supernatural kingdom where the ideals are actual only because
actuality itself is ideal, then the closure of the syllogism-systems must
indeed end in a formal eschatology. But then an encroachment is after all
being made upon theology's well-earned rights; then the logician does
after all come at last to institute such considerations in the upper world
as only a speculative dogmatician, a mystically profound genius who in an
otherworldly ``cloister life'' had devoted himself solely and alone to
``eschatological considerations'', would be able to carry through in the
underworld!''

In this respect we can at once reassure criticism by referring to the
problem: ``Two Worlds''.\footnote{Forelæsninger over ``Phil.\ Propæd.''
fra Universitetsaar 1861--62. p.~123--137.}
{} The discussion of this problem ends precisely with an entirely negative
result. ``If we hold fast'', it is said, ``to the likeness between Here
and Yonder, then we are surprised by the unlikeness; if we hold fast to
the unlikeness, then the likeness again forces itself upon us. If we keep
to actuality, then it comes into conflict with the ideal; if the ideal is
itself to be actual, then we get two kinds of ideals: ideals which we
know, but which never become actual; ideals which we do not know, and of
whose actuality we can form no representation, let alone any clear
concept. But when I am not in a position to see into an ideal's
possibility, how should
% --- p. 459 ---
\opage{459}I then be able to assume its actuality? \textit{A non posse esse
ad non esse valet consequentia}''.

What a precarious undertaking it is in science to have to depict for
oneself, by analogies from this present world, the condition in the life
to come, can be seen from many examples; the theologians' attempts with
the objective doctrine of the resurrection are instructive in particular,
also in a logical respect. The thought objectively ``redeemed'' in
dogmatics has, namely, in union with the equally objectively ``redeemed''
imagination, taken many and various and characteristic runs at a
scientific explanation of the nature transfigured in the resurrection. The
attempts, carried out now with scholastic boldness, now with the
smoothing-over caution of a modern consciousness, now inaesthetically and
tastelessly, now again with tact and taste, all move on the basis of one
and the same contradiction, the contradiction of \emph{a natureless
nature}.

That philosophical criticism, every time it has occupied itself with
exposing what is insipid, flat and self-contradictory in the dogmaticians'
attempts, has been able not only to render \textit{præstanda} but to
celebrate triumphs, is understandable. But neither ought it to be passed
over in silence that criticism, dazzled by sundry cheap triumphs, has not
seldom criticized itself so far into finite contradictions of the
understanding that it has \textit{eo ipso} criticized itself out of the
idea. By proving the impossibility of making the condition in the world to
come intuitable, negative criticism has, in accordance with an easily
explicable self-deception, been brought to assume that it thereby proved
the impossibility of a world to come. But this is a precipitancy.

If the world to come is posited as identical with the ideal real world and
this as identical with the middle of the idea-syllogism, then no doubt can
any longer be raised about its validity; its necessity is given. Insofar
as what is at stake is then not to make its actuality intuitable in the
form of a condition --- for this is, as has been said, impossible --- but
to
% --- p. 460 ---
\opage{460}determine its concept as the necessary form for the completion
of the idea-syllogism, reflection can in this case rest precisely on the
familiar proposition of formal logic: \textit{a necesse esse ad posse esse
valet consequentia}.

There can as little be talk of any leap or outward transition from this
world to that as there can be talk of a leap or outward transition from a
syllogism's extreme to its middle. There is indeed a difference of sphere
between Here and Yonder, but in the same way there is also a pervasive
opposition and qualitative difference between middle and extreme. Just
because the opposition of extreme and middle reaches its highest point in
the idea-syllogism, the identity is all the stronger. It can be seen at
all principal points in the extreme that the one-sidedness, albeit in a
one-sided manner, is integrated from the middle: through the middle there
runs a continued movement from extreme to extreme.

If we now regard the three regions, the ideal and the real system as
extremes and the ideal real system as middle, then it is at bottom the
idea, the one and the same idea, which with all particular ideas runs
through all three regions. In the ideal system the idea has as idea its
pre-existence, in the real system its immediate real existence; but each
of these forms of existence is one-sided and in its one-sidedness
unsatisfying. Pre-existence lacks reality, real existence lacks ideality;
if these lacks could not be remedied, the idea would have to remain in
eternal strife with itself: strife and contradiction would then, in
meaningless meaning, be the principle of existence (πόλεμος πάντων πατήρ).

But in each of the extremes the middle is itself preformed. The ideal
disjunction by which the particular idea obtains its pre-existence
coincides precisely, as has been shown above, with the real disjunction by
which real existence obtains its disposition. The real originally works
upon the ideal in order that the ideal may work upon the real; there is in
the idea nothing
% --- p. 461 ---
\opage{461}so ideal that there must not be in its energy a moment, a
reflex of real immediacy, and there is in the idea nothing so real that it
does not have an ideality in its possibility. By the doubling that in this
way obtains in each of the extremes, the middle can gain power; for what
now is the middle? The all-sidedly determined existence of the idea, that
completion of existence which alone can give the extremes their
significance. The middle is goal, the goal is closure, closure is the
truth of all syllogisms; from this standpoint it can then be seen that the
functions of the idea-syllogism are essentially modal functions of truth.

\emph{The middle is truth.} The extremes indeed also have their truth, but
only in union with the middle. If the middle is thought away, then
otherness is on both sides inadequate, and the thought-syllogism E--A--S
(as movement A--(E, S)--A) is not in a position to satisfy the idea. In
the ideal extreme otherness is abstractly ideal, the concretion in E and S
lacks reality; but to lack reality is to lack truth. In the real system
otherness is one-sidedly elementary, the concretion in E and S lacks
ideality; but to lack ideality is to lack truth. In the middle, by
contrast, otherness is adequate, and the idea of truth satisfied.

\emph{The middle is beauty}; for beauty is the intuitable harmony of the
ideal and the real. In the ideal system, as regards the
intuition-syllogism E--S--A (as movement S--(E, A)--S), E is with all its
transparency so feeble that S precisely thereby comes to lack the striking
immediacy which the beautiful may never be without. In the real extreme E
has immediacy enough; but since otherness is elementary and, by its
elementary immediacy, contingent, it is in rupture with A and
non-transparent for S. From the intuitable disproportion between E, S and
A there arises the unbeautiful. In the middle, by contrast, otherness is
adequate; E has transparency enough to make A intuitable and immediacy
enough
% --- p. 462 ---
\opage{462}to take up, individualizingly, the shape S: the idea of beauty
is satisfied.

\emph{The middle is the good}. In the ideal extreme the will lacks its
necessary barrier, the actual opposition by whose sublation it proves its
energy. The will-syllogism S--E--A (as movement E--(S, A)--E) is empty,
for E is a shadow, otherness is only a reflex of selfhood. In the real
system otherness is strong enough to offer the will resistance; but now
the resistance is so elementary that the will and the barrier find
themselves in unceasing struggle. In unity with the middle, by contrast,
both extremes are good; the all-sided energy with which the will works in
and by the equilibrium of ideal and real otherness that characterizes the
middle is precisely conditioned by the one-sidedness with which it works
in the extremes.

What holds of the ideal ideas, the true, the beautiful, the good, in all
generality, holds naturally in a peculiar sense of the real ideas, the
foundation-ideas and the selfhood-ideas ascending from them. For the
individual existence of these ideas elementary otherness acquires a double
significance: first, in general, through the real extreme to which it
belongs; next, in particular, through the substrate it furnishes for the
individual idea's own real development. Thus e.g. the idea of the
world-structure has not merely a subsistence reflected in the subjective
extreme, a pre-existence belonging under the ideal system, but in the real
extreme, the objective extreme, it has at the same time a material and
immediate subsistence as this determinate world-structure (this solar
system). What now does the middle signify? The world-structure developed
and transfigured into all-sided existence! Here it is a matter of
attending to the double-sided role of the real extreme.

If the real extreme is apprehended individually in individual opposition
to the ideal, hence this existing world-structure in individual opposition
to its pre-existing idea, then the middle's significance is unmistakable.
If the middle were not,
% --- p. 463 ---
\opage{463}the idea in its pre-existence would lack reality, and immediate
real existence would lack ideality; between such extremes no reciprocal
relation could then be thought: both extremes would be without truth. That
the middle is the truth of the extremes thus means that the
world-structure can have an immediate and material existence only on the
presupposition that it at the same time has a real existence reflected in
its closure and idealized by the reflection. There can thus, to put it
plainly, exist a sensuous world-structure only on the presupposition that
a supersensuous one at the same time comes into existence.

But as pre-existing idea the world-structure has a beginning in eternity,
as real existence it has a beginning in time; the ideal extreme bears the
stamp of eternity, the real extreme that of temporality: how does the
eternal stand, in the closure, the middle of the syllogism, to the
temporal? If the beginning in eternity did not originally reflect a
determination of time, there could here be no talk of relative
pre-existence. Relative pre-existence indicates precisely a becoming of
the idea out of its ideal foundation, a coming-to-be marked by the
disjunction's \textit{punctum saliens}, a development of reflection
ascending in the principle, hence a succession within the eternal.
Conversely, real existence beginning in time could not possibly begin as
an actual disposition if there were not something eternal in the temporal
beginning. In the middle itself, ideal real existence, the temporal will
thus in an infinite manner convert itself into the eternal, and the latter
into the former.

But at this point an ambiguity enters. That the middle subsists as long as
the extremes subsist is a matter of course; but if now one of the extremes
perishes, what then becomes of the middle? That a supersensuous world
exists as long as this material world-structure exists must, if
% --- p. 464 ---
\opage{464}it is true that by the supersensuous there is here thought not
a new outward and immediate, but an essential, inwardly reflected
existence, indeed be conceded; but everything that arises in time must
perish with time; when, then, this sensuous world-structure one day
perishes: what then becomes of the supersensuous one? For the answering of
such questions it is necessary to set clearly before one's eye the
relation between selfhood and otherness that is decisive for the idea's
pregnancy and for infinite productivity.

In all regions, under all shapes, \emph{otherness has the double-sided
significance: to be at once both coordinated with and subordinated to
selfhood}. That otherness is \emph{coordinated} with selfhood, that there
could be no subjectivity if the objective were lacking; that the knowing
self could have no energy either in thinking, intuiting or willing if that
manifestation of power which gives otherness its reality were thought
away --- all this has been set forth in our foregoing. But that otherness
with all this is nonetheless, and must be, \emph{subordinated} to
selfhood; that subjectivity, which through its real opposition always
relates itself to itself, expresses its acknowledgment of the objective
precisely by impressing its seal upon it in infinite productivity, and
thus sealing its objectivity --- this too must then be assumed to have
been sufficiently set forth.

The contradiction concealed in otherness, that it is at once both
coordinated with and subordinated to selfhood, is resolved only by the
double role which the otherness-elements evidently play in the real
extreme.

That a real extreme is in all generality necessary in order to form an
opposition to the ideal extreme means that the ideal world would have to
vanish if this actual, material world-order could vanish as a whole; in
this
% --- p. 465 ---
\opage{465}respect the immediately real stands on a par with the ideal:
otherness is coordinated with selfhood. But as regards the individual real
existence corresponding to the determinate idea, elementary otherness is,
despite all its apparent independence, essentially subordinated to
selfhood and subjectivity.

Every idea, even the foundation-idea most deeply sunk in otherness, is
nevertheless according to its original essence a selfhood. No world-idea
could thus come to actual existence in the actual world-structure if the
idea in its real self-development were not in a position to master the
elements of otherness.

In the real extreme the elements of otherness can never --- this expressly
conflicts with the character of the real extreme --- be transformed into
real idea-moments. But without developing an indwelling system of real
idea-moments the idea could not possibly develop itself as idea. Since now
the real idea-moments, although from the beginning they must be
articulated by the material elements of otherness, are by no means one
with these: the otherness-elements can fall back again to the universal
real system, and the existing idea nevertheless preserve the ideal-real
result of the development and with this result be absorbed into the
closing middle. In this middle real otherness --- even if its reality is
otherwise ever so great, even if in the new system it should be able to be
released anew into elementary extremes, indeed, though it should thunder
in the midst of ``eternity's summer'' --- is nevertheless in every way so
merged with the ideal, so absolutely subjected to selfhood in every
respect, that idea-existence must with full right be said to be satisfied:
the supersensuous world is incorruptible.

\begin{center}---\end{center}

\end{document}
